\input{preamble}

% OK, start here.
%
\begin{document}

\title{Catégories dérivées des espaces algébriques}


\maketitle

\phantomsection
\label{section-phantom}

\tableofcontents

\section{Introduction}
\label{section-introduction}

\noindent
Dans ce chapitre, nous étudions les catégories dérivées de modules sur les
espaces algébriques. Il ne semble pas exister de bonne référence introductive
consacrée à ce sujet ; la littérature l'aborde en renvoyant à des articles sur
les catégories dérivées de modules sur les champs algébriques ; voir par exemple
\cite{olsson_sheaves}.



\section{Conventions}
\label{section-conventions}

\noindent
Si $\mathcal{A}$ est une catégorie abélienne et $M$ un objet de
$\mathcal{A}$, nous notons encore $M$ l'objet de $K(\mathcal{A})$
et/ou de $D(\mathcal{A})$ correspondant au complexe égal à
$M$ en degré $0$ et nul en tout autre degré.

\medskip\noindent
Si $A$ est un anneau, $K(A)$ désigne la catégorie homotopique
des complexes de $A$-modules et $D(A)$ la catégorie dérivée associée.
De même, si $(X, \mathcal{O}_X)$ est un espace annelé, le symbole
$K(\mathcal{O}_X)$ désigne la catégorie homotopique des complexes de
$\mathcal{O}_X$-modules et $D(\mathcal{O}_X)$ la catégorie dérivée
associée.





\section{Généralités}
\label{section-generalities}

\noindent
Dans cette section, nous rassemblons quelques résultats généraux sur la
cohomologie des complexes non bornés de modules sur les espaces algébriques.

\begin{lemma}
\label{lemma-restrict-direct-image-open}
Soit $S$ un schéma. Soit $f : X \to Y$ un morphisme d'espaces algébriques
au-dessus de $S$. Étant donné un morphisme étale $V \to Y$, posons $U = V \times_Y X$
et notons $g : U \to V$ le morphisme de projection. Alors
$(Rf_*E)|_V = Rg_*(E|_U)$ pour tout $E$ de $D(\mathcal{O}_X)$.
\end{lemma}

\begin{proof}
Représentons $E$ par un complexe K-injectif $\mathcal{I}^\bullet$ de
$\mathcal{O}_X$-modules. Alors $Rf_*(E) = f_*\mathcal{I}^\bullet$
et $Rg_*(E|_U) = g_*(\mathcal{I}^\bullet|_U)$ d'après le lemme
de Cohomology on Sites
\ref{sites-cohomology-lemma-restrict-K-injective-to-open}.
Le résultat découle donc du lemme
de Properties of Spaces
\ref{spaces-properties-lemma-pushforward-etale-base-change-modules}.
\end{proof}

\begin{definition}
\label{definition-supported-on}
Soit $S$ un schéma. Soit $X$ un espace algébrique au-dessus de $S$.
Soit $E$ un objet de $D(\mathcal{O}_X)$.
Soit $T \subset |X|$ une partie fermée.
Nous disons que $E$ est {\it à support dans $T$} si les
faisceaux de cohomologie $H^i(E)$ sont à support dans $T$.
\end{definition}








\section{Catégorie dérivée des modules quasi-cohérents sur le petit site \'etale}
\label{section-derived-quasi-coherent-etale}

\noindent
Soit $X$ un schéma. Dans cette section, nous montrons que
$D_\QCoh(\mathcal{O}_X)$
peut être définie en termes du petit site étale $X_\etale$ de $X$.
Notons $\mathcal{O}_\etale$ le faisceau structural sur
$X_\etale$. 
Considérons le morphisme de sites annelés
\begin{equation}
\label{equation-epsilon}
\epsilon :
(X_\etale, \mathcal{O}_\etale)
\longrightarrow
(X_{Zar}, \mathcal{O}_X).
\end{equation}
noté $\text{id}_{small, \etale, Zar}$ dans
Descent, Lemma \ref{descent-lemma-compare-sites}.

\begin{lemma}
\label{lemma-epsilon-flat}
Le morphisme $\epsilon$ de (\ref{equation-epsilon})
est un morphisme plat de sites annelés. En particulier, le foncteur
$\epsilon^* : \textit{Mod}(\mathcal{O}_X) \to
\textit{Mod}(\mathcal{O}_\etale)$ est exact.
De plus, si $\epsilon^*\mathcal{F} = 0$, alors $\mathcal{F} = 0$.
\end{lemma}

\begin{proof}
La platitude du morphisme $\epsilon$ résulte de
Descent, Lemma \ref{descent-lemma-compare-etale-zariski-flat}.
Voici une autre preuve. Il suffit de montrer que
$\mathcal{O}_\etale$ est un $\epsilon^{-1}\mathcal{O}_X$-module plat.
Pour cela, il suffit de vérifier que
$\mathcal{O}_{X, x} \to \mathcal{O}_{\etale, \overline{x}}$
est plat pour tout point géométrique $\overline{x}$ de $X$, voir
Modules on Sites, Lemma
\ref{sites-modules-lemma-check-flat-stalks},
Sites, Lemma
\ref{sites-lemma-point-morphism-sites},
et
\'Etale Cohomology, Remarques
\ref{etale-cohomology-remarks-enough-points}.
Par \'Etale Cohomology, Lemme
\ref{etale-cohomology-lemma-describe-etale-local-ring}
on voit que $\mathcal{O}_{\etale, \overline{x}}$ est l'hensélisation
stricte de $\mathcal{O}_{X, x}$. Ainsi
$\mathcal{O}_{X, x} \to \mathcal{O}_{\etale, \overline{x}}$
est fidèlement plat d'après More on Algebra,
Lemme \ref{more-algebra-lemma-dumb-properties-henselization}.

\medskip\noindent
L'exactitude de $\epsilon^*$ découle de la platitude
de $\epsilon$ par
Modules on Sites, Lemma \ref{sites-modules-lemma-flat-pullback-exact}.

\medskip\noindent
Soit $\mathcal{F}$ un $\mathcal{O}_X$-module. Si
$\epsilon^*\mathcal{F} = 0$, alors, avec les notations ci-dessus,
$$
0 = \epsilon^*\mathcal{F}_{\overline{x}} =
\mathcal{F}_x \otimes_{\mathcal{O}_{X, x}} \mathcal{O}_{\etale, \overline{x}}
$$
(Modules on Sites, Lemme \ref{sites-modules-lemma-pullback-stalk})
pour tout point géométrique $\overline{x}$. Par la fidélité plate de
$\mathcal{O}_{X, x} \to \mathcal{O}_{\etale, \overline{x}}$
nous concluons que $\mathcal{F}_x = 0$ pour tout $x \in X$.
\end{proof}

\noindent
Soit $X$ un schéma. Nous conservons les notations de (\ref{equation-epsilon}).
Rappelons que $\epsilon^* : \QCoh(\mathcal{O}_X)
\to \QCoh(\mathcal{O}_\etale)$
est une équivalence d'après
Descent, Proposition \ref{descent-proposition-equivalence-quasi-coherent} et
Remark \ref{descent-remark-change-topologies-ringed-sites}.
De plus, $\QCoh(\mathcal{O}_\etale)$ est une sous-catégorie de Serre de
$\textit{Mod}(\mathcal{O}_\etale)$ d'après
Descent, Lemma \ref{descent-lemma-equivalence-quasi-coherent-limits}.
Nous pouvons donc définir $D_\QCoh(\mathcal{O}_\etale)$ comme la sous-catégorie
triangulée de $D(\mathcal{O}_\etale)$ dont les objets sont les complexes
dont les faisceaux de cohomologie sont quasi-cohérents, voir
Derived Categories, Section \ref{derived-section-triangulated-sub}.
Le foncteur $\epsilon^*$ est exact (Lemme \ref{lemma-epsilon-flat})
et induit donc
$\epsilon^* :  D(\mathcal{O}_X) \to D(\mathcal{O}_\etale)$
et, puisque les images réciproques de modules quasi-cohérents sont quasi-cohérentes,
également $\epsilon^* : D_\QCoh(\mathcal{O}_X) \to
D_\QCoh(\mathcal{O}_\etale)$.

\begin{lemma}
\label{lemma-derived-quasi-coherent-small-etale-site}
Soit $X$ un schéma. Le foncteur
$\epsilon^* : D_\QCoh(\mathcal{O}_X) \to
D_\QCoh(\mathcal{O}_\etale)$
défini ci-dessus est une équivalence.
\end{lemma}

\begin{proof}
Nous démontrerons ceci en montrant que le foncteur
$R\epsilon_* : D(\mathcal{O}_\etale) \to D(\mathcal{O}_X)$
induit un quasi-inverse. Nous utiliserons librement le fait que $\epsilon_*$
est donné par restriction à $X_{Zar} \subset X_\etale$ et par la description de
$\epsilon^* = \text{id}_{small, \etale, Zar}^*$
dans Descent, Lemme \ref{descent-lemma-compare-sites}.

\medskip\noindent
Pour un $\mathcal{O}_X$-module quasi-cohérent $\mathcal{F}$, le morphisme d'adjonction
$\mathcal{F} \to \epsilon_*\epsilon^*\mathcal{F}$ est un isomorphisme, car
le faisceau associé $\mathcal{F}^a$
(Descent, Définition \ref{descent-definition-structure-sheaf})
est un faisceau, comme le montre
Descent, Lemma \ref{descent-lemma-sheaf-condition-holds}.
Réciproquement, tout $\mathcal{O}_\etale$-module quasi-cohérent
$\mathcal{H}$ est de la forme $\epsilon^*\mathcal{F}$ pour un certain
$\mathcal{O}_X$-module quasi-cohérent $\mathcal{F}$, voir
Descent, Proposition \ref{descent-proposition-equivalence-quasi-coherent}.
Alors $\mathcal{F} = \epsilon_*\mathcal{H}$ d'après ce qui précède, et
nous concluons que le morphisme d'adjonction
$\epsilon^*\epsilon_*\mathcal{H} \to \mathcal{H}$ est un isomorphisme pour tout
$\mathcal{O}_\etale$-module quasi-cohérent $\mathcal{H}$.

\medskip\noindent
Soit $E$ un objet de $D_\QCoh(\mathcal{O}_\etale)$
et notons $\mathcal{H}^q = H^q(E)$ son $q$-ième faisceau de cohomologie.
Soit $\mathcal{B}$ l'ensemble des objets affines de $X_\etale$.
Alors $H^p(U, \mathcal{H}^q) = 0$ pour tout $p > 0$, tout $q \in \mathbf{Z}$,
et tout $U \in \mathcal{B}$, voir
Descent, Proposition \ref{descent-proposition-same-cohomology-quasi-coherent}
et
Cohomology of Schemes, Lemma
\ref{coherent-lemma-quasi-coherent-affine-cohomology-zero}.
D'après Cohomology on Sites, Lemme
\ref{sites-cohomology-lemma-cohomology-over-U-trivial}
cela signifie que
$$
H^q(U, E) = H^0(U, \mathcal{H}^q)
$$
pour tout $U \in \mathcal{B}$. En particulier, cela vaut pour les ouverts affines
$U \subset X$. Il s'ensuit que le $q$-ième faisceau de cohomologie de
$R\epsilon_*E$ sur $U$ est la valeur du faisceau $\epsilon_*\mathcal{H}^q$
sur $U$. En faisceautisant, nous obtenons
$$
H^q(R\epsilon_*E) = \epsilon_*\mathcal{H}^q
$$
ce qui montre en particulier que $R\epsilon_*$ induit un foncteur
$D_\QCoh(\mathcal{O}_\etale) \to D_\QCoh(\mathcal{O}_X)$.
Comme $\epsilon^*$ est exact, nous obtenons alors
$H^q(\epsilon^*R\epsilon_*E) = \epsilon^*\epsilon_*\mathcal{H}^q =
\mathcal{H}^q$ (d'après la discussion ci-dessus). Le morphisme d'adjonction
$\epsilon^*R\epsilon_*E \to E$ est donc un isomorphisme.

\medskip\noindent
Réciproquement, pour $F \in D_\QCoh(\mathcal{O}_X)$, le
morphisme d'adjonction $F \to R\epsilon_*\epsilon^*F$
est un isomorphisme pour la même raison, c'est-à-dire parce que
les faisceaux de cohomologie de $R\epsilon_*\epsilon^*F$
sont isomorphes à
$\epsilon_*H^m(\epsilon^*F) = \epsilon_*\epsilon^*H^m(F) = H^m(F)$.
\end{proof}










\section{Catégorie dérivée des modules quasi-cohérents}
\label{section-derived-quasi-coherent}

\noindent
Soit $S$ un schéma. Le lemme
\ref{lemma-derived-quasi-coherent-small-etale-site}
montre que la catégorie $D_\QCoh(\mathcal{O}_S)$ peut être définie
à l'aide de complexes de $\mathcal{O}_S$-modules sur le schéma $S$
ou de complexes de $\mathcal{O}$-modules sur le petit site étale
de $S$. La définition suivante est donc compatible avec celle
du cas des schémas.

\begin{definition}
\label{definition-derived-quasi-coherent}
Soit $S$ un schéma. Soit $X$ un espace algébrique au-dessus de $S$.
La {\it catégorie dérivée des $\mathcal{O}_X$-modules dont les
faisceaux de cohomologie sont quasi-cohérents} est notée
$D_\QCoh(\mathcal{O}_X)$.
\end{definition}

\noindent
Cette définition a un sens d'après
Properties of Spaces, Lemme
\ref{spaces-properties-lemma-properties-quasi-coherent}
et
Derived Categories, Section \ref{derived-section-triangulated-sub}.
Nous obtenons ainsi un foncteur canonique
\begin{equation}
\label{equation-compare}
D(\QCoh(\mathcal{O}_X))
\longrightarrow
D_\QCoh(\mathcal{O}_X)
\end{equation}
voir Derived Categories, Équation (\ref{derived-equation-compare}).

\medskip\noindent
Observons qu'un morphisme plat $f : Y \to X$ d'espaces algébriques
induit un foncteur exact
$f^* : \textit{Mod}(\mathcal{O}_X) \to \textit{Mod}(\mathcal{O}_Y)$,
voir
Morphisms of Spaces, Lemma \ref{spaces-morphisms-lemma-flat-morphism-sites}
et
Modules on Sites, Lemma \ref{sites-modules-lemma-flat-pullback-exact}.
En particulier, $Lf^* : D(\mathcal{O}_X) \to D(\mathcal{O}_Y)$
se calcule sur n'importe quel complexe représentant
(Derived Categories, Lemma \ref{derived-lemma-right-derived-exact-functor}).
Nous écrirons $Lf^* = f^*$ lorsque $f$ est plat ; dans ce cas,
$H^i(f^*E) = f^*H^i(E)$ pour $E$ dans $D(\mathcal{O}_X)$.
Nous utiliserons souvent cette convention lorsque $f$ est étale. Dans le cas étale,
le foncteur image réciproque n'est autre que la restriction à $Y_\etale$,
voir Properties of Spaces, Équation
(\ref{spaces-properties-equation-restrict-modules}).

\begin{lemma}
\label{lemma-check-quasi-coherence-on-covering}
Soit $S$ un schéma. Soit $X$ un espace algébrique au-dessus de $S$.
Soit $E$ un objet de $D(\mathcal{O}_X)$. Les conditions suivantes sont équivalentes :
\begin{enumerate}
\item $E$ appartient à $D_\QCoh(\mathcal{O}_X)$ ;
\item pour tout morphisme étale $\varphi : U \to X$ où $U$ est un
schéma affine, $\varphi^*E$ est un objet de
$D_\QCoh(\mathcal{O}_U)$,
\item pour tout morphisme étale $\varphi : U \to X$ où $U$ est un schéma,
$\varphi^*E$ est un objet de
$D_\QCoh(\mathcal{O}_U)$,
\item il existe un morphisme étale surjectif $\varphi : U \to X$
où $U$ est un schéma tel que $\varphi^*E$ soit un objet de
$D_\QCoh(\mathcal{O}_U)$, et
\item il existe un morphisme étale surjectif d'espaces algébriques
$f : Y \to X$ tel que $Lf^*E$ soit un objet de
$D_\QCoh(\mathcal{O}_Y)$.
\end{enumerate}
\end{lemma}

\begin{proof}
Cela découle immédiatement de la discussion qui précède le lemme et de
Properties of Spaces, Lemma
\ref{spaces-properties-lemma-characterize-quasi-coherent}.
\end{proof}

\begin{lemma}
\label{lemma-quasi-coherence-direct-sums}
Soit $S$ un schéma. Soit $X$ un espace algébrique au-dessus de $S$.
Alors $D_\QCoh(\mathcal{O}_X)$ admet des sommes directes.
\end{lemma}

\begin{proof}
Par Injectives, Lemme \ref{injectives-lemma-derived-products},
la catégorie dérivée $D(\mathcal{O}_X)$ admet des sommes directes,
qui se calculent en prenant terme à terme les sommes directes de représentants quelconques.
Il est donc clair que le faisceau de cohomologie d'une somme directe est la
somme directe des faisceaux de cohomologie, puisque la prise de sommes directes
est un foncteur exact (dans toute catégorie abélienne de Grothendieck). Le lemme
résulte alors du fait que la somme directe de faisceaux quasi-cohérents est quasi-cohérente, voir
Properties of Spaces, Lemma
\ref{spaces-properties-lemma-properties-quasi-coherent}.
\end{proof}

\noindent
Nous aurons besoin de quelques informations sur les limites dérivées. Nous avertissons le lecteur
que, dans le lemme ci-dessous, la limite dérivée ne sera en général pas
un objet de $D_\QCoh$.

\begin{lemma}
\label{lemma-Rlim-quasi-coherent}
Soit $S$ un schéma. Soit $X$ un espace algébrique au-dessus de $S$.
Soit $(K_n)$ un système projectif de
$D_\QCoh(\mathcal{O}_X)$ dont la limite dérivée
$K = R\lim K_n$ dans $D(\mathcal{O}_X)$. Supposons que $H^q(K_{n + 1}) \to H^q(K_n)$
soit surjectif pour tout $q \in \mathbf{Z}$ et tout $n \geq 1$.
Alors
\begin{enumerate}
\item $H^q(K) = \lim H^q(K_n)$ ;
\item $R\lim H^q(K_n) = \lim H^q(K_n)$ ;
\item pour tout ouvert affine $U \subset X$, on a
$H^p(U, \lim H^q(K_n)) = 0$ pour $p > 0$.
\end{enumerate}
\end{lemma}

\begin{proof}
Soit $\mathcal{B} \subset \Ob(X_\etale)$ l'ensemble des objets affines.
Comme $H^q(K_n)$ est quasi-cohérent, on a $H^p(U, H^q(K_n)) = 0$
pour $U \in \mathcal{B}$ d'après la discussion dans
Cohomology of Spaces, Section
\ref{spaces-cohomology-section-higher-direct-image}
et
Cohomology of Schemes, Lemme
\ref{coherent-lemma-quasi-coherent-affine-cohomology-zero}.
De plus, les applications $H^0(U, H^q(K_{n + 1})) \to H^0(U, H^q(K_n))$
sont surjectives pour $U \in \mathcal{B}$ par un raisonnement analogue.
La partie (1) découle de Cohomology on Sites, Lemme
\ref{sites-cohomology-lemma-derived-limit-suitable-system}
dont nous venons de vérifier les conditions.
Les parties (2) et (3) découlent de
Cohomology on Sites, Lemme
\ref{sites-cohomology-lemma-inverse-limit-is-derived-limit}.
\end{proof}

\begin{lemma}
\label{lemma-quasi-coherence-pullback}
Soit $S$ un schéma.
Soit $f : Y \to X$ un morphisme d'espaces algébriques au-dessus de $S$.
Le foncteur $Lf^*$ envoie $D_\QCoh(\mathcal{O}_X)$
dans $D_\QCoh(\mathcal{O}_Y)$.
\end{lemma}

\begin{proof}
Choisissons un diagramme
$$
\xymatrix{
U \ar[d]_a \ar[r]_h & V \ar[d]^b \\
X \ar[r]^f & Y
}
$$
où $U$ et $V$ sont des schémas, les flèches verticales sont étales et
$a$ est surjectif. Comme $a^* \circ Lf^* = Lh^* \circ b^*$, le résultat
découle de
Lemma \ref{lemma-check-quasi-coherence-on-covering}
et du cas des schémas, qui est traité dans
Derived Categories of Schemes, Lemma
\ref{perfect-lemma-quasi-coherence-pullback}.
\end{proof}

\begin{lemma}
\label{lemma-quasi-coherence-tensor-product}
Soit $S$ un schéma. Soit $X$ un espace algébrique au-dessus de $S$.
Pour des objets $K, L$ de $D_\QCoh(\mathcal{O}_X)$,
le produit tensoriel dérivé $K \otimes^\mathbf{L} L$ appartient à
$D_\QCoh(\mathcal{O}_X)$.
\end{lemma}

\begin{proof}
Soit $\varphi : U \to X$ un morphisme étale surjectif provenant d'un schéma $U$.
Comme
$\varphi^*(K \otimes_{\mathcal{O}_X}^\mathbf{L} L) =
\varphi^*K \otimes_{\mathcal{O}_U}^\mathbf{L} \varphi^*L$
il résulte de
Lemma \ref{lemma-check-quasi-coherence-on-covering}
que cela découle du cas des schémas, traité dans
Derived Categories of Schemes, Lemma
\ref{perfect-lemma-quasi-coherence-tensor-product}.
\end{proof}

\noindent
Le lemme suivant nous aidera à « calculer » un foncteur dérivé à droite
sur un objet de $D_\QCoh(\mathcal{O}_X)$.

\begin{lemma}
\label{lemma-nice-K-injective}
Soit $S$ un schéma. Soit $X$ un espace algébrique au-dessus de $S$. Soit $E$ un
objet de $D_\QCoh(\mathcal{O}_X)$. Alors le morphisme
$E \to R\lim \tau_{\geq -n}E$ of
Derived Categories, Remark
\ref{derived-remark-map-into-derived-limit-truncations}
est un isomorphisme\footnote{En particulier,
$E$ admet un représentant K-injectif, voir
Derived Categories, Lemma \ref{derived-lemma-difficulty-K-injectives}.}.
\end{lemma}

\begin{proof}
Notons $\mathcal{H}^i = H^i(E)$ le $i$-ième faisceau de cohomologie de $E$.
Soit $\mathcal{B}$ l'ensemble des objets affines de $X_\etale$.
Alors $H^p(U, \mathcal{H}^i) = 0$ pour tout $p > 0$, tout $i \in \mathbf{Z}$,
et tout $U \in \mathcal{B}$, puisque $U$ est un schéma affine.
Voir la discussion dans
Cohomology of Spaces, Section
\ref{spaces-cohomology-section-higher-direct-image}
and
Cohomology of Schemes, Lemma
\ref{coherent-lemma-quasi-coherent-affine-cohomology-zero}.
Le lemme découle donc de
Cohomology on Sites, Lemma \ref{sites-cohomology-lemma-is-limit-dimension}
with $d = 0$.
\end{proof}

\begin{lemma}
\label{lemma-application-nice-K-injective}
Soit $S$ un schéma. Soit $X$ un espace algébrique au-dessus de $S$.
Soit $F : \textit{Mod}(\mathcal{O}_X) \to \textit{Ab}$
un foncteur et $N \geq 0$ un entier. Supposons que
\begin{enumerate}
\item $F$ est exact à gauche ;
\item $F$ commute aux produits directs dénombrables ;
\item $R^pF(\mathcal{F}) = 0$ pour tout $p \geq N$ et tout $\mathcal{F}$
quasi-cohérent.
\end{enumerate}
Alors, pour $E \in D_\QCoh(\mathcal{O}_X)$,
\begin{enumerate}
\item $H^i(RF(\tau_{\leq a}E)) \to H^i(RF(E))$ est un isomorphisme
pour $i \leq a$ ;
\item $H^i(RF(E)) \to H^i(RF(\tau_{\geq b - N + 1}E))$ est un isomorphisme
pour $i \geq b$ ;
\item si $H^i(E) = 0$ pour $i \not \in [a, b]$ pour certains
$-\infty \leq a \leq b \leq \infty$, alors $H^i(RF(E)) = 0$
pour $i \not \in [a, b + N - 1]$.
\end{enumerate}
\end{lemma}

\begin{proof}
L'assertion (1) est
Derived Categories, Lemma \ref{derived-lemma-negative-vanishing}.

\medskip\noindent
Preuve de l'assertion (2). Écrivons $E_n = \tau_{\geq -n}E$. Nous avons
$E = R\lim E_n$, voir le lemme \ref{lemma-nice-K-injective}. Ainsi,
$RF(E) = R\lim RF(E_n)$ dans $D(\textit{Ab})$ d'après Injectives, Lemme
\ref{injectives-lemma-RF-commutes-with-Rlim}. Ainsi, pour tout $i \in \mathbf{Z}$,
nous avons une suite exacte courte
$$
0 \to R^1\lim H^{i - 1}(RF(E_n)) \to H^i(RF(E)) \to \lim H^i(RF(E_n)) \to 0
$$
voir More on Algebra, Remarque
\ref{more-algebra-remark-compare-derived-limit}.
Pour démontrer (2), nous montrerons que le terme de gauche est nul
et que le terme de droite est égal à $H^i(RF(E_{-b + N - 1}))$
pour tout $b$ tel que $i \geq b$.

\medskip\noindent
Pour tout $n$, nous avons un triangle distingué
$$
H^{-n}(E)[n] \to E_n \to E_{n - 1} \to H^{-n}(E)[n + 1]
$$
(Derived Categories, Remarque
\ref{derived-remark-truncation-distinguished-triangle})
dans $D(\mathcal{O}_X)$. Comme $H^{-n}(E)$ est quasi-cohérent, nous avons
$$
H^i(RF(H^{-n}(E)[n])) = R^{i + n}F(H^{-n}(E)) = 0
$$
pour $i + n \geq N$, et
$$
H^i(RF(H^{-n}(E)[n + 1])) = R^{i + n + 1}F(H^{-n}(E)) = 0
$$
pour $i + n + 1 \geq N$. Nous en concluons que
$$
H^i(RF(E_n)) \to H^i(RF(E_{n - 1}))
$$
est un isomorphisme pour $n \geq N - i$. Ainsi, les systèmes $H^i(RF(E_n))$
satisfont tous la condition ML et le terme $R^1\lim$ de notre suite exacte
courte est nul (voir la discussion dans
More on Algebra, Section \ref{more-algebra-section-Rlim}).
De plus, le système $H^i(RF(E_n))$ est constant à partir de
$n = N - i - 1$, comme voulu.

\medskip\noindent
Preuve de (3). Sous l'hypothèse portant sur $E$, nous avons
$\tau_{\leq a - 1}E = 0$ et (1) donne l'annulation
de $H^i(RF(E))$ pour $i \leq a - 1$.
De même, $\tau_{\geq b + 1}E = 0$ et donc
nous obtenons l'annulation de $H^i(RF(E))$ pour $i \geq b + N$ d'après
la partie (2).
\end{proof}









\section{Image directe totale}
\label{section-total-direct-image}

\noindent
Le lemme suivant est l'analogue de
Cohomology of Spaces, Lemma
\ref{spaces-cohomology-lemma-vanishing-higher-direct-images}.

\begin{lemma}
\label{lemma-quasi-coherence-direct-image}
Soit $S$ un schéma. Soit $f : X \to Y$ un morphisme quasi-séparé et quasi-compact
d'espaces algébriques au-dessus de $S$.
\begin{enumerate}
\item Le foncteur $Rf_*$ envoie $D_\QCoh(\mathcal{O}_X)$
dans $D_\QCoh(\mathcal{O}_Y)$.
\item Si $Y$ est quasi-compact, il existe un entier $N = N(X, Y, f)$
tel que, pour un objet $E$ de $D_\QCoh(\mathcal{O}_X)$
tel que $H^m(E) = 0$ pour $m > 0$, on ait
$H^m(Rf_*E) = 0$ pour $m \geq N$.
\item En fait, si $Y$ est quasi-compact, on peut trouver $N = N(X, Y, f)$
tel que, pour tout morphisme d'espaces algébriques $Y' \to Y$,
la même conclusion vaille pour le foncteur $R(f')_*$,
où $f' : X' \to Y'$ est le changement de base de $f$.
\end{enumerate}
\end{lemma}

\begin{proof}
Soit $E$ un objet de $D_\QCoh(\mathcal{O}_X)$.
Pour démontrer (1), nous devons montrer que $Rf_*E$ a des faisceaux
de cohomologie quasi-cohérents. Cette question est locale sur $Y$ ; nous pouvons
donc supposer $Y$ quasi-compact. Choisissons $N = N(X, Y, f)$ comme dans
Cohomology of Spaces, Lemma
\ref{spaces-cohomology-lemma-vanishing-higher-direct-images}.
Ainsi $R^pf_*\mathcal{F} = 0$ pour tout $\mathcal{O}_X$-module quasi-cohérent
$\mathcal{F}$ et tout $p \geq N$. De plus, $R^pf_*\mathcal{F}$
est quasi-cohérent pour tout $p$ d'après
Cohomology of Spaces, Lemma \ref{spaces-cohomology-lemma-higher-direct-image}.
Ces assertions restent vraies après changement de base.

\medskip\noindent
Commençons par supposer $E$ borné inférieurement. Nous allons montrer que (1), (2) et (3)
valent pour un tel $E$ avec notre choix de $N$. Dans ce cas, nous pouvons par exemple utiliser la
spectral sequence
$$
R^pf_*H^q(E) \Rightarrow R^{p + q}f_*E
$$
(Derived Categories, Lemma \ref{derived-lemma-two-ss-complex-functor}),
la quasi-cohérence de $R^pf_*H^q(E)$ et l'annulation de $R^pf_*H^q(E)$
pour $p \geq N$ pour voir que (1), (2) et (3) valent dans ce cas.

\medskip\noindent
Passons à la preuve de (2) et (3). Supposons $H^m(E) = 0$ pour $m > 0$.
Soit $V$ un objet affine de $Y_\etale$.
On a $H^p(V \times_Y X, \mathcal{F}) = 0$ pour $p \geq N$, voir
Cohomology of Spaces, Lemma
\ref{spaces-cohomology-lemma-quasi-coherence-higher-direct-images-application}.
Nous pouvons donc appliquer le lemme \ref{lemma-application-nice-K-injective}
au foncteur $\Gamma(V \times_Y X, -)$ pour obtenir
$$
R\Gamma(V, Rf_*E) = R\Gamma(V \times_Y X, E)
$$
dont la cohomologie s'annule en degrés $\geq N$. Comme ceci vaut pour
tout $V$ affine de $Y_\etale$, nous concluons que $H^m(Rf_*E) = 0$
pour $m \geq N$.

\medskip\noindent
Prouvons maintenant (1) dans le cas général. Rappelons qu'il existe un
triangle distingué
$$
\tau_{\leq -n - 1}E \to E \to \tau_{\geq -n}E \to
(\tau_{\leq -n - 1}E)[1]
$$
dans $D(\mathcal{O}_X)$, voir Derived Categories, Remarque
\ref{derived-remark-truncation-distinguished-triangle}.
D'après (2), nous voyons que $Rf_*\tau_{\leq -n - 1}E$
a des faisceaux de cohomologie nuls en degrés $\geq -n + N$.
Ainsi, pour un entier $q$ donné, $R^qf_*E$ est égal
à $R^qf_*\tau_{\geq -n}E$ pour un certain $n$, et le résultat
ci-dessus s'applique.
\end{proof}

\begin{lemma}
\label{lemma-quasi-coherence-pushforward-direct-sums}
Soit $S$ un schéma. Soit $f : X \to Y$ un morphisme quasi-séparé et
quasi-compact d'espaces algébriques au-dessus de $S$. Alors
$Rf_* : D_\QCoh(\mathcal{O}_X) \to D_\QCoh(\mathcal{O}_Y)$
commute aux sommes directes.
\end{lemma}

\begin{proof}
Soit $(E_i)$ une famille d'objets de $D_\QCoh(\mathcal{O}_X)$
et posons $E = \bigoplus E_i$. Nous voulons montrer que le morphisme
$$
\bigoplus Rf_*E_i \longrightarrow Rf_*E
$$
est un isomorphisme. Nous montrerons qu'il induit un isomorphisme sur les
faisceaux de cohomologie en degré $0$, ce qui entraînera le lemme.
Choisissons un entier $N$ comme dans le lemme \ref{lemma-quasi-coherence-direct-image}.
Alors $R^0f_*E = R^0f_*\tau_{\geq -N}E$ et
$R^0f_*E_i = R^0f_*\tau_{\geq -N}E_i$ d'après le lemme cité. Observons que
$\tau_{\geq -N}E = \bigoplus \tau_{\geq -N}E_i$.
Nous pouvons donc supposer que tous les $E_i$ ont des faisceaux de cohomologie
nuls en degrés $< -N$. Nous utilisons ensuite les suites spectrales
$$
R^pf_*H^q(E) \Rightarrow R^{p + q}f_*E
\quad\text{et}\quad
R^pf_*H^q(E_i) \Rightarrow R^{p + q}f_*E_i
$$
(Derived Categories, Lemme \ref{derived-lemma-two-ss-complex-functor})
pour nous ramener au cas d'une somme directe de faisceaux quasi-cohérents.
Ce cas est traité par
Cohomology of Spaces, Lemma \ref{spaces-cohomology-lemma-colimit-cohomology}.
\end{proof}

\begin{remark}
\label{remark-match-total-direct-images}
Soit $S$ un schéma. Soit $f : X \to Y$ un morphisme d'espaces algébriques
représentables $X$ et $Y$ au-dessus de $S$. Soit $f_0 : X_0 \to Y_0$ un
morphisme de schémas représentant $f$ (notation maladroite mais provisoire).
Alors le diagramme
$$
\xymatrix{
D_\QCoh(\mathcal{O}_{X_0})
\ar@{=}[rrrrrr]_{\text{Lemme
\ref{lemma-derived-quasi-coherent-small-etale-site}}}
& & & & & &
D_\QCoh(\mathcal{O}_X) \\
D_\QCoh(\mathcal{O}_{Y_0})
\ar[u]^{Lf^*_0}
\ar@{=}[rrrrrr]^{\text{Lemme
\ref{lemma-derived-quasi-coherent-small-etale-site}}}
& & & & & &
D_\QCoh(\mathcal{O}_Y) \ar[u]_{Lf^*}
}
$$
(Lemme \ref{lemma-quasi-coherence-pullback} et
Derived Categories of Schemes, Lemma
\ref{perfect-lemma-quasi-coherence-pullback})
est commutatif. Cela résulte du fait que les
équivalences
$D_\QCoh(\mathcal{O}_{X_0}) \to D_\QCoh(\mathcal{O}_X)$
et
$D_\QCoh(\mathcal{O}_{Y_0}) \to D_\QCoh(\mathcal{O}_Y)$
du lemme \ref{lemma-derived-quasi-coherent-small-etale-site}
proviennent de l'image réciproque par les morphismes (plats) de sites annelés
$\epsilon : X_\etale \to X_{0, Zar}$ et
$\epsilon : Y_\etale \to Y_{0, Zar}$
et que le diagramme de sites annelés
$$
\xymatrix{
X_{0, Zar} \ar[d]_{f_0} & X_\etale \ar[l]^\epsilon \ar[d]^f \\
Y_{0, Zar} & Y_\etale \ar[l]_\epsilon
}
$$
est commutatif (nous omettons les détails). Si $f$ est quasi-compact et
quasi-séparé, ce qui équivaut à ce que $f_0$ soit quasi-compact et
quasi-séparé, alors nous affirmons que
$$
\xymatrix{
D_\QCoh(\mathcal{O}_{X_0})
\ar[d]_{Rf_{0, *}} \ar@{=}[rrrrrr]_{\text{Lemma
\ref{lemma-derived-quasi-coherent-small-etale-site}}}
& & & & & &
D_\QCoh(\mathcal{O}_X) \ar[d]^{Rf_*} \\
D_\QCoh(\mathcal{O}_{Y_0})
\ar@{=}[rrrrrr]^{\text{Lemma
\ref{lemma-derived-quasi-coherent-small-etale-site}}}
& & & & & &
D_\QCoh(\mathcal{O}_Y)
}
$$
(Lemme \ref{lemma-quasi-coherence-direct-image} et
Derived Categories of Schemes, Lemma
\ref{perfect-lemma-quasi-coherence-direct-image})
est également commutatif. Cela découle encore du
diagramme commutatif de sites affiché ci-dessus, puisque la preuve du lemme
\ref{lemma-derived-quasi-coherent-small-etale-site}
montre que le foncteur $R\epsilon_*$ fournit les équivalences
$D_\QCoh(\mathcal{O}_X) \to D_\QCoh(\mathcal{O}_{X_0})$
et
$D_\QCoh(\mathcal{O}_Y) \to D_\QCoh(\mathcal{O}_{Y_0})$.
\end{remark}

\begin{lemma}
\label{lemma-affine-morphism}
Soit $S$ un schéma. Soit $f : X \to Y$ un morphisme affine d'espaces algébriques
au-dessus de $S$. Alors
$Rf_* : D_\QCoh(\mathcal{O}_X) \to D_\QCoh(\mathcal{O}_Y)$
réfléchit les isomorphismes.
\end{lemma}

\begin{proof}
L'assertion signifie qu'un morphisme $\alpha : E \to F$ de
$D_\QCoh(\mathcal{O}_X)$ est un isomorphisme dès que
$Rf_*\alpha$ est un isomorphisme. Nous pouvons le vérifier sur les faisceaux de cohomologie.
En particulier, la question est locale pour la topologie étale sur $Y$. Nous pouvons donc supposer
$Y$, et par conséquent $X$, affines. Dans ce cas, le problème se ramène au
cas des schémas
(Derived Categories of Schemes, Lemme \ref{perfect-lemma-affine-morphism})
par le Lemme \ref{lemma-derived-quasi-coherent-small-etale-site} et la
Remarque \ref{remark-match-total-direct-images}.
\end{proof}

\begin{lemma}
\label{lemma-affine-morphism-pull-push}
Soit $S$ un schéma. Soit $f : X \to Y$ un morphisme affine d'espaces algébriques
au-dessus de $S$. Pour $E$ dans $D_\QCoh(\mathcal{O}_Y)$, on a
$Rf_* Lf^* E = E \otimes^\mathbf{L}_{\mathcal{O}_Y} f_*\mathcal{O}_X$.
\end{lemma}

\begin{proof}
Comme $f$ est affine, le morphisme $f_*\mathcal{O}_X \to Rf_*\mathcal{O}_X$
est un isomorphisme (Cohomology of Spaces, Lemme
\ref{spaces-cohomology-lemma-affine-vanishing-higher-direct-images}).
Il existe un morphisme canonique
$E \otimes^\mathbf{L} f_*\mathcal{O}_X =
E \otimes^\mathbf{L} Rf_*\mathcal{O}_X \to Rf_* Lf^* E$
adjoint au morphisme
$$
Lf^*(E \otimes^\mathbf{L} Rf_*\mathcal{O}_X) =
Lf^*E \otimes^\mathbf{L} Lf^*Rf_*\mathcal{O}_X \longrightarrow
Lf^* E \otimes^\mathbf{L} \mathcal{O}_X = Lf^* E
$$
provenant de $1 : Lf^*E \to Lf^*E$ et du morphisme canonique
$Lf^*Rf_*\mathcal{O}_X \to \mathcal{O}_X$. Pour vérifier que le morphisme ainsi construit
est un isomorphisme, nous pouvons travailler localement sur $Y$. Nous pouvons donc supposer
$Y$, et par conséquent $X$, affine. Dans ce cas, le problème se ramène au
cas des schémas
(Derived Categories of Schemes, Lemme
\ref{perfect-lemma-affine-morphism-pull-push})
par le Lemme \ref{lemma-derived-quasi-coherent-small-etale-site} et la
Remarque \ref{remark-match-total-direct-images}.
\end{proof}








\section{Être propre au-dessus d'une base}
\label{section-proper-over-base}

\noindent
Cette section est l'analogue de la section de Cohomologie des schémas
\ref{coherent-section-proper-over-base}.
Comme toujours pour les résultats concernant la topologie des ensembles de points,
il faut être prudent en transposant le matériel aux espaces algébriques.

\begin{lemma}
\label{lemma-closed-proper-over-base}
Soit $S$ un schéma. Soit $f : X \to Y$ un morphisme d'espaces algébriques
sur $S$, localement de type fini. Soit $T \subset |X|$ un sous-ensemble fermé.
Les conditions suivantes sont équivalentes
\begin{enumerate}
\item le morphisme $Z \to Y$ est propre si $Z$ est la structure d'espace algébrique réduite
induite sur $T$
(Propriétés des espaces, définition
\ref{spaces-properties-definition-reduced-induced-space}),
\item pour un sous-espace fermé $Z \subset X$ tel que $|Z| = T$,
le morphisme $Z \to Y$ est propre, et
\item pour tout sous-espace fermé $Z \subset X$ tel que $|Z| = T$, le morphisme
$Z \to Y$ est propre.
\end{enumerate}
\end{lemma}

\begin{proof}
Les implications (3) $\Rightarrow$ (1) et (1) $\Rightarrow$ (2)
sont immédiates. Il suffit donc de démontrer que (2) implique (3).
Nous invitons le lecteur à trouver sa propre démonstration de ce fait.
Soient $Z'$ et $Z''$ des sous-espaces fermés tels que $T = |Z'| = |Z''|$
et tels que $Z' \to Y$ soit un morphisme propre d'espaces algébriques.
Nous devons montrer que $Z'' \to Y$ est également propre.
Soit $Z''' = Z' \cup Z''$ l'union schématique, voir
Morphismes des espaces, définition
\ref{spaces-morphisms-definition-scheme-theoretic-intersection-union}.
Alors $Z'''$ est un autre sous-espace fermé tel que $|Z'''| = T$.
Cela découle par exemple de la description des unions schématiques dans
Morphismes des espaces, lemme \ref{spaces-morphisms-lemma-scheme-theoretic-union}.
Comme $Z'' \to Z'''$ est une immersion fermée, il suffit de démontrer
que $Z''' \to Y$ est propre (voir
Morphismes des espaces, lemmes
\ref{spaces-morphisms-lemma-closed-immersion-proper} et
\ref{spaces-morphisms-lemma-composition-proper}).
Le morphisme $Z' \to Z'''$ est une immersion fermée bijective
et donc, en particulier, surjectif et universellement fermé.
Le fait que $Z' \to Y$ soit séparé implique alors que
$Z''' \to Y$ est séparé, voir
Morphismes des espaces, lemme
\ref{spaces-morphisms-lemma-image-universally-closed-separated}.
De plus, $Z''' \to Y$ est localement de type fini
puisque $X \to Y$ est localement de type fini
(Morphismes des espaces, lemmes
\ref{spaces-morphisms-lemma-immersion-locally-finite-type} et
\ref{spaces-morphisms-lemma-composition-finite-type}).
Comme $Z' \to Y$ est quasi-compact et $Z' \to Z'''$ est un
homéomorphisme universel, $Z''' \to Y$ est quasi-compact.
Enfin, puisque $Z' \to Y$ est universellement fermé,
il en va de même pour $Z''' \to Y$ d'après
Morphismes des espaces, lemme \ref{spaces-morphisms-lemma-image-proper-is-proper}.
Ceci achève la démonstration.
\end{proof}

\begin{definition}
\label{definition-proper-over-base}
Soit $S$ un schéma.
Soit $f : X \to Y$ un morphisme d'espaces algébriques sur $S$
localement de type fini.
Soit $T \subset |X|$ un sous-ensemble fermé.
Nous disons que {\it $T$ est propre au-dessus de $Y$}
si les conditions équivalentes du lemme \ref{lemma-closed-proper-over-base}
sont satisfaites.
\end{definition}

\noindent
Le lemme utilisé dans la définition ci-dessus est faux si le morphisme
$f : X \to Y$ n'est pas localement de type fini. Nous invitons donc
le lecteur à ne pas utiliser cette terminologie si $f$ n'est pas de
type fini localement.

\begin{lemma}
\label{lemma-closed-closed-proper-over-base}
Soit $S$ un schéma.
Soit $f : X \to Y$ un morphisme d'espaces algébriques sur $S$
localement de type fini.
Soit $T' \subset T \subset |X|$ des sous-ensembles fermés.
Si $T$ est propre au-dessus de $Y$, il en va de même pour $T'$.
\end{lemma}

\begin{proof}
Omis.
\end{proof}

\begin{lemma}
\label{lemma-base-change-closed-proper-over-base}
Soit $S$ un schéma.
Considérons un diagramme cartésien d'espaces algébriques sur $S$
$$
\xymatrix{
X' \ar[d]_{f'} \ar[r]_{g'} & X \ar[d]^f \\
Y' \ar[r]^g & Y
}
$$
tel que $f$ soit localement de type fini.
Si $T$ est un sous-ensemble fermé de $|X|$ propre au-dessus de $Y$, alors
$|g'|^{-1}(T)$ est un sous-ensemble fermé de $|X'|$ propre au-dessus de $Y'$.
\end{lemma}

\begin{proof}
Observons que l'énoncé a un sens puisque $f'$ est localement de
type fini, d'après Morphismes des espaces, lemme
\ref{spaces-morphisms-lemma-base-change-finite-type}.
Soit $Z \subset X$ la structure de sous-espace fermé réduite induite sur $T$.
Notons $Z' = (g')^{-1}(Z)$ l'image inverse schématique.
Alors $Z' = X' \times_X Z = (Y' \times_Y X) \times_X Z = Y' \times_Y Z$
est propre au-dessus de $Y'$ comme changement de base de $Z$ au-dessus de $Y$
(Morphismes des espaces, lemme \ref{spaces-morphisms-lemma-base-change-proper}).
D'autre part, $T' = |Z'|$. Le lemme en découle.
\end{proof}

\begin{lemma}
\label{lemma-functoriality-closed-proper-over-base}
Soit $S$ un schéma. Soit $B$ un espace algébrique sur $S$.
Soit $f : X \to Y$ un morphisme d'espaces algébriques qui
sont localement de type fini sur $B$.
\begin{enumerate}
\item Si $Y$ est séparé sur $B$ et si $T \subset |X|$ est un sous-ensemble fermé
propre au-dessus de $B$, alors $|f|(T)$ est un sous-ensemble fermé de $|Y|$ propre au-dessus de $B$.
\item Si $f$ est universellement fermé et si $T \subset |X|$ est un
sous-ensemble fermé propre au-dessus de $B$, alors $|f|(T)$ est un sous-ensemble fermé
de $Y$ propre au-dessus de $B$.
\item Si $f$ est propre et si $T \subset |Y|$ est un sous-ensemble fermé
propre au-dessus de $B$, alors $|f|^{-1}(T)$ est un sous-ensemble fermé de $|X|$
propre au-dessus de $B$.
\end{enumerate}
\end{lemma}

\begin{proof}
Démonstration de (1). Supposons que $Y$ soit séparé sur $B$ et que $T \subset |X|$
soit un sous-ensemble fermé propre au-dessus de $B$. Soit $Z$ la structure réduite induite
de sous-espace fermé sur $T$ et appliquons
Morphismes des espaces, lemme
\ref{spaces-morphisms-lemma-scheme-theoretic-image-is-proper}
à $Z \to Y$ au-dessus de $B$ pour conclure.

\medskip\noindent
Démonstration de (2). Supposons que $f$ soit universellement fermé et que $T \subset |X|$ soit un
sous-ensemble fermé propre au-dessus de $B$. Soit $Z$ la structure réduite induite
de sous-espace fermé sur $T$ et soit $Z'$ la structure réduite
de sous-espace fermé induite sur $|f|(T)$. Nous obtenons un morphisme
induit $Z \to Z'$.
Notons $Z'' = f^{-1}(Z')$ l'image inverse schématique.
Alors $Z'' \to Z'$ est universellement fermé comme changement de base de $f$
(Morphismes des espaces, lemme \ref{spaces-morphisms-lemma-base-change-proper}).
Par conséquent, $Z \to Z'$ est universellement fermé comme composée de
l'immersion fermée $Z \to Z''$ et de $Z'' \to Z'$
(Morphismes des espaces, lemmes
\ref{spaces-morphisms-lemma-closed-immersion-proper} et
\ref{spaces-morphisms-lemma-composition-proper}).
Nous concluons que $Z' \to B$ est séparé d'après
Morphismes des espaces, lemme
\ref{spaces-morphisms-lemma-image-universally-closed-separated}.
Comme $Z \to B$ est quasi-compact et $Z \to Z'$ surjectif,
on voit que $Z' \to B$ est quasi-compact.
Comme $Z' \to B$ est la composée de $Z' \to Y$ et de $Y \to B$,
on voit que $Z' \to B$ est localement de type fini
(Morphismes des espaces, lemmes
\ref{spaces-morphisms-lemma-immersion-locally-finite-type} et
\ref{spaces-morphisms-lemma-composition-finite-type}).
Enfin, puisque $Z \to B$ est universellement fermé,
il en va de même pour $Z' \to B$ d'après
Morphismes des espaces, lemme
\ref{spaces-morphisms-lemma-image-proper-is-proper}.
Ceci achève la démonstration.

\medskip\noindent
Démonstration de (3). Supposons $f$ propre et $T \subset |Y|$ un sous-ensemble fermé
propre au-dessus de $B$. Soit $Z$ la structure de sous-espace fermé réduite induite
sur $T$. Notons $Z' = f^{-1}(Z)$ l'image inverse schématique.
Alors $Z' \to Z$ est propre comme changement de base de $f$
(Morphismes des espaces, lemme \ref{spaces-morphisms-lemma-base-change-proper}).
Par conséquent, $Z' \to B$ est propre comme composée de $Z' \to Z$
et de $Z \to B$
(Morphismes des espaces, lemme \ref{spaces-morphisms-lemma-composition-proper}).
Ceci achève la démonstration.
\end{proof}

\begin{lemma}
\label{lemma-union-closed-proper-over-base}
Soit $S$ un schéma.
Soit $f : X \to Y$ un morphisme d'espaces algébriques sur $S$
localement de type fini.
Soient $T_i \subset |X|$, $i = 1, \ldots, n$, des sous-ensembles fermés.
Si les $T_i$, $i = 1, \ldots, n$, sont propres au-dessus de $Y$, il en va de même
pour $T_1 \cup \ldots \cup T_n$.
\end{lemma}

\begin{proof}
Soit $Z_i$ la structure de sous-schéma fermé réduite induite sur $T_i$.
Le morphisme
$$
Z_1 \amalg \ldots \amalg Z_n \longrightarrow X
$$
est fini d'après Morphisms of Spaces, lemmes
\ref{spaces-morphisms-lemma-closed-immersion-finite} et
\ref{spaces-morphisms-lemma-finite-union-finite}.
Comme les morphismes finis sont universellement fermés
(Morphisms of Spaces, Lemme \ref{spaces-morphisms-lemma-finite-proper})
et que $Z_1 \amalg \ldots \amalg Z_n$ est propre au-dessus de $S$,
nous concluons par la partie (2) du lemme
\ref{lemma-functoriality-closed-proper-over-base}
que l'image $Z_1 \cup \ldots \cup Z_n$ est propre au-dessus de $S$.
\end{proof}

\noindent
Soit $S$ un schéma.
Soit $f : X \to Y$ un morphisme d'espaces algébriques au-dessus de $S$
localement de type fini. Soit $\mathcal{F}$ un
$\mathcal{O}_X$-module de type fini et quasi-cohérent. Alors le support
$\text{Supp}(\mathcal{F})$ de $\mathcal{F}$ est un sous-ensemble fermé de $|X|$, voir
Morphisms of Spaces, Lemma \ref{spaces-morphisms-lemma-support-finite-type}.
Il est donc sensé de dire que
« le support de $\mathcal{F}$ est propre au-dessus de $Y$ ».

\begin{lemma}
\label{lemma-module-support-proper-over-base}
Soit $S$ un schéma. Soit $f : X \to Y$ un morphisme d'espaces
algébriques au-dessus de $S$, localement de type fini.
Soit $\mathcal{F}$ un $\mathcal{O}_X$-module de type fini et
quasi-cohérent. Les conditions suivantes sont équivalentes :
\begin{enumerate}
\item le support de $\mathcal{F}$ est propre au-dessus de $Y$ ;
\item le support schématique de $\mathcal{F}$
(Morphisms of Spaces, Definition
\ref{spaces-morphisms-definition-scheme-theoretic-support})
est propre au-dessus de $Y$ ;
\item il existe un sous-espace fermé $Z \subset X$ et un
$\mathcal{O}_Z$-module $\mathcal{G}$ de type fini et quasi-cohérent
tel que (a) $Z \to Y$ soit propre et (b)
$(Z \to X)_*\mathcal{G} = \mathcal{F}$.
\end{enumerate}
\end{lemma}

\begin{proof}
Le support $\text{Supp}(\mathcal{F})$ de $\mathcal{F}$ est un sous-ensemble fermé
de $|X|$, voir Morphisms of Spaces, Lemme
\ref{spaces-morphisms-lemma-support-finite-type}.
Nous pouvons donc appliquer la définition \ref{definition-proper-over-base}.
Comme le support schématique de $\mathcal{F}$ est un sous-espace fermé
dont le sous-ensemble fermé sous-jacent est $\text{Supp}(\mathcal{F})$,
les conditions (1) et (2) sont équivalentes d'après la
définition \ref{definition-proper-over-base}.
Il est clair que (2) implique (3).
Réciproquement, si (3) est vraie, alors
$\text{Supp}(\mathcal{F}) \subset |Z|$
et donc $\text{Supp}(\mathcal{F})$
est propre au-dessus de $Y$, par exemple d'après le
Lemma \ref{lemma-closed-closed-proper-over-base}.
\end{proof}

\begin{lemma}
\label{lemma-base-change-module-support-proper-over-base}
Soit $S$ un schéma.
Considérons un diagramme cartésien d'espaces algébriques au-dessus de $S$
$$
\xymatrix{
X' \ar[d]_{f'} \ar[r]_{g'} & X \ar[d]^f \\
Y' \ar[r]^g & Y
}
$$
tel que $f$ soit localement de type fini. Soit $\mathcal{F}$ un
$\mathcal{O}_X$-module de type fini et quasi-cohérent.
Si le support de $\mathcal{F}$ est propre au-dessus de $Y$, alors
le support de $(g')^*\mathcal{F}$ est propre au-dessus de $Y'$.
\end{lemma}

\begin{proof}
Observons que l'énoncé a un sens puisque
$(g')^*\mathcal{F}$ est de type fini d'après
Modules on Sites, Lemma \ref{sites-modules-lemma-local-pullback}.
On a $\text{Supp}((g')^*\mathcal{F}) = |g'|^{-1}(\text{Supp}(\mathcal{F}))$
d'après Morphisms of Spaces, Lemme \ref{spaces-morphisms-lemma-support-finite-type}.
Le lemme découle donc du
Lemma \ref{lemma-base-change-closed-proper-over-base}.
\end{proof}

\begin{lemma}
\label{lemma-cat-module-support-proper-over-base}
Soit $S$ un schéma.
Soit $f : X \to Y$ un morphisme d'espaces algébriques au-dessus de $S$
localement de type fini. Soient $\mathcal{F}$ et $\mathcal{G}$
des $\mathcal{O}_X$-modules de type fini et quasi-cohérents.
\begin{enumerate}
\item Si les supports de $\mathcal{F}$ et $\mathcal{G}$
sont propres au-dessus de $Y$, il en va de même
pour $\mathcal{F} \oplus \mathcal{G}$, toute extension
de $\mathcal{G}$ par $\mathcal{F}$, $\Im(u)$ et $\Coker(u)$
pour tout morphisme de $\mathcal{O}_X$-modules $u : \mathcal{F} \to \mathcal{G}$,
ainsi que pour tout quotient quasi-cohérent de $\mathcal{F}$ ou de $\mathcal{G}$.
\item Si $Y$ est localement noethérien, la catégorie des
$\mathcal{O}_X$-modules cohérents dont le support est propre au-dessus de
$Y$ est une sous-catégorie de Serre (Homology, Définition
\ref{homology-definition-serre-subcategory})
de la catégorie abélienne des $\mathcal{O}_X$-modules cohérents.
\end{enumerate}
\end{lemma}

\begin{proof}
Preuve de (1). Soient $T$ et $T'$ les supports de $\mathcal{F}$
et de $\mathcal{G}$. Tous les faisceaux mentionnés en (1)
ont alors un support contenu dans $T \cup T'$. L'assertion découle donc
des lemmes \ref{lemma-closed-closed-proper-over-base} et
\ref{lemma-union-closed-proper-over-base}
pourvu que nous vérifiions que ces faisceaux sont de type fini
et quasi-cohérents.
Pour la quasi-cohérence, nous renvoyons le lecteur à
Properties of Spaces, Section \ref{spaces-properties-section-quasi-coherent}.
Pour la propriété « de type fini », nous renvoyons le lecteur à
Properties of Spaces, Section
\ref{spaces-properties-section-properties-modules}.

\medskip\noindent
Preuve de (2). La preuve est la même que celle de (1). Notons que
les assertions ont un sens puisque $X$ est localement noethérien d'après
Morphisms of Spaces, Lemma
\ref{spaces-morphisms-lemma-locally-finite-type-locally-noetherian}
et d'après la description de la catégorie des modules cohérents dans
Cohomology of Spaces, Section \ref{spaces-cohomology-section-coherent}.
\end{proof}

\begin{lemma}
\label{lemma-support-proper-over-base-pushforward}
Soit $S$ un schéma. Soit $f : X \to Y$ un morphisme d'espaces algébriques
au-dessus de $S$. Supposons $f$ localement de type fini et $Y$ localement noethérien.
Soit $\mathcal{F}$ un $\mathcal{O}_X$-module cohérent dont le support
est propre au-dessus de $Y$. Alors $R^pf_*\mathcal{F}$ est un
$\mathcal{O}_Y$-module cohérent pour tout $p \geq 0$.
\end{lemma}

\begin{proof}
D'après le lemme \ref{lemma-module-support-proper-over-base},
il existe une immersion fermée $i : Z \to X$ telle que
$g = f \circ i : Z \to Y$ soit propre et
$\mathcal{F} = i_*\mathcal{G}$ pour un certain module cohérent $\mathcal{G}$
sur $Z$. On voit que $R^pg_*\mathcal{G}$
est cohérent sur $S$ d'après Cohomology of Spaces, Lemme
\ref{spaces-cohomology-lemma-proper-pushforward-coherent}.
D'autre part, $R^qi_*\mathcal{G} = 0$ pour $q > 0$
(Cohomology of Spaces, Lemme
\ref{spaces-cohomology-lemma-finite-pushforward-coherent}).
Par Cohomology on Sites, Lemme \ref{sites-cohomology-lemma-relative-Leray},
nous obtenons $R^pf_*\mathcal{F} = R^pg_*\mathcal{G}$, d'où le lemme.
\end{proof}






\section{Catégorie dérivée des modules cohérents}
\label{section-derived-coherent}

\noindent
Soit $S$ un schéma. Soit $X$ un espace algébrique localement noethérien au-dessus de $S$.
Dans ce cas, la catégorie
$\textit{Coh}(\mathcal{O}_X) \subset \textit{Mod}(\mathcal{O}_X)$
des $\mathcal{O}_X$-modules cohérents est une sous-catégorie de Serre faible, voir
Homology, Section \ref{homology-section-serre-subcategories}
et
Cohomology of Spaces, Lemme
\ref{spaces-cohomology-lemma-coherent-abelian-Noetherian}.
Denote
$$
D_{\textit{Coh}}(\mathcal{O}_X) \subset D(\mathcal{O}_X)
$$
la sous-catégorie des complexes dont les faisceaux de cohomologie sont cohérents, voir
Derived Categories, Section \ref{derived-section-triangulated-sub}.
Nous obtenons ainsi un foncteur canonique
\begin{equation}
\label{equation-compare-coherent}
D(\textit{Coh}(\mathcal{O}_X))
\longrightarrow
D_{\textit{Coh}}(\mathcal{O}_X)
\end{equation}
voir Derived Categories, Équation (\ref{derived-equation-compare}).

\begin{lemma}
\label{lemma-direct-image-coherent}
Soit $S$ un schéma. Soit $f : X \to Y$ un morphisme d'espaces algébriques
au-dessus de $S$. Supposons $f$ localement de type fini et $Y$ noethérien.
Soit $E$ un objet de $D^b_{\textit{Coh}}(\mathcal{O}_X)$ tel que le
support de $H^i(E)$ soit propre au-dessus de $Y$ pour tout $i$.
Alors $Rf_*E$ est un objet de $D^b_{\textit{Coh}}(\mathcal{O}_Y)$.
\end{lemma}

\begin{proof}
Considérons la suite spectrale
$$
R^pf_*H^q(E) \Rightarrow R^{p + q}f_*E
$$
voir Derived Categories, Lemme \ref{derived-lemma-two-ss-complex-functor}.
D'après l'hypothèse et le lemme \ref{lemma-support-proper-over-base-pushforward},
les faisceaux $R^pf_*H^q(E)$ sont cohérents. Donc
$R^{p + q}f_*E$ est cohérent, c'est-à-dire $E \in D_{\textit{Coh}}(\mathcal{O}_Y)$.
La bornitude inférieure est triviale. La bornitude supérieure
découle de
Cohomology of Spaces, Lemma
\ref{spaces-cohomology-lemma-vanishing-higher-direct-images}
ou du
lemme \ref{lemma-quasi-coherence-direct-image}.
\end{proof}

\begin{lemma}
\label{lemma-direct-image-coherent-bdd-below}
Soit $S$ un schéma. Soit $f : X \to Y$ un morphisme d'espaces algébriques
au-dessus de $S$. Supposons $f$ localement de type fini et $Y$ noethérien.
Soit $E$ un objet de $D^+_{\textit{Coh}}(\mathcal{O}_X)$ tel que le support de $H^i(E)$
soit propre au-dessus de $S$ pour tout $i$.
Alors $Rf_*E$ est un objet de $D^+_{\textit{Coh}}(\mathcal{O}_Y)$.
\end{lemma}

\begin{proof}
La preuve est la même que celle du
Lemma \ref{lemma-direct-image-coherent}.
On peut aussi le déduire du
Lemma \ref{lemma-direct-image-coherent}
en considérant l'action du foncteur exact $Rf_*$ sur les
triangles distingués
$\tau_{\leq a}E \to E \to \tau_{\geq a + 1}E \to \tau_{\leq a}E[1]$.
\end{proof}

\begin{lemma}
\label{lemma-coherent-internal-hom}
Soit $S$ un schéma. Soit $X$ un espace algébrique localement noethérien au-dessus de $S$.
Si $L$ appartient à $D^+_{\textit{Coh}}(\mathcal{O}_X)$
et $K$ à $D^-_{\textit{Coh}}(\mathcal{O}_X)$, alors
$R\SheafHom(K, L)$ appartient à $D^+_{\textit{Coh}}(\mathcal{O}_X)$.
\end{lemma}

\begin{proof}
Nous pouvons vérifier qu'un objet de $D(\mathcal{O}_X)$ appartient à
$D_{\textit{Coh}}(\mathcal{O}_X)$ localement pour la topologie étale sur $X$, voir
Cohomology of Spaces, Lemme \ref{spaces-cohomology-lemma-coherent-Noetherian}.
Le lemme découle donc du cas des schémas, voir
Derived Categories of Schemes, Lemme \ref{perfect-lemma-coherent-internal-hom}.
\end{proof}

\begin{lemma}
\label{lemma-ext-finite}
Soit $A$ un anneau noethérien. Soit $X$ un espace algébrique propre au-dessus de $A$.
Pour $L$ dans $D^+_{\textit{Coh}}(\mathcal{O}_X)$ et $K$ dans
$D^-_{\textit{Coh}}(\mathcal{O}_X)$, les $A$-modules
$\Ext_{\mathcal{O}_X}^n(K, L)$ sont de type fini.
\end{lemma}

\begin{proof}
Rappelons que
$$
\Ext_{\mathcal{O}_X}^n(K, L) =
H^n(X, R\SheafHom_{\mathcal{O}_X}(K, L)) =
H^n(\Spec(A), Rf_*R\SheafHom_{\mathcal{O}_X}(K, L))
$$
voir Cohomology on Sites, Lemme \ref{sites-cohomology-lemma-section-RHom-over-U}
et Cohomology on Sites, Section \ref{sites-cohomology-section-leray}.
Le résultat découle donc des lemmes \ref{lemma-coherent-internal-hom} et
\ref{lemma-direct-image-coherent-bdd-below}.
\end{proof}




\section{Principe de récurrence}
\label{section-induction}

\noindent
Dans cette section, nous discutons d'un principe de récurrence pour les espaces algébriques,
analogue à celui qui est établi dans
Cohomology of Schemes, Lemma \ref{coherent-lemma-induction-principle}
pour les schémas. Pour le formuler, nous introduisons la notion de
{\it carré distingué élémentaire} ; cette terminologie est empruntée
à \cite{MV}.
Le principe sous cette forme est implicite dans l'article \cite{GruRay}
de Raynaud et Gruson.
Un principe apparenté pour les champs algébriques est donné dans
\cite[Théorème D]{rydh_etale_devissage} par David Rydh.

\begin{definition}
\label{definition-elementary-distinguished-square}
Soit $S$ un schéma. Un diagramme commutatif
$$
\xymatrix{
U \times_W V \ar[r] \ar[d] & V \ar[d]^f \\
U \ar[r]^j & W
}
$$
d'espaces algébriques au-dessus de $S$ est appelé un {\it carré distingué élémentaire}
si
\begin{enumerate}
\item $U$ est un sous-espace ouvert de $W$ et $j$ est le morphisme d'inclusion ;
\item $f$ est étale ;
\item en posant $T = W \setminus U$ (muni de la structure réduite induite),
le morphisme $f^{-1}(T) \to T$ est un isomorphisme.
\end{enumerate}
Nous l'indiquerons en disant : « Soit $(U \subset W, f : V \to W)$
un carré distingué élémentaire. »
\end{definition}

\noindent
Notons que si $(U \subset W, f : V \to W)$ est un carré distingué
élémentaire, alors $W = U \cup f(V)$. Ainsi, $\{U \to W, V \to W\}$ est
un recouvrement étale de $W$. Ces recouvrements étales possèdent
de bonnes propriétés et sont, en un certain sens,
« suffisamment nombreux ».

\begin{lemma}
\label{lemma-make-more-elementary-distinguished-squares}
Soit $S$ un schéma. Soit $(U \subset W, f : V \to W)$ un carré
distingué élémentaire d'espaces algébriques au-dessus de $S$.
\begin{enumerate}
\item Si $V' \subset V$ et si
$U \subset U' \subset W$ sont des sous-espaces ouverts et $W' = U' \cup f(V')$,
alors $(U' \subset W', f|_{V'} : V' \to W')$ est un carré distingué
élémentaire.
\item Si $p : W' \to W$ est un morphisme d'espaces algébriques, alors
$(p^{-1}(U) \subset W', V \times_W W' \to W')$ est un carré distingué
élémentaire.
\item Si $S' \to S$ est un morphisme de schémas, alors
$(S' \times_S U \subset S' \times_S W, S' \times_S V \to S' \times_S W)$
est un carré distingué élémentaire.
\end{enumerate}
\end{lemma}

\begin{proof}
Omis.
\end{proof}

\begin{lemma}
\label{lemma-induction-principle}
Soit $S$ un schéma. Soit $X$ un espace algébrique quasi-compact et quasi-séparé
au-dessus de $S$. Soit $P$ une propriété des objets quasi-compacts
et quasi-séparés de $X_{spaces, \etale}$. Supposons que
\begin{enumerate}
\item tout objet affine de $X_{spaces, \etale}$ v\'erifie $P$ ;
\item pour tout carré distingué élémentaire $(U \subset W, f : V \to W)$
tel que
\begin{enumerate}
\item $W$ soit un objet quasi-compact et quasi-séparé de
$X_{spaces, \etale}$ ;
\item $U$ soit quasi-compact ;
\item $V$ soit affine ;
\item $U$, $V$ et $U \times_W V$ v\'erifient $P$,
\end{enumerate}
alors $W$ v\'erifie $P$.
\end{enumerate}
Alors tout objet quasi-compact et quasi-séparé de
$X_{spaces, \etale}$ v\'erifie $P$, et en particulier $X$.
\end{lemma}

\begin{proof}
Nous affirmons d'abord que tout objet représentable, quasi-compact et
quasi-séparé de $X_{spaces, \etale}$ v\'erifie $P$.
Supposons en effet que $U \to X$ soit étale et que $U$ soit un schéma
quasi-compact et quasi-séparé. Par l'hypothèse (1), tout ouvert affine de
$U$ v\'erifie $P$. De plus, si $W, V \subset U$ sont des ouverts
quasi-compacts, avec $V$ affine, et si $W$, $V$ et $W \cap V$ v\'erifient
$P$, alors $W \cup V$ v\'erifie $P$ d'après (2)
(car le couple $(W \subset W \cup V, V \to W \cup V)$ est un carré
distingué élémentaire). Le principe de récurrence pour les schémas montre donc
que $U$ v\'erifie $P$; voir
Cohomology of Schemes, Lemma \ref{coherent-lemma-induction-principle}.

\medskip\noindent
Pour achever la preuve, il suffit de montrer que $X$ v\'erifie $P$
(car nous pouvons remplacer $X$ par n'importe quel objet quasi-compact et quasi-séparé
de $X_{spaces, \etale}$ auquel nous voulons appliquer le résultat).
Nous utiliserons la filtration
$$
\emptyset = U_{n + 1} \subset
U_n \subset U_{n - 1} \subset \ldots \subset U_1 = X
$$
et les morphismes $f_p : V_p \to U_p$ du
Decent Spaces, Lemma
\ref{decent-spaces-lemma-filter-quasi-compact-quasi-separated}.
Nous allons montrer par récurrence descendante sur $p$ que $U_p$ v\'erifie
$P$. Notons que $U_{n + 1}$ v\'erifie $P$ d'après (1), puisque l'espace
algébrique vide est affine. Supposons que $U_{p + 1}$ v\'erifie $P$.
Le couple $(U_{p + 1} \subset U_p, f_p : V_p \to U_p)$ est un carré distingué
élémentaire, mais (2) ne s'applique pas nécessairement, car $V_p$ peut ne pas être affine.
Toutefois, puisque $V_p$ est un schéma quasi-compact, nous pouvons choisir un recouvrement
fini par des ouverts affines $V_p = V_{p, 1} \cup \ldots \cup V_{p, m}$.
Posons $W_{p, 0} = U_{p + 1}$ et
$$
W_{p, i} = U_{p + 1} \cup f_p(V_{p, 1} \cup \ldots \cup V_{p, i})
$$
pour $i = 1, \ldots, m$. Ce sont des sous-espaces ouverts quasi-compacts de $X$.
Nous avons alors
$$
U_{p + 1} = W_{p, 0} \subset
W_{p, 1} \subset \ldots \subset
W_{p, m} = U_p
$$
et les couples
$$
(W_{p, 0} \subset W_{p, 1}, f_p|_{V_{p, 1}}),
(W_{p, 1} \subset W_{p, 2}, f_p|_{V_{p, 2}}),\ldots,
(W_{p, m - 1} \subset W_{p, m}, f_p|_{V_{p, m}})
$$
sont des carrés distingués élémentaires d'après le
lemme \ref{lemma-make-more-elementary-distinguished-squares}.
Notons que chaque $V_{p, i}$ v\'erifie $P$ (c'est un schéma affine), de même
que $W_{p, i} \times_{W_{p, i + 1}} V_{p, i + 1}$, car celui-ci est un ouvert
quasi-compact de $V_{p, i + 1}$ et vérifie donc $P$ d'après le premier
paragraphe de cette preuve. Ainsi, (2) s'applique à chacun de ces carrés et,
par récurrence, on conclut que
$W_{p, 1}, \ldots, W_{p, m} = U_p$ vérifient $P$.
\end{proof}

\begin{lemma}
\label{lemma-induction-principle-separated}
Soit $S$ un schéma. Soit $X$ un espace algébrique quasi-compact et quasi-séparé
au-dessus de $S$. Soit
$\mathcal{B} \subset \Ob(X_{spaces, \etale})$.
Soit $P$ une propriété des éléments de $\mathcal{B}$.
Supposons que
\begin{enumerate}
\item tout $W \in \mathcal{B}$ soit quasi-compact et quasi-séparé ;
\item si $W \in \mathcal{B}$ et si $U \subset W$ est un ouvert quasi-compact, alors
$U \in \mathcal{B}$ ;
\item si $V \in \Ob(X_{spaces, \etale})$ est affine, alors
(a) $V \in \mathcal{B}$ et (b) $V$ v\'erifie $P$ ;
\item pour tout carré distingué élémentaire $(U \subset W, f : V \to W)$
tel que
\begin{enumerate}
\item $W \in \mathcal{B}$ ;
\item $U$ soit quasi-compact ;
\item $V$ soit affine ;
\item $U$, $V$ et $U \times_W V$ v\'erifient $P$,
\end{enumerate}
alors $W$ v\'erifie $P$.
\end{enumerate}
Alors tout $W \in \mathcal{B}$ v\'erifie $P$.
\end{lemma}

\begin{proof}
La preuve est exactement la même que celle du
lemme \ref{lemma-induction-principle}.
(Notons que (4)(d) a un sens puisque $U \times_W V$ est un ouvert
quasi-compact de $V$, donc un élément de $\mathcal{B}$ d'après les conditions
(2) et (3).)
\end{proof}

\begin{remark}
\label{remark-how-to}
Comment choisir la collection $\mathcal{B}$ dans le
lemme \ref{lemma-induction-principle-separated} ?
Voici quelques exemples :
\begin{enumerate}
\item Si $X$ est quasi-compact et séparé, nous pouvons choisir pour
$\mathcal{B}$ l'ensemble des objets quasi-compacts et séparés
de $X_{spaces, \etale}$. Alors $X \in \mathcal{B}$ et $\mathcal{B}$
satisfait (1), (2) et (3)(a). Avec ce choix de $\mathcal{B}$, le
lemme \ref{lemma-induction-principle-separated} redonne le
lemme \ref{lemma-induction-principle}.
\item Si $X$ est quasi-compact et à diagonale affine
sur $\mathbf{Z}$ (comme dans Properties of Spaces, Définition
\ref{spaces-properties-definition-separated}), nous pouvons choisir pour
$\mathcal{B}$ l'ensemble des objets de $X_{spaces, \etale}$ qui sont
quasi-compacts et dont la diagonale sur $\mathbf{Z}$ est affine. Là encore,
$X \in \mathcal{B}$ et $\mathcal{B}$ satisfait (1), (2) et (3)(a).
\item Si $X$ est quasi-compact et quasi-séparé, le plus petit sous-ensemble
$\mathcal{B}$ qui contient $X$ et satisfait (1), (2) et (3)(a) est donné par
la règle suivante : $W \in \mathcal{B}$ si et seulement si $W$ est soit un
sous-espace ouvert quasi-compact de $X$, soit un ouvert quasi-compact d'un
objet affine de $X_{spaces, \etale}$.
\end{enumerate}
\end{remark}

\noindent
Voici une variante dans laquelle on propage la validit\'e d'une propri\'et\'e
d'un ouvert \`a des ouverts plus grands.

\begin{lemma}
\label{lemma-induction-principle-enlarge}
Soit $S$ un sch\'ema. Soit $X$ un espace alg\'ebrique quasi-compact et
quasi-s\'epar\'e sur $S$. Soit $W \subset X$ un sous-espace ouvert
quasi-compact. Soit $P$ une propri\'et\'e des sous-espaces ouverts
quasi-compacts de $X$. Supposons que
\begin{enumerate}
\item $W$ v\'erifie $P$, et
\item pour tout carr\'e distingu\'e \`el\'ementaire
$(W_1 \subset W_2, f : V \to W_2)$ satisfaisant aux conditions suivantes :
\begin{enumerate}
\item $W_1$ et $W_2$ soient des sous-espaces ouverts quasi-compacts de $X$,
\item $W \subset W_1$,
\item $V$ soit affine, et
\item $W_1$ v\'erifie $P$,
\end{enumerate}
alors $W_2$ v\'erifie $P$.
\end{enumerate}
Alors $X$ v\'erifie $P$.
\end{lemma}

\begin{proof}
On peut d\'eduire ce r\'esultat du
lemme \ref{lemma-induction-principle-separated}, mais nous allons plut\^ot
donner un argument direct en reprenant explicitement la preuve du
lemme \ref{lemma-induction-principle}. Nous utiliserons la filtration
$$
\emptyset = U_{n + 1} \subset
U_n \subset U_{n - 1} \subset \ldots \subset U_1 = X
$$
et les morphismes $f_p : V_p \to U_p$ du lemme
\ref{decent-spaces-lemma-filter-quasi-compact-quasi-separated} de
Decent Spaces. Nous allons montrer, par r\'ecurrence descendante sur $p$,
que $W_p = W \cup U_p$ v\'erifie $P$. Cela ach\`evera la preuve puisque
$W_1 = X$. Notons que $W_{n + 1} = W \cup U_{n + 1} = W$ v\'erifie $P$
d'apr\`es (1). Supposons que $W_{p + 1}$ v\'erifie $P$. Le sous-espace
$W_p \setminus W_{p + 1}$, muni de la structure r\'eduite induite, est un
sous-espace ferm\'e de $U_p \setminus U_{p + 1}$.
Comme $(U_{p + 1} \subset U_p, f_p : V_p \to U_p)$ est un carr\'e distingu\'e
\'el\'ementaire, il en est de m\^eme de
$(W_{p + 1} \subset W_p, f_p : V_p \to W_p)$.
Il se peut toutefois que (2) ne s'applique pas, car $V_p$ n'est pas
n\'ecessairement affine. Puisque $V_p$ est un sch\'ema quasi-compact, nous
pouvons choisir un recouvrement ouvert affine fini
$V_p = V_{p, 1} \cup \ldots \cup V_{p, m}$. Posons
$W_{p, 0} = W_{p + 1}$ et
$$
W_{p, i} = W_{p + 1} \cup f_p(V_{p, 1} \cup \ldots \cup V_{p, i})
$$
pour $i = 1, \ldots, m$. Ce sont des sous-espaces ouverts quasi-compacts
de $X$ contenant $W$. On a alors
$$
W_{p + 1} = W_{p, 0} \subset
W_{p, 1} \subset \ldots \subset
W_{p, m} = W_p
$$
et les couples
$$
(W_{p, 0} \subset W_{p, 1}, f_p|_{V_{p, 1}}),
(W_{p, 1} \subset W_{p, 2}, f_p|_{V_{p, 2}}),\ldots,
(W_{p, m - 1} \subset W_{p, m}, f_p|_{V_{p, m}})
$$
sont des carr\'es distingu\'es \`el\'ementaires d'apr\`es le
lemme \ref{lemma-make-more-elementary-distinguished-squares}.
La condition (2) s'applique donc \`a chacun d'eux et, par r\'ecurrence sur
$i$, on conclut que $W_{p, 1}, \ldots, W_{p, m} = W_p$ v\'erifient $P$.
\end{proof}



\section{Mayer-Vietoris}
\label{section-mayer-vietoris}

\noindent
Dans cette section, nous montrons qu'un carr\'e distingu\'e \`el\'ementaire
donne naissance \`a diverses suites de Mayer--Vietoris.

\medskip\noindent
Soit $S$ un sch\'ema et soit $U \to X$ un morphisme \'etale d'espaces
alg\'ebriques sur $S$. Dans la section
\ref{spaces-properties-section-localize} de Propri\'et\'es des espaces,
on a montr\'e que
$U_{spaces, \etale} = X_{spaces, \etale}/U$
de fa\c{c}on compatible avec les faisceaux structuraux. Dans cette situation,
on consid\`ere donc souvent le morphisme $j_U : U \to X$ comme un morphisme
de localisation (voir Modules sur les sites, d\'efinition
\ref{sites-modules-definition-localize-ringed-site}).
En particulier, on consid\`ere l'image inverse $j_U^*$ comme la restriction
\`a $U$ et on la note souvent ${}|_U$; cette convention est compatible avec
l'\'equation
(\ref{spaces-properties-equation-restrict-modules}).
de Propri\'et\'es des espaces. On voit notamment que
\begin{equation}
\label{equation-stalk-restriction}
(\mathcal{F}|_U)_{\overline{u}} = \mathcal{F}_{\overline{x}}
\end{equation}
si $\overline{u}$ est un point g\'eom\'etrique de $U$ et $\overline{x}$ son
image dans $X$. En outre, le foncteur de restriction admet un adjoint \`a
gauche exact $j_{U!}$; voir Modules sur les sites, lemmes
\ref{sites-modules-lemma-extension-by-zero} et
\ref{sites-modules-lemma-extension-by-zero-exact}.
Enfin, rappelons que, si $\mathcal{G}$ est un $\mathcal{O}_X$-module, alors
\begin{equation}
\label{equation-stalk-j-shriek}
(j_{U!}\mathcal{G})_{\overline{x}} =
\bigoplus\nolimits_{\overline{u}} \mathcal{G}_{\overline{u}}
\end{equation}
pour tout point g\'eom\'etrique $\overline{x} : \Spec(k) \to X$, o\`u la
somme directe porte sur les morphismes $\overline{u} : \Spec(k) \to U$
tels que $j_U \circ \overline{u} = \overline{x}$; voir Modules sur les sites,
lemme \ref{sites-modules-lemma-stalk-j-shriek}, et Propri\'et\'es des espaces,
lemme
\ref{spaces-properties-lemma-points-small-etale-site}.

\begin{lemma}
\label{lemma-exact-sequence-lower-shriek}
Soit $S$ un sch\'ema et soit $(U \subset X, V \to X)$ un carr\'e distingu\'e
\'el\'ementaire d'espaces alg\'ebriques sur $S$.
\begin{enumerate}
\item Pour tout faisceau de $\mathcal{O}_X$-modules $\mathcal{F}$, on a une
suite exacte courte
$$
0 \to j_{U \times_X V!}\mathcal{F}|_{U \times_X V} \to
j_{U!}\mathcal{F}|_U \oplus j_{V!}\mathcal{F}|_V \to \mathcal{F} \to 0
$$
\item Pour tout objet $E$ de $D(\mathcal{O}_X)$, on a un triangle distingu\'e
$$
j_{U \times_X V!}E|_{U \times_X V} \to
j_{U!}E|_U \oplus j_{V!}E|_V \to E \to 
j_{U \times_X V!}E|_{U \times_X V}[1]
$$
dans $D(\mathcal{O}_X)$.
\end{enumerate}
\end{lemma}

\begin{proof}
Pour montrer que la suite de (1) est exacte, il suffit de le v\'erifier sur
les fibres aux points g\'eom\'etriques, d'apr\`es Propri\'et\'es des espaces,
th\'eor\`eme
\ref{spaces-properties-theorem-exactness-stalks}.
Soit $\overline{x}$ un point g\'eom\'etrique de $X$. En prenant les fibres en
$\overline{x}$ et en utilisant les \`equations
(\ref{equation-stalk-restriction}) et (\ref{equation-stalk-j-shriek}), on
obtient la suite
$$
0 \to
\bigoplus\nolimits_{(\overline{u}, \overline{v})} \mathcal{F}_{\overline{x}}
\to
\bigoplus\nolimits_{\overline{u}} \mathcal{F}_{\overline{x}}
\oplus
\bigoplus\nolimits_{\overline{v}} \mathcal{F}_{\overline{x}}
\to
\mathcal{F}_{\overline{x}} \to 0
$$
Cette suite est exacte car, pour tout $\overline{x}$, ou bien il existe
exactement un $\overline{u}$ qui s'envoie sur $\overline{x}$, ou bien il
n'en existe aucun et il existe exactement un $\overline{v}$ qui s'envoie
sur $\overline{x}$.

\medskip\noindent
Preuve de (2). Nous avons vu dans Cohomologie sur les sites, section
\ref{sites-cohomology-section-properties-K-injective}, que les foncteurs de
restriction et de prolongement par z\'ero entre cat\'egories d\'eriv\'ees se
calculent en appliquant simplement le foncteur terme \`a terme \`a un
complexe quelconque. Soit $\mathcal{E}^\bullet$ un complexe de
$\mathcal{O}_X$-modules repr\'esentant $E$. Le triangle distingu\'e du lemme
est celui associ\'e, d'apr\`es Cat\'egories d\'eriv\'ees, section
\ref{derived-section-canonical-delta-functor}, et plus particuli\`erement le
lemme \ref{derived-lemma-derived-canonical-delta-functor}, \`a la suite exacte
courte de complexes de $\mathcal{O}_X$-modules
$$
0 \to j_{U \times_X V!}\mathcal{E}^\bullet|_{U \times_X V} \to
j_{U!}\mathcal{E}^\bullet|_U \oplus j_{V!}\mathcal{E}^\bullet|_V
\to \mathcal{E}^\bullet \to 0
$$
qui est exacte courte d'apr\`es (1).
\end{proof}

\begin{lemma}
\label{lemma-exact-sequence-j-star}
Soit $S$ un sch\'ema et soit $(U \subset X, V \to X)$ un carr\'e distingu\'e
\'el\'ementaire d'espaces alg\'ebriques sur $S$.
\begin{enumerate}
\item Pour tout faisceau de $\mathcal{O}_X$-modules $\mathcal{F}$, on a une
suite exacte courte
$$
0 \to \mathcal{F} \to
j_{U, *}\mathcal{F}|_U \oplus j_{V, *}\mathcal{F}|_V \to
j_{U \times_X V, *}\mathcal{F}|_{U \times_X V} \to 0
$$
\item Pour tout objet $E$ de $D(\mathcal{O}_X)$, on a un triangle distingu\'e
$$
E \to 
Rj_{U, *}E|_U \oplus Rj_{V, *}E|_V \to
Rj_{U \times_X V, *}E|_{U \times_X V} \to
E[1]
$$
dans $D(\mathcal{O}_X)$.
\end{enumerate}
\end{lemma}

\begin{proof}
Soit $W$ un objet de $X_\etale$. Nous affirmons que la suite
$$
0 \to
\mathcal{F}(W) \to
\mathcal{F}(W \times_X U) \oplus \mathcal{F}(W \times_X V) \to
\mathcal{F}(W \times_X U \times_X V)
$$
est exacte et que, localement sur $W$, tout \`el\'ement du dernier groupe
se rel\`eve dans celui du milieu. D'apr\`es le
lemme \ref{lemma-make-more-elementary-distinguished-squares}, les donn\'ees
$(W \times_X U \subset W, V \times_X W \to W)$ forment un carr\'e distingu\'e
\'el\'ementaire. Nous pouvons donc supposer que $W = X$, et il suffit
d'\'etablir le m\^eme fait pour
$$
0 \to
\mathcal{F}(X) \to
\mathcal{F}(U) \oplus \mathcal{F}(V) \to
\mathcal{F}(U \times_X V)
$$
Nous avons vu que
$$
0 \to j_{U \times_X V!}\mathcal{O}_{U \times_X V}
\to j_{U!}\mathcal{O}_U \oplus
j_{V!}\mathcal{O}_V \to
\mathcal{O}_X \to 0
$$
est une suite exacte de $\mathcal{O}_X$-modules, d'apr\`es le
lemme \ref{lemma-exact-sequence-lower-shriek}. L'application du foncteur
contravariant $\Hom_{\mathcal{O}_X}(- , \mathcal{F})$ donne alors la suite
ci-dessus. Cela signifie aussi que l'obstruction au rel\`evement de
$s \in \mathcal{F}(U \times_X V)$ en un \`el\'ement de
$\mathcal{F}(U) \oplus \mathcal{F}(V)$ appartient \`a
$\Ext^1_{\mathcal{O}_X}(\mathcal{O}_X, \mathcal{F}) =
H^1(X, \mathcal{F})$. Par localit\'e de la cohomologie (Cohomologie sur les
sites, lemme
\ref{sites-cohomology-lemma-kill-cohomology-class-on-covering})
cette obstruction s'annule localement pour la topologie \'etale sur $X$, ce
qui ach\`eve la preuve de (1).

\medskip\noindent
Preuve de (2). Choisissons un complexe K-injectif
$\mathcal{I}^\bullet$ repr\'esentant $E$ dont les termes $\mathcal{I}^n$
sont des objets injectifs de $\textit{Mod}(\mathcal{O}_X)$; voir Injectifs,
th\'eor\`eme
\ref{injectives-theorem-K-injective-embedding-grothendieck}.
Alors $\mathcal{I}^\bullet|U$ est un complexe K-injectif (Cohomologie sur
les sites, lemme
\ref{sites-cohomology-lemma-restrict-K-injective-to-open}).
Ainsi $Rj_{U, *}E|_U$ est repr\'esent\'e par
$j_{U, *}\mathcal{I}^\bullet|_U$, et de m\^eme pour $V$ et $U \times_X V$.
Le triangle distingu\'e du lemme est donc celui associ\'e, d'apr\`es
Cat\'egories d\'eriv\'ees, section
\ref{derived-section-canonical-delta-functor}, et plus particuli\`erement le
lemme \ref{derived-lemma-derived-canonical-delta-functor}, \`a la suite exacte
courte de complexes
$$
0 \to
\mathcal{I}^\bullet \to
j_{U, *}\mathcal{I}^\bullet|_U \oplus j_{V, *}\mathcal{I}^\bullet|_V \to
j_{U \times_X V, *}\mathcal{I}^\bullet|_{U \times_X V} \to
0.
$$
Cette suite est exacte d'apr\`es (1).
\end{proof}

\begin{lemma}
\label{lemma-unbounded-relative-mayer-vietoris}
Soit $S$ un sch\'ema. Soit $f : X \to Y$ un morphisme d'espaces alg\'ebriques
sur $S$, et soit $(U \subset X, V \to X)$ un carr\'e distingu\'e
\'el\'ementaire. Notons $a = f|_U : U \to Y$, $b = f|_V : V \to Y$ et
$c = f|_{U \times_X V} : U \times_X V \to Y$ les restrictions.
Pour tout objet $E$ de $D(\mathcal{O}_X)$, il existe un triangle distingu\'e
$$
Rf_*E \to
Ra_*(E|_U) \oplus Rb_*(E|_V) \to
Rc_*(E|_{U \times_X V}) \to
Rf_*E[1]
$$
dans $D(\mathcal{O}_Y)$. Ce triangle est fonctoriel en $E$.
\end{lemma}

\begin{proof}
Choisissons un complexe K-injectif $\mathcal{I}^\bullet$ repr\'esentant $E$.
Nous pouvons supposer que $\mathcal{I}^n$ est un objet injectif de
$\textit{Mod}(\mathcal{O}_X)$ pour tout $n$; voir Injectifs, th\'eor\`eme
\ref{injectives-theorem-K-injective-embedding-grothendieck}.
Alors $Rf_*E$ est calcul\'e par $f_*\mathcal{I}^\bullet$. Il en est de m\^eme
sur $U$, $V$ et $U \times_X V$, d'apr\`es Cohomologie sur les sites, lemme
\ref{sites-cohomology-lemma-restrict-K-injective-to-open}. Le triangle
distingu\'e du lemme est donc celui associ\'e, d'apr\`es Cat\'egories
d\'eriv\'ees, section \ref{derived-section-canonical-delta-functor}, et plus
particuli\`erement le lemme
\ref{derived-lemma-derived-canonical-delta-functor}, \`a la suite exacte
courte de complexes
$$
0 \to
f_*\mathcal{I}^\bullet \to
a_*\mathcal{I}^\bullet|_U \oplus b_*\mathcal{I}^\bullet|_V \to
c_*\mathcal{I}^\bullet|_{U \times_X V} \to
0.
$$
Pour voir qu'il s'agit d'une suite exacte courte de complexes, raisonnons
comme suit. Prenons un objet injectif $\mathcal{I}$ de
$\textit{Mod}(\mathcal{O}_X)$. Appliquons $f_*$ \`a la suite exacte courte
$$
0 \to \mathcal{I} \to
j_{U, *}\mathcal{I}|_U \oplus j_{V, *}\mathcal{I}|_V \to
j_{U \times_X V, *}\mathcal{I}|_{U \times_X V} \to 0
$$
du lemme \ref{lemma-exact-sequence-j-star}. Comme
$R^1f_*\mathcal{I} = 0$, on obtient la suite exacte courte
$$
0 \to f_*\mathcal{I} \to
f_*j_{U, *}\mathcal{I}|_U \oplus f_*j_{V, *}\mathcal{I}|_V \to
f_*j_{U \times_X V, *}\mathcal{I}|_{U \times_X V} \to 0
$$
Pour conclure, il suffit d'observer que $a_* = f_*j_{U, *}$, et de m\^eme
pour $b_*$ et $c_*$.
\end{proof}

\begin{lemma}
\label{lemma-mayer-vietoris-hom}
Soit $S$ un sch\'ema et soit $(U \subset X, V \to X)$ un carr\'e distingu\'e
\'el\'ementaire d'espaces alg\'ebriques sur $S$. Pour tous objets $E$ et $F$
de $D(\mathcal{O}_X)$, on a une suite de Mayer--Vietoris
$$
\xymatrix{
& \ldots \ar[r] &
\Ext^{-1}(E_{U \times_X V}, F_{U \times_X V}) \ar[lld] \\
\Hom(E, F) \ar[r] &
\Hom(E_U, F_U) \oplus
\Hom(E_V, F_V) \ar[r] &
\Hom(E_{U \times_X V}, F_{U \times_X V})
}
$$
o\`u les indices d\'esignent les restrictions aux ouverts correspondants et
les $\Hom$ sont pris dans les cat\'egories d\'eriv\'ees correspondantes.
\end{lemma}

\begin{proof}
Le triangle distingu\'e du lemme
\ref{lemma-exact-sequence-lower-shriek} donne une suite exacte longue de
$\Hom$ (Cat\'egories d\'eriv\'ees, lemme
\ref{derived-lemma-representable-homological}). On utilise ensuite l'\'egalit\'e
$\Hom(j_{U!}E|_U, F) = \Hom(E|_U, F|_U)$, qui r\'esulte de Cohomologie sur les
sites, lemme
\ref{sites-cohomology-lemma-adjoint-lower-shriek-restrict}.
\end{proof}

\begin{lemma}
\label{lemma-unbounded-mayer-vietoris}
Soit $S$ un sch\'ema et soit $(U \subset X, V \to X)$ un carr\'e distingu\'e
\'el\'ementaire d'espaces alg\'ebriques sur $S$. Pour tout objet $E$ de
$D(\mathcal{O}_X)$, on a un triangle distingu\'e
$$
R\Gamma(X, E) \to R\Gamma(U, E) \oplus R\Gamma(V, E) \to
R\Gamma(U \times_X V, E) \to R\Gamma(X, E)[1]
$$
et, en particulier, une suite exacte longue de cohomologie
$$
\ldots \to
H^n(X, E) \to
H^n(U, E) \oplus H^n(V, E) \to
H^n(U \times_X V, E) \to
H^{n + 1}(X, E) \to \ldots
$$
La construction du triangle distingu\'e et de la suite exacte longue est
fonctorielle en $E$.
\end{lemma}

\begin{proof}
Choisissons un complexe K-injectif $\mathcal{I}^\bullet$ repr\'esentant $E$
dont les termes $\mathcal{I}^n$ sont des objets injectifs de
$\textit{Mod}(\mathcal{O}_X)$; voir Injectifs, th\'eor\`eme
\ref{injectives-theorem-K-injective-embedding-grothendieck}.
Dans la preuve du lemme \ref{lemma-exact-sequence-j-star}, nous avons obtenu
une suite exacte courte de complexes
$$
0 \to \mathcal{I}^\bullet \to
j_{U, *}\mathcal{I}^\bullet|_U \oplus j_{V, *}\mathcal{I}^\bullet|_V \to
j_{U \times_X V, *}\mathcal{I}^\bullet|_{U \times_X V} \to 0
$$
Comme $H^1(X, \mathcal{I}^n) = 0$, le passage aux sections globales donne
une suite exacte de complexes
$$
0 \to \Gamma(X, \mathcal{I}^\bullet) \to
\Gamma(U, \mathcal{I}^\bullet) \oplus
\Gamma(V, \mathcal{I}^\bullet) \to
\Gamma(U \times_X V, \mathcal{I}^\bullet) \to 0
$$
Puisque ces complexes repr\'esentent $R\Gamma(X, E)$, $R\Gamma(U, E)$,
$R\Gamma(V, E)$ et $R\Gamma(U \times_X V, E)$, on obtient un triangle
distingu\'e d'apr\`es Cat\'egories d\'eriv\'ees, section
\ref{derived-section-canonical-delta-functor}, et plus particuli\`erement le
lemme \ref{derived-lemma-derived-canonical-delta-functor}.
\end{proof}

\begin{lemma}
\label{lemma-restrict-lower-shriek}
Soit $S$ un sch\'ema. Soit $j : U \to X$ un morphisme \'etale d'espaces
alg\'ebriques sur $S$. Pour tout morphisme \'etale $V \to X$, posons
$W = V \times_X U$ et notons $j_W : W \to V$ le morphisme de projection.
Alors $(j_!E)|_V = j_{W!}(E|_W)$ pour tout $E$ de $D(\mathcal{O}_U)$.
\end{lemma}

\begin{proof}
Cela r\'esulte de l'\'egalit\'e
$(j_!\mathcal{F})|_V = j_{W!}(\mathcal{F}|_W)$ pour tout
$\mathcal{O}_X$-module $\mathcal{F}$, laquelle d\'ecoule imm\'ediatement de la
construction des foncteurs $j_!$ et $j_{W!}$; voir Modules sur les sites,
lemme \ref{sites-modules-lemma-extension-by-zero}.
\end{proof}

\begin{lemma}
\label{lemma-pushforward-with-support-in-open}
Soit $S$ un sch\'ema et soit $(U \subset X, j : V \to X)$ un carr\'e
distingu\'e \`el\'ementaire d'espaces alg\'ebriques sur $S$. Posons
$T = |X| \setminus |U|$.
\begin{enumerate}
\item Si $E$ est un objet de $D(\mathcal{O}_X)$ \`a support dans $T$, alors
(a) $E \to Rj_*(E|_V)$ et (b) $j_!(E|_V) \to E$ sont des isomorphismes.
\item Si $F$ est un objet de $D(\mathcal{O}_V)$ \`a support dans $j^{-1}T$,
alors (a) $F \to (j_!F)|_V$, (b) $(Rj_*F)|_V \to F$ et (c)
$j_!F \to Rj_*F$ sont des isomorphismes.
\end{enumerate}
\end{lemma}

\begin{proof}
Soit $E$ un objet de $D(\mathcal{O}_X)$ dont les faisceaux de cohomologie
sont \`a support dans $T$. Alors $E|_U = 0$ et
$E|_{U \times_X V} = 0$, puisque $T$ ne rencontre pas $U$ et que $j^{-1}T$
ne rencontre pas $U \times_X V$. L'assertion (1)(a) r\'esulte donc du
lemme \ref{lemma-exact-sequence-j-star}. De la m\^eme mani\`ere, (1)(b)
r\'esulte du lemme \ref{lemma-exact-sequence-lower-shriek}.

\medskip\noindent
Soit $F$ un objet de $D(\mathcal{O}_V)$ dont les faisceaux de cohomologie
sont \`a support dans $j^{-1}T$. D'apr\`es le
lemme \ref{lemma-restrict-direct-image-open}, on a
$(Rj_*F)|_U = Rj_{W, *}(F|_W) = 0$, puisque $F|_W = 0$ par hypoth\`ese.
De m\^eme, $(j_!F)|_U = j_{W!}(F|_W) = 0$ d'apr\`es le
lemme \ref{lemma-restrict-lower-shriek}. Ainsi $j_!F$ et $Rj_*F$ sont \`a
support dans $T$, tandis que $(j_!F)|_V$ et $(Rj_*F)|_V$ sont \`a support
dans $j^{-1}(T)$. Pour v\'erifier que les morphismes (2)(a), (b) et (c) sont
des isomorphismes dans la cat\'egorie d\'eriv\'ee, il suffit de v\'erifier
qu'ils induisent des isomorphismes sur les fibres des faisceaux de
cohomologie aux points g\'eom\'etriques de $T$ et de $j^{-1}(T)$, d'apr\`es
Propri\'et\'es des espaces, th\'eor\`eme
\ref{spaces-properties-theorem-exactness-stalks}.
Nous pouvons le faire apr\`es avoir remplac\'e $X$ par $V$, $U$ par
$U \times_X V$, $V$ par $V \times_X V$ et $F$ par $F|_{V \times_X V}$
(restriction suivant la premi\`ere projection); voir les lemmes
\ref{lemma-restrict-direct-image-open},
\ref{lemma-restrict-lower-shriek} et
\ref{lemma-make-more-elementary-distinguished-squares}.
Comme $V \times_X V \to V$ admet une section, cela ram\`ene (2) au cas o\`u
$j : V \to X$ admet une section.

\medskip\noindent
Supposons que $j$ admette une section $\sigma : X \to V$. Posons
$V' = \sigma(X)$; c'est un sous-espace ouvert de $V$. Posons aussi
$U' = j^{-1}(U)$; c'est un autre sous-espace ouvert de $V$. Alors
$(U' \subset V, V' \to V)$ est un carr\'e distingu\'e \`el\'ementaire.
Comme $F$ est \`a support dans $j^{-1}(T)$, on a $F|_{U'} = 0$ et
$F|_{V' \cap U'} = 0$. Notons $j' : V' \to V$ l'immersion ouverte et
$j_{V'} : V' \to X$ le compos\'e $V' \to V \to X$, qui est l'inverse de
$\sigma$. Posons $F' = \sigma^*F$. Les triangles distingu\'es des lemmes
\ref{lemma-exact-sequence-lower-shriek} et
\ref{lemma-exact-sequence-j-star} montrent que
$F = j'_!(F|_{V'})$ et $F = Rj'_*(F|_{V'})$. Il s'ensuit que
$j_!F = j_!j'_!(F|_{V'}) = j_{V'!}F = F'$, car
$j_{V'} : V' \to X$ est un isomorphisme inverse de $\sigma$. De m\^eme,
$Rj_*F = Rj_*Rj'_*F = Rj_{V', *}F = F'$. Cela prouve (2)(c). Pour prouver
(2)(a) et (2)(b), il suffit de montrer que $F = F'|_V$. C'est clair, car
$F$ et $F'|_V$ ont tous deux pour restriction z\'ero sur $U'$ et sur
$U' \cap V'$, et le m\^eme objet sur $V'$.
\end{proof}

\noindent
On peut recoller des complexes !

\begin{lemma}
\label{lemma-glue}
Soit $S$ un sch\'ema et soit $(U \subset X, V \to X)$ un carr\'e distingu\'e
\'el\'ementaire d'espaces alg\'ebriques sur $S$. Supposons donn\'es
\begin{enumerate}
\item un objet $A$ de $D(\mathcal{O}_U)$,
\item un objet $B$ de $D(\mathcal{O}_V)$, et
\item un isomorphisme $c : A|_{U \times_X V} \to B|_{U \times_X V}$.
\end{enumerate}
Alors il existe un objet $F$ de $D(\mathcal{O}_X)$ et des isomorphismes
$f : F|_U \to A$ et $g : F|_V \to B$ tels que
$c = g|_{U \times_X V} \circ f^{-1}|_{U \times_X V}$. En outre, soient
\begin{enumerate}
\item un objet $E$ de $D(\mathcal{O}_X)$,
\item un morphisme $a : A \to E|_U$ dans $D(\mathcal{O}_U)$,
\item un morphisme $b : B \to E|_V$ dans $D(\mathcal{O}_V)$,
\end{enumerate}
tels que
$$
a|_{U \times_X V}  = b|_{U \times_X V} \circ c.
$$
Il existe alors un morphisme $F \to E$ dans $D(\mathcal{O}_X)$ dont la
restriction \`a $U$ est $a \circ f$ et la restriction \`a $V$ est
$b \circ g$.
\end{lemma}

\begin{proof}
Notons $j_U$, $j_V$ et $j_{U \times_X V}$ les morphismes correspondants
vers $X$. Choisissons un triangle distingu\'e
$$
F \to Rj_{U, *}A \oplus Rj_{V, *}B \to
Rj_{U \times_X V, *}(B|_{U \times_X V}) \to F[1]
$$
Ici, le morphisme
$Rj_{V, *}B \to Rj_{U \times_X V, *}(B|_{U \times_X V})$ est le morphisme
naturel. Le morphisme
$Rj_{U, *}A \to Rj_{U \times_X V, *}(B|_{U \times_X V})$
est le compos\'e de
$Rj_{U, *}A \to Rj_{U \times_X V, *}(A|_{U \times_X V})$
avec $Rj_{U \times_X V, *}c$. En restreignant \`a $U$, on obtient
$$
F|_U \to A \oplus (Rj_{V, *}B)|_U \to
(Rj_{U \times_X V, *}(B|_{U \times_X V}))|_U \to F|_U[1]
$$
Notons $j : U \times_X V \to U$. La compatibilit\'e entre restriction et
image directe totale (lemme \ref{lemma-restrict-direct-image-open}) montre
que $(Rj_{V, *}B)|_U$ et
$(Rj_{U \times_X V, *}(B|_{U \times_X V}))|_U$
sont tous deux canoniquement isomorphes \`a
$Rj_*(B|_{U \times_X V})$. Le second morphisme de la derni\`ere formule
affich\'ee admet donc une section, d'o\`u l'on conclut que
$F|_U \to A$ est un isomorphisme.

\medskip\noindent
Pour voir que $F|_V \to B$ est un isomorphisme, nous emploierons l'astuce
suivante. Choisissons un triangle distingu\'e
$$
F|_V \to B \to B' \to F[1]|_V
$$
dans $D(\mathcal{O}_V)$. Puisque $F|_U \to A$ est un isomorphisme et que
l'on dispose de l'isomorphisme
$c : A|_{U \times_X V} \to B|_{U \times_X V}$, la restriction de
$F|_V \to B$ \`a $U \times_X V$ est un isomorphisme. Ainsi $B'$ est \`a
support dans $j_V^{-1}(T)$, o\`u $T = |X| \setminus |U|$. D'autre part, on a
un morphisme de triangles distingu\'es
$$
\xymatrix{
F \ar[r] \ar[d] &
Rj_{U, *}F|_U \oplus Rj_{V, *}F|_V \ar[r] \ar[d] &
Rj_{U \times_X V, *}F|_{U \times_X V} \ar[r] \ar[d] &
F[1] \ar[d] \\
F \ar[r] &
Rj_{U, *}A \oplus Rj_{V, *}B \ar[r] &
Rj_{U \times_X V, *}(B|_{U \times_X V}) \ar[r] &
F[1]
}
$$
Tous les morphismes verticaux de ce diagramme, sauf peut-\^etre
$Rj_{V, *}F|_V \to Rj_{V, *}B$, sont des isomorphismes; ce dernier en est
donc aussi un (Cat\'egories d\'eriv\'ees, lemme
\ref{derived-lemma-third-isomorphism-triangle}). Il en r\'esulte que
$Rj_{V, *}B' = 0$, puis que $B' = 0$ d'apr\`es le
lemme \ref{lemma-pushforward-with-support-in-open}.

\medskip\noindent
L'existence du morphisme $F \to E$ r\'esulte de la suite de Mayer--Vietoris
pour $\Hom$; voir le lemme \ref{lemma-mayer-vietoris-hom}.
\end{proof}







\section{Le coh\'erateur}
\label{section-coherator}

\noindent
Soit $S$ un sch\'ema et soit $X$ un espace alg\'ebrique sur $S$.
Le {\it coh\'erateur} est un foncteur
$$
Q_X :
\textit{Mod}(\mathcal{O}_X)
\longrightarrow
\QCoh(\mathcal{O}_X)
$$
adjoint \`a droite au foncteur d'inclusion
$\QCoh(\mathcal{O}_X) \to \textit{Mod}(\mathcal{O}_X)$.
Il existe pour tout espace alg\'ebrique $X$ et, de plus, le morphisme
d'adjonction $Q_X(\mathcal{F}) \to \mathcal{F}$ est un isomorphisme pour tout
module quasi-coh\'erent $\mathcal{F}$; voir Propri\'et\'es des espaces,
proposition
\ref{spaces-properties-proposition-coherator}.
Comme $Q_X$ est exact \`a gauche (en tant qu'adjoint \`a droite), on peut
consid\'erer son extension d\'eriv\'ee \`a droite
$$
RQ_X :
D(\mathcal{O}_X)
\longrightarrow
D(\QCoh(\mathcal{O}_X)).
$$
Puisque $Q_X$ est adjoint \`a droite au foncteur d'inclusion
$\QCoh(\mathcal{O}_X) \to \textit{Mod}(\mathcal{O}_X)$, le foncteur $RQ_X$
est adjoint \`a droite au foncteur canonique
$D(\QCoh(\mathcal{O}_X)) \to D(\mathcal{O}_X)$, d'apr\`es Cat\'egories
d\'eriv\'ees, lemme \ref{derived-lemma-derived-adjoint-functors}.

\medskip\noindent
Dans cette section, nous \'etudierons le foncteur $RQ_X$. Dans la
section \ref{section-better-coherator}, nous \'etudierons l'adjoint \`a droite,
\'etroitement apparent\'e, du foncteur d'inclusion
$D_\QCoh(\mathcal{O}_X) \to D(\mathcal{O}_X)$, lorsqu'il existe.

\begin{lemma}
\label{lemma-affine-pushforward}
Soit $S$ un sch\'ema et soit $f : X \to Y$ un morphisme affine d'espaces
alg\'ebriques sur $S$. Alors $f_*$ d\'efinit un foncteur d\'eriv\'e
$f_* : D(\QCoh(\mathcal{O}_X)) \to D(\QCoh(\mathcal{O}_Y))$ tel que le
diagramme
$$
\xymatrix{
D(\QCoh(\mathcal{O}_X)) \ar[d]_{f_*} \ar[r] &
D_\QCoh(\mathcal{O}_X) \ar[d]^{Rf_*} \\
D(\QCoh(\mathcal{O}_Y)) \ar[r] &
D_\QCoh(\mathcal{O}_Y)
}
$$
est commutatif.
\end{lemma}

\begin{proof}
Le foncteur
$f_* : \QCoh(\mathcal{O}_X) \to \QCoh(\mathcal{O}_Y)$
est exact; voir Cohomologie des espaces, lemme
\ref{spaces-cohomology-lemma-affine-vanishing-higher-direct-images}.
Ainsi $f_*$ d\'efinit un foncteur d\'eriv\'e
$f_* : D(\QCoh(\mathcal{O}_X)) \to D(\QCoh(\mathcal{O}_Y))$ en appliquant
simplement $f_*$ \`a un complexe repr\'esentant l'objet consid\'er\'e; voir
Cat\'egories d\'eriv\'ees, lemme
\ref{derived-lemma-right-derived-exact-functor}. Pour tout complexe
$\mathcal{F}^\bullet$ de $\mathcal{O}_X$-modules, on dispose d'un morphisme
canonique
$f_*\mathcal{F}^\bullet \to Rf_*\mathcal{F}^\bullet$.
Pour achever la preuve, montrons que c'est un quasi-isomorphisme lorsque
chaque terme $\mathcal{F}^n$ du complexe $\mathcal{F}^\bullet$ est
quasi-coh\'erent. L'assertion est locale sur $Y$ pour la topologie \'etale;
nous pouvons donc supposer $Y$ affine. Tout morphisme affine \'etant
repr\'esentable, la compatibilit\'e de la
remarque \ref{remark-match-total-direct-images} nous ram\`ene au cas des
sch\'emas. Ce cas est trait\'e dans Cat\'egories d\'eriv\'ees des sch\'emas,
lemme \ref{perfect-lemma-affine-pushforward}.
\end{proof}

\begin{lemma}
\label{lemma-flat-pushforward-coherator}
Soit $S$ un sch\'ema et soit $f : X \to Y$ un morphisme d'espaces alg\'ebriques
sur $S$. Supposons $f$ quasi-compact, quasi-s\'epar\'e et plat. Si l'on note
$$
\Phi : D(\QCoh(\mathcal{O}_X)) \to D(\QCoh(\mathcal{O}_Y))
$$
le foncteur d\'eriv\'e droit de
$f_* : \QCoh(\mathcal{O}_X) \to \QCoh(\mathcal{O}_Y)$
alors $RQ_Y \circ Rf_* = \Phi \circ RQ_X$.
\end{lemma}

\begin{proof}
Nous allons le prouver en montrant que $RQ_Y \circ Rf_*$ et
$\Phi \circ RQ_X$ sont adjoints \`a droite au m\^eme foncteur
$D(\QCoh(\mathcal{O}_Y)) \to D(\mathcal{O}_X)$.

\medskip\noindent
Comme $f$ est quasi-compact et quasi-s\'epar\'e, $f_*$ pr\'eserve la
quasi-coh\'erence; voir Morphismes d'espaces, lemme
\ref{spaces-morphisms-lemma-pushforward}. Rappelons que
$\QCoh(\mathcal{O}_X)$ est une cat\'egorie ab\'elienne de Grothendieck
(Propri\'et\'es des espaces, proposition
\ref{spaces-properties-proposition-coherator}).
Ainsi tout $K$ de $D(\QCoh(\mathcal{O}_X))$ peut \^etre repr\'esent\'e par un
complexe K-injectif $\mathcal{I}^\bullet$ de $\QCoh(\mathcal{O}_X)$; voir
Injectifs, th\'eor\`eme
\ref{injectives-theorem-K-injective-embedding-grothendieck}.
On peut alors poser $\Phi(K) = f_*\mathcal{I}^\bullet$.

\medskip\noindent
Comme $f$ est plat, le foncteur $f^*$ est exact. Il d\'efinit donc les
foncteurs
$f^* : D(\mathcal{O}_Y) \to D(\mathcal{O}_X)$ et
$f^* : D(\QCoh(\mathcal{O}_Y)) \to D(\QCoh(\mathcal{O}_X))$.
Le foncteur $f^* = Lf^* : D(\mathcal{O}_Y) \to D(\mathcal{O}_X)$ est adjoint
\`a gauche \`a
$Rf_* : D(\mathcal{O}_X) \to D(\mathcal{O}_Y)$,
voir Cohomologie sur les sites, lemme
\ref{sites-cohomology-lemma-adjoint}. De m\^eme, le foncteur
$f^* : D(\QCoh(\mathcal{O}_Y)) \to D(\QCoh(\mathcal{O}_X))$
est adjoint \`a gauche \`a
$\Phi : D(\QCoh(\mathcal{O}_X)) \to D(\QCoh(\mathcal{O}_Y))$, d'apr\`es
Cat\'egories d\'eriv\'ees, lemme
\ref{derived-lemma-derived-adjoint-functors}.

\medskip\noindent
Soient $A$ un objet de $D(\QCoh(\mathcal{O}_Y))$ et $E$ un objet de
$D(\mathcal{O}_X)$. Alors
\begin{align*}
\Hom_{D(\QCoh(\mathcal{O}_Y))}(A, RQ_Y(Rf_*E))
& =
\Hom_{D(\mathcal{O}_Y)}(A, Rf_*E) \\
& =
\Hom_{D(\mathcal{O}_X)}(f^*A, E) \\
& =
\Hom_{D(\QCoh(\mathcal{O}_X))}(f^*A, RQ_X(E)) \\
& =
\Hom_{D(\QCoh(\mathcal{O}_Y))}(A, \Phi(RQ_X(E)))
\end{align*}
Cela donne le r\'esultat voulu.
\end{proof}

\begin{lemma}
\label{lemma-affine-coherator}
Soit $S$ un sch\'ema et soit $X$ un espace alg\'ebrique affine sur $S$.
Posons $A = \Gamma(X, \mathcal{O}_X)$. Alors
\begin{enumerate}
\item $Q_X : \textit{Mod}(\mathcal{O}_X) \to \QCoh(\mathcal{O}_X)$
est le foncteur qui associe \`a $\mathcal{F}$ le
$\mathcal{O}_X$-module quasi-coh\'erent associ\'e au $A$-module
$\Gamma(X, \mathcal{F})$,
\item $RQ_X : D(\mathcal{O}_X) \to D(\QCoh(\mathcal{O}_X))$
est le foncteur qui associe \`a $E$ le complexe de
$\mathcal{O}_X$-modules quasi-coh\'erents associ\'e \`a l'objet
$R\Gamma(X, E)$ de $D(A)$,
\item la restriction de $RQ_X$ \`a $D_\QCoh(\mathcal{O}_X)$ d\'efinit un
quasi-inverse de (\ref{equation-compare}).
\end{enumerate}
\end{lemma}

\begin{proof}
Soit $X_0 = \Spec(A)$ le sch\'ema affine qui repr\'esente $X$. Rappelons qu'il
existe un morphisme de sites annel\'es
$\epsilon : X_\etale \to X_{0, Zar}$
qui induit des \'equivalences
$$
\xymatrix{
\QCoh(\mathcal{O}_X) \ar@<1ex>[r]^{{\epsilon_*}} &
\QCoh(\mathcal{O}_{X_0}) \ar@<1ex>[l]^{{\epsilon^*}}
}
$$
voir le lemme
\ref{lemma-derived-quasi-coherent-small-etale-site}.
Par unicit\'e des foncteurs adjoints, on a donc
$Q_X = \epsilon^* \circ Q_{X_0} \circ \epsilon_*$. L'assertion (1) r\'esulte
alors de la description de $Q_{X_0}$ dans Cat\'egories d\'eriv\'ees des
sch\'emas, lemme \ref{perfect-lemma-affine-coherator}, et de l'\'egalit\'e
$\Gamma(X_0, \epsilon_*\mathcal{F}) = \Gamma(X, \mathcal{F})$.
L'assertion (2) r\'esulte de (1) et de l'exactitude du foncteur des
$A$-modules vers les $\mathcal{O}_X$-modules quasi-coh\'erents. La troisi\`eme
assertion d\'ecoule alors du r\'esultat pour les sch\'emas (Cat\'egories
d\'eriv\'ees des sch\'emas, lemme \ref{perfect-lemma-affine-coherator}) et du
lemme
\ref{lemma-derived-quasi-coherent-small-etale-site}.
\end{proof}

\noindent
Nous \'etablissons maintenant un crit\`ere assurant que le foncteur
$D(\QCoh(\mathcal{O}_X)) \to D_\QCoh(\mathcal{O}_X)$
est une \'equivalence.

\begin{lemma}
\label{lemma-argument-proves}
Soit $S$ un sch\'ema et soit $X$ un espace alg\'ebrique quasi-compact et
quasi-s\'epar\'e sur $S$. Supposons que, pour tout morphisme \'etale
$j : V \to W$, o\`u $W \subset X$ est ouvert quasi-compact et $V$ affine,
le foncteur d\'eriv\'e droit
$$
\Phi : D(\QCoh(\mathcal{O}_V)) \to D(\QCoh(\mathcal{O}_W))
$$
du foncteur exact \`a gauche
$j_* : \QCoh(\mathcal{O}_V) \to \QCoh(\mathcal{O}_W)$
s'inscrive dans le diagramme commutatif
$$
\xymatrix{
D(\QCoh(\mathcal{O}_V)) \ar[d]_\Phi \ar[r]_{i_V} &
D_\QCoh(\mathcal{O}_V) \ar[d]^{Rj_*} \\
D(\QCoh(\mathcal{O}_W)) \ar[r]^{i_W} &
D_\QCoh(\mathcal{O}_W)
}
$$
Alors le foncteur (\ref{equation-compare})
$$
D(\QCoh(\mathcal{O}_X))
\longrightarrow
D_\QCoh(\mathcal{O}_X)
$$
est une \'equivalence, dont $RQ_X$ est un quasi-inverse.
\end{lemma}

\begin{proof}
Nous utilisons d'abord le principe de r\'ecurrence pour montrer que $i_X$ est
pleinement fid\`ele. Plus pr\'ecis\'ement, nous appliquerons le
lemme \ref{lemma-induction-principle-enlarge}. Soit
$(U \subset W, V \to W)$ un carr\'e distingu\'e \`el\'ementaire, avec $V$
affine et $U,W$ ouverts quasi-compacts de $X$. Supposons $i_U$ pleinement
fid\`ele. Il faut montrer que $i_W$ est pleinement fid\`ele. Nous pouvons
remplacer $X$ par $W$, c'est-\`a-dire supposer que $W = X$; cette op\'eration
ne sert qu'\`a simplifier les notations et pr\'eserve la condition de l'\'enonc\'e.

\medskip\noindent
Soient $A$ et $B$ des objets de $D(\QCoh(\mathcal{O}_X))$. Nous voulons
montrer que
$$
\Hom_{D(\QCoh(\mathcal{O}_X))}(A, B)
\longrightarrow
\Hom_{D(\mathcal{O}_X)}(i_X(A), i_X(B))
$$
est bijectif. Posons $T = |X| \setminus |U|$.

\medskip\noindent
Supposons d'abord que $i_X(B)$ soit \`a support dans $T$. Dans ce cas, le
morphisme
$$
i_X(B) \to Rj_{V, *}(i_X(B)|_V) = Rj_{V, *}(i_V(B|_V))
$$
est un quasi-isomorphisme (lemme
\ref{lemma-pushforward-with-support-in-open}). Par hypoth\`ese, on a dans
$D(\mathcal{O}_X)$ un isomorphisme
$i_X(\Phi(B|_V)) \to Rj_{V, *}(i_V(B|_V))$. En outre, $\Phi$ et ${-}|_V$
sont des foncteurs adjoints entre les cat\'egories d\'eriv\'ees de modules
quasi-coh\'erents (Cat\'egories d\'eriv\'ees, lemme
\ref{derived-lemma-derived-adjoint-functors}). Le morphisme d'adjonction
$B \to \Phi(B|_V)$ devient un isomorphisme apr\`es application de $i_X$;
c'est donc un isomorphisme dans $D(\QCoh(\mathcal{O}_X))$. Par cons\'equent,
\begin{align*}
\Mor_{D(\QCoh(\mathcal{O}_X))}(A, B)
& =
\Mor_{D(\QCoh(\mathcal{O}_X))}(A, \Phi(B|_V)) \\
& =
\Mor_{D(\QCoh(\mathcal{O}_V))}(A|_V, B|_V) \\
& =
\Mor_{D(\mathcal{O}_V)}(i_V(A|_V), i_V(B|_V)) \\
& =
\Mor_{D(\mathcal{O}_X)}(i_X(A), Rj_{V, *}(i_V(B|_V))) \\
& =
\Mor_{D(\mathcal{O}_X)}(i_X(A), i_X(B))
\end{align*}
comme voulu. Nous avons utilis\'e ici le fait que $i_V$ est pleinement fid\`ele
(lemme \ref{lemma-affine-coherator}).

\medskip\noindent
Dans le cas g\'en\'eral, choisissons un complexe quelconque
$\mathcal{B}^\bullet$ de $\mathcal{O}_X$-modules quasi-coh\'erents qui
repr\'esente $B$, puis un quasi-isomorphisme quelconque
$s : \mathcal{B}^\bullet|_U \to \mathcal{C}^\bullet$ de complexes de modules
quasi-coh\'erents sur $U$. Comme $j_U : U \to X$ est quasi-compact et
quasi-s\'epar\'e, le foncteur $j_{U, *}$ envoie les modules quasi-coh\'erents
sur des modules quasi-coh\'erents (Morphismes d'espaces, lemme
\ref{spaces-morphisms-lemma-pushforward}). On dispose donc d'un morphisme
canonique
$\mathcal{B}^\bullet \to j_{U, *}(\mathcal{B}^\bullet|_U) \to
j_{U, *}\mathcal{C}^\bullet$
de complexes de modules quasi-coh\'erents sur $X$. Posons
$B'' = j_{U, *}\mathcal{C}^\bullet$ dans $D(\QCoh(\mathcal{O}_X))$ et
choisissons un triangle distingu\'e
$$
B \to B'' \to B' \to B[1]
$$
dans $D(\QCoh(\mathcal{O}_X))$. Comme le premier morphisme du triangle se
restreint en un isomorphisme sur $U$, l'objet $B'$ est \`a support dans $T$.
Dans le diagramme
$$
\xymatrix{
\Hom_{D(\QCoh(\mathcal{O}_X))}(A, B'[-1]) \ar[r] \ar[d] &
\Hom_{D(\mathcal{O}_X)}(i_X(A), i_X(B')[-1]) \ar[d] \\
\Hom_{D(\QCoh(\mathcal{O}_X))}(A, B) \ar[r] \ar[d] &
\Hom_{D(\mathcal{O}_X)}(i_X(A), i_X(B)) \ar[d] \\
\Hom_{D(\QCoh(\mathcal{O}_X))}(A, B'') \ar[r] \ar[d] &
\Hom_{D(\mathcal{O}_X)}(i_X(A), i_X(B'')) \ar[d] \\
\Hom_{D(\QCoh(\mathcal{O}_X))}(A, B') \ar[r] &
\Hom_{D(\mathcal{O}_X)}(i_X(A), i_X(B'))
}
$$
les colonnes sont exactes et les fl\`eches horizontales sup\'erieure et
inf\'erieure sont bijectives. Choisissons enfin un complexe
$\mathcal{A}^\bullet$ de modules quasi-coh\'erents repr\'esentant $A$.

\medskip\noindent
Soit $\alpha : i_X(A) \to i_X(B)$ un morphisme dans
$D(\mathcal{O}_X)$. Puisque $i_U$ est pleinement fid\`ele, la restriction
$\alpha|_U$ provient d'un morphisme de $D(\QCoh(\mathcal{O}_U))$. On peut
donc choisir $s : \mathcal{B}^\bullet|_U \to \mathcal{C}^\bullet$ comme
ci-dessus de telle sorte que $\alpha|_U$ soit repr\'esent\'e par un v\'eritable
morphisme de complexes
$\mathcal{A}^\bullet|_U \to \mathcal{C}^\bullet$. Celui-ci correspond \`a un
morphisme de complexes
$\mathcal{A}^\bullet \to j_{U, *}\mathcal{C}^\bullet$. Autrement dit, l'image
de $\alpha$ dans $\Hom_{D(\mathcal{O}_X)}(i_X(A), i_X(B''))$ provient d'un
\'el\'ement de $\Hom_{D(\QCoh(\mathcal{O}_X))}(A, B'')$. Une chasse au
diagramme montre alors que $\alpha$ provient d'un morphisme $A \to B$ dans
$D(\QCoh(\mathcal{O}_X))$.

Supposons enfin que $a : A \to B$ soit un morphisme de
$D(\QCoh(\mathcal{O}_X))$ qui devienne nul dans $D(\mathcal{O}_X)$. En
choisissant convenablement $\mathcal{B}^\bullet$, nous pouvons supposer que
$a$ est repr\'esent\'e par un morphisme de complexes
$a^\bullet : \mathcal{A}^\bullet \to \mathcal{B}^\bullet$. Comme $i_U$ est
pleinement fid\`ele, la restriction $a^\bullet|_U$ est nulle dans
$D(\QCoh(\mathcal{O}_U))$. On peut donc choisir $s$ de sorte que
$s \circ a^\bullet|_U : \mathcal{A}^\bullet|_U \to \mathcal{C}^\bullet$
soit homotope \`a z\'ero. En appliquant le foncteur $j_{U, *}$, on en d\'eduit
que $\mathcal{A}^\bullet \to j_{U, *}\mathcal{C}^\bullet$ est homotope \`a
z\'ero. Ainsi $a$ s'envoie sur z\'ero dans
$\Hom_{D(\QCoh(\mathcal{O}_X))}(A, B'')$.
Nous pouvons donc supposer que $a$ est l'image d'un \`el\'ement
$b \in \Hom_{D(\QCoh(\mathcal{O}_X))}(A, B'[-1])$. L'image de $b$ dans
$\Hom_{D(\mathcal{O}_X)}(i_X(A), i_X(B')[-1])$ provient d'un
$\gamma \in \Hom_{D(\mathcal{O}_X)}(A, B''[-1])$, puisque $a$ s'envoie sur
z\'ero dans le groupe de droite. Les fl\`eches horizontales \'etant
surjectives, comme nous l'avons vu plus haut, $\gamma$ provient d'un
$c$ de
$\Hom_{D(\QCoh(\mathcal{O}_X))}(A, B''[-1])$
et il en r\'esulte que $a = 0$, comme voulu.

\medskip\noindent
Nous savons maintenant que $i_X$ est pleinement fid\`ele pour l'espace $X$
initial. Comme $RQ_X$ est son adjoint \`a droite, on a
$RQ_X \circ i_X = \text{id}$ (Cat\'egories, lemme
\ref{categories-lemma-adjoint-fully-faithful}).
Pour achever la preuve, montrons que, pour tout
$E$ de $D_\QCoh(\mathcal{O}_X)$, le morphisme
$i_X(RQ_X(E)) \to E$ est un isomorphisme. Choisissons un triangle distingu\'e
$$
i_X(RQ_X(E)) \to E \to E' \to i_X(RQ_X(E))[1]
$$
dans $D_\QCoh(\mathcal{O}_X)$. Un argument formel fond\'e sur ce qui pr\'ec\`ede
montre que $i_X(RQ_X(E')) = 0$. Il suffit donc de prouver que, pour
$E \in D_\QCoh(\mathcal{O}_X)$, la condition $i_X(RQ_X(E)) = 0$ entra\^ine
$E = 0$. Consid\'erons un morphisme \'etale $j : V \to X$, avec $V$ affine.
D'apr\`es les lemmes \ref{lemma-affine-coherator} et
\ref{lemma-flat-pushforward-coherator}, ainsi que notre hypoth\`ese, on a
$$
Rj_*(E|_V) = Rj_*(i_V(RQ_V(E|_V))) = i_X(\Phi(RQ_V(E|_V))) =
i_X(RQ_X(Rj_*(E|_V)))
$$
Choisissons un triangle distingu\'e
$$
E \to Rj_*(E|_V) \to E' \to E[1]
$$
En appliquant $RQ_X$, on obtient un triangle distingu\'e
$$
0 \to RQ_X(Rj_*(E|_V)) \to RQ_X(E') \to 0[1]
$$
c'est-\`a-dire que le morphisme du milieu est un isomorphisme. Avec la suite
d'\'egalit\'es ci-dessus, on voit que notre premier triangle distingu\'e devient
un triangle distingu\'e
$$
E \to i_X(RQ_X(E')) \to E' \to E[1]
$$
o\`u le morphisme du milieu est le morphisme d'adjonction. Or le compos\'e
$E \to E'$ est nul; par adjonction, $E \to i_X(RQ_X(E'))$ est donc nul.
Ce morphisme \'etant isomorphe au morphisme
$E \to Rj_*(E|_V)$ adjoint \`a $\text{id} : E|_V \to E|_V$, on en d\'eduit
que $E|_V$ est nul. Comme cela vaut pour tout $V$ affine et \'etale sur $X$,
on conclut que $E$ est nul, comme voulu.
\end{proof}

\begin{proposition}
\label{proposition-quasi-compact-affine-diagonal}
Soit $S$ un sch\'ema et soit $X$ un espace alg\'ebrique quasi-compact sur $S$
dont la diagonale sur $\mathbf{Z}$ est affine (au sens de Propri\'et\'es des
espaces, d\'efinition \ref{spaces-properties-definition-separated}). Alors le
foncteur
(\ref{equation-compare})
$$
D(\QCoh(\mathcal{O}_X))
\longrightarrow
D_\QCoh(\mathcal{O}_X)
$$
est une \'equivalence, dont $RQ_X$ est un quasi-inverse.
\end{proposition}

\begin{proof}
Soit $V \to W$ un morphisme \'etale, avec $V$ affine et $W$ un sous-espace
ouvert quasi-compact de $X$. Le morphisme $V \to W$ est affine, car la
diagonale de $W$ sur $\mathbf{Z}$ est affine et $V$ est affine (Morphismes
d'espaces, lemme \ref{spaces-morphisms-lemma-affine-permanence}). Le
lemme \ref{lemma-affine-pushforward} assure alors que l'hypoth\`ese du
lemme \ref{lemma-argument-proves} est satisfaite, ce qui permet de conclure.
\end{proof}

\begin{lemma}
\label{lemma-direct-image-coherator}
Soit $S$ un sch\'ema et soit $f : X \to Y$ un morphisme d'espaces alg\'ebriques
sur $S$. Supposons $X$ et $Y$ quasi-compacts et \`a diagonale affine sur
$\mathbf{Z}$ (au sens de Propri\'et\'es des espaces, d\'efinition
\ref{spaces-properties-definition-separated}). Si l'on note
$$
\Phi : D(\QCoh(\mathcal{O}_X)) \to D(\QCoh(\mathcal{O}_Y))
$$
le foncteur d\'eriv\'e droit de
$f_* : \QCoh(\mathcal{O}_X) \to \QCoh(\mathcal{O}_Y)$
alors le diagramme
$$
\xymatrix{
D(\QCoh(\mathcal{O}_X)) \ar[d]_\Phi \ar[r] &
D_\QCoh(\mathcal{O}_X) \ar[d]^{Rf_*} \\
D(\QCoh(\mathcal{O}_Y)) \ar[r] &
D_\QCoh(\mathcal{O}_Y)
}
$$
est commutatif.
\end{lemma}

\begin{proof}
Les fl\`eches horizontales du diagramme sont des \'equivalences de cat\'egories
d'apr\`es la proposition
\ref{proposition-quasi-compact-affine-diagonal}. Nous pouvons donc identifier
ces cat\'egories, et faire de m\^eme pour tout autre espace alg\'ebrique
quasi-compact \`a diagonale affine. L'\'enonc\'e du lemme affirme alors que le
morphisme canonique $\Phi(K) \to Rf_*(K)$ est un isomorphisme pour tout
$K$ de $D(\QCoh(\mathcal{O}_X))$. Remarquons que, si
$K_1 \to K_2 \to K_3 \to K_1[1]$ est un triangle distingu\'e dans
$D(\QCoh(\mathcal{O}_X))$ et que l'assertion est vraie pour deux des trois
objets, elle est vraie pour le troisi\`eme.

\medskip\noindent
Soit $\mathcal{B} \subset \Ob(X_{spaces, \etale})$ l'ensemble des objets
quasi-compacts \`a diagonale affine. Pour $U \in \mathcal{B}$ et tout
morphisme $g : U \to Z$, o\`u $Z$ est un espace alg\'ebrique quasi-compact
sur $S$ \`a diagonale affine, notons
$$
\Phi_g : D(\QCoh(\mathcal{O}_U)) \to D(\QCoh(\mathcal{O}_Z))
$$
le foncteur d\'eriv\'e droit de $g_*$. D\'efinissons $P(U)$ comme l'assertion :
``pour tout $K$ de $D(\QCoh(\mathcal{O}_U))$ et tout $g : U \to Z$ comme
ci-dessus, le morphisme $\Phi_g(K) \to Rg_*K$ est un isomorphisme''.
D'apr\`es la remarque \ref{remark-how-to}, les conditions (1), (2) et (3)(a)
du lemme \ref{lemma-induction-principle-separated} sont satisfaites; il reste
\`a prouver (3)(b) et (4).

\medskip\noindent
V\'erifions la condition (3)(b). Soit $U$ un sch\'ema affine \'etale sur $X$
et soit $g : U \to Z$ comme ci-dessus. Puisque la diagonale de $Z$ est
affine, le morphisme $g : U \to Z$ est affine (Morphismes d'espaces, lemme
\ref{spaces-morphisms-lemma-affine-permanence}).
Le lemme \ref{lemma-affine-pushforward} montre donc que $U$ v\'erifie $P$.

\medskip\noindent
V\'erifions la condition (4). Soit $(U \subset W, V \to W)$ un carr\'e
distingu\'e \`el\'ementaire de $X_{spaces, \etale}$, avec
$U,W,V \in \mathcal{B}$ et $V$ affine. Supposons que $U$, $V$ et
$U \times_W V$ v\'erifient $P$. Il faut montrer que $W$ v\'erifie $P$.
Soit $g : W \to Z$ un morphisme vers un espace alg\'ebrique quasi-compact
\`a diagonale affine, et soit $K$ un objet de
$D(\QCoh(\mathcal{O}_W))$. Consid\'erons le triangle distingu\'e
$$
K \to Rj_{U, *}K|_U \oplus Rj_{V, *}K|_V \to
Rj_{U \times_W V, *}K|_{U \times_W V} \to K[1]
$$
dans $D(\mathcal{O}_W)$. Par la propri\'et\'e des deux sur trois rappel\'ee
ci-dessus, il suffit de montrer que
$\Phi_g(Rj_{U, *}K|_U) \to Rg_*(Rj_{U, *}K|_U)$ est un isomorphisme, et de
m\^eme pour $V$ et $U \times_W V$. C'est l'objet du paragraphe suivant.

\medskip\noindent
Soit $j : U \to W$ un morphisme de $X_{spaces, \etale}$, avec
$U,W \in \mathcal{B}$ et $U$ v\'erifiant $P$. Soit $g : W \to Z$ un
morphisme vers un espace alg\'ebrique quasi-compact \`a diagonale affine.
Pour achever la preuve, il faut montrer que
$\Phi_g(Rj_*K) \to Rg_*(Rj_*K)$
soit un isomorphisme pour tout $K$ de $D(\QCoh(\mathcal{O}_U))$.
Soit $\mathcal{I}^\bullet$ un complexe K-injectif de
$\QCoh(\mathcal{O}_U)$ repr\'esentant $K$. En appliquant $P(U)$ \`a $j$, on
voit que $j_*\mathcal{I}^\bullet$ repr\'esente $Rj_*K$. Comme
$j_* : \QCoh(\mathcal{O}_U) \to \QCoh(\mathcal{O}_W)$ admet l'adjoint \`a
gauche exact $j^* : \QCoh(\mathcal{O}_W) \to \QCoh(\mathcal{O}_U)$, le
complexe $j_*\mathcal{I}^\bullet$ est K-injectif dans
$\QCoh(\mathcal{O}_W)$; voir Cat\'egories d\'eriv\'ees, lemme
\ref{derived-lemma-adjoint-preserve-K-injectives}. Ainsi
$\Phi_g(Rj_*K)$ est repr\'esent\'e par
$g_*j_*\mathcal{I}^\bullet = (g \circ j)_*\mathcal{I}^\bullet$.
En appliquant $P(U)$ \`a $g \circ j$, on voit que ce complexe repr\'esente
$R_{g \circ j, *}(K) = Rg_*(Rj_*K)$, ce qui ach\`eve la preuve.
\end{proof}









\section{Le coh\'erateur pour les espaces noeth\'eriens}
\label{section-coherator-Noetherian}

\noindent
Il nous faut encore quelques r\'esultats sur les modules injectifs pour traiter
le cas d'un espace alg\'ebrique noeth\'erien.

\begin{lemma}
\label{lemma-affine-injective-colimit-direct-sum-pushforwards-artin}
Soit $S$ un sch\'ema affine noeth\'erien. Tout objet injectif de
$\QCoh(\mathcal{O}_S)$ est une limite inductive filtrante
$\colim_i \mathcal{F}_i$ de faisceaux quasi-coh\'erents de la forme
$$
\mathcal{F}_i = (Z_i \to S)_*\mathcal{G}_i
$$
o\`u $Z_i$ est le spectre d'un anneau artinien et $\mathcal{G}_i$ un module
coh\'erent sur $Z_i$.
\end{lemma}

\begin{proof}
Posons $S = \Spec(A)$. Soit $\mathcal{J}$ un objet injectif de
$\QCoh(\mathcal{O}_S)$. Comme $\QCoh(\mathcal{O}_S)$ est
\'equivalente \`a la cat\'egorie des $A$-modules, on a
$\mathcal{J}=\widetilde{J}$ pour un $A$-module injectif $J$.
D'apr\`es Complexes dualisants, proposition
\ref{dualizing-proposition-structure-injectives-noetherian}
on peut \'ecrire $J = \bigoplus E_\alpha$, avec $E_\alpha$ ind\'ecomposable,
donc isomorphe \`a l'enveloppe injective du corps r\'esiduel en un point.
Ainsi (puisqu'une r\'eunion disjointe finie de sch\'emas artiniens est
artinienne), on peut supposer que $J$ est l'enveloppe injective de
$\kappa(\mathfrak p)$ pour un id\'eal premier $\mathfrak p$ de $A$.
Alors $J = \bigcup J[\mathfrak p^n]$, o\`u $J[\mathfrak p^n]$ est
l'enveloppe injective de $\kappa(\mathfrak p)$ sur
$A_\mathfrak p/\mathfrak p^nA_\mathfrak p$; voir
Complexes dualisants, lemme \ref{dualizing-lemma-union-artinian}.
Ainsi $\widetilde{J}$ est la limite inductive des faisceaux
$(Z_n \to S)_*\mathcal{G}_n$, o\`u
$Z_n = \Spec(A_\mathfrak p/\mathfrak p^nA_\mathfrak p)$ et
$\mathcal{G}_n$ est le faisceau coh\'erent associ\'e au
$A_\mathfrak p/\mathfrak p^nA_\mathfrak p$-module fini
$J[\mathfrak p^n]$. La finitude r\'esulte de
Complexes dualisants, lemme \ref{dualizing-lemma-finite}.
\end{proof}

\begin{lemma}
\label{lemma-injective-colimit-direct-sum-pushforwards-artin}
Soit $S$ un sch\'ema affine et soit $X$ un espace alg\'ebrique noeth\'erien
sur $S$. Tout objet injectif de $\QCoh(\mathcal{O}_X)$ est facteur direct
d'une limite inductive filtrante $\colim_i \mathcal{F}_i$ de faisceaux
quasi-coh\'erents de la forme
$$
\mathcal{F}_i = (Z_i \to X)_*\mathcal{G}_i
$$
o\`u $Z_i$ est le spectre d'un anneau artinien et $\mathcal{G}_i$ un module
coh\'erent sur $Z_i$.
\end{lemma}

\begin{proof}
Choisissons un sch\'ema affine $U$ et un morphisme \'etale surjectif
$j : U \to X$ (Propri\'et\'es des espaces, lemme
\ref{spaces-properties-lemma-quasi-compact-affine-cover}).
Alors $U$ est un sch\'ema affine noeth\'erien. Choisissons un objet injectif
$\mathcal{J}'$ de $\QCoh(\mathcal{O}_U)$ et une injection
$\mathcal{J}|_U \to \mathcal{J}'$. Le morphisme
$$
\mathcal{J} \to j_*\mathcal{J}'
$$
est alors injectif dans $\QCoh(\mathcal{O}_X)$; il identifie donc
$\mathcal{J}$ \`a un facteur direct de $j_*\mathcal{J}'$.
Le r\'esultat d\'ecoule donc du r\'esultat correspondant pour
$\mathcal{J}'$, d\'emontr\'e dans le lemme
\ref{lemma-affine-injective-colimit-direct-sum-pushforwards-artin}.
\end{proof}

\begin{lemma}
\label{lemma-flat-pullback-injective-quasi-coherent}
Soit $S$ un sch\'ema. Soit $f : X \to Y$ un morphisme plat,
quasi-compact et quasi-s\'epar\'e d'espaces alg\'ebriques sur $S$. Si
$\mathcal{J}$ est un objet injectif de $\QCoh(\mathcal{O}_X)$,
alors $f_*\mathcal{J}$ est un objet injectif de
$\QCoh(\mathcal{O}_Y)$.
\end{lemma}

\begin{proof}
Puisque $f$ est quasi-compact et quasi-s\'epar\'e, le foncteur
$f_*$ transforme les faisceaux quasi-coh\'erents en faisceaux
quasi-coh\'erents (Morphismes d'espaces, lemme
\ref{spaces-morphisms-lemma-pushforward}). Le foncteur $f^*$ est un adjoint
\`a gauche de $f_*$ qui transforme les injections en injections.
Le r\'esultat d\'ecoule donc d'Homologie, lemme
\ref{homology-lemma-adjoint-preserve-injectives}.
\end{proof}

\begin{lemma}
\label{lemma-injective-pushforward}
Soit $S$ un sch\'ema et soit $X$ un espace alg\'ebrique noeth\'erien sur $S$.
Si $\mathcal{J}$ est un objet injectif de $\QCoh(\mathcal{O}_X)$, alors
\begin{enumerate}
\item $H^p(U, \mathcal{J}|_U) = 0$ pour $p > 0$ et pour tout espace
alg\'ebrique $U$ quasi-compact et quasi-s\'epar\'e, \'etale sur $X$;
\item pour tout morphisme $f : X \to Y$ d'espaces alg\'ebriques sur $S$,
avec $Y$ quasi-s\'epar\'e, on a $R^pf_*\mathcal{J} = 0$ pour $p > 0$.
\end{enumerate}
\end{lemma}

\begin{proof}
D\'emonstration de (1). \'Ecrivons $\mathcal{J}$ comme facteur direct de
$\colim \mathcal{F}_i$, avec
$\mathcal{F}_i = (Z_i \to X)_*\mathcal{G}_i$, comme dans le lemme
\ref{lemma-injective-colimit-direct-sum-pushforwards-artin}.
Il suffit manifestement de d\'emontrer l'annulation pour
$\colim \mathcal{F}_i$. Comme l'image inverse commute aux limites inductives
et que $U$ est quasi-compact et quasi-s\'epar\'e, il suffit de d\'emontrer
$H^p(U, \mathcal{F}_i|_U) = 0$ pour $p > 0$; voir Cohomologie des espaces,
lemme \ref{spaces-cohomology-lemma-colimits}. Observons que
$Z_i \to X$ est un morphisme affine; voir Morphismes d'espaces, lemme
\ref{spaces-morphisms-lemma-Artinian-affine}. Par suite,
$$
\mathcal{F}_i|_U = (Z_i \times_X U \to U)_*\mathcal{G}'_i =
R(Z_i \times_X U \to U)_*\mathcal{G}'_i
$$
o\`u $\mathcal{G}'_i$ est l'image inverse de $\mathcal{G}_i$ sur
$Z_i \times_X U$; voir Cohomologie des espaces, lemme
\ref{spaces-cohomology-lemma-affine-base-change}.
Comme $Z_i \times_X U$ est affine, on en conclut que $\mathcal{G}'_i$
n'a pas de cohomologie sup\'erieure sur $Z_i \times_X U$.
La suite spectrale de Leray donne le m\^eme r\'esultat pour
$\mathcal{F}_i|_U$ (Cohomologie sur les sites, lemme
\ref{sites-cohomology-lemma-apply-Leray}).

\medskip\noindent
D\'emonstration de (2). Soit $f : X \to Y$ un morphisme d'espaces
alg\'ebriques sur $S$, et soit $V \to Y$ un morphisme \'etale avec $V$
affine. Alors $V \times_Y X \to X$ est un morphisme \'etale et
$V \times_Y X$ est un espace alg\'ebrique quasi-compact et quasi-s\'epar\'e,
\'etale sur $X$ (d\'etails omis). Par (1),
$H^p(V \times_Y X, \mathcal{J})$ est donc nul. Comme
$R^pf_*\mathcal{J}$ est le faisceau associ\'e au pr\'efaisceau
$V \mapsto H^p(V \times_Y X, \mathcal{J})$, la d\'emonstration est achev\'ee.
\end{proof}

\begin{lemma}
\label{lemma-Noetherian-pushforward}
Soit $S$ un sch\'ema et soit $f : X \to Y$ un morphisme d'espaces
alg\'ebriques noeth\'eriens sur $S$. Le foncteur $f_*$ sur les faisceaux
quasi-coh\'erents admet alors une extension d\'eriv\'ee \`a droite
$\Phi : D(\QCoh(\mathcal{O}_X)) \to D(\QCoh(\mathcal{O}_Y))$
telle que le diagramme
$$
\xymatrix{
D(\QCoh(\mathcal{O}_X)) \ar[d]_{\Phi} \ar[r] &
D_\QCoh(\mathcal{O}_X) \ar[d]^{Rf_*} \\
D(\QCoh(\mathcal{O}_Y)) \ar[r] &
D_\QCoh(\mathcal{O}_Y)
}
$$
soit commutatif.
\end{lemma}

\begin{proof}
Comme $X$ et $Y$ sont noeth\'eriens, le morphisme est quasi-compact et
quasi-s\'epar\'e (voir Morphismes d'espaces, lemme
\ref{spaces-morphisms-lemma-quasi-compact-quasi-separated-permanence}).
Ainsi $f_*$ pr\'eserve la quasi-coh\'erence; voir Morphismes d'espaces,
lemme \ref{spaces-morphisms-lemma-pushforward}.
Soit maintenant $K$ un objet de $D(\QCoh(\mathcal{O}_X))$.
Comme $\QCoh(\mathcal{O}_X)$ est une cat\'egorie ab\'elienne de Grothendieck
(Propri\'et\'es des espaces, proposition
\ref{spaces-properties-proposition-coherator}), on peut repr\'esenter $K$
par un complexe K-injectif $\mathcal{I}^\bullet$ tel que chaque
$\mathcal{I}^n$ soit un objet injectif de $\QCoh(\mathcal{O}_X)$; voir
Injectifs, th\'eor\`eme
\ref{injectives-theorem-K-injective-embedding-grothendieck}.
On d\'efinit alors le foncteur $\Phi$ en posant
$$
\Phi(K) = f_*\mathcal{I}^\bullet
$$
o\`u le membre de droite est consid\'er\'e comme un objet de
$D(\QCoh(\mathcal{O}_Y))$. Pour achever la preuve du lemme, il suffit de
montrer que le morphisme canonique
$$
f_*\mathcal{I}^\bullet \longrightarrow Rf_*\mathcal{I}^\bullet
$$
est un isomorphisme dans $D(\mathcal{O}_Y)$. Pour cela, il suffit de montrer
qu'il induit un isomorphisme sur les faisceaux de cohomologie. Fixons
$m \in \mathbf{Z}$ et soit $N = N(X, Y, f)$ comme dans le lemme
\ref{lemma-quasi-coherence-direct-image}. Consid\'erons la suite exacte courte
$$
0 \to \sigma_{\geq m - N - 1}\mathcal{I}^\bullet \to
\mathcal{I}^\bullet \to \sigma_{\leq m - N - 2}\mathcal{I}^\bullet \to 0
$$
de complexes de faisceaux quasi-coh\'erents sur $X$. D'apr\`es le lemme
\ref{lemma-quasi-coherence-direct-image}, les faisceaux de cohomologie de
$Rf_*\sigma_{\leq m - N - 2}\mathcal{I}^\bullet$
sont nuls en degr\'es $\geq m - 1$. On en d\'eduit que
$R^mf_*\mathcal{I}^\bullet$ est isomorphe \`a
$R^mf_*\sigma_{\geq m - N - 1}\mathcal{I}^\bullet$.
Autrement dit, on peut supposer que $\mathcal{I}^\bullet$ est un complexe
born\'e inf\'erieurement d'objets injectifs de $\QCoh(\mathcal{O}_X)$.
Ce cas r\'esulte du lemme d'acyclicit\'e de Leray (Cat\'egories d\'eriv\'ees,
lemme \ref{derived-lemma-leray-acyclicity}), l'annulation requise \'etant
fournie par le lemme \ref{lemma-injective-pushforward}.
\end{proof}

\begin{proposition}
\label{proposition-Noetherian}
Soit $S$ un sch\'ema et soit $X$ un espace alg\'ebrique noeth\'erien sur $S$.
Alors le foncteur (\ref{equation-compare})
$$
D(\QCoh(\mathcal{O}_X))
\longrightarrow
D_\QCoh(\mathcal{O}_X)
$$
est une \'equivalence dont un quasi-inverse est donn\'e par $RQ_X$.
\end{proposition}

\begin{proof}
Cela r\'esulte imm\'ediatement des lemmes
\ref{lemma-Noetherian-pushforward} et \ref{lemma-argument-proves}.
\end{proof}













\section{Complexes pseudo-coh\'erents et parfaits}
\label{section-spell-out}

\noindent
Dans cette section, nous \'etudions les notions g\'en\'erales d\'efinies dans
Cohomologie sur les sites, sections
\ref{sites-cohomology-section-strictly-perfect},
\ref{sites-cohomology-section-pseudo-coherent},
\ref{sites-cohomology-section-tor}, et
\ref{sites-cohomology-section-perfect}
pour le site \'etale d'un espace alg\'ebrique. Nous les comparons en particulier
aux notions correspondantes pour les sch\'emas.

\medskip\noindent
Comparons tout d'abord la notion de complexe pseudo-coh\'erent sur un sch\'ema
et sur le petit site \'etale qui lui est associ\'e.

\begin{lemma}
\label{lemma-descend-finite-type}
Soit $X$ un sch\'ema et soit $\mathcal{F}$ un $\mathcal{O}_X$-module.
Les conditions suivantes sont \'equivalentes :
\begin{enumerate}
\item $\mathcal{F}$ est de type fini comme $\mathcal{O}_X$-module;
\item $\epsilon^*\mathcal{F}$ est de type fini comme
$\mathcal{O}_\etale$-module sur le petit site \'etale de $X$.
\end{enumerate}
Ici $\epsilon$ est comme dans (\ref{equation-epsilon}).
\end{lemma}

\begin{proof}
L'implication (1) $\Rightarrow$ (2) est un fait g\'en\'eral; voir Modules sur
les sites, lemme \ref{sites-modules-lemma-local-pullback}.
Supposons (2). Il existe par hypoth\`ese une famille couvrante
$\{f_i : X_i \to X\}$ pour la topologie \'etale telle que
$\epsilon^*\mathcal{F}|_{(X_i)_\etale}$ soit engendr\'e par un nombre fini de
sections. Soit $x \in X$. Nous allons montrer que $\mathcal{F}$ est engendr\'e
par un nombre fini de sections dans un voisinage de $x$.
Choisissons $i$ tel que $x$ appartienne \`a l'image de $X_i \to X$ et posons
$X' = X_i$. Soient
$s_1, \ldots, s_n \in
\Gamma(X', \epsilon^*\mathcal{F}|_{X'_\etale})$
des sections g\'en\'eratrices. Comme
$\epsilon^*\mathcal{F} =
\epsilon^{-1}\mathcal{F} \otimes_{\epsilon^{-1}\mathcal{O}_X}
\mathcal{O}_\etale$
on peut trouver un morphisme \'etale $X'' \to X'$ tel que $x$ appartienne
\`a l'image de $X'' \to X$ et que
$s_i|_{X''} = \sum s_{ij} \otimes a_{ij}$, pour des sections
$s_{ij} \in \epsilon^{-1}\mathcal{F}(X'')$ et
$a_{ij} \in \mathcal{O}_\etale(X'')$. Notons $U \subset X$ l'image de
$X'' \to X$. C'est un sous-sch\'ema ouvert, car $f'' : X'' \to X$ est \'etale
(Morphismes, lemme \ref{morphisms-lemma-etale-open}). Quitte \`a r\'etr\'ecir
encore $X''$, on peut supposer que les $s_{ij}$ proviennent d'\'el\'ements
$t_{ij} \in \mathcal{F}(U)$, d'apr\`es la construction du foncteur image
inverse $\epsilon^{-1}$. Nous affirmons que les $t_{ij}$ engendrent
$\mathcal{F}|_U$, ce qui ach\`eve la preuve du lemme. En effet, l'image inverse
sur $X''$ du morphisme correspondant
$\mathcal{O}_U^{\oplus N} \to \mathcal{F}|_U$ est surjective. Comme
$f'' : X'' \to U$ est un morphisme plat surjectif de sch\'emas, on en d\'eduit
que $\mathcal{O}_U^{\oplus N} \to \mathcal{F}|_U$ est surjectif en passant
aux fibres et en utilisant le fait que
$\mathcal{O}_{U, f''(z)} \to \mathcal{O}_{X'', z}$
est fid\`element plat pour tout $z \in X''$.
\end{proof}

\noindent
Dans la situation ci-dessus, le morphisme de sites $\epsilon$ est plat et
d\'efinit donc une image inverse sur les complexes de modules.

\begin{lemma}
\label{lemma-descend-pseudo-coherent}
Soit $X$ un sch\'ema et soit $E$ un objet de $D(\mathcal{O}_X)$.
Les conditions suivantes sont \'equivalentes :
\begin{enumerate}
\item $E$ est $m$-pseudo-coh\'erent;
\item $\epsilon^*E$ est $m$-pseudo-coh\'erent sur le petit site \'etale de $X$.
\end{enumerate}
Ici $\epsilon$ est comme dans (\ref{equation-epsilon}).
\end{lemma}

\begin{proof}
L'implication (1) $\Rightarrow$ (2) est un fait g\'en\'eral; voir Cohomologie
sur les sites, lemme
\ref{sites-cohomology-lemma-pseudo-coherent-pullback}.
Supposons $\epsilon^*E$ $m$-pseudo-coh\'erent. Nous utiliserons sans plus le
mentionner le fait que $\epsilon^*$ est un foncteur exact et que, par suite,
$$
\epsilon^*H^i(E) = H^i(\epsilon^*E).
$$
Pour montrer que $E$ est $m$-pseudo-coh\'erent, on peut travailler localement
sur $X$ et donc supposer $X$ quasi-compact (par exemple affine).
Comme $X$ est quasi-compact, toute famille couvrante
$\{U_i \to X\}$ pour la topologie \'etale admet un raffinement fini. Ainsi
$\epsilon^*E$ est un objet de
$D^{-}(\mathcal{O}_\etale)$; voir les commentaires qui suivent Cohomologie
sur les sites, d\'efinition
\ref{sites-cohomology-definition-pseudo-coherent}.
D'apr\`es le lemme \ref{lemma-epsilon-flat}, il s'ensuit que $E$ est un objet
de $D^-(\mathcal{O}_X)$.

\medskip\noindent
Soit $n \in \mathbf{Z}$ le plus grand entier tel que $H^n(E)$ soit non nul;
c'est aussi le plus grand entier tel que $H^n(\epsilon^*E)$ soit non nul.
Nous d\'emontrerons le lemme par r\'ecurrence sur $n-m$.
Si $n < m$, le lemme est clair. Si $n \geq m$, alors
$H^n(\epsilon^*E)$ est un $\mathcal{O}_\etale$-module de type fini; voir
Cohomologie sur les sites, lemme
\ref{sites-cohomology-lemma-finite-cohomology}.
Le $\mathcal{O}_X$-module $H^n(E)$ est donc de type fini d'apr\`es le lemme
\ref{lemma-descend-finite-type}. En rempla\c{c}ant $X$ par les membres d'un
recouvrement ouvert, on peut supposer qu'il existe une surjection
$\mathcal{O}_X^{\oplus t} \to H^n(E)$. Localement sur $X$, on peut la relever
en un morphisme de complexes
$\alpha : \mathcal{O}_X^{\oplus t}[-n] \to E$ (d\'etails omis).
Choisissons un triangle distingu\'e
$$
\mathcal{O}_X^{\oplus t}[-n] \to E \to C \to \mathcal{O}_X^{\oplus t}[-n + 1]
$$
La cohomologie de $C$ est alors nulle en degr\'es $\geq n$. D'autre part, le
complexe $\epsilon^*C$ est $m$-pseudo-coh\'erent; voir Cohomologie sur les
sites, lemme \ref{sites-cohomology-lemma-cone-pseudo-coherent}.
Par l'hypoth\`ese de r\'ecurrence, $C$ est donc $m$-pseudo-coh\'erent. Une
nouvelle application de Cohomologie sur les sites, lemme
\ref{sites-cohomology-lemma-cone-pseudo-coherent}, permet de conclure.
\end{proof}

\begin{lemma}
\label{lemma-descend-tor-amplitude}
Soit $X$ un sch\'ema et soit $E$ un objet de $D(\mathcal{O}_X)$. Alors :
\begin{enumerate}
\item $E$ est de Tor-amplitude contenue dans $[a,b]$ si et seulement si
$\epsilon^*E$ est de Tor-amplitude contenue dans $[a,b]$;
\item $E$ est de Tor-dimension finie si et seulement si $\epsilon^*E$ est de
Tor-dimension finie.
\end{enumerate}
Ici $\epsilon$ est comme dans (\ref{equation-epsilon}).
\end{lemma}

\begin{proof}
L'implication directe r\'esulte de Cohomologie sur les sites, lemme
\ref{sites-cohomology-lemma-tor-amplitude-pullback}.
R\'eciproquement, supposons $\epsilon^*E$ de Tor-amplitude contenue dans
$[a,b]$.
Soit $\mathcal{F}$ un $\mathcal{O}_X$-module. Comme $\epsilon$ est un
morphisme plat de sites annel\'es (lemme \ref{lemma-epsilon-flat}), on a
$$
\epsilon^*(E \otimes^\mathbf{L}_{\mathcal{O}_X} \mathcal{F})
=
\epsilon^*E
\otimes^\mathbf{L}_{\mathcal{O}_\etale}
\epsilon^*\mathcal{F}
$$
L'annulation suppos\'ee des faisceaux de cohomologie du membre de droite
entra\^ine donc, par le lemme \ref{lemma-epsilon-flat}, l'annulation voulue
des faisceaux de cohomologie de
$E \otimes^\mathbf{L}_{\mathcal{O}_X} \mathcal{F}$.
\end{proof}

\begin{lemma}
\label{lemma-tor-dimension-rel}
Soit $f : X \to Y$ un morphisme de sch\'emas et soit $E$ un objet de
$D(\mathcal{O}_X)$. Alors :
\begin{enumerate}
\item $E$, consid\'er\'e comme objet de $D(f^{-1}\mathcal{O}_Y)$, est
de Tor-amplitude contenue dans $[a,b]$ si et seulement si $\epsilon^*E$,
consid\'er\'e
comme objet de $D(f_{small}^{-1}\mathcal{O}_{Y_\etale})$, est d'amplitude de
Tor dans $[a,b]$;
\item $E$, consid\'er\'e comme objet de $D(f^{-1}\mathcal{O}_Y)$, est
localement de Tor-dimension finie si et seulement si $\epsilon^*E$,
consid\'er\'e comme objet de
$D(f_{small}^{-1}\mathcal{O}_{Y_\etale})$, est localement de Tor-dimension
finie.
\end{enumerate}
Ici $\epsilon$ est comme dans (\ref{equation-epsilon}).
\end{lemma}

\begin{proof}
L'implication directe de (1) r\'esulte de Cohomologie sur les sites, lemme
\ref{sites-cohomology-lemma-tor-amplitude-pullback}.
Soit $x \in X$ un point et soit $\overline{x}$ un point g\'eom\'etrique
au-dessus de $x$. Posons $y=f(x)$ et notons $\overline{y}$ le point
g\'eom\'etrique de $Y$ provenant de $\overline{x}$. Alors
$(f^{-1}\mathcal{O}_Y)_x = \mathcal{O}_{Y,y}$ (Faisceaux, lemme
\ref{sheaves-lemma-stalk-pullback}) et
$$
(f_{small}^{-1}\mathcal{O}_{Y_\etale})_{\overline{x}} =
\mathcal{O}_{Y_\etale, \overline{y}} =
\mathcal{O}_{Y, y}^{sh}
$$
est le hens\'elis\'e strict (d'apr\`es Cohomologie \'etale, lemmes
\ref{etale-cohomology-lemma-stalk-pullback} et
\ref{etale-cohomology-lemma-describe-etale-local-ring}).
Comme la fibre de $\mathcal{O}_{X_\etale}$ en $\overline{x}$ est
$\mathcal{O}_{X,x}^{sh}$, on obtient
$$
(\epsilon^*E)_{\overline{x}} =
E_x \otimes_{\mathcal{O}_{X, x}}^\mathbf{L} \mathcal{O}_{X, x}^{sh}
$$
par transitivit\'e des images inverses. Si $\epsilon^*E$ est de Tor-amplitude
contenue dans $[a,b]$ comme complexe de
$f_{small}^{-1}\mathcal{O}_{Y_\etale}$-modules, alors
$(\epsilon^*E)_{\overline{x}}$ est de Tor-amplitude contenue dans $[a,b]$
comme
complexe de $\mathcal{O}_{Y,y}^{sh}$-modules (car le passage aux fibres est
une image inverse, d'apr\`es le lemme cit\'e ci-dessus). Par Autour de la
platitude, lemme \ref{flat-lemma-tor-amplitude-up-down-henselization},
la Tor-amplitude de
$(\epsilon^*E)_{\overline{x}}$
comme complexe de $\mathcal{O}_{Y,y}$-modules est contenue dans $[a,b]$.
Comme $\mathcal{O}_{X,x} \to \mathcal{O}_{X,x}^{sh}$ est fid\`element plat
(Autour de l'alg\`ebre, lemme
\ref{more-algebra-lemma-dumb-properties-henselization}) et que
$(\epsilon^*E)_{\overline{x}} =
E_x \otimes_{\mathcal{O}_{X, x}}^\mathbf{L} \mathcal{O}_{X, x}^{sh}$
on peut appliquer Autour de l'alg\`ebre, lemme
\ref{more-algebra-lemma-no-change-tor-amplitude}, et conclure que la
Tor-amplitude
de Tor de $E_x$, comme complexe de $\mathcal{O}_{Y,y}$-modules, est contenue
dans $[a,b]$. D'apr\`es Cohomologie, lemme
\ref{cohomology-lemma-tor-amplitude-stalk}, $E$, consid\'er\'e comme objet de
$D(f^{-1}\mathcal{O}_Y)$, est donc de Tor-amplitude contenue dans $[a,b]$.
Ceci d\'emontre l'implication r\'eciproque de (1). L'assertion (2) d\'ecoule
formellement de (1).
\end{proof}

\begin{lemma}
\label{lemma-descend-perfect}
Soit $X$ un sch\'ema et soit $E$ un objet de $D(\mathcal{O}_X)$. Alors $E$ est
un objet parfait de $D(\mathcal{O}_X)$ si et seulement si $\epsilon^*E$ est
un objet parfait de $D(\mathcal{O}_\etale)$. Ici $\epsilon$ est comme dans
(\ref{equation-epsilon}).
\end{lemma}

\begin{proof}
L'implication directe r\'esulte du r\'esultat g\'en\'eral de Cohomologie sur les
sites, lemme \ref{sites-cohomology-lemma-perfect-pullback}.
Pour la r\'eciproque, on utilise l'\'equivalence de Cohomologie sur les sites,
lemme \ref{sites-cohomology-lemma-perfect}, et les r\'esultats correspondants
pour les complexes pseudo-coh\'erents et les complexes de Tor-dimension finie,
\`a savoir les lemmes \ref{lemma-descend-pseudo-coherent} et
\ref{lemma-descend-tor-amplitude}. Quelques d\'etails sont omis.
\end{proof}

\begin{lemma}
\label{lemma-pseudo-coherent}
Soit $S$ un sch\'ema et soit $X$ un espace alg\'ebrique sur $S$.
Si $E$ est un objet $m$-pseudo-coh\'erent de $D(\mathcal{O}_X)$, alors
$H^i(E)$ est un $\mathcal{O}_X$-module quasi-coh\'erent pour $i>m$.
Si $E$ est pseudo-coh\'erent, alors $E$ est un objet de
$D_\QCoh(\mathcal{O}_X)$.
\end{lemma}

\begin{proof}
Localement, $H^i(E)$ est isomorphe \`a $H^i(\mathcal{E}^\bullet)$, o\`u
$\mathcal{E}^\bullet$ est strictement parfait. Les faisceaux
$\mathcal{E}^i$ sont facteurs directs de modules libres de type fini, donc
quasi-coh\'erents. Le lemme en r\'esulte.
\end{proof}

\begin{lemma}
\label{lemma-identify-pseudo-coherent-noetherian}
Soit $S$ un sch\'ema, soit $X$ un espace alg\'ebrique noeth\'erien sur $S$ et
soit $E$ un objet de $D_\QCoh(\mathcal{O}_X)$. Pour
$m \in \mathbf{Z}$, les conditions suivantes sont \'equivalentes :
\begin{enumerate}
\item $H^i(E)$ est coh\'erent pour $i \geq m$ et nul pour $i \gg 0$;
\item $E$ est $m$-pseudo-coh\'erent.
\end{enumerate}
En particulier, $E$ est pseudo-coh\'erent si et seulement si $E$ est un objet
de $D^-_{\textit{Coh}}(\mathcal{O}_X)$.
\end{lemma}

\begin{proof}
Comme $X$ est quasi-compact, on peut trouver un sch\'ema affine $U$ et un
morphisme \'etale surjectif $U \to X$ (Propri\'et\'es des espaces, lemme
\ref{spaces-properties-lemma-quasi-compact-affine-cover}).
Observons que $U$ est noeth\'erien. L'objet $E$ est $m$-pseudo-coh\'erent si et
seulement si $E|_U$ l'est (cela r\'esulte de la d\'efinition ou de Cohomologie
sur les sites, lemme
\ref{sites-cohomology-lemma-pseudo-coherent-independent-representative}).
De m\^eme, $H^i(E)$ est coh\'erent si et seulement si
$H^i(E)|_U=H^i(E|_U)$ est coh\'erent (voir Cohomologie des espaces, lemme
\ref{spaces-cohomology-lemma-coherent-Noetherian}).
On peut donc supposer $X$ repr\'esentable.

\medskip\noindent
Si $X$ est repr\'esent\'e par un sch\'ema $X_0$, on peut (lemme
\ref{lemma-derived-quasi-coherent-small-etale-site}) \'ecrire
$E=\epsilon^*E_0$, o\`u $E_0$ est un objet de
$D_\QCoh(\mathcal{O}_{X_0})$ et o\`u
$\epsilon : X_\etale \to (X_0)_{Zar}$ est comme dans
(\ref{equation-epsilon}). Dans ce cas, le lemme
\ref{lemma-descend-pseudo-coherent} montre que $E$ est
$m$-pseudo-coh\'erent si et seulement si $E_0$ l'est. De m\^eme, d'apr\`es le
lemme \ref{lemma-descend-finite-type}, $H^i(E_0)$ est de type fini (c'est-\`a-dire
coh\'erent) si et seulement si $H^i(E)$ l'est. Enfin, le lemme
\ref{lemma-epsilon-flat} montre que $H^i(E_0)=0$ si et seulement si
$H^i(E)=0$. On est donc ramen\'e au cas des sch\'emas, qui est le lemme de
Cat\'egories d\'eriv\'ees des sch\'emas
\ref{perfect-lemma-identify-pseudo-coherent-noetherian}.
\end{proof}

\begin{lemma}
\label{lemma-tor-qc-qs}
Soit $S$ un sch\'ema, soit $X$ un espace alg\'ebrique quasi-s\'epar\'e sur $S$,
soit $E$ un objet de $D_\QCoh(\mathcal{O}_X)$ et soient $a \leq b$.
Les conditions suivantes sont \'equivalentes :
\begin{enumerate}
\item $E$ est de Tor-amplitude contenue dans $[a,b]$;
\item pour tout $\mathcal{F}$ de $\QCoh(\mathcal{O}_X)$, on a
$H^i(E \otimes_{\mathcal{O}_X}^\mathbf{L} \mathcal{F})=0$ pour
$i \notin [a,b]$.
\end{enumerate}
\end{lemma}

\begin{proof}
Il est clair que (1) implique (2). Supposons (2). Soit $j : U \to X$ un
morphisme \'etale avec $U$ affine. Comme $X$ est quasi-s\'epar\'e, $j$ est \`a
la fois quasi-compact et s\'epar\'e. Par cons\'equent, $j_*$ transforme les
modules quasi-coh\'erents en modules quasi-coh\'erents (Morphismes d'espaces,
lemme
\ref{spaces-morphisms-lemma-pushforward}).
Soit $\mathcal{G}$ un module quasi-coh\'erent quelconque sur $U$.
Par hypoth\`ese,
$H^i(E \otimes_{\mathcal{O}_X}^\mathbf{L} j_*\mathcal{G})$ est nul pour
$i \notin [a,b]$. En prenant l'image inverse par le morphisme plat $j$, on
voit que
$H^i(E|_U \otimes_{\mathcal{O}_U}^\mathbf{L} j^*j_*\mathcal{G})$ est nul
pour $i \notin [a,b]$. D'apr\`es Cohomologie des espaces, lemme
\ref{spaces-cohomology-lemma-etale-pull-push-split}, $\mathcal{G}$ est un
facteur direct de $j^*j_*\mathcal{G}$. La condition (2) entra\^ine donc
l'annulation de
$H^i(E|_U \otimes_{\mathcal{O}_U}^\mathbf{L} \mathcal{G})$
pour $i \notin [a,b]$ et pour tout $\mathcal{O}_U$-module quasi-coh\'erent
$\mathcal{G}$. Puisqu'il suffit de d\'emontrer que $E|_U$ est d'amplitude de
Tor dans $[a,b]$, on est ramen\'e au cas o\`u $X$ est repr\'esentable.

\medskip\noindent
Si $X$ est repr\'esent\'e par un sch\'ema $X_0$, on peut (lemme
\ref{lemma-derived-quasi-coherent-small-etale-site}) \'ecrire
$E=\epsilon^*E_0$, o\`u $E_0$ est un objet de
$D_\QCoh(\mathcal{O}_{X_0})$ et o\`u
$\epsilon : X_\etale \to (X_0)_{Zar}$ est comme dans
(\ref{equation-epsilon}). Pour tout module quasi-coh\'erent
$\mathcal{F}_0$ sur $X_0$, le module $\epsilon^*\mathcal{F}_0$ est
quasi-coh\'erent sur $X$ et
$$
H^i(E \otimes_{\mathcal{O}_X}^\mathbf{L} \epsilon^*\mathcal{F}_0)
=
\epsilon^*H^i(E_0 \otimes_{\mathcal{O}_{X_0}}^\mathbf{L} \mathcal{F}_0)
$$
car $\epsilon$ est plat (lemme \ref{lemma-epsilon-flat}). En outre,
l'annulation de ces faisceaux pour $i \notin [a,b]$ entra\^ine celle de
$H^i(E_0 \otimes_{\mathcal{O}_{X_0}}^\mathbf{L} \mathcal{F}_0)$
par le m\^eme lemme. Le probl\`eme est donc ramen\'e au cas des sch\'emas,
trait\'e dans Cat\'egories d\'eriv\'ees des sch\'emas, lemme
\ref{perfect-lemma-tor-qc-qs}.
\end{proof}

\begin{lemma}
\label{lemma-descend-RHom}
Soit $X$ un sch\'ema et soient $E,F$ des objets de $D(\mathcal{O}_X)$.
Supposons que l'une des conditions suivantes soit satisfaite :
\begin{enumerate}
\item $E$ est pseudo-coh\'erent et $F$ appartient \`a
$D^+(\mathcal{O}_X)$;
\item $E$ est parfait et $F$ est quelconque.
\end{enumerate}
Alors il existe un isomorphisme canonique
$$
\epsilon^*R\SheafHom(E, F) \longrightarrow R\SheafHom(\epsilon^*E, \epsilon^*F)
$$
Ici $\epsilon$ est comme dans (\ref{equation-epsilon}).
\end{lemma}

\begin{proof}
Rappelons que $\epsilon$ est plat (lemme \ref{lemma-epsilon-flat}), donc que
$\epsilon^*=L\epsilon^*$. Il existe un morphisme canonique du membre de gauche
vers celui de droite d'apr\`es Cohomologie sur les sites, remarque
\ref{sites-cohomology-remark-prepare-fancy-base-change}.
Pour voir que c'est un isomorphisme, on peut travailler localement, c'est-\`a-dire
supposer que $X$ est un sch\'ema affine.

\medskip\noindent
Dans le cas (1), on peut repr\'esenter $E$ par un complexe born\'e
sup\'erieurement $\mathcal{E}^\bullet$ de $\mathcal{O}_X$-modules libres de
type fini; voir Cat\'egories d\'eriv\'ees des sch\'emas, lemme
\ref{perfect-lemma-lift-pseudo-coherent}. On peut aussi repr\'esenter $F$ par
un complexe born\'e inf\'erieurement $\mathcal{F}^\bullet$ de
$\mathcal{O}_X$-modules. En appliquant Cohomologie, lemme
\ref{cohomology-lemma-Rhom-complex-of-direct-summands-finite-free}
on voit que $R\SheafHom(E,F)$ est repr\'esent\'e par le complexe de termes
$$
\bigoplus\nolimits_{n = - p + q}
\SheafHom_{\mathcal{O}_X}(\mathcal{E}^p, \mathcal{F}^q)
$$
En appliquant Cohomologie sur les sites, lemme
\ref{sites-cohomology-lemma-Rhom-complex-of-direct-summands-finite-free}
on voit que $R\SheafHom(\epsilon^*E,\epsilon^*F)$ est repr\'esent\'e par le
complexe de termes
$$
\bigoplus\nolimits_{n = - p + q}
\SheafHom_{\mathcal{O}_\etale}
(\epsilon^*\mathcal{E}^p, \epsilon^*\mathcal{F}^q)
$$
L'assertion du lemme se ram\`ene donc au fait, vrai, que le morphisme canonique
$$
\epsilon^*\SheafHom_{\mathcal{O}_X}(\mathcal{E}, \mathcal{F})
\longrightarrow
\SheafHom_{\mathcal{O}_\etale}
(\epsilon^*\mathcal{E}, \epsilon^*\mathcal{F})
$$
est un isomorphisme pour tout $\mathcal{O}_X$-module $\mathcal{F}$ et tout
$\mathcal{O}_X$-module libre de type fini $\mathcal{E}$.

\medskip\noindent
Dans le cas (2), on peut repr\'esenter $E$ par un complexe strictement parfait
$\mathcal{E}^\bullet$ de $\mathcal{O}_X$-modules; on utilise Cat\'egories
d\'eriv\'ees des sch\'emas, lemmes
\ref{perfect-lemma-affine-compare-bounded} et
\ref{perfect-lemma-perfect-affine}, ainsi que le fait qu'un complexe parfait
de modules est repr\'esent\'e par un complexe fini de modules projectifs de
type fini. On peut donc reprendre exactement la d\'emonstration ci-dessus en
rempla\c{c}ant la r\'ef\'erence \`a Cohomologie, lemme
\ref{cohomology-lemma-Rhom-complex-of-direct-summands-finite-free}
par une r\'ef\'erence \`a Cohomologie, lemme
\ref{cohomology-lemma-Rhom-strictly-perfect}.
\end{proof}

\begin{lemma}
\label{lemma-quasi-coherence-internal-hom}
Soit $S$ un sch\'ema, soit $X$ un espace alg\'ebrique sur $S$ et soient $L,K$
des objets de $D(\mathcal{O}_X)$. Si l'une des conditions suivantes est
satisfaite :
\begin{enumerate}
\item $L$ appartient \`a $D^+_\QCoh(\mathcal{O}_X)$ et $K$ est
pseudo-coh\'erent;
\item $L$ appartient \`a $D_\QCoh(\mathcal{O}_X)$ et $K$ est parfait;
\end{enumerate}
alors $R\SheafHom(K,L)$ appartient \`a $D_\QCoh(\mathcal{O}_X)$.
\end{lemma}

\begin{proof}
Cela r\'esulte de l'analogue pour les sch\'emas (Cat\'egories d\'eriv\'ees des
sch\'emas, lemme \ref{perfect-lemma-quasi-coherence-internal-hom}), du crit\`ere
du lemme \ref{lemma-check-quasi-coherence-on-covering}, des crit\`eres des
lemmes \ref{lemma-descend-pseudo-coherent} et \ref{lemma-descend-perfect},
et du r\'esultat du lemme \ref{lemma-descend-RHom}.
\end{proof}

\begin{lemma}
\label{lemma-internal-hom-evaluate-tensor-isomorphism}
Soit $S$ un sch\'ema, soit $X$ un espace alg\'ebrique sur $S$ et soient
$K,L,M$ des objets de $D_\QCoh(\mathcal{O}_X)$. Le morphisme
$$
K \otimes_{\mathcal{O}_X}^\mathbf{L} R\SheafHom(M, L)
\longrightarrow
R\SheafHom(M, K \otimes_{\mathcal{O}_X}^\mathbf{L} L)
$$
de Cohomologie sur les sites, lemme
\ref{sites-cohomology-lemma-internal-hom-diagonal-better}
est un isomorphisme dans les cas suivants :
\begin{enumerate}
\item $M$ est parfait;
\item $K$ est parfait;
\item $M$ est pseudo-coh\'erent, $L \in D^+(\mathcal{O}_X)$ et $K$ est de
Tor-dimension finie.
\end{enumerate}
\end{lemma}

\begin{proof}
La propri\'et\'e pour ce morphisme d'\^etre un isomorphisme se v\'erifie
localement pour la topologie \'etale; on peut donc supposer que $X$ est un
sch\'ema affine. D'apr\`es le lemme
\ref{lemma-derived-quasi-coherent-small-etale-site}, on peut alors \'ecrire
$K,L,M$ sous la forme $\epsilon^*K_0,\epsilon^*L_0,\epsilon^*M_0$, pour des
$K_0,L_0,M_0$ de $D_\QCoh(\mathcal{O}_X)$. Le lemme
\ref{lemma-descend-RHom} ram\`ene les cas (1) et (3) au cas des sch\'emas,
qui est Cat\'egories d\'eriv\'ees des sch\'emas, lemme
\ref{perfect-lemma-internal-hom-evaluate-tensor-isomorphism}.
Si $K$ est parfait mais qu'aucune autre hypoth\`ese n'est faite, on ne sait pas
si l'un ou l'autre membre de la fl\`eche appartient \`a
$D_\QCoh(\mathcal{O}_X)$; le r\'esultat reste toutefois vrai, car, apr\`es une
localisation suppl\'ementaire, $K$ sera repr\'esent\'e par un complexe fini de
modules libres de type fini, auquel cas l'assertion est claire.
\end{proof}









\section{Approximation par des complexes parfaits}
\label{section-approximation}

\noindent
Dans cette section, nous poursuivons l'\'etude commenc\'ee dans Cat\'egories
d\'eriv\'ees des sch\'emas, section \ref{perfect-section-approximation}.

\begin{definition}
\label{definition-approximation-holds}
Soit $S$ un sch\'ema et soit $X$ un espace alg\'ebrique sur $S$.
Consid\'erons les triplets $(T,E,m)$ tels que
\begin{enumerate}
\item $T \subset |X|$ est une partie ferm\'ee;
\item $E$ est un objet de $D_\QCoh(\mathcal{O}_X)$;
\item $m \in \mathbf{Z}$.
\end{enumerate}
On dit que {\it la propri\'et\'e d'approximation est satisfaite pour le triplet}
$(T,E,m)$ s'il existe un objet parfait $P$ de $D(\mathcal{O}_X)$ \`a support
dans $T$ et un morphisme $\alpha : P \to E$ qui induit des isomorphismes
$H^i(P) \to H^i(E)$ pour $i>m$ et une surjection
$H^m(P) \to H^m(E)$.
\end{definition}

\noindent
La propri\'et\'e d'approximation ne peut pas \^etre satisfaite pour tout
triplet. Les remarques qui suivent Cat\'egories d\'eriv\'ees des sch\'emas,
d\'efinition \ref{perfect-definition-approximation-holds}, en expliquent la
raison.

\begin{definition}
\label{definition-approximation}
Soit $S$ un sch\'ema et soit $X$ un espace alg\'ebrique sur $S$.
On dit que {\it la propri\'et\'e d'approximation par des complexes parfaits est
satisfaite} sur $X$ si, pour toute partie ferm\'ee $T \subset |X|$ telle que le
morphisme $X \setminus T \to X$ soit quasi-compact, il existe un entier $r$
tel que, pour tout triplet $(T,E,m)$ comme dans la d\'efinition
\ref{definition-approximation-holds} et v\'erifiant
\begin{enumerate}
\item $E$ est $(m-r)$-pseudo-coh\'erent;
\item $H^i(E)$ est \`a support dans $T$ pour $i \geq m-r$,
\end{enumerate}
la propri\'et\'e d'approximation soit satisfaite.
\end{definition}

\begin{lemma}
\label{lemma-pushforward-perfect}
Soit $S$ un sch\'ema et soit $(U \subset X, j : V \to X)$ un carr\'e
distingu\'e \'el\'ementaire d'espaces alg\'ebriques sur $S$.
Soit $E$ un objet parfait de $D(\mathcal{O}_V)$ \`a support dans
$j^{-1}(T)$, o\`u $T=|X|\setminus|U|$. Alors $Rj_*E$ est un objet parfait de
$D(\mathcal{O}_X)$.
\end{lemma}

\begin{proof}
La propri\'et\'e d'\^etre parfait est locale sur $X_\etale$. Il suffit donc de
v\'erifier que les restrictions de $Rj_*E$ \`a $U$ et \`a $V$ sont parfaites.
On a $Rj_*E|_V=E$ d'apr\`es le lemme
\ref{lemma-pushforward-with-support-in-open}, et cet objet est parfait.
On a $Rj_*E|_U=0$, car $E|_{V\setminus j^{-1}(T)}=0$ (utiliser le lemme
\ref{lemma-restrict-direct-image-open}).
\end{proof}

\begin{lemma}
\label{lemma-open}
Soit $S$ un sch\'ema et soit $(U \subset X,j : V \to X)$ un carr\'e distingu\'e
\'el\'ementaire d'espaces alg\'ebriques sur $S$. Soit $T$ une partie ferm\'ee de
$|X|\setminus|U|$ et soit $(T,E,m)$ un triplet comme dans la d\'efinition
\ref{definition-approximation-holds}. Si
\begin{enumerate}
\item la propri\'et\'e d'approximation est satisfaite pour
$(j^{-1}T,E|_V,m)$;
\item les faisceaux $H^i(E)$ sont \`a support dans $T$ pour $i\geq m$,
\end{enumerate}
alors la propri\'et\'e d'approximation est satisfaite pour $(T,E,m)$.
\end{lemma}

\begin{proof}
Soit $P \to E|_V$ une approximation du triplet $(j^{-1}T,E|_V,m)$ sur $V$.
D'apr\`es le lemme \ref{lemma-pushforward-perfect}, $Rj_*P$ est un objet
parfait de $D(\mathcal{O}_X)$. D'autre part, le lemme
\ref{lemma-pushforward-with-support-in-open} donne $Rj_*P=j_!P$.
L'objet $j_!P$ est \`a support dans $T$, comme le montre par exemple
(\ref{equation-stalk-j-shriek}). On obtient donc une approximation
$Rj_*P=j_!P \to j_!(E|_V) \to E$.
\end{proof}

\begin{lemma}
\label{lemma-approximation-affine}
Soit $S$ un sch\'ema et soit $X$ un espace alg\'ebrique sur $S$ repr\'esentable
par un sch\'ema affine. Alors la propri\'et\'e d'approximation est satisfaite
pour tout triplet $(T,E,m)$ comme dans la d\'efinition
\ref{definition-approximation-holds} pour lequel il existe un entier
$r\geq0$ tel que
\begin{enumerate}
\item $E$ est $m$-pseudo-coh\'erent;
\item $H^i(E)$ est \`a support dans $T$ pour $i\geq m-r+1$;
\item $X\setminus T$ est la r\'eunion de $r$ ouverts affines.
\end{enumerate}
En particulier, la propri\'et\'e d'approximation par des complexes parfaits est
satisfaite pour les sch\'emas affines.
\end{lemma}

\begin{proof}
Soit $X_0$ un sch\'ema affine repr\'esentant $X$ et soit $T_0\subset X_0$ la
partie ferm\'ee correspondant \`a $T$. Soit
$\epsilon : X_\etale \to X_{0,Zar}$ le morphisme
(\ref{equation-epsilon}). On peut \'ecrire $E=\epsilon^*E_0$ pour un objet
$E_0$ de $D_\QCoh(\mathcal{O}_{X_0})$; voir le lemme
\ref{lemma-derived-quasi-coherent-small-etale-site}.
Le lemme \ref{lemma-descend-pseudo-coherent} montre que $E_0$ est
$m$-pseudo-coh\'erent. En comparant les fibres des faisceaux de cohomologie
(voir la preuve du lemme \ref{lemma-epsilon-flat}), on voit que $H^i(E_0)$ est
\`a support dans $T_0$ pour $i\geq m-r+1$. D'apr\`es Cat\'egories d\'eriv\'ees
des sch\'emas, lemme \ref{perfect-lemma-approximation-affine}, il existe une
approximation $P_0\to E_0$ de $(T_0,E_0,m)$. Le lemme
\ref{lemma-descend-perfect} montre que $P=\epsilon^*P_0$ est un objet parfait
de $D(\mathcal{O}_X)$. En prenant l'image inverse, on obtient l'approximation
voulue $P=\epsilon^*P_0\to\epsilon^*E_0=E$.
\end{proof}

\begin{lemma}
\label{lemma-induction-step}
Soit $S$ un sch\'ema et soit $(U \subset X,j : V \to X)$ un carr\'e distingu\'e
\'el\'ementaire d'espaces alg\'ebriques sur $S$. Supposons $U$ quasi-compact,
$V$ affine et $U\times_XV$ quasi-compact. Si la propri\'et\'e d'approximation
par des complexes parfaits est satisfaite sur $U$, elle l'est sur $X$.
\end{lemma}

\begin{proof}
Soit $T\subset|X|$ une partie ferm\'ee telle que
$X\setminus T\to X$ soit quasi-compact. Soit $r_U$ l'entier de la d\'efinition
\ref{definition-approximation} associ\'e au couple $(U,T\cap|U|)$.
Posons $T'=T\setminus|U|$ et munissons $T'$ de la structure r\'eduite induite
de sous-espace. Comme $|T'|$ est contenu dans $|X|\setminus|U|$, le morphisme
$j^{-1}(T')\to T'$ est un isomorphisme. De plus,
$V\setminus j^{-1}(T')$ est quasi-compact : c'est le produit fibr\'e de
$U\times_XV$ et de $X\setminus T$ au-dessus de $X$, et nous avons suppos\'e
$U\times_XV$ quasi-compact ainsi que $X\setminus T\to X$ quasi-compact.
Soit $r'$ le nombre d'ouverts affines n\'ecessaires pour recouvrir
$V\setminus j^{-1}(T')$. Nous affirmons que $r=\max(r_U,r')$ convient au
couple $(X,T)$.

\medskip\noindent
Pour le voir, choisissons un triplet $(T,E,m)$ tel que $E$ soit
$(m-r)$-pseudo-coh\'erent et que $H^i(E)$ soit \`a support dans $T$ pour
$i\geq m-r$. Soit $t$ le plus grand entier tel que $H^t(E)|_U$ soit non nul.
(Un tel entier existe, car $U$ est quasi-compact et $E|_U$ est
$(m-r)$-pseudo-coh\'erent.) Nous montrerons par r\'ecurrence sur $t$ que $E$
admet une approximation.

\medskip\noindent
Initialisation : $t\leq m-r'$. Cela signifie que $H^i(E)$ est \`a support
dans $T'$ pour $i\geq m-r'$. Le lemme
\ref{lemma-approximation-affine} assure donc l'existence sur $V$ d'une
approximation $P\to E|_V$ de $(T',E|_V,m)$. En appliquant le lemme
\ref{lemma-open}, on voit que $(T',E,m)$ admet une approximation. Celle-ci est
aussi une approximation de $(T,E,m)$.

\medskip\noindent
H\'er\'edit\'e. Choisissons une approximation $P\to E|_U$ de
$(T\cap|U|,E|_U,m)$. Elle fournit en particulier une surjection
$H^t(P)\to H^t(E|_U)$. Dans la suite de la preuve, nous utiliserons
l'\'equivalence du lemme
\ref{lemma-derived-quasi-coherent-small-etale-site} (et les compatibilit\'es de
la remarque \ref{remark-match-total-direct-images}) pour les espaces
alg\'ebriques repr\'esentables $V$ et $U\times_XV$. Nous utiliserons aussi le
fait que les notions de $(m-r)$-pseudo-coh\'erence, respectivement de
perfection, co\"incident sur les sites de Zariski et \'etale; voir les lemmes
\ref{lemma-descend-pseudo-coherent} et \ref{lemma-descend-perfect}.
On peut donc appliquer les r\'esultats de Cat\'egories d\'eriv\'ees des sch\'emas,
section \ref{perfect-section-lift}, \`a l'immersion ouverte
$U\times_XV\subset V$. Ainsi, Cat\'egories d\'eriv\'ees des sch\'emas, lemme
\ref{perfect-lemma-lift-perfect-complex-plus-shift-support}, montre qu'il
existe un objet parfait $Q$ de $D(\mathcal{O}_V)$, \`a support dans
$j^{-1}(T)$, et un isomorphisme
$Q|_{U\times_XV}\to(P\oplus P[1])|_{U\times_XV}$. D'apr\`es Cat\'egories
d\'eriv\'ees des sch\'emas, lemme \ref{perfect-lemma-lift-map}, on peut
remplacer $Q$ par $Q\otimes^\mathbf{L}I$ et supposer que le morphisme
$$
Q|_{U \times_X V} \longrightarrow
(P \oplus P[1])|_{U \times_X V} \longrightarrow
P|_{U \times_X V} \longrightarrow E|_{U \times_X V}
$$
se rel\`eve en $Q\to E|_V$. Le lemme \ref{lemma-glue} fournit alors un
morphisme $a:R\to E$ de $D(\mathcal{O}_X)$ tel que $a|_U$ soit isomorphe \`a
$P\oplus P[1]\to E|_U$ et que $a|_V$ soit isomorphe \`a $Q\to E|_V$.
Ainsi $R$ est parfait et \`a support dans $T$, et le morphisme
$H^t(R)\to H^t(E)$ est surjectif apr\`es restriction \`a $U$.
Choisissons un triangle distingu\'e
$$
R \to E \to E' \to R[1]
$$
Alors $E'$ est $(m-r)$-pseudo-coh\'erent (Cohomologie sur les sites, lemme
\ref{sites-cohomology-lemma-cone-pseudo-coherent}), $H^i(E')|_U=0$ pour
$i\geq t$, et $H^i(E')$ est \`a support dans $T$ pour $i\geq m-r$.
Par l'hypoth\`ese de r\'ecurrence, on trouve une approximation $R'\to E'$ de
$(T,E',m)$. Ins\'erons le compos\'e $R'\to E'\to R[1]$ dans un triangle
distingu\'e $R\to R''\to R'\to R[1]$, puis prolongeons les morphismes
$R'\to E'$ et $R[1]\to R[1]$ en un morphisme de triangles distingu\'es
$$
\xymatrix{
R \ar[r] \ar[d] & R'' \ar[d] \ar[r] & R' \ar[d] \ar[r] & R[1] \ar[d] \\
R \ar[r] &  E \ar[r] & E' \ar[r] & R[1]
}
$$
en utilisant TR3. Alors $R''$ est un complexe parfait (Cohomologie sur les
sites, lemme \ref{sites-cohomology-lemma-two-out-of-three-perfect}) \`a support
dans $T$. Une chasse au diagramme imm\'ediate montre que $R''\to E$ est
l'approximation cherch\'ee.
\end{proof}

\begin{theorem}
\label{theorem-approximation}
Soit $S$ un sch\'ema et soit $X$ un espace alg\'ebrique quasi-compact et
quasi-s\'epar\'e sur $S$. Alors la propri\'et\'e d'approximation par des
complexes parfaits est satisfaite sur $X$.
\end{theorem}

\begin{proof}
Cela r\'esulte du principe de r\'ecurrence du lemme
\ref{lemma-induction-principle} et des lemmes \ref{lemma-induction-step} et
\ref{lemma-approximation-affine}.
\end{proof}






\section{Engendrement des cat\'egories d\'eriv\'ees}
\label{section-generating}

\noindent
Cette section est l'analogue de Cat\'egories d\'eriv\'ees des sch\'emas,
section \ref{perfect-section-generating}. Nous commen\c{c}ons toutefois par
d\'emontrer le lemme suivant, analogue de Cat\'egories d\'eriv\'ees des
sch\'emas, lemme
\ref{perfect-lemma-lift-map-from-perfect-complex-with-support}.

\begin{lemma}
\label{lemma-lift-map-from-perfect-complex-with-support}
Soit $S$ un sch\'ema. Soit $X$ un espace alg\'ebrique quasi-compact et
quasi-s\'epar\'e sur $S$. Soit $W\subset X$ un ouvert quasi-compact. Soit
$T\subset|X|$ une partie ferm\'ee telle que le morphisme
$X\setminus T\to X$ soit quasi-compact. Soit $E$ un objet de
$D_\QCoh(\mathcal{O}_X)$. Soit $\alpha:P\to E|_W$ un morphisme, o\`u $P$ est
un objet parfait de $D(\mathcal{O}_W)$ \`a support dans $T\cap W$. Il existe
alors un morphisme $\beta:R\to E$, o\`u $R$ est un objet parfait de
$D(\mathcal{O}_X)$ \`a support dans $T$, tel que $P$ soit un facteur direct
de $R|_W$ dans $D(\mathcal{O}_W)$, de fa\c{c}on compatible avec $\alpha$ et
$\beta|_W$.
\end{lemma}

\begin{proof}
Nous allons utiliser le principe de r\'ecurrence du lemme
\ref{lemma-induction-principle-enlarge}. On se ram\`ene donc imm\'ediatement
au cas d'un carr\'e distingu\'e \'el\'ementaire
$(W\subset X,f:V\to X)$, avec $V$ affine et $P\to E|_W$ comme dans l'\'enonc\'e
du lemme. Dans la suite de la d\'emonstration, nous utiliserons le lemme
\ref{lemma-derived-quasi-coherent-small-etale-site} (ainsi que les
compatibilit\'es de la remarque \ref{remark-match-total-direct-images}) pour
les espaces alg\'ebriques repr\'esentables $V$ et $W\times_XV$. Nous
utiliserons aussi le fait que la perfection sur le site de Zariski et sur le
site \'etale co\"incident; voir le lemme \ref{lemma-descend-perfect}.

\medskip\noindent
D'apr\`es Cat\'egories d\'eriv\'ees des sch\'emas, lemme
\ref{perfect-lemma-lift-perfect-complex-plus-shift-support}, on peut choisir un
objet parfait $Q$ de $D(\mathcal{O}_V)$, \`a support dans $f^{-1}T$, et un
isomorphisme
$Q|_{W\times_XV}\to(P\oplus P[1])|_{W\times_XV}$. D'apr\`es Cat\'egories
d\'eriv\'ees des sch\'emas, lemme \ref{perfect-lemma-lift-map}, on peut
remplacer $Q$ par $Q\otimes^\mathbf{L}I$ (qui est encore \`a support dans
$f^{-1}T$) et supposer que le morphisme
$$
Q|_{W \times_X V} \to (P \oplus P[1])|_{W \times_X V}
\longrightarrow P|_{W \times_X V}
\longrightarrow
E|_{W \times_X V}
$$
se rel\`eve en $Q\to E|_V$. Le lemme \ref{lemma-glue} fournit un morphisme
$a:R\to E$ de $D(\mathcal{O}_X)$ tel que $a|_W$ soit isomorphe \`a
$P\oplus P[1]\to E|_W$ et que $a|_V$ soit isomorphe \`a $Q\to E|_V$.
Ainsi, $R$ est parfait et \`a support dans $T$, comme souhait\'e.
\end{proof}

\begin{remark}
\label{remark-addendum}
La d\'emonstration du lemme
\ref{lemma-lift-map-from-perfect-complex-with-support} montre que
$$
R|_W = P \oplus P^{\oplus n_1}[1] \oplus \ldots \oplus P^{\oplus n_m}[m]
$$
pour certains $m\geq0$ et $n_j\geq0$. Le faisceau de cohomologie de plus haut
degr\'e de $R|_W$ est donc \`egal \`a celui de $P$. En r\'ep\'etant la
construction pour le morphisme
$P^{\oplus n_1}[1]\oplus\ldots\oplus P^{\oplus n_m}[m]\to R|_W$, en prenant
des c\^ones et en raisonnant par r\'ecurrence, on peut obtenir l'\'egalit\'e des
faisceaux de cohomologie de $R|_W$ et de $P$ au-dessus de tout degr\'e fix\'e.
\end{remark}

\begin{lemma}
\label{lemma-direct-summand-of-a-restriction}
Soit $S$ un sch\'ema. Soit $X$ un espace alg\'ebrique quasi-compact et
quasi-s\'epar\'e sur $S$. Soit $W$ un sous-espace ouvert quasi-compact de $X$.
Soit $P$ un objet parfait de $D(\mathcal{O}_W)$. Alors $P$ est un facteur
direct de la restriction d'un objet parfait de $D(\mathcal{O}_X)$.
\end{lemma}

\begin{proof}
C'est un cas particulier du lemme
\ref{lemma-lift-map-from-perfect-complex-with-support}.
\end{proof}

\begin{theorem}
\label{theorem-bondal-van-den-Bergh}
Soit $S$ un sch\'ema. Soit $X$ un espace alg\'ebrique quasi-compact et
quasi-s\'epar\'e sur $S$. La cat\'egorie $D_\QCoh(\mathcal{O}_X)$ peut \^etre
engendr\'ee par un seul objet parfait. Plus pr\'ecis\'ement, il existe un objet
parfait $P$ de $D(\mathcal{O}_X)$ tel que, pour
$E\in D_\QCoh(\mathcal{O}_X)$, les conditions suivantes soient \'equivalentes :
\begin{enumerate}
\item $E = 0$;
\item $\Hom_{D(\mathcal{O}_X)}(P[n], E) = 0$ pour tout $n \in \mathbf{Z}$.
\end{enumerate}
\end{theorem}

\begin{proof}
Nous allons le d\'emontrer au moyen du principe de r\'ecurrence du lemme
\ref{lemma-induction-principle}.

\medskip\noindent
Si $X$ est affine, alors $\mathcal{O}_X$ est un g\'en\'erateur parfait. Cela
r\'esulte du lemme \ref{lemma-derived-quasi-coherent-small-etale-site} et de
Cat\'egories d\'eriv\'ees des sch\'emas, lemme
\ref{perfect-lemma-affine-compare-bounded}.

\medskip\noindent
Supposons que $(U\subset X,f:V\to X)$ soit un carr\'e distingu\'e
\'el\'ementaire, avec $U$ quasi-compact, tel que le th\'eor\`eme soit vrai pour
$U$ et que $V$ soit un sch\'ema affine. Soit $P$ un objet parfait de
$D(\mathcal{O}_U)$ qui engendre $D_\QCoh(\mathcal{O}_U)$. En utilisant le
lemme \ref{lemma-direct-summand-of-a-restriction}, on peut choisir un objet
parfait $Q$ de $D(\mathcal{O}_X)$ dont la restriction \`a $U$ est une somme
directe ayant $P$ pour facteur direct. Posons $V=\Spec(A)$. Soit $Z\subset V$
le sous-sch\'ema ferm\'e r\'eduit, image inverse de $X\setminus U$, qui s'envoie
isomorphiquement sur celui-ci (voir la d\'efinition
\ref{definition-elementary-distinguished-square}). C'est une partie ferm\'ee
r\'etrocompacte de $V$. Choisissons $f_1,\ldots,f_r\in A$ tels que
$Z=V(f_1,\ldots,f_r)$. Soit $K\in D(\mathcal{O}_V)$ l'objet parfait
correspondant au complexe de Koszul associ\'e \`a $f_1,\ldots,f_r$ sur $A$.
Comme $K$ est \`a support dans $Z$, son image directe $K'=Rf_*K$ est un objet
parfait de $D(\mathcal{O}_X)$ dont la restriction \`a $V$ est $K$ (voir les
lemmes \ref{lemma-pushforward-perfect} et
\ref{lemma-pushforward-with-support-in-open}). Nous affirmons que
$Q\oplus K'$ engendre $D_\QCoh(\mathcal{O}_X)$.

\medskip\noindent
Soit $E$ un objet de $D_\QCoh(\mathcal{O}_X)$ tel qu'il n'existe aucun
morphisme non nul d'un translat\'e de $Q\oplus K'$ vers $E$. D'apr\`es le lemme
\ref{lemma-pushforward-with-support-in-open}, on a $K'=f_!K$, donc
$$
\Hom_{D(\mathcal{O}_X)}(K'[n], E) = \Hom_{D(\mathcal{O}_V)}(K[n], E|_V)
$$
Par suite, Cat\'egories d\'eriv\'ees des sch\'emas, lemme
\ref{perfect-lemma-orthogonal-koszul-complex} (en utilisant aussi le lemme
\ref{lemma-derived-quasi-coherent-small-etale-site}), montre que l'annulation
de ces groupes entra\^ine que $E|_V$ est isomorphe \`a
$R(U\times_XV\to V)_*E|_{U\times_XV}$. Il en r\'esulte que
$E=R(U\to X)_*E|_U$ (nous omettons un d\'etail mineur). Dans ce cas,
$$
\Hom_{D(\mathcal{O}_X)}(Q[n], E) = \Hom_{D(\mathcal{O}_U)}(Q|_U[n], E|_U)
$$
contient $\Hom_{D(\mathcal{O}_U)}(P[n],E|_U)$ comme facteur direct. Le choix
de $P$ montre donc que l'annulation de ces groupes entra\^ine $E|_U=0$, d'o\`u
$E=0$.
\end{proof}

\begin{remark}
\label{remark-pullback-generator}
Soit $S$ un sch\'ema. Soit $f:X\to Y$ un morphisme d'espaces alg\'ebriques
quasi-compacts et quasi-s\'epar\'es sur $S$. Soit
$E\in D_\QCoh(\mathcal{O}_Y)$ un g\'en\'erateur (voir le th\'eor\`eme
\ref{theorem-bondal-van-den-Bergh}). Les conditions suivantes sont alors
\'equivalentes :
\begin{enumerate}
\item pour $K\in D_\QCoh(\mathcal{O}_X)$, on a $Rf_*K=0$ si et seulement si
$K=0$;
\item $Rf_* : D_\QCoh(\mathcal{O}_X) \to D_\QCoh(\mathcal{O}_Y)$ est
conservatif;
\item $Lf^*E$ engendre $D_\QCoh(\mathcal{O}_X)$.
\end{enumerate}
L'\'equivalence de (1) et (2) est une cons\'equence formelle du fait que
$Rf_*:D_\QCoh(\mathcal{O}_X)\to D_\QCoh(\mathcal{O}_Y)$ est un foncteur exact
de cat\'egories triangul\'ees. De m\^eme, l'\'equivalence de (1) et (3) r\'esulte
formellement du fait que $Lf^*$ est adjoint \`a gauche de $Rf_*$. Ces
conditions sont satisfaites si $f$ est affine (lemme
\ref{lemma-affine-morphism}), si $f$ est une immersion ouverte, ou si $f$ est
une compos\'ee de morphismes de ces types.
\end{remark}

\noindent
Le r\'esultat suivant renforce le th\'eor\`eme
\ref{theorem-bondal-van-den-Bergh}; il se d\'emontre par les m\^emes m\'ethodes.
Soit $T\subset|X|$ une partie ferm\'ee, o\`u $X$ est un espace alg\'ebrique.
Notons $D_T(\mathcal{O}_X)$ la sous-cat\'egorie triangul\'ee strictement pleine
et satur\'ee form\'ee des complexes dont les faisceaux de cohomologie sont \`a
support dans $T$.

\begin{lemma}
\label{lemma-generator-with-support}
Soit $S$ un sch\'ema. Soit $X$ un espace alg\'ebrique quasi-compact et
quasi-s\'epar\'e sur $S$. Soit $T\subset|X|$ une partie ferm\'ee telle que
$|X|\setminus T$ soit quasi-compact. Avec les notations ci-dessus, la
cat\'egorie $D_{\QCoh,T}(\mathcal{O}_X)$ est engendr\'ee par un seul objet
parfait.
\end{lemma}

\begin{proof}
Nous allons le d\'emontrer au moyen du principe de r\'ecurrence du lemme
\ref{lemma-induction-principle}. D'apr\`es Cat\'egories d\'eriv\'ees des
sch\'emas, lemme \ref{perfect-lemma-generator-with-support}, la propri\'et\'e est
vraie pour les objets repr\'esentables, quasi-compacts et quasi-s\'epar\'es du
site $X_{spaces,\etale}$.

\medskip\noindent
Supposons que $(U\subset X,f:V\to X)$ soit un carr\'e distingu\'e
\'el\'ementaire tel que le lemme soit vrai pour $U$ et que $V$ soit affine.
Pour achever la d\'emonstration, il faut \'etablir le r\'esultat pour $X$. Soit
$P$ un objet parfait de $D(\mathcal{O}_U)$, \`a support dans $T\cap U$, qui
engendre $D_{\QCoh,T\cap U}(\mathcal{O}_U)$. D'apr\`es le lemme
\ref{lemma-lift-map-from-perfect-complex-with-support}, on peut choisir un
objet parfait $Q$ de $D(\mathcal{O}_X)$, \`a support dans $T$, dont la
restriction \`a $U$ est une somme directe ayant $P$ pour facteur direct.
\'Ecrivons $V=\Spec(B)$ et posons $Z=X\setminus U$. Alors $f^{-1}Z$ est une
partie ferm\'ee de $V$ telle que $V\setminus f^{-1}Z$ soit quasi-compact.
Comme $X$ est quasi-s\'epar\'e, il s'ensuit que
$f^{-1}Z\cap f^{-1}T=f^{-1}(Z\cap T)$ est une partie ferm\'ee de $V$ telle que
$W=V\setminus f^{-1}(Z\cap T)$ soit quasi-compact. On peut donc choisir
$g_1,\ldots,g_s\in B$ tels que
$f^{-1}(Z\cap T)=V(g_1,\ldots,g_s)$. Soit $K\in D(\mathcal{O}_V)$ l'objet
parfait correspondant au complexe de Koszul associ\'e \`a
$g_1,\ldots,g_s$ sur $B$. Comme $K$ est \`a support dans la partie ferm\'ee
$f^{-1}(Z\cap T)\subset V$, son image directe $K'=R(V\to X)_*K$ est un objet
parfait de $D(\mathcal{O}_X)$ dont la restriction \`a $V$ est $K$ (voir les
lemmes \ref{lemma-pushforward-perfect} et
\ref{lemma-pushforward-with-support-in-open}). Nous affirmons que
$Q\oplus K'$ engendre $D_{\QCoh,T}(\mathcal{O}_X)$.

\medskip\noindent
Soit $E$ un objet de $D_{\QCoh,T}(\mathcal{O}_X)$ tel qu'il n'existe aucun
morphisme non nul d'un translat\'e de $Q\oplus K'$ vers $E$. D'apr\`es le lemme
\ref{lemma-pushforward-with-support-in-open}, on a
$K'=R(V\to X)_!K$, donc
$$
\Hom_{D(\mathcal{O}_X)}(K'[n], E) = \Hom_{D(\mathcal{O}_V)}(K[n], E|_V)
$$
Cat\'egories d\'eriv\'ees des sch\'emas, lemme
\ref{perfect-lemma-orthogonal-koszul-complex}, donne donc
$E|_V=Rj_*E|_W$, o\`u $j:W\to V$ est l'inclusion. On a le diagramme
$$
\xymatrix{
W \ar[r]_j & V & f^{-1}(Z \cap T) \ar[l] \ar[d] \\
V \setminus f^{-1}Z \ar[u]^{j'} \ar[ru]_{j''} & & f^{-1}Z \ar[lu]
}
$$
Comme $E$ est \`a support dans $T$, on voit que $E|_W$ est \`a support dans
$f^{-1}T\cap W=f^{-1}T\cap(V\setminus f^{-1}Z)$, qui est ferm\'e dans $W$.
On en conclut que
$$
E|_V = Rj_*(E|_W) = Rj_*(Rj'_*(E|_{U \times_X V})) =
Rj''_*(E|_{U \times_X V})
$$
La seconde \'egalit\'e est l'assertion (1) de Cohomologie, lemme
\ref{cohomology-lemma-pushforward-restriction}. Celui-ci s'applique parce que
$V$ est un sch\'ema et que $E$ a des faisceaux de cohomologie quasi-coh\'erents;
l'image directe par l'immersion ouverte quasi-compacte $j'$ co\"incide donc
avec l'image directe sur les sch\'emas sous-jacents; voir la remarque
\ref{remark-match-total-direct-images}. Il en r\'esulte que
$E=R(U\to X)_*E|_U$ (nous omettons un d\'etail mineur). Dans ce cas,
$$
\Hom_{D(\mathcal{O}_X)}(Q[n], E) = \Hom_{D(\mathcal{O}_U)}(Q|_U[n], E|_U)
$$
contient $\Hom_{D(\mathcal{O}_U)}(P[n],E|_U)$ comme facteur direct. Le choix
de $P$ montre donc que l'annulation de ces groupes entra\^ine $E|_U=0$, d'o\`u
$E=0$.
\end{proof}










\section{Objets compacts et parfaits}
\label{section-compact}

\noindent
Cette section est l'analogue de
Cat\'egories d\'eriv\'ees des sch\'emas, section \ref{perfect-section-compact}.

\begin{proposition}
\label{proposition-compact-is-perfect}
Soit $S$ un sch\'ema.
Soit $X$ un espace alg\'ebrique quasi-compact et quasi-s\'epar\'e sur $S$.
Un objet de $D_\QCoh(\mathcal{O}_X)$ est compact
si et seulement s'il est parfait.
\end{proposition}

\begin{proof}
Si $K$ est un objet parfait de $D(\mathcal{O}_X)$, de dual
$K^\vee$ (Cohomologie sur les sites, lemme
\ref{sites-cohomology-lemma-dual-perfect-complex}),
on a
$$
\Hom_{D(\mathcal{O}_X)}(K, M) =
H^0(X, K^\vee \otimes_{\mathcal{O}_X}^\mathbf{L} M)
$$
fonctoriellement en $M$. Comme $K^\vee \otimes_{\mathcal{O}_X}^\mathbf{L} -$
commute aux sommes directes et que $H^0(X, -)$ commute aux sommes
directes dans $D_\QCoh(\mathcal{O}_X)$ d'apr\`es le
lemme \ref{lemma-quasi-coherence-pushforward-direct-sums},
on en conclut que $K$ est compact dans $D_\QCoh(\mathcal{O}_X)$.

\medskip\noindent
R\'eciproquement, soit $K$ un objet compact de $D_\QCoh(\mathcal{O}_X)$.
Pour montrer que $K$ est parfait, il suffit de montrer que
$K|_U$ est parfait pour tout sch\'ema affine $U$ \'etale sur $X$; voir
Cohomologie sur les sites, lemme
\ref{sites-cohomology-lemma-perfect-independent-representative}.
Observons que $j : U \to X$ est un morphisme quasi-compact et s\'epar\'e.
Par cons\'equent,
$Rj_* : D_\QCoh(\mathcal{O}_U) \to D_\QCoh(\mathcal{O}_X)$
commute aux sommes directes; voir le
lemme \ref{lemma-quasi-coherence-pushforward-direct-sums}.
L'adjonction entre la restriction \`a $U$ et $Rj_*$ montre donc que
$K|_U$ est un objet compact de $D_\QCoh(\mathcal{O}_U)$.
On se ram\`ene ainsi au cas o\`u $X$ est affine, et en particulier \`a celui
d'un sch\'ema quasi-compact et quasi-s\'epar\'e. Au moyen des
lemmes \ref{lemma-derived-quasi-coherent-small-etale-site} et
\ref{lemma-descend-perfect},
on se ram\`ene au cas des sch\'emas, c'est-\`a-dire \`a
Cat\'egories d\'eriv\'ees des sch\'emas, proposition
\ref{perfect-proposition-compact-is-perfect}.
\end{proof}

\begin{remark}
\label{remark-classical-generator}
Soit $S$ un sch\'ema.
Soit $X$ un espace alg\'ebrique quasi-compact et quasi-s\'epar\'e sur $S$.
Soit $G$ un
objet parfait de $D(\mathcal{O}_X)$ qui engendre
$D_\QCoh(\mathcal{O}_X)$. D'apr\`es le th\'eor\`eme
\ref{theorem-bondal-van-den-Bergh}, il en existe au moins un. En combinant
le lemme \ref{lemma-quasi-coherence-direct-sums}, la
proposition \ref{proposition-compact-is-perfect} et
Cat\'egories d\'eriv\'ees, proposition
\ref{derived-proposition-generator-versus-classical-generator},
on voit que $G$ est un g\'en\'erateur classique de
$D_{perf}(\mathcal{O}_X)$.
\end{remark}

\noindent
Le r\'esultat suivant renforce la
proposition \ref{proposition-compact-is-perfect}.
Soit $T \subset |X|$ une partie ferm\'ee, o\`u $X$ est un espace alg\'ebrique.
Comme pr\'ec\'edemment, $D_T(\mathcal{O}_X)$ d\'esigne la sous-cat\'egorie
triangul\'ee strictement pleine et satur\'ee form\'ee des complexes dont les
faisceaux de cohomologie sont \`a support dans $T$. Comme la formation des
sommes directes commute \`a celle des faisceaux de cohomologie, il s'ensuit
que $D_T(\mathcal{O}_X)$ admet des sommes directes et que celles-ci co\"incident
avec les sommes directes dans $D(\mathcal{O}_X)$.

\begin{lemma}
\label{lemma-compact-is-perfect-with-support}
Soit $S$ un sch\'ema.
Soit $X$ un espace alg\'ebrique quasi-compact et quasi-s\'epar\'e sur $S$.
Soit $T \subset |X|$ une partie ferm\'ee telle que $|X| \setminus T$
soit quasi-compact. Un objet de $D_{\QCoh, T}(\mathcal{O}_X)$ est compact
si et seulement s'il est parfait en tant qu'objet de $D(\mathcal{O}_X)$.
\end{lemma}

\begin{proof}
La remarque qui pr\'ec\`ede le lemme montre que
$D_{\QCoh, T}(\mathcal{O}_X)$ est une cat\'egorie triangul\'ee
admettant des sommes directes.
D'apr\`es la proposition \ref{proposition-compact-is-perfect},
les objets parfaits donnent des objets compacts de
$D_\QCoh(\mathcal{O}_X)$, et donc a fortiori de toute sous-cat\'egorie
stable par sommes directes. Pour la r\'eciproque, nous utiliserons le fait
qu'il existe un g\'en\'erateur
$E \in D_{\QCoh, T}(\mathcal{O}_X)$ qui est un complexe parfait
de $\mathcal{O}_X$-modules; voir le
lemme \ref{lemma-generator-with-support}.
D'apr\`es ce qui pr\'ec\`ede, $E$ est donc compact. Il r\'esulte alors de
Cat\'egories d\'eriv\'ees, proposition
\ref{derived-proposition-generator-versus-classical-generator},
que $E$ est un g\'en\'erateur classique de la sous-cat\'egorie pleine
des objets compacts de $D_{\QCoh, T}(\mathcal{O}_X)$.
Ainsi, tout objet compact se construit \`a partir de $E$ par
une suite finie d'op\'erations consistant \`a
(a) prendre des translat\'es, (b) prendre des sommes directes finies,
(c) prendre des c\^ones et
(d) prendre des facteurs directs. Chacune de ces op\'erations pr\'eserve
la propri\'et\'e d'\^etre parfait, ce qui prouve le r\'esultat.
\end{proof}

\begin{remark}
\label{remark-classical-generator-with-support}
Soit $S$ un sch\'ema.
Soit $X$ un espace alg\'ebrique quasi-compact et quasi-s\'epar\'e sur $S$.
Soit $T \subset |X|$ une partie ferm\'ee telle que $|X| \setminus T$
soit quasi-compact. Soit $G$ un
objet parfait de $D_{\QCoh, T}(\mathcal{O}_X)$ qui engendre
$D_{\QCoh, T}(\mathcal{O}_X)$. D'apr\`es le lemme
\ref{lemma-generator-with-support}, il en existe au moins un. En combinant
le fait que $D_{\QCoh, T}(\mathcal{O}_X)$ admet des sommes directes, le
lemme \ref{lemma-compact-is-perfect-with-support} et
Cat\'egories d\'eriv\'ees, proposition
\ref{derived-proposition-generator-versus-classical-generator},
on voit que $G$ est un g\'en\'erateur classique de
$D_{perf, T}(\mathcal{O}_X)$.
\end{remark}

\noindent
Le lemme suivant applique les id\'ees qui interviennent dans
la d\'emonstration du lemme pr\'ec\'edent.

\begin{lemma}
\label{lemma-map-from-pseudo-coherent-to-complex-with-support}
Soit $S$ un sch\'ema. Soit $X$ un espace alg\'ebrique quasi-compact et
quasi-s\'epar\'e sur $S$. Soit $T \subset |X|$
une partie ferm\'ee dont le compl\'ementaire $U \subset X$ est quasi-compact.
Soit $\alpha : P \to E$ un morphisme de $D_\QCoh(\mathcal{O}_X)$ se trouvant
dans l'un des deux cas suivants :
\begin{enumerate}
\item $P$ est parfait et $E$ est \`a support dans $T$;
\item $P$ est pseudo-coh\'erent, $E$ est \`a support dans $T$ et $E$ est
born\'e inf\'erieurement.
\end{enumerate}
Alors il existe un complexe parfait de $\mathcal{O}_X$-modules $I$
et un morphisme $I \to \mathcal{O}_X[0]$ tels que
$I \otimes^\mathbf{L} P \to E$ soit nul et que
$I|_U \to \mathcal{O}_U[0]$ soit un
isomorphisme.
\end{lemma}

\begin{proof}
Posons $\mathcal{D} = D_{\QCoh, T}(\mathcal{O}_X)$. Dans les deux cas, le complexe
$K = R\SheafHom(P, E)$ est un objet de $\mathcal{D}$. Voir le
lemme \ref{lemma-quasi-coherence-internal-hom} pour la quasi-coh\'erence.
Il est clair que $K$ est \`a support dans $T$, puisque la formation de
$R\SheafHom$ commute aux restrictions aux ouverts.
Le morphisme $\alpha$ d\'efinit un \'el\'ement de
$H^0(K) = \Hom_{D(\mathcal{O}_X)}(\mathcal{O}_X[0], K)$.
Il suffit donc de d\'emontrer le r\'esultat pour le morphisme
$\alpha : \mathcal{O}_X[0] \to K$.

\medskip\noindent
Soit $E \in \mathcal{D}$ un g\'en\'erateur parfait; voir le
lemme \ref{lemma-generator-with-support}. \'Ecrivons
$$
K = \text{hocolim} K_n
$$
comme dans Cat\'egories d\'eriv\'ees, lemme
\ref{derived-lemma-write-as-colimit}, en utilisant le g\'en\'erateur $E$.
Comme le foncteur $\mathcal{D} \to D(\mathcal{O}_X)$ commute aux sommes
directes, l'\'egalit\'e $K = \text{hocolim} K_n$ vaut aussi dans
$D(\mathcal{O}_X)$. Puisque $\mathcal{O}_X$ est un objet compact de
$D(\mathcal{O}_X)$, il existe un entier $n$ et un morphisme
$\alpha_n : \mathcal{O}_X \to K_n$ qui induit $\alpha$; voir
Cat\'egories d\'eriv\'ees, lemme
\ref{derived-lemma-commutes-with-countable-sums}.
En appliquant Cat\'egories d\'eriv\'ees, lemme
\ref{derived-lemma-factor-through}, au morphisme
$\mathcal{O}_X[0] \to K_n$ de la cat\'egorie ambiante
$D(\mathcal{O}_X)$, on voit que $\alpha_n$ se factorise sous la forme
$\mathcal{O}_X[0] \to Q \to K_n$, o\`u $Q$ est un objet
de $\langle E \rangle$. On en conclut que $Q$ est un complexe parfait
\`a support dans $T$.

\medskip\noindent
Choisissons un triangle distingu\'e
$$
I \to \mathcal{O}_X[0] \to Q \to I[1]
$$
Par construction, $I$ est parfait, le morphisme $I \to \mathcal{O}_X[0]$
se restreint \`a un isomorphisme sur $U$, et le compos\'e
$I \to K$ est nul puisque $\alpha$ se factorise par $Q$.
Cela prouve le lemme.
\end{proof}

\section{Les cat\'egories d\'eriv\'ees comme cat\'egories de modules}
\label{section-derived-is-dga}

\noindent
Cette section est l'analogue de Cat\'egories d\'eriv\'ees des sch\'emas,
section \ref{perfect-section-derived-is-dga}.

\begin{lemma}
\label{lemma-tensor-with-QCoh-complex}
Soit $S$ un sch\'ema et soit $X$ un espace alg\'ebrique sur $S$. Soit
$K^\bullet$ un complexe de $\mathcal{O}_X$-modules dont les faisceaux de
cohomologie sont quasi-coh\'erents. Soit
$(E,d)=\Hom_{\text{Comp}^{dg}(\mathcal{O}_X)}(K^\bullet,K^\bullet)$ l'alg\`ebre
diff\'erentielle gradu\'ee des endomorphismes. Alors le foncteur
$$
- \otimes_E^\mathbf{L} K^\bullet :
D(E, \text{d}) \longrightarrow D(\mathcal{O}_X)
$$
du chapitre Alg\`ebres diff\'erentielles gradu\'ees, lemme
\ref{dga-lemma-tensor-with-complex-derived}, prend ses valeurs dans
$D_\QCoh(\mathcal{O}_X)$.
\end{lemma}

\begin{proof}
Soit $P$ un $E$-module diff\'erentiel gradu\'e poss\'edant la propri\'et\'e P.
Soit $F_\bullet$ une filtration de $P$ comme dans Alg\`ebres diff\'erentielles
gradu\'ees, section \ref{dga-section-P-resolutions}. On a alors
$$
P \otimes_E K^\bullet = \text{hocolim}\ F_iP \otimes_E K^\bullet.
$$
Chacun des $F_iP$ poss\`ede une filtration finie dont les gradu\'es associ\'es
sont des sommes directes de $E[k]$. Le r\'esultat s'en d\'eduit ais\'ement.
\end{proof}

\noindent
On peut renforcer le lemme suivant : l'annulation est uniforme lorsque $L$
parcourt les complexes dont les faisceaux de cohomologie non nuls sont
concentr\'es dans un intervalle fix\'e.

\begin{lemma}
\label{lemma-ext-from-perfect-into-bounded-QCoh}
Soit $S$ un sch\'ema. Soit $X$ un espace alg\'ebrique quasi-compact et
quasi-s\'epar\'e sur $S$. Soit $K$ un objet parfait de
$D(\mathcal{O}_X)$. Alors
\begin{enumerate}
\item il existe des entiers $a\leq b$ tels que
$\Hom_{D(\mathcal{O}_X)}(K,L)=0$ pour tout
$L\in D_\QCoh(\mathcal{O}_X)$ tel que $H^i(L)=0$ pour $i\in[a,b]$;
\item si $L$ est born\'e, alors
$\Ext^n_{D(\mathcal{O}_X)}(K,L)$ est nul sauf pour un nombre fini de valeurs de
$n$.
\end{enumerate}
\end{lemma}

\begin{proof}
L'assertion (2) r\'esulte de (1), puisque
$\Ext^n_{D(\mathcal{O}_X)}(K,L)=
\Hom_{D(\mathcal{O}_X)}(K,L[n])$. D\'emontrons (1). Comme $K$ est parfait, on a
$$
\Ext^i_{D(\mathcal{O}_X)}(K, L) =
H^i(X, K^\vee \otimes_{\mathcal{O}_X}^\mathbf{L} L),
$$
o\`u $K^\vee$ est le complexe parfait dual de $K$; voir Cohomologie sur les
sites, lemme \ref{sites-cohomology-lemma-dual-perfect-complex}. Remarquons que
$P=K^\vee\otimes_{\mathcal{O}_X}^\mathbf{L}L$ appartient \`a $D_\QCoh(X)$
d'apr\`es les lemmes \ref{lemma-quasi-coherence-tensor-product} et
\ref{lemma-pseudo-coherent} (ce dernier montre qu'un complexe parfait a des
faisceaux de cohomologie quasi-coh\'erents). Supposons que $K^\vee$ soit de
Tor-amplitude contenue dans $[a,b]$. La suite spectrale
$$
E_1^{p, q} = H^p(K^\vee \otimes_{\mathcal{O}_X}^\mathbf{L} H^q(L))
\Rightarrow
H^{p + q}(K^\vee \otimes_{\mathcal{O}_X}^\mathbf{L} L)
$$
montre que $H^j(K^\vee\otimes_{\mathcal{O}_X}^\mathbf{L}L)$ est nul si
$H^q(L)=0$ pour $q\in[j-b,j-a]$. Soit $N$ l'entier
$\max(d_p+p)$ du lemme de Cohomologie des espaces
\ref{spaces-cohomology-lemma-vanishing-quasi-separated}. Alors
$H^0(X,K^\vee\otimes_{\mathcal{O}_X}^\mathbf{L}L)$ s'annule si les faisceaux de
cohomologie
$$
H^{-N}(K^\vee \otimes_{\mathcal{O}_X}^\mathbf{L} L),
\ H^{-N + 1}(K^\vee \otimes_{\mathcal{O}_X}^\mathbf{L} L),
\ \ldots,
\ H^0(K^\vee \otimes_{\mathcal{O}_X}^\mathbf{L} L)
$$
sont nuls. En effet, le lemme cit\'e et le lemme
\ref{lemma-application-nice-K-injective} donnent
$$
H^0(X, K^\vee \otimes_{\mathcal{O}_X}^\mathbf{L} L) =
H^0(X, \tau_{\geq -N}(K^\vee \otimes_{\mathcal{O}_X}^\mathbf{L} L));
$$
l'annulation des faisceaux de cohomologie montre que ce dernier groupe est
\'egal \`a
$H^0(X,\tau_{\geq1}(K^\vee\otimes_{\mathcal{O}_X}^\mathbf{L}L))$, lequel est
nul d'apr\`es Cat\'egories d\'eriv\'ees, lemme
\ref{derived-lemma-negative-vanishing}. Par cons\'equent,
$\Hom_{D(\mathcal{O}_X)}(K,L)$ est nul si $H^i(L)=0$ pour
$i\in[-b-N,-a]$.
\end{proof}

\noindent
Le th\'eor\`eme suivant est l'analogue de Cat\'egories d\'eriv\'ees des sch\'emas,
th\'eor\`eme \ref{perfect-theorem-DQCoh-is-Ddga}.

\begin{theorem}
\label{theorem-DQCoh-is-Ddga}
Soit $S$ un sch\'ema. Soit $X$ un espace alg\'ebrique quasi-compact et
quasi-s\'epar\'e sur $S$. Il existe alors une alg\`ebre diff\'erentielle gradu\'ee
$(E,\text{d})$ dont seuls un nombre fini de groupes de cohomologie $H^i(E)$
sont non nuls, et telle que $D_\QCoh(\mathcal{O}_X)$ soit \'equivalente \`a
$D(E,\text{d})$.
\end{theorem}

\begin{proof}
Soit $K^\bullet$ un complexe K-injectif de $\mathcal{O}$-modules qui soit
parfait et qui engendre $D_\QCoh(\mathcal{O}_X)$. Un tel complexe existe par le
th\'eor\`eme \ref{theorem-bondal-van-den-Bergh} et par l'existence des
r\'esolutions K-injectives. Nous allons montrer que le th\'eor\`eme est vrai pour
$$
(E, \text{d}) = \Hom_{\text{Comp}^{dg}(\mathcal{O}_X)}(K^\bullet, K^\bullet),
$$
o\`u $\text{Comp}^{dg}(\mathcal{O}_X)$ est la cat\'egorie diff\'erentielle
gradu\'ee des complexes de $\mathcal{O}$-modules. Voir Alg\`ebres
diff\'erentielles gradu\'ees, section
\ref{dga-section-variant-base-change}. Comme $K^\bullet$ est K-injectif, on a
\begin{equation}
\label{equation-E-is-OK}
H^n(E) = \Ext^n_{D(\mathcal{O}_X)}(K^\bullet, K^\bullet)
\end{equation}
pour tout $n\in\mathbf{Z}$. Seul un nombre fini de ces groupes Ext sont non
nuls, d'apr\`es le lemme \ref{lemma-ext-from-perfect-into-bounded-QCoh}.
Consid\'erons le foncteur
$$
- \otimes_E^\mathbf{L} K^\bullet :
D(E, \text{d}) \longrightarrow D(\mathcal{O}_X)
$$
d'Alg\`ebres diff\'erentielles gradu\'ees, lemme
\ref{dga-lemma-tensor-with-complex-derived}. Comme $K^\bullet$ est parfait, il
d\'efinit un objet compact de $D(\mathcal{O}_X)$; voir la proposition
\ref{proposition-compact-is-perfect}. Avec (\ref{equation-E-is-OK}), le lemme
d'Alg\`ebres diff\'erentielles gradu\'ees
\ref{dga-lemma-fully-faithful-in-compact-case} montre que le foncteur ci-dessus
est pleinement fid\`ele. Il admet pour adjoint \`a droite
$$
R\Hom(K^\bullet, - ) : D(\mathcal{O}_X) \longrightarrow D(E, \text{d})
$$
d'apr\`es Alg\`ebres diff\'erentielles gradu\'ees, lemme
\ref{dga-lemma-tensor-with-complex-hom-adjoint}; par les propri\'et\'es g\'en\'erales
des foncteurs adjoints, cet adjoint est un quasi-inverse \`a gauche. D'autre
part, le lemme \ref{lemma-tensor-with-QCoh-complex} donne un foncteur
$$
- \otimes_E^\mathbf{L} K^\bullet :
D(E, \text{d}) \longrightarrow D_\QCoh(\mathcal{O}_X),
$$
et, puisque $K^\bullet$ engendre $D_\QCoh(\mathcal{O}_X)$, le noyau de
l'adjoint restreint \`a $D_\QCoh(\mathcal{O}_X)$ est nul. Un argument formel
fournit alors l'\'equivalence cherch\'ee; voir Cat\'egories d\'eriv\'ees, lemme
\ref{derived-lemma-fully-faithful-adjoint-kernel-zero}.
\end{proof}

\begin{remark}[Variante avec support]
\label{remark-DQCoh-is-Ddga-with-support}
Soit $S$ un sch\'ema. Soit $X$ un espace alg\'ebrique quasi-compact et
quasi-s\'epar\'e sur $S$. Soit $T\subset|X|$ une partie ferm\'ee telle que
$|X|\setminus T$ soit quasi-compact. L'analogue du th\'eor\`eme
\ref{theorem-DQCoh-is-Ddga} est vrai pour
$D_{\QCoh,T}(\mathcal{O}_X)$. Cela r\'esulte exactement du m\^eme argument que
dans la d\'emonstration du th\'eor\`eme, en utilisant les lemmes
\ref{lemma-generator-with-support} et
\ref{lemma-compact-is-perfect-with-support}, ainsi qu'une variante avec
supports du lemme \ref{lemma-tensor-with-QCoh-complex}. Si ce r\'esultat devient
n\'ecessaire, nous en donnerons ici un \'enonc\'e pr\'ecis et une d\'emonstration
d\'etaill\'ee.
\end{remark}

\begin{remark}[Unicit\'e de l'alg\`ebre diff\'erentielle gradu\'ee]
\label{remark-independence-choice}
Soit $X$ un espace alg\'ebrique quasi-compact et quasi-s\'epar\'e sur un anneau
$R$. La construction de la d\'emonstration du th\'eor\`eme
\ref{theorem-DQCoh-is-Ddga} fournit une alg\`ebre diff\'erentielle gradu\'ee
$(A,\text{d})$ sur $R$ telle que $D_\QCoh(X)$ soit $R$-lin\'eairement
\'equivalente \`a $D(A,\text{d})$ comme cat\'egorie triangul\'ee. On peut se
demander dans quelle mesure $(A,\text{d})$ est unique. La r\'eponse est
seulement un peu plus pr\'ecise que l'affirmation suivant laquelle
$(A,\text{d})$ est d\'efinie \`a \'equivalence d\'eriv\'ee pr\`es. Supposons en
effet que $(B,\text{d})$ soit une seconde paire de ce type. On a alors
$$
(A, \text{d}) = \Hom_{\text{Comp}^{dg}(\mathcal{O}_X)}(K^\bullet, K^\bullet)
$$
et
$$
(B, \text{d}) = \Hom_{\text{Comp}^{dg}(\mathcal{O}_X)}(L^\bullet, L^\bullet)
$$
pour des complexes K-injectifs $K^\bullet$ et $L^\bullet$ de
$\mathcal{O}_X$-modules correspondant \`a des g\'en\'erateurs parfaits de
$D_\QCoh(\mathcal{O}_X)$. Posons
$$
\Omega = \Hom_{\text{Comp}^{dg}(\mathcal{O}_X)}(K^\bullet, L^\bullet)
\quad
\Omega' = \Hom_{\text{Comp}^{dg}(\mathcal{O}_X)}(L^\bullet, K^\bullet).
$$
Alors $\Omega$ est un $B^{opp}\otimes_RA$-module diff\'erentiel gradu\'e, et
$\Omega'$ un $A^{opp}\otimes_RB$-module diff\'erentiel gradu\'e. De plus,
l'\'equivalence
$$
D(A, \text{d}) \to D_\QCoh(\mathcal{O}_X) \to
D(B, \text{d})
$$
est donn\'ee par le foncteur $-\otimes_A^\mathbf{L}\Omega'$, et de m\^eme pour
le quasi-inverse. Nous sommes donc dans la situation d'Alg\`ebres
diff\'erentielles gradu\'ees, remarque
\ref{dga-remark-hochschild-cohomology}. Si cette remarque devient n\'ecessaire,
nous en donnerons ici un \'enonc\'e pr\'ecis et une d\'emonstration d\'etaill\'ee.
\end{remark}

\section{Caract\'erisation des complexes pseudo-coh\'erents, I}
\label{section-pseudo-coherent-hocolim}

\noindent
Cette \'etude se poursuivra dans Compl\'ements sur les morphismes d'espaces,
section
\ref{spaces-more-morphisms-section-characterize-pseudo-coherent}.
On peut caract\'eriser les objets pseudo-coh\'erents
comme des colimites homotopiques d'approximations par des objets parfaits.

\begin{lemma}
\label{lemma-pseudo-coherent-hocolim}
Soit $S$ un sch\'ema.
Soit $X$ un espace alg\'ebrique quasi-compact et quasi-s\'epar\'e sur $S$.
Soit $K \in D(\mathcal{O}_X)$. Les conditions suivantes sont \'equivalentes :
\begin{enumerate}
\item $K$ est pseudo-coh\'erent;
\item $K = \text{hocolim} K_n$, o\`u
$K_n$ est parfait et $\tau_{\geq -n}K_n \to \tau_{\geq -n}K$
est un isomorphisme pour tout $n$.
\end{enumerate}
\end{lemma}

\begin{proof}
L'implication (2) $\Rightarrow$ (1) est valable sur tout site annel\'e.
Supposons en effet que (2) soit satisfaite. Rappelons qu'un objet parfait de la
cat\'egorie d\'eriv\'ee est pseudo-coh\'erent; voir
Cohomologie sur les sites, lemme \ref{sites-cohomology-lemma-perfect}.
Il r\'esulte alors des d\'efinitions que
$\tau_{\geq -n}K_n$ est $(-n + 1)$-pseudo-coh\'erent,
donc que $\tau_{\geq -n}K$ est $(-n + 1)$-pseudo-coh\'erent,
et par suite que $K$ est $(-n + 1)$-pseudo-coh\'erent. Comme ceci vaut pour
tout $n$, le complexe $K$ est pseudo-coh\'erent; voir
Cohomologie sur les sites, d\'efinition
\ref{sites-cohomology-definition-pseudo-coherent}.

\medskip\noindent
Supposons (1). Commen\c{c}ons par choisir une approximation
$K_1 \to K$ de $(X, K, -2)$ par un complexe parfait $K_1$; voir les
d\'efinitions \ref{definition-approximation-holds} et
\ref{definition-approximation}, ainsi que le
th\'eor\`eme \ref{theorem-approximation}.
Supposons par r\'ecurrence que l'on ait
$$
K_1 \to K_2 \to \ldots \to K_n \to K
$$
avec les $K_i$ parfaits et tels que
$\tau_{\geq -i}K_i \to \tau_{\geq -i}K$ soit un isomorphisme
pour tout $1 \leq i \leq n$. Choisissons alors $a \leq b$ comme dans le
lemme \ref{lemma-ext-from-perfect-into-bounded-QCoh}
pour l'objet parfait $K_n$. Choisissons une approximation
$K_{n + 1} \to K$ de $(X, K, \min(a - 1, -n - 1))$.
Choisissons un triangle distingu\'e
$$
K_{n + 1} \to K \to C \to K_{n + 1}[1]
$$
On voit alors que $C \in D_\QCoh(\mathcal{O}_X)$ v\'erifie
$H^i(C) = 0$ pour $i \geq a$. Le choix de $a, b$ donne donc
$\Hom_{D(\mathcal{O}_X)}(K_n, C) = 0$.
Le compos\'e $K_n \to K \to C$ est par cons\'equent nul. D'apr\`es
Cat\'egories d\'eriv\'ees, lemme
\ref{derived-lemma-representable-homological},
le morphisme $K_n \to K$ se factorise par $K_{n + 1}$,
ce qui \'etablit l'h\'er\'edit\'e.

\medskip\noindent
Il reste \`a d\'emontrer que $K = \text{hocolim} K_n$.
Cela r\'esulte de l'application de
Cat\'egories d\'eriv\'ees, lemme \ref{derived-lemma-cohomology-of-hocolim},
aux foncteurs
$H^i( - ) : D(\mathcal{O}_X) \to \textit{Mod}(\mathcal{O}_X)$
et du choix des $K_n$.
\end{proof}

\begin{lemma}
\label{lemma-pseudo-coherent-hocolim-with-support}
Soit $X$ un sch\'ema quasi-compact et quasi-s\'epar\'e.
Soit $T \subset X$ une partie ferm\'ee telle que $X \setminus T$
soit quasi-compact. Soit $K \in D(\mathcal{O}_X)$ \`a support dans $T$.
Les conditions suivantes sont \'equivalentes :
\begin{enumerate}
\item $K$ est pseudo-coh\'erent;
\item $K = \text{hocolim} K_n$, o\`u
$K_n$ est parfait, \`a support dans $T$, et
$\tau_{\geq -n}K_n \to \tau_{\geq -n}K$ est un isomorphisme pour tout $n$.
\end{enumerate}
\end{lemma}

\begin{proof}
La d\'emonstration est exactement la m\^eme que celle du
lemme \ref{lemma-pseudo-coherent-hocolim},
si ce n'est que, pour choisir les approximations, on utilise
les triplets $(T, K, m)$.
\end{proof}

\section{Retour sur le cohérateur}
\label{section-better-coherator}

\noindent
Dans la section \ref{section-coherator}, nous avons construit et étudié
l'adjoint à droite $RQ_X$ du foncteur canonique
$D(\QCoh(\mathcal{O}_X)) \to D(\mathcal{O}_X)$.
On l'a construit comme le foncteur dérivé droit du cohérateur
$Q_X : \textit{Mod}(\mathcal{O}_X) \to \QCoh(\mathcal{O}_X)$.
Dans la présente section, nous étudions les conditions sous lesquelles le foncteur d'inclusion
$$
D_\QCoh(\mathcal{O}_X) \longrightarrow D(\mathcal{O}_X)
$$
admet un adjoint à droite. S'il existe, nous le
noterons\footnote{Cette notation n'est probablement pas standard. Nous avons toutefois déjà
utilisé $Q_X$ pour le cohérateur et $RQ_X$ pour son foncteur dérivé.}
$$
DQ_X :
D(\mathcal{O}_X) \longrightarrow D_\QCoh(\mathcal{O}_X)
$$
Il s'avère que les espaces algébriques quasi-compacts et quasi-séparés
admettent un tel adjoint à droite.

\begin{lemma}
\label{lemma-better-coherator}
Soit $S$ un schéma.
Soit $X$ un espace algébrique quasi-compact et quasi-séparé sur $S$.
Le foncteur d'inclusion $D_\QCoh(\mathcal{O}_X) \to D(\mathcal{O}_X)$
admet un adjoint à droite.
\end{lemma}

\begin{proof}[Première démonstration]
Nous allons le démontrer à l'aide du principe de récurrence du
lemme \ref{lemma-induction-principle}.
Si $D(\QCoh(\mathcal{O}_X)) \to D_\QCoh(\mathcal{O}_X)$
est une équivalence, le lemme est vrai, car le foncteur
$RQ_X$ de la section \ref{section-coherator} est adjoint à droite au foncteur
$D(\QCoh(\mathcal{O}_X)) \to D(\mathcal{O}_X)$.
En particulier, notre lemme est vrai pour les espaces algébriques affines ; voir le
lemme \ref{lemma-affine-coherator}.
Il suffit donc de montrer ceci : si $(U \subset X, f : V \to X)$
est un carré distingué élémentaire, avec $U$ quasi-compact
et $V$ affine, et si le lemme vaut pour $U$, $V$ et $U \times_X V$,
alors il vaut pour $X$.

\medskip\noindent
L'adjoint existe si et seulement si, pour tout objet $E$ de
$D(\mathcal{O}_X)$, on peut trouver un triangle distingué
$$
E' \to E \to K \to E'[1]
$$
dans $D(\mathcal{O}_X)$
tel que $E'$ appartienne à $D_\QCoh(\mathcal{O}_X)$ et que
$\Hom(M, K) = 0$ pour tout $M$ de $D_\QCoh(\mathcal{O}_X)$. Voir
Catégories dérivées, lemme \ref{derived-lemma-right-adjoint}.
Considérons le triangle distingué
$$
E \to Rj_{U, *}E|_U \oplus Rj_{V, *}E|_V \to
Rj_{U \times_X V, *}E|_{U \times_X V} \to E[1]
$$
dans $D(\mathcal{O}_X)$ du lemme \ref{lemma-exact-sequence-j-star}.
D'après Catégories dérivées, lemme \ref{derived-lemma-prepare-adjoint},
il suffit de construire les triangles distingués voulus
pour $Rj_{U, *}E|_U$, $Rj_{V, *}E|_V$ et
$Rj_{U \times_X V, *}E|_{U \times_X V}$. Nous sommes ainsi ramenés à l'assertion
étudiée au paragraphe suivant.

\medskip\noindent
Soit $j : U \to X$ un morphisme étale, où
$U$ est quasi-compact et quasi-séparé et où le lemme est vrai pour $U$.
Soit $L$ un objet de $D(\mathcal{O}_U)$.
Il existe alors un triangle distingué
$$
E' \to Rj_*L \to K \to E'[1]
$$
dans $D(\mathcal{O}_X)$
tel que $E'$ appartienne à $D_\QCoh(\mathcal{O}_X)$ et que
$\Hom(M, K) = 0$ pour tout $M$ de $D_\QCoh(\mathcal{O}_X)$.
Pour le voir, choisissons un triangle distingué
$$
L' \to L \to Q \to L'[1]
$$
dans $D(\mathcal{O}_U)$ tel que $L'$ appartienne à $D_\QCoh(\mathcal{O}_U)$
et que $\Hom(N, Q) = 0$ pour tout $N$ de $D_\QCoh(\mathcal{O}_U)$.
Cela est possible puisque l'énoncé de
Catégories dérivées, lemme \ref{derived-lemma-right-adjoint},
est une équivalence.
Nous obtenons un triangle distingué
$$
Rj_*L' \to Rj_*L \to Rj_*Q \to Rj_*L'[1]
$$
dans $D(\mathcal{O}_X)$. Notons que $Rj_*L'$ appartient à $D_\QCoh(\mathcal{O}_X)$
d'après le lemme \ref{lemma-quasi-coherence-direct-image}.
D'autre part, si $M$ appartient à $D_\QCoh(\mathcal{O}_X)$, alors
$$
\Hom(M, Rj_*Q) = \Hom(Lj^*M, Q) = 0
$$
car $Lj^*M$ appartient à $D_\QCoh(\mathcal{O}_U)$ d'après le
lemme \ref{lemma-quasi-coherence-pullback}.
Cela achève la démonstration.
\end{proof}

\begin{proof}[Seconde démonstration]
L'adjoint existe d'après Catégories dérivées, proposition
\ref{derived-proposition-brown}. Les hypothèses sont satisfaites :
d'abord, $D_\QCoh(\mathcal{O}_X)$ admet les sommes directes,
et le foncteur d'inclusion commute aux sommes directes
(lemme \ref{lemma-quasi-coherence-direct-sums}).
D'autre part, $D_\QCoh(\mathcal{O}_X)$
est compactement engendrée, puisqu'elle possède un
générateur parfait d'après le théorème \ref{theorem-bondal-van-den-Bergh},
et puisque les objets parfaits sont compacts d'après la
proposition \ref{proposition-compact-is-perfect}.
\end{proof}

\begin{lemma}
\label{lemma-pushforward-better-coherator}
Soit $S$ un schéma.
Soit $f : X \to Y$ un morphisme quasi-compact et quasi-séparé
d'espaces algébriques sur $S$.
Si les adjoints à droite $DQ_X$ et $DQ_Y$
des foncteurs d'inclusion $D_\QCoh \to D$ existent pour $X$ et $Y$, alors
$$
Rf_* \circ DQ_X = DQ_Y \circ Rf_*
$$
\end{lemma}

\begin{proof}
L'énoncé a un sens parce que $Rf_*$ envoie
$D_\QCoh(\mathcal{O}_X)$ dans $D_\QCoh(\mathcal{O}_Y)$ d'après le
lemme \ref{lemma-quasi-coherence-direct-image}.
Il est vrai parce que $Lf^*$ envoie de même
$D_\QCoh(\mathcal{O}_Y)$ dans $D_\QCoh(\mathcal{O}_X)$
(lemme \ref{lemma-quasi-coherence-pullback}) ;
ainsi, $Rf_* \circ DQ_X$ et $DQ_Y \circ Rf_*$ sont tous deux
adjoints à droite de $Lf^* : D_\QCoh(\mathcal{O}_Y) \to D(\mathcal{O}_X)$.
\end{proof}

\begin{remark}
\label{remark-explain-consequence}
Soit $S$ un schéma. Soit $(U \subset X, f : V \to X)$ un
carré distingué élémentaire d'espaces algébriques sur $S$.
Supposons $X$, $U$ et $V$ quasi-compacts et quasi-séparés.
D'après le lemme \ref{lemma-better-coherator}, les foncteurs
$DQ_X$, $DQ_U$, $DQ_V$, $DQ_{U \times_X V}$ existent. De plus, on a un
triangle distingué canonique
$$
DQ_X(K) \to Rj_{U, *}DQ_U(K|_U) \oplus Rj_{V, *}DQ_V(K|_V)
\to Rj_{U \times_X V, *}DQ_{U \times_X V}(K|_{U \times_X V})
\to DQ_X(K)[1]
$$
pour tout $K \in D(\mathcal{O}_X)$. Cela résulte de l'application du
foncteur exact $DQ_X$ au triangle distingué du
lemme \ref{lemma-exact-sequence-j-star}
et de trois applications du lemme \ref{lemma-pushforward-better-coherator}.
\end{remark}

\begin{lemma}
\label{lemma-boundedness-better-coherator}
Soit $S$ un schéma.
Soit $X$ un espace algébrique quasi-compact et quasi-séparé sur $S$.
Le foncteur $DQ_X$ du lemme \ref{lemma-better-coherator}
possède la propriété de bornitude suivante :
il existe un entier $N = N(X)$ tel que, si
$K$ appartient à $D(\mathcal{O}_X)$ et si
$H^i(U, K) = 0$ pour tout $U$ affine étale sur $X$ et tout $i \not \in [a, b]$, alors
les faisceaux de cohomologie $H^i(DQ_X(K))$ sont nuls pour
$i \not \in [a, b + N]$.
\end{lemma}

\begin{proof}
Nous allons le démontrer à l'aide du principe de récurrence du
lemme \ref{lemma-induction-principle}.

\medskip\noindent
Si $X$ est affine, le lemme est vrai avec $N = 0$, car
$RQ_X = DQ_X$ s'obtient alors en prenant le complexe de
faisceaux quasi-cohérents associé à $R\Gamma(X, K)$.
Voir le lemme \ref{lemma-affine-coherator}.

\medskip\noindent
Soit $(U \subset W, f : V \to W)$ un carré distingué élémentaire,
où $W$ est quasi-compact et quasi-séparé, $U \subset W$
est un ouvert quasi-compact et $V$ est affine, tel que
le lemme vaille pour $U$, $V$ et $U \times_W V$.
Notons $N(U)$, $N(V)$ et $N(U \times_W V)$ les entiers correspondants.
Supposons maintenant que $K$ appartienne à $D(\mathcal{O}_W)$ et que
$H^i(T, K) = 0$ pour tout $T$ affine étale sur $W$ et tout $i \not \in [a, b]$.
Alors $K|_U$, $K|_V$ et $K|_{U \times_W V}$ ont la même propriété.
Il en résulte que $DQ_U(K|_U)$, $DQ_V(K|_V)$ et
$DQ_{U \times_W V}(K|_{U \times_W V})$ ont leurs faisceaux de cohomologie nuls
hors de l'intervalle $[a, b + \max(N(U), N(V), N(U \times_W V))]$.
Puisque les foncteurs $Rj_{U, *}$, $Rj_{V, *}$ et $Rj_{U \times_W V, *}$
sont de dimension cohomologique finie sur $D_\QCoh$ d'après le
lemme \ref{lemma-quasi-coherence-direct-image},
il existe un entier $N$ tel que
$Rj_{U, *}DQ_U(K|_U)$, $Rj_{V, *}DQ_V(K|_V)$ et
$Rj_{U \times_W V, *}DQ_{U \times_W V}(K|_{U \times_W V})$ aient leurs
faisceaux de cohomologie nuls hors de l'intervalle $[a, b + N]$.
On conclut enfin au moyen du triangle distingué de la
remarque \ref{remark-explain-consequence}.
\end{proof}

\begin{example}
\label{example-inverse-limit-quasi-coherent}
Soit $S$ un schéma.
Soit $X$ un espace algébrique quasi-compact et quasi-séparé sur $S$.
Soit $(\mathcal{F}_n)$
un système projectif de faisceaux quasi-cohérents sur $X$.
Comme $DQ_X$ est adjoint à droite, il commute aux
produits, donc aux limites projectives dérivées.
Il en résulte que
$$
DQ_X(R\lim \mathcal{F}_n) =
(R\lim\text{ dans }D_\QCoh(\mathcal{O}_X))(\mathcal{F}_n)
$$
où la première $R\lim$ est calculée dans $D(\mathcal{O}_X)$.
En fait, posons $K = R\lim \mathcal{F}_n$.
Pour tout $U$ affine étale sur $X$, on a
$$
H^i(U, K) =
H^i(R\Gamma(U, R\lim \mathcal{F}_n)) =
H^i(R\lim R\Gamma(U, \mathcal{F}_n)) =
H^i(R\lim \Gamma(U, \mathcal{F}_n))
$$
En effet, la cohomologie commute aux limites projectives dérivées et les
faisceaux quasi-cohérents
$\mathcal{F}_n$ n'ont pas de cohomologie supérieure sur les espaces affines.
Le calcul de $R\lim$ dans la catégorie des
groupes abéliens montre que $H^i(U, K) = 0$
sauf si $i \in [0, 1]$. On en conclut que
le $R\lim$ dans $D_\QCoh(\mathcal{O}_X)$, qui est
$DQ_X(K)$ d'après ce qui précède, appartient à $D^b_\QCoh(\mathcal{O}_X)$
et que ses faisceaux de cohomologie sont nuls en degrés négatifs
d'après le lemme \ref{lemma-boundedness-better-coherator}.
\end{example}

\section{Cohomologie et changement de base, IV}
\label{section-cohomology-and-base-change-perfect}

\noindent
Cette section est l'analogue de Cat\'egories d\'eriv\'ees des sch\'emas,
section \ref{perfect-section-cohomology-and-base-change-perfect}.

\begin{lemma}
\label{lemma-cohomology-base-change}
Soit $S$ un sch\'ema. Soit $f:X\to Y$ un morphisme quasi-compact et
quasi-s\'epar\'e d'espaces alg\'ebriques sur $S$. Pour
$E\in D_\QCoh(\mathcal{O}_X)$ et $K\in D_\QCoh(\mathcal{O}_Y)$, on a
$$
Rf_*(E) \otimes_{\mathcal{O}_Y}^\mathbf{L} K =
Rf_*(E \otimes_{\mathcal{O}_X}^\mathbf{L} Lf^*K).
$$
\end{lemma}

\begin{proof}
Sans aucune hypoth\`ese, il existe un morphisme
$Rf_*(E)\otimes_{\mathcal{O}_Y}^\mathbf{L}K\to
Rf_*(E\otimes_{\mathcal{O}_X}^\mathbf{L}Lf^*K)$. C'est en effet l'adjoint du
morphisme canonique
$$
Lf^*(Rf_*(E) \otimes_{\mathcal{O}_Y}^\mathbf{L} K) =
Lf^*(Rf_*(E)) \otimes_{\mathcal{O}_X}^\mathbf{L} Lf^*K
\longrightarrow
E \otimes_{\mathcal{O}_X}^\mathbf{L} Lf^*K
$$
provenant du morphisme $Lf^*Rf_*E\to E$; voir Cohomologie sur les sites,
lemmes \ref{sites-cohomology-lemma-pullback-tensor-product} et
\ref{sites-cohomology-lemma-adjoint}. Pour v\'erifier qu'il s'agit d'un
isomorphisme, on peut travailler localement pour la topologie \'etale sur $Y$.
On se ram\`ene donc au cas o\`u $Y$ est un sch\'ema affine.

\medskip\noindent
Supposons que $K=\bigoplus K_i$ soit une somme directe de complexes
$K_i\in D_\QCoh(\mathcal{O}_Y)$. Si l'assertion est vraie pour chacun des
$K_i$, elle l'est pour $K$. En effet, les foncteurs $Lf^*$ et
$\otimes^\mathbf{L}$ commutent aux sommes directes par construction, et $Rf_*$
commute aux sommes directes pour les complexes \`a faisceaux de cohomologie
quasi-coh\'erents, d'apr\`es le lemme
\ref{lemma-quasi-coherence-pushforward-direct-sums}. De plus, si
$K\to L\to M\to K[1]$ est un triangle distingu\'e de $D_\QCoh(Y)$ et si
l'assertion du lemme est vraie pour deux des trois objets $K,L,M$, elle l'est
pour le troisi\`eme, car les foncteurs en jeu sont des foncteurs exacts de
cat\'egories triangul\'ees.

\medskip\noindent
Supposons $Y$ affine, disons $Y=\Spec(A)$. Le foncteur
$\widetilde{\ }:D(A)\to D_\QCoh(\mathcal{O}_Y)$ est une \'equivalence d'apr\`es
le lemme \ref{lemma-derived-quasi-coherent-small-etale-site} et Cat\'egories
d\'eriv\'ees des sch\'emas, lemme
\ref{perfect-lemma-affine-compare-bounded}. Notons $T$ la propri\'et\'e, pour
$K\in D(A)$, selon laquelle l'assertion du lemme est vraie pour
$\widetilde K$. La discussion pr\'ec\'edente et Compl\'ements d'alg\`ebre,
remarque \ref{more-algebra-remark-P-resolution}, montrent qu'il suffit de
prouver $T$ pour $A[k]$. Cela ach\`eve la d\'emonstration, puisque l'assertion
du lemme est claire pour les translat\'es du faisceau structural.
\end{proof}

\begin{definition}
\label{definition-tor-independent}
Soit $S$ un sch\'ema. Soit $B$ un espace alg\'ebrique sur $S$. Soient $X$ et
$Y$ des espaces alg\'ebriques sur $B$. On dit que $X$ et $Y$ sont
{\it Tor-ind\'ependants sur $B$} si, pour tout diagramme commutatif
$$
\xymatrix{
\Spec(k) \ar[d]_{\overline{y}} \ar[dr]_{\overline{b}} \ar[r]_-{\overline{x}} &
X \ar[d] \\
Y \ar[r] & B
}
$$
de points g\'eom\'etriques, les anneaux
$\mathcal{O}_{X,\overline{x}}$ et $\mathcal{O}_{Y,\overline{y}}$ sont
Tor-ind\'ependants sur $\mathcal{O}_{B,\overline{b}}$ (voir Compl\'ements
d'alg\`ebre, d\'efinition \ref{more-algebra-definition-tor-independent}).
\end{definition}

\noindent
Le lemme suivant montre en particulier que cette d\'efinition co\"incide avec
la d\'efinition donn\'ee pour les espaces alg\'ebriques repr\'esentables.

\begin{lemma}
\label{lemma-tor-independent}
Soit $S$ un sch\'ema. Soit $B$ un espace alg\'ebrique sur $S$. Soient $X$ et
$Y$ des espaces alg\'ebriques sur $B$. Les conditions suivantes sont
\'equivalentes :
\begin{enumerate}
\item $X$ et $Y$ sont Tor-ind\'ependants sur $B$;
\item pour tout diagramme commutatif
$$
\xymatrix{
U \ar[d] \ar[r] & W \ar[d] & V \ar[d] \ar[l] \\
X \ar[r] & B & Y \ar[l]
}
$$
dont les fl\`eches verticales sont \'etales, $U$ et $V$ sont
Tor-ind\'ependants sur $W$;
\item il existe un diagramme commutatif comme en (2) tel que (a) $W\to B$
soit \'etale surjectif, (b) $U\to X\times_BW$ soit \'etale surjectif,
(c) $V\to Y\times_BW$ soit \'etale surjectif, et tel que $U$ et $V$ soient
Tor-ind\'ependants sur $W$;
\item il existe un diagramme commutatif comme en (3), avec $U,V,W$ des
sch\'emas, tel que $U$ et $V$ soient Tor-ind\'ependants sur $W$ au sens de
Cat\'egories d\'eriv\'ees des sch\'emas, d\'efinition
\ref{perfect-definition-tor-independent}.
\end{enumerate}
\end{lemma}

\begin{proof}
Pour un morphisme \'etale $\varphi:U\to X$ d'espaces alg\'ebriques et un point
g\'eom\'etrique $\overline u$, le morphisme d'anneaux locaux
$\mathcal{O}_{X,\varphi(\overline u)}\to\mathcal{O}_{U,\overline u}$ est un
isomorphisme. On en d\'eduit l'\'equivalence de (1) et (2), ainsi que
l'implication (1) $\Rightarrow$ (3). Supposons (3) et choisissons un diagramme
de points g\'eom\'etriques comme dans la d\'efinition
\ref{definition-tor-independent}. Les hypoth\`eses permettent d'abord de
relever $\overline b$ en un point g\'eom\'etrique $\overline w$ de $W$, puis de
relever le point g\'eom\'etrique $(\overline x,\overline b)$ en un point
g\'eom\'etrique $\overline u$ de $U$, et enfin de relever
$(\overline y,\overline b)$ en un point g\'eom\'etrique $\overline v$ de $V$.
Pour trouver ces rel\`evements, on utilise Propri\'et\'es des espaces, lemme
\ref{spaces-properties-lemma-geometric-lift-to-usual}. La remarque pr\'ec\'edente
sur les anneaux locaux montre alors que la condition de la d\'efinition est
satisfaite pour le diagramme donn\'e.

\medskip\noindent
Ces points pr\'eliminaires \'etant fix\'es, l'assertion (4) revient 
manifestement \`a dire que la d\'efinition \ref{definition-tor-independent}
co\"incide avec Cat\'egories d\'eriv\'ees des sch\'emas, d\'efinition
\ref{perfect-definition-tor-independent}, lorsque $X,Y,B$ sont des sch\'emas.

\medskip\noindent
Soient $\overline x,\overline b,\overline y$ comme dans la d\'efinition
\ref{definition-tor-independent}, au-dessus des points $x,y,b$. Rappelons que
$\mathcal{O}_{X,\overline{x}}=\mathcal{O}_{X,x}^{sh}$ (Propri\'et\'es des
espaces, lemme \ref{spaces-properties-lemma-describe-etale-local-ring}), et de
m\^eme pour les deux autres. D'apr\`es Alg\`ebre, lemme
\ref{algebra-lemma-strictly-henselian-functorial-improve},
$\mathcal{O}_{X,\overline x}$ est un hens\'elis\'e strict de
$\mathcal{O}_{X,x}\otimes_{\mathcal{O}_{B,b}}\mathcal{O}_{B,\overline b}$.
En particulier, le morphisme d'anneaux
$$
\mathcal{O}_{X, x} \otimes_{\mathcal{O}_{B, b}} \mathcal{O}_{B, \overline{b}}
\longrightarrow
\mathcal{O}_{X, \overline{x}}
$$
est plat (Compl\'ements d'alg\`ebre, lemme
\ref{more-algebra-lemma-dumb-properties-henselization}). D'apr\`es Compl\'ements
d'alg\`ebre, lemme \ref{more-algebra-lemma-tor-independent-flat}, on a
$$
\text{Tor}_i^{\mathcal{O}_{B, b}}(\mathcal{O}_{X, x}, \mathcal{O}_{Y, y})
\otimes_{\mathcal{O}_{X, x} \otimes_{\mathcal{O}_{B, b}} \mathcal{O}_{Y, y}}
(\mathcal{O}_{X, \overline{x}} \otimes_{\mathcal{O}_{B, \overline{b}}}
\mathcal{O}_{Y, \overline y})
=
\text{Tor}_i^{\mathcal{O}_{B, \overline{b}}}(
\mathcal{O}_{X, \overline{x}}, \mathcal{O}_{Y, \overline{y}}).
$$
Par cons\'equent, si $X$ et $Y$ sont Tor-ind\'ependants sur $B$ comme
sch\'emas, ils le sont aussi comme espaces alg\'ebriques sur $B$.

\medskip\noindent
R\'eciproquement, on peut supposer $X,Y,B$ affines. Le morphisme d'anneaux
$$
\mathcal{O}_{X, x} \otimes_{\mathcal{O}_{B, b}} \mathcal{O}_{Y, y}
\longrightarrow
\mathcal{O}_{X, \overline{x}} \otimes_{\mathcal{O}_{B, \overline{b}}}
\mathcal{O}_{Y, \overline y}
$$
est plat d'apr\`es les observations pr\'ec\'edentes. De plus, l'image du
morphisme sur les spectres contient tous les id\'eaux premiers
$\mathfrak s\subset
\mathcal{O}_{X,x}\otimes_{\mathcal{O}_{B,b}}\mathcal{O}_{Y,y}$ situ\'es
au-dessus de $\mathfrak m_x$ et $\mathfrak m_y$. Cette remarque et la formule
pr\'ec\'edente pour les Tor montrent que, si $X$ et $Y$ sont
Tor-ind\'ependants sur $B$ comme espaces alg\'ebriques, alors
$$
\text{Tor}_i^{\mathcal{O}_{B, b}}
(\mathcal{O}_{X, x}, \mathcal{O}_{Y, y})_\mathfrak s = 0
$$
pour tout $i>0$ et tout $\mathfrak s$ comme ci-dessus. Le lemme de Compl\'ements
d'alg\`ebre \ref{more-algebra-lemma-tor-independent}, appliqu\'e aux morphismes
d'anneaux
$\Gamma(B,\mathcal{O}_B)\to\Gamma(X,\mathcal{O}_X)$ et
$\Gamma(B,\mathcal{O}_B)\to\Gamma(Y,\mathcal{O}_Y)$, montre alors que $X$ et
$Y$ sont Tor-ind\'ependants sur $B$.
\end{proof}

\begin{lemma}
\label{lemma-compare-base-change}
Soit $S$ un sch\'ema. Soit $g:Y'\to Y$ un morphisme d'espaces alg\'ebriques
sur $S$. Soit $f:X\to Y$ un morphisme quasi-compact et quasi-s\'epar\'e
d'espaces alg\'ebriques sur $S$. Consid\'erons le diagramme de changement de
base
$$
\xymatrix{
X' \ar[r]_{g'} \ar[d]_{f'} &
X \ar[d]^f \\
Y' \ar[r]^g &
Y
}
$$
Si $X$ et $Y'$ sont Tor-ind\'ependants sur $Y$, alors, pour tout
$E\in D_\QCoh(\mathcal{O}_X)$, on a
$Rf'_*L(g')^*E=Lg^*Rf_*E$.
\end{lemma}

\begin{proof}
Pour tout objet $E$ de $D(\mathcal{O}_X)$, la remarque de Cohomologie sur les
sites \ref{sites-cohomology-remark-base-change} fournit un morphisme canonique
de changement de base $Lg^*Rf_*E\to Rf'_*L(g')^*E$. Pour v\'erifier que c'est
un isomorphisme, on peut travailler localement pour la topologie \'etale sur
$Y'$. On peut donc supposer que $g:Y'\to Y$ est un morphisme de sch\'emas
affines. En particulier, $g$ est affine et il suffit de montrer que
$$
Rg_*Lg^*Rf_*E \to Rg_*Rf'_*L(g')^*E = Rf_*(Rg'_* L(g')^* E)
$$
est un isomorphisme; voir le lemme \ref{lemma-affine-morphism} (et utiliser les
lemmes \ref{lemma-quasi-coherence-pullback},
\ref{lemma-quasi-coherence-tensor-product} et
\ref{lemma-quasi-coherence-direct-image} pour voir que les objets
$Rf'_*L(g')^*E$ et $Lg^*Rf_*E$ ont des faisceaux de cohomologie
quasi-coh\'erents). Le morphisme $g'$ est lui aussi affine (Morphismes
d'espaces, lemme \ref{spaces-morphisms-lemma-base-change-affine}). D'apr\`es le
lemme \ref{lemma-affine-morphism-pull-push}, le morphisme devient
$$
Rf_*E \otimes_{\mathcal{O}_Y}^\mathbf{L} g_*\mathcal{O}_{Y'}
\longrightarrow
Rf_*(E \otimes_{\mathcal{O}_X}^\mathbf{L} g'_*\mathcal{O}_{X'}).
$$
On a $g'_*\mathcal{O}_{X'}=f^*g_*\mathcal{O}_{Y'}$. Le lemme
\ref{lemma-cohomology-base-change} ram\`ene donc la d\'emonstration \`a
$Lf^*g_*\mathcal{O}_{Y'}=f^*g_*\mathcal{O}_{Y'}$. Cette \'egalit\'e r\'esulte de
l'hypoth\`ese selon laquelle $X$ et $Y'$ sont Tor-ind\'ependants sur $Y$. Pour
la v\'erifier, on peut travailler localement pour la topologie \'etale sur $X$
et supposer aussi $X$ affine. \'Ecrivons $X=\Spec(A)$, $Y=\Spec(R)$ et
$Y'=\Spec(R')$. L'hypoth\`ese entra\^ine que $A$ et $R'$ sont
Tor-ind\'ependants sur $R$ (voir le lemme \ref{lemma-tor-independent} et
Compl\'ements d'alg\`ebre, lemme \ref{more-algebra-lemma-tor-independent}),
c'est-\`a-dire que $\text{Tor}_i^R(A,R')=0$ pour $i>0$. Autrement dit,
$A\otimes_R^\mathbf{L}R'=A\otimes_RR'$, ce qui signifie exactement que
$Lf^*g_*\mathcal{O}_{Y'}=f^*g_*\mathcal{O}_{Y'}$.
\end{proof}

\noindent
Le lemme suivant sera utilis\'e dans le chapitre consacr\'e aux complexes
dualisants.

\begin{lemma}
\label{lemma-affine-morphism-and-hom-out-of-perfect}
Soit $g:S'\to S$ un morphisme de sch\'emas affines. Consid\'erons un carr\'e
cart\'esien
$$
\xymatrix{
X' \ar[r]_{g'} \ar[d]_{f'} & X \ar[d]^f \\
S' \ar[r]^g & S
}
$$
d'espaces alg\'ebriques quasi-compacts et quasi-s\'epar\'es. Supposons $g$ et
$f$ Tor-ind\'ependants. \'Ecrivons $S=\Spec(R)$ et $S'=\Spec(R')$. Pour
$M,K\in D(\mathcal{O}_X)$, le morphisme canonique
$$
R\Hom_X(M, K) \otimes^\mathbf{L}_R R'
\longrightarrow
R\Hom_{X'}(L(g')^*M, L(g')^*K)
$$
de $D(R')$ est un isomorphisme dans les deux cas suivants :
\begin{enumerate}
\item $M\in D(\mathcal{O}_X)$ est parfait et $K\in D_\QCoh(X)$;
\item $M\in D(\mathcal{O}_X)$ est pseudo-coh\'erent,
$K\in D_\QCoh^+(X)$ et $R'$ est de Tor-dimension finie sur $R$.
\end{enumerate}
\end{lemma}

\begin{proof}
Il existe dans $D(\Gamma(X,\mathcal{O}_X))$ un morphisme canonique
$R\Hom_X(M,K)\to R\Hom_{X'}(L(g')^*M,L(g')^*K)$ entre complexes
d'homomorphismes globaux; voir Cohomologie sur les sites, section
\ref{sites-cohomology-section-global-RHom}. Par restriction des scalaires, on
peut le regarder comme un morphisme de $D(R)$. L'adjonction entre la
restriction et $-\otimes_R^\mathbf{L}R'$ donne alors le morphisme affich\'e dans
l'\'enonc\'e. Ce morphisme \'etant d\'efini, il suffit de montrer que c'est un
isomorphisme dans la cat\'egorie d\'eriv\'ee des groupes ab\'eliens.

\medskip\noindent
Le membre de droite est \'egal \`a
$$
R\Hom_X(M, R(g')_*L(g')^*K) =
R\Hom_X(M, K \otimes_{\mathcal{O}_X}^\mathbf{L} g'_*\mathcal{O}_{X'})
$$
d'apr\`es le lemme \ref{lemma-affine-morphism-pull-push}. Dans les deux cas,
le complexe $R\SheafHom(M,K)$ est un objet de $D_\QCoh(\mathcal{O}_X)$
d'apr\`es le lemme \ref{lemma-quasi-coherence-internal-hom}. Il existe un
morphisme naturel
$$
R\SheafHom(M, K) \otimes_{\mathcal{O}_X}^\mathbf{L} g'_*\mathcal{O}_{X'}
\longrightarrow
R\SheafHom(M, K \otimes_{\mathcal{O}_X}^\mathbf{L} g'_*\mathcal{O}_{X'})
$$
qui est un isomorphisme dans les deux cas, d'apr\`es le lemme
\ref{lemma-internal-hom-evaluate-tensor-isomorphism}. Pour voir que ce lemme
s'applique dans le cas (2), remarquons que
$g'_*\mathcal{O}_{X'}=Rg'_*\mathcal{O}_{X'}=
Lf^*g_*\mathcal{O}_{S'}$, la seconde \'egalit\'e r\'esultant du lemme
\ref{lemma-compare-base-change}. Les lemmes Cat\'egories d\'eriv\'ees des
sch\'emas \ref{perfect-lemma-tor-dimension-affine},
\ref{lemma-descend-tor-amplitude} et Cohomologie sur les sites
\ref{sites-cohomology-lemma-tor-amplitude-pullback} montrent alors que
$g'_*\mathcal{O}_{X'}$ est de Tor-dimension finie. Dans les deux cas, en
rempla\c{c}ant $K$ par $R\SheafHom(M,K)$, on se ram\`ene donc \`a montrer que
$$
R\Gamma(X, K) \otimes^\mathbf{L}_R R' \longrightarrow
R\Gamma(X, K \otimes^\mathbf{L}_{\mathcal{O}_X} g'_*\mathcal{O}_{X'})
$$
est un isomorphisme. Le membre de gauche est \'egal \`a
$R\Gamma(X',L(g')^*K)$ d'apr\`es le lemme
\ref{lemma-affine-morphism-pull-push}. Le r\'esultat d\'ecoule donc du lemme
\ref{lemma-compare-base-change}.
\end{proof}

\begin{remark}
\label{remark-multiplication-map}
Avec les notations du lemme
\ref{lemma-affine-morphism-and-hom-out-of-perfect}, le diagramme
$$
\xymatrix{
R\Hom_X(M, Rg'_*L) \otimes_R^\mathbf{L} R' \ar[r] \ar[d]_\mu &
R\Hom_{X'}(L(g')^*M, L(g')^*Rg'_*L) \ar[d]^a \\
R\Hom_X(M, R(g')_*L) \ar@{=}[r] &
R\Hom_{X'}(L(g')^*M, L)
}
$$
est commutatif; la fl\`eche horizontale sup\'erieure est le morphisme du lemme,
$\mu$ est le morphisme de multiplication, et $a$ provient du morphisme
d'adjonction $L(g')^*Rg'_*L\to L$. Pour tout $K'\in D(R')$, le morphisme de
multiplication est le morphisme d'adjonction
$K'\otimes_R^\mathbf{L}R'\to K'$.
\end{remark}

\begin{lemma}
\label{lemma-tor-independence-and-tor-amplitude}
Soit $S$ un sch\'ema. Consid\'erons un carr\'e cart\'esien d'espaces alg\'ebriques
$$
\xymatrix{
X' \ar[r]_{g'} \ar[d]_{f'} & X \ar[d]^f \\
Y' \ar[r]^g & Y
}
$$
sur $S$. Supposons $g$ et $f$ Tor-ind\'ependants.
\begin{enumerate}
\item Si $E\in D(\mathcal{O}_X)$ est de Tor-amplitude contenue dans $[a,b]$
comme complexe de $f^{-1}\mathcal{O}_Y$-modules, alors $L(g')^*E$ est de
Tor-amplitude contenue dans $[a,b]$ comme complexe de
$(f')^{-1}\mathcal{O}_{Y'}$-modules.
\item Si $\mathcal{G}$ est un $\mathcal{O}_X$-module plat sur $Y$, alors
$L(g')^*\mathcal{G}=(g')^*\mathcal{G}$.
\end{enumerate}
\end{lemma}

\begin{proof}
On peut calculer la Tor-dimension sur les fibres; voir Cohomologie sur les
sites, lemme \ref{sites-cohomology-lemma-tor-amplitude-stalk}, et Propri\'et\'es
des espaces, th\'eor\`eme \ref{spaces-properties-theorem-exactness-stalks}.
Si $\overline{x}'$ est un point g\'eom\'etrique de $X'$ d'image $\overline x$
dans $X$, alors
$$
(L(g')^*E)_{\overline{x}'} =
E_{\overline{x}}
\otimes_{\mathcal{O}_{X, \overline{x}}}^\mathbf{L}
\mathcal{O}_{X', \overline{x}'}.
$$
Soient $\overline y'$ dans $Y'$ et $\overline y$ dans $Y$ les images de
$\overline x'$ et $\overline x$. Comme $X$ et $Y'$ sont Tor-ind\'ependants sur
$Y$, Compl\'ements d'alg\`ebre, lemme
\ref{more-algebra-lemma-base-change-comparison}, montre que le membre de droite
de la formule pr\'ec\'edente est \'egal \`a
$E_{\overline{x}}
\otimes_{\mathcal{O}_{Y, \overline{y}}}^\mathbf{L}
\mathcal{O}_{Y', \overline{y}'}$
dans $D(\mathcal{O}_{Y',\overline y'})$. L'assertion (1) r\'esulte donc de
Compl\'ements d'alg\`ebre, lemme
\ref{more-algebra-lemma-pull-tor-amplitude}. Pour (2), il suffit d'observer que
la platitude de $\mathcal G$ \'equivaut \`a ce que $\mathcal G[0]$ soit de
Tor-amplitude contenue dans $[0,0]$, puis d'appliquer (1).
\end{proof}

\section{Cohomologie et changement de base, V}
\label{section-cohomology-base-change}

\noindent
Cette section est l'analogue de Catégories dérivées des schémas, section
\ref{perfect-section-cohomology-base-change}. Dans la
section \ref{section-cohomology-and-base-change-perfect},
nous avons vu qu'un théorème de changement de base vaut lorsque les morphismes sont Tor-indépendants.
Même dans le cas affine, il ne peut exister de théorème de changement de base sans une telle
condition ; voir
Compléments d'algèbre, section \ref{more-algebra-section-tor-independence}.
Dans cette section, nous cherchons quand on peut obtenir un résultat de changement de base
``complexe par complexe''.

\medskip\noindent
Pour mettre cela en œuvre, soit $S$ un schéma de base et
supposons donné un diagramme commutatif
$$
\xymatrix{
X' \ar[r]_{g'} \ar[d]_{f'} &
X \ar[d]^f \\
Y' \ar[r]^g &
Y
}
$$
d'espaces algébriques sur $S$ (nous supposerons généralement qu'il est cartésien).
Soit $K \in D_\QCoh(\mathcal{O}_X)$
et soit $L(g')^*K \to K'$ un morphisme dans $D_\QCoh(\mathcal{O}_{X'})$.
Pour un point géométrique $\overline{x}'$ de $X'$, considérons les points
géométriques
$\overline{x} = g'(\overline{x}')$,
$\overline{y}' = f'(\overline{x}')$,
$\overline{y} = f(\overline{x}) = g(\overline{y}')$
de $X$, $Y'$ et $Y$.
Nous pouvons alors considérer les morphismes
$$
K_{\overline{x}}
\otimes_{\mathcal{O}_{Y, \overline{y}}}^\mathbf{L}
\mathcal{O}_{Y', \overline{y}'}
\to
K_{\overline{x}}
\otimes_{\mathcal{O}_{X, \overline{x}}}^\mathbf{L}
\mathcal{O}_{X', \overline{x}'} \to
K'_{\overline{x}'}
$$
où le premier morphisme est celui de Compléments d'algèbre,
équation (\ref{more-algebra-equation-comparison-map}),
et le second provient de
$(L(g')^*K)_{\overline{x}'} =
K_{\overline{x}} \otimes_{\mathcal{O}_{X, \overline{x}}}^\mathbf{L}
\mathcal{O}_{X', \overline{x}'}$
et du morphisme donné $L(g')^*K \to K'$.
Pour tout $i \in \mathbf{Z}$, on obtient une structure de
$\mathcal{O}_{X, \overline{x}} \otimes_{\mathcal{O}_{Y, \overline{y}}}
\mathcal{O}_{Y', \overline{y}'}$-module sur
$H^i(K_{\overline{x}}
\otimes_{\mathcal{O}_{Y, \overline{y}}}^\mathbf{L}
\mathcal{O}_{Y', \overline{y}'})$.
En regroupant ces données, on obtient des morphismes canoniques
\begin{equation}
\label{equation-bc}
H^i(K_{\overline{x}}
\otimes_{\mathcal{O}_{Y, \overline{y}}}^\mathbf{L}
\mathcal{O}_{Y', \overline{y}'})
\otimes_{(\mathcal{O}_{X, \overline{x}}
\otimes_{\mathcal{O}_{Y, \overline{y}}}
\mathcal{O}_{Y', \overline{y}'})}
\mathcal{O}_{X', \overline{x}'}
\longrightarrow
H^i(K'_{\overline{x}'})
\end{equation}
de $\mathcal{O}_{X', \overline{x}'}$-modules.

\begin{lemma}
\label{lemma-single-complex-base-change-condition}
Soit $S$ un schéma. Soit
$$
\xymatrix{
X' \ar[r]_{g'} \ar[d]_{f'} &
X \ar[d]^f \\
Y' \ar[r]^g &
Y
}
$$
un diagramme cartésien d'espaces algébriques sur $S$.
Soit $K \in D_\QCoh(\mathcal{O}_X)$ et soit $L(g')^*K \to K'$
un morphisme dans $D_\QCoh(\mathcal{O}_{X'})$. Les conditions suivantes sont équivalentes :
\begin{enumerate}
\item pour tout $x' \in X'$ et tout $i \in \mathbf{Z}$, le morphisme (\ref{equation-bc})
est un isomorphisme ;
\item pour tout diagramme commutatif
$$
\xymatrix{
& U \ar[d] \ar[rd]^a \\
V' \ar[r] \ar[rd]^c & V \ar[rd]^b & X \ar[d]^f \\
& Y' \ar[r]^g & Y
}
$$
où $a, b, c$ sont étales, où $U, V, V'$ sont des schémas et où $U' = V' \times_V U$,
les conditions équivalentes de
Catégories dérivées des schémas, lemme
\ref{perfect-lemma-single-complex-base-change-condition},
sont satisfaites pour $(U \to X)^*K$ et $(U' \to X')^*K'$ ;
\item il existe un diagramme comme en (2) tel que $U' \to X'$ soit surjectif.
\end{enumerate}
\end{lemma}

\begin{proof}
Notons que la condition (1) est locale pour la topologie étale sur $X'$.
En explicitant les implications formelles des résultats connus, on voit qu'il suffit de prouver
que la condition (1) du présent lemme équivaut à la condition
(1) de Catégories dérivées des schémas, lemme
\ref{perfect-lemma-single-complex-base-change-condition},
lorsque $X, Y, Y', X'$ sont représentables par des schémas
$X_0, Y_0, Y'_0, X'_0$. Notons $f_0, g_0, g'_0, f'_0$ les
morphismes entre ces schémas qui correspondent à $f, g, g', f'$.
On peut supposer que $K = \epsilon^*K_0$ et $K' = \epsilon^*K'_0$
pour certains objets $K_0 \in D_\QCoh(\mathcal{O}_{X_0})$
et $K'_0 \in D_\QCoh(\mathcal{O}_{X'_0})$ ; voir le
lemme \ref{lemma-derived-quasi-coherent-small-etale-site}.
De plus, le morphisme $L(g')^*K \to K'$ est l'image inverse
d'un morphisme $L(g_0)^*K_0 \to K'_0$, avec les notations de la
remarque \ref{remark-match-total-direct-images}.
Rappelons que $\mathcal{O}_{X, \overline{x}}$ est l'hensélisé
strict de $\mathcal{O}_{X, x}$
(Propriétés des espaces, lemme
\ref{spaces-properties-lemma-describe-etale-local-ring})
et que l'on a
$$
K_{\overline{x}} =
K_{0, x} \otimes_{\mathcal{O}_{X, x}}^\mathbf{L} \mathcal{O}_{X, \overline{x}}
\quad\text{et}\quad
K'_{\overline{x}'} =
K'_{0, x'} \otimes_{\mathcal{O}_{X', x'}}^\mathbf{L}
\mathcal{O}_{X', \overline{x}'}
$$
(de façon analogue à Propriétés des espaces, lemme
\ref{spaces-properties-lemma-stalk-quasi-coherent}).
Considérons le diagramme commutatif
$$
\xymatrix{
H^i(K_{\overline{x}}
\otimes_{\mathcal{O}_{Y, \overline{y}}}^\mathbf{L}
\mathcal{O}_{Y', \overline{y}'})
\otimes_{(\mathcal{O}_{X, \overline{x}}
\otimes_{\mathcal{O}_{Y, \overline{y}}}
\mathcal{O}_{Y', \overline{y}'})}
\mathcal{O}_{X', \overline{x}'}
\ar[r] &
H^i(K'_{\overline{x}'}) \\
H^i(K_{0, x} \otimes_{\mathcal{O}_{Y, y}}^\mathbf{L} \mathcal{O}_{Y', y'})
\otimes_{(\mathcal{O}_{X, x} \otimes_{\mathcal{O}_{Y, y}} \mathcal{O}_{Y', y'})}
\mathcal{O}_{X', x'}
\ar[u] \ar[r] &
H^i(K'_{0, x'}) \ar[u]
}
$$
Il faut montrer que le morphisme horizontal inférieur est un isomorphisme si et seulement
si le morphisme horizontal supérieur en est un. Comme
$\mathcal{O}_{X', x'} \to \mathcal{O}_{X', \overline{x}'}$
est fidèlement plat (Compléments d'algèbre, lemme
\ref{more-algebra-lemma-dumb-properties-henselization}),
il suffit de montrer que le morphisme supérieur s'obtient du morphisme
inférieur par le changement de base donné par ce morphisme.
Cela résulte immédiatement des relations entre les fibres
données plus haut pour les objets de droite.
Pour les objets de gauche, il suffit de montrer que
\begin{align*}
&
H^i\left(
(K_{0, x}
\otimes_{\mathcal{O}_{X, x}}^\mathbf{L} \mathcal{O}_{X, \overline{x}})
\otimes_{\mathcal{O}_{Y, \overline{y}}}^\mathbf{L}
\mathcal{O}_{Y', \overline{y}'}\right) \\
& =
H^i(K_{0, x} \otimes_{\mathcal{O}_{Y, y}}^\mathbf{L} \mathcal{O}_{Y', y'})
\otimes_{(\mathcal{O}_{X, x} \otimes_{\mathcal{O}_{Y, y}} \mathcal{O}_{Y', y'})}
(\mathcal{O}_{X, \overline{x}}
\otimes_{\mathcal{O}_{Y, \overline{y}}}
\mathcal{O}_{Y', \overline{y}'})
\end{align*}
Cela résulte de Compléments d'algèbre, lemme
\ref{more-algebra-lemma-lemma-tor-independent-flat-compare}.
Les hypothèses de platitude de ce lemme découlent de ce qui précède
ainsi que d'Algèbre, lemme
\ref{algebra-lemma-strictly-henselian-functorial-improve},
qui entraîne que
$\mathcal{O}_{X, \overline{x}}$ est l'hensélisé strict de
$\mathcal{O}_{X, x} \otimes_{\mathcal{O}_{Y, y}}
\mathcal{O}_{Y, \overline{y}}$
et que
$\mathcal{O}_{Y', \overline{y}'}$ est l'hensélisé strict de
$\mathcal{O}_{Y', y'} \otimes_{\mathcal{O}_{Y, y}}
\mathcal{O}_{Y, \overline{y}}$.
\end{proof}

\begin{lemma}
\label{lemma-single-complex-base-change}
Soit $S$ un schéma. Soit
$$
\xymatrix{
X' \ar[r]_{g'} \ar[d]_{f'} &
X \ar[d]^f \\
Y' \ar[r]^g &
Y
}
$$
un diagramme cartésien d'espaces algébriques sur $S$.
Soit $K \in D_\QCoh(\mathcal{O}_X)$ et soit $L(g')^*K \to K'$
un morphisme dans $D_\QCoh(\mathcal{O}_{X'})$. Si
\begin{enumerate}
\item les conditions équivalentes du
lemme \ref{lemma-single-complex-base-change-condition} sont satisfaites ; et
\item $f$ est quasi-compact et quasi-séparé,
\end{enumerate}
alors le composé $Lg^*Rf_*K \to Rf'_*L(g')^*K \to Rf'_*K'$
est un isomorphisme.
\end{lemma}

\begin{proof}
Pour vérifier que le morphisme est un isomorphisme, on peut raisonner localement pour la topologie étale sur $Y'$.
On peut donc supposer que $g : Y' \to Y$ est un morphisme de schémas affines.
Dans ce cas, nous utiliserons le principe de récurrence du
lemme \ref{lemma-induction-principle}
pour montrer que, pour tout espace algébrique quasi-compact et quasi-séparé
$U$ étale sur $X$,
le morphisme construit de même
$Lg^*R(U \to Y)_*K|_U \to R(U' \to Y')_*K'|_{U'}$
est un isomorphisme. Ici, $U' = X' \times_{g', X} U = Y' \times_{g, Y} U$.

\medskip\noindent
Si $U$ est un schéma (par exemple s'il est affine), le résultat est vrai.
En effet, $Y, Y', U, U'$ sont alors des schémas, et $K$ et $K'$ proviennent
d'objets des catégories dérivées des schémas sous-jacents d'après le
lemme \ref{lemma-derived-quasi-coherent-small-etale-site},
et la condition de
Catégories dérivées des schémas,
lemme \ref{perfect-lemma-single-complex-base-change-condition},
est satisfaite pour ces complexes d'après le
lemme \ref{lemma-single-complex-base-change-condition}.
Ainsi, en vertu des compatibilités expliquées dans la
remarque \ref{remark-match-total-direct-images},
on peut appliquer le résultat dans le cas des schémas, à savoir
Catégories dérivées des schémas, lemme
\ref{perfect-lemma-single-complex-base-change}.

\medskip\noindent
Passons à l'étape de récurrence. Soit $(U \subset W, V \to W)$ un carré
distingué élémentaire, où $W$ est un espace algébrique quasi-compact
et quasi-séparé, étale sur $X$, où $U$ est quasi-compact, où $V$ est affine,
et pour lequel le résultat vaut pour $U$, $V$ et $U \times_W V$.
Pour simplifier les notations, remplaçons $W$ par $X$ (ce qui est permis à ce stade).
Notons $a : U \to Y$, $b : V \to Y$ et $c : U \times_X V \to Y$
les morphismes évidents. Soient $a' : U' \to Y'$, $b' : V' \to Y'$
et $c' : U' \times_{X'} V' \to Y'$ les changements de base de $a$, $b$ et $c$.
À l'aide des triangles distingués de Mayer--Vietoris relatif
(lemme \ref{lemma-unbounded-relative-mayer-vietoris}),
on obtient un diagramme commutatif
$$
\xymatrix{
Lg^*Rf_*K \ar[r] \ar[d] & Rf'_*K' \ar[d] \\
Lg^*Ra_*K|_U \oplus Lg^*Rb_*K|_V \ar[r] \ar[d] &
Ra'_* K'|_{U'} \oplus Rb'_* K'|_{V'} \ar[d] \\
Lg^*Rc_*K|_{U \times_X V} \ar[r] \ar[d] &
Rc'_*K'|_{U' \times_{X'} V'} \ar[d] \\
Lg^*Rf_* K[1] \ar[r] &
Rf'_* K'[1]
}
$$
Comme les deuxième et troisième morphismes horizontaux sont des isomorphismes, le premier
l'est aussi (Catégories dérivées, lemme \ref{derived-lemma-third-isomorphism-triangle}),
ce qui achève la démonstration du lemme.
\end{proof}

\begin{lemma}
\label{lemma-single-complex-base-change-condition-inherited}
Soit $S$ un schéma. Soit
$$
\xymatrix{
X' \ar[r]_{g'} \ar[d]_{f'} &
X \ar[d]^f \\
S' \ar[r]^g &
S
}
$$
un diagramme cartésien d'espaces algébriques sur $S$.
Soit $K \in D_\QCoh(\mathcal{O}_X)$ et soit $L(g')^*K \to K'$
un morphisme dans $D_\QCoh(\mathcal{O}_{X'})$. Si les conditions équivalentes du
lemme \ref{lemma-single-complex-base-change-condition} sont satisfaites, alors :
\begin{enumerate}
\item pour $E \in D_\QCoh(\mathcal{O}_X)$, les conditions
équivalentes du lemme \ref{lemma-single-complex-base-change-condition} sont satisfaites
pour $L(g')^*(E \otimes^\mathbf{L} K) \to L(g')^*E \otimes^\mathbf{L} K'$ ;
\item si $E$ est un objet parfait de $D(\mathcal{O}_X)$, les conditions équivalentes du
lemme \ref{lemma-single-complex-base-change-condition} sont satisfaites pour
$L(g')^*R\SheafHom(E, K) \to R\SheafHom(L(g')^*E, K')$ ;
\item si $K$ est borné inférieurement et si $E$ est un objet pseudo-cohérent de
$D(\mathcal{O}_X)$, les conditions équivalentes du
lemme \ref{lemma-single-complex-base-change-condition} sont satisfaites pour
$L(g')^*R\SheafHom(E, K) \to R\SheafHom(L(g')^*E, K')$.
\end{enumerate}
\end{lemma}

\begin{proof}
L'énoncé a un sens, car les complexes en jeu ont des faisceaux de cohomologie
quasi-cohérents d'après les lemmes
\ref{lemma-quasi-coherence-pullback},
\ref{lemma-quasi-coherence-tensor-product} et
\ref{lemma-quasi-coherence-internal-hom}, ainsi que
Cohomologie sur les sites, lemmes
\ref{sites-cohomology-lemma-pseudo-coherent-pullback} et
\ref{sites-cohomology-lemma-perfect-pullback}.
Cela étant dit, dans le cas (1), on peut vérifier que les morphismes (\ref{equation-bc})
sont des isomorphismes en calculant la source et le but
de (\ref{equation-bc}) à l'aide de l'associativité du produit tensoriel ; voir
Compléments d'algèbre, lemme \ref{more-algebra-lemma-triple-tensor-product}.
Le morphisme des cas (2) et (3) est le composé
$$
L(g')^*R\SheafHom(E, K) \to R\SheafHom(L(g')^*E, L(g')^*K)
\to R\SheafHom(L(g')^*E, K')
$$
où le premier morphisme est celui de
Cohomologie sur les sites, remarque
\ref{sites-cohomology-remark-prepare-fancy-base-change},
et où le second provient du morphisme donné $L(g')^*K \to K'$.
Pour prouver que les morphismes (\ref{equation-bc}) sont des isomorphismes, on représente
$E_x$ par un complexe borné de $\mathcal{O}_{X, x}$-modules projectifs de type fini
dans le cas (2), ou par un complexe borné supérieurement de modules libres de type fini dans le cas (3),
puis on calcule la source et le but du morphisme.
Nous omettons quelques détails.
\end{proof}

\begin{lemma}
\label{lemma-base-change-tensor}
Soit $S$ un schéma. Soit $f : X \to Y$ un morphisme quasi-compact et
quasi-séparé d'espaces algébriques sur $S$.
Soit $E \in D_\QCoh(\mathcal{O}_X)$.
Soit $\mathcal{G}^\bullet$ un complexe borné supérieurement de
$\mathcal{O}_X$-modules quasi-cohérents plats sur $Y$.
Alors la formation de
$$
Rf_*(E \otimes^\mathbf{L}_{\mathcal{O}_X} \mathcal{G}^\bullet)
$$
commute à tout changement de base (voir la démonstration pour l'énoncé précis).
\end{lemma}

\begin{proof}
L'énoncé signifie ce qui suit. Soit $g : Y' \to Y$ un morphisme d'espaces
algébriques et considérons le diagramme de changement de base
$$
\xymatrix{
X' \ar[r]_{g'} \ar[d]_{f'} &
X \ar[d]^f \\
Y' \ar[r]^g &
Y
}
$$
autrement dit, $X' = Y' \times_Y X$. Le lemme affirme que
$$
Lg^*Rf_*(E \otimes^\mathbf{L}_{\mathcal{O}_X} \mathcal{G}^\bullet)
\longrightarrow
Rf'_*(L(g')^*E \otimes^\mathbf{L}_{\mathcal{O}_{X'}} (g')^*\mathcal{G}^\bullet)
$$
est un isomorphisme. Notons qu'au second membre nous n'utilisons {\bf pas}
l'image inverse dérivée de $\mathcal{G}^\bullet$.
Pour le démontrer, appliquons les lemmes \ref{lemma-single-complex-base-change} et
\ref{lemma-single-complex-base-change-condition-inherited} : il
suffit de prouver que le morphisme canonique
$$
L(g')^*\mathcal{G}^\bullet \to (g')^*\mathcal{G}^\bullet
$$
satisfait aux conditions équivalentes du
lemme \ref{lemma-single-complex-base-change-condition}.
Cela résulte de la vérification de la condition sur les fibres, où le résultat
découle immédiatement du fait que
$\mathcal{G}^\bullet_{\overline{x}}
\otimes_{\mathcal{O}_{Y, \overline{y}}}
\mathcal{O}_{Y', \overline{y}'}$
calcule le produit tensoriel dérivé en vertu de nos hypothèses sur le complexe
$\mathcal{G}^\bullet$.
\end{proof}

\begin{lemma}
\label{lemma-base-change-RHom}
Soit $S$ un schéma. Soit $f : X \to Y$ un morphisme quasi-compact et
quasi-séparé d'espaces algébriques sur $S$. Soit $E$
un objet de $D(\mathcal{O}_X)$. Soit $\mathcal{G}^\bullet$
un complexe de $\mathcal{O}_X$-modules quasi-cohérents. Si
\begin{enumerate}
\item $E$ est parfait, $\mathcal{G}^\bullet$ est borné supérieurement
et $\mathcal{G}^n$ est plat sur $Y$ ; ou si
\item $E$ est pseudo-cohérent, $\mathcal{G}^\bullet$ est borné
et $\mathcal{G}^n$ est plat sur $Y$,
\end{enumerate}
alors la formation de
$$
Rf_*R\SheafHom(E, \mathcal{G}^\bullet)
$$
commute à tout changement de base (voir la démonstration pour l'énoncé précis).
\end{lemma}

\begin{proof}
L'énoncé signifie ce qui suit. Soit $g : Y' \to Y$ un morphisme d'espaces
algébriques et considérons le diagramme de changement de base
$$
\xymatrix{
X' \ar[r]_{g'} \ar[d]_{f'} &
X \ar[d]^f \\
Y' \ar[r]^g &
Y
}
$$
autrement dit, $X' = Y' \times_Y X$. Le lemme affirme que
$$
Lg^*Rf_*R\SheafHom(E, \mathcal{G}^\bullet)
\longrightarrow
R(f')_*R\SheafHom(L(g')^*E, (g')^*\mathcal{G}^\bullet)
$$
est un isomorphisme. Notons qu'au second membre nous n'utilisons {\bf pas}
l'image inverse dérivée de $\mathcal{G}^\bullet$. Pour le démontrer, appliquons les
lemmes \ref{lemma-single-complex-base-change} et
\ref{lemma-single-complex-base-change-condition-inherited} : il
suffit de prouver que le morphisme canonique
$$
L(g')^*\mathcal{G}^\bullet \to (g')^*\mathcal{G}^\bullet
$$
satisfait aux conditions équivalentes du
lemme \ref{lemma-single-complex-base-change-condition}.
Cela a été montré dans la démonstration du lemme \ref{lemma-base-change-tensor}.
\end{proof}

\section{Construction de complexes parfaits}
\label{section-producing-perfect}

\noindent
Le lemme suivant est notre principal outil technique pour construire des
complexes parfaits. Des versions ult\'erieures de ce r\'esultat s'y ram\`eneront
par approximation noeth\'erienne.

\begin{lemma}
\label{lemma-perfect-direct-image}
Soit $S$ un sch\'ema. Soit $Y$ un espace alg\'ebrique noeth\'erien sur $S$.
Soit $f:X\to Y$ un morphisme d'espaces alg\'ebriques localement de type fini
et quasi-s\'epar\'e. Soit $E\in D(\mathcal{O}_X)$ tel que
\begin{enumerate}
\item $E\in D^b_{\textit{Coh}}(\mathcal{O}_X)$ ;
\item le support de $H^i(E)$ est propre sur $Y$ pour tout $i$ ;
\item $E$ est de Tor-dimension finie comme objet de
$D(f^{-1}\mathcal{O}_Y)$.
\end{enumerate}
Alors $Rf_*E$ est un objet parfait de $D(\mathcal{O}_Y)$.
\end{lemma}

\begin{proof}
D'apr\`es le lemme \ref{lemma-direct-image-coherent}, $Rf_*E$ est un objet de
$D^b_{\textit{Coh}}(\mathcal{O}_Y)$. Il est donc pseudo-coh\'erent
(lemme \ref{lemma-identify-pseudo-coherent-noetherian}). Il suffit par
cons\'equent de montrer que $Rf_*E$ est de Tor-dimension finie ; voir
Cohomologie sur les sites, lemme \ref{sites-cohomology-lemma-perfect}.
D'apr\`es le lemme \ref{lemma-tor-qc-qs}, il suffit de v\'erifier que
$Rf_*(E)\otimes_{\mathcal{O}_Y}^\mathbf{L}\mathcal{F}$ a une cohomologie
uniform\'ement born\'ee pour tout faisceau quasi-coh\'erent $\mathcal{F}$ sur
$Y$. La propri\'et\'e d'\^etre born\'e sup\'erieurement est imm\'ediate, puisque
$Rf_*(E)$ l'est. Soit $T\subset |X|$ la r\'eunion des supports de $H^i(E)$ pour
tout $i$. Les hypoth\`eses (1) et (2), ainsi que le lemme
\ref{lemma-union-closed-proper-over-base}, montrent que $T$ est propre
sur $Y$. Il existe donc un sous-espace ouvert quasi-compact
$X'\subset X$ contenant $T$. En posant $f'=f|_{X'}$, on a
$Rf_*(E)=Rf'_*(E|_{X'})$, car la restriction de $E$ \`a $X\setminus T$ est
nulle. On peut ainsi remplacer $X$ par $X'$ et supposer $f$ quasi-compact.
Nous avons suppos\'e $f$ quasi-s\'epar\'e. Par cons\'equent,
$$
Rf_*(E) \otimes_{\mathcal{O}_Y}^\mathbf{L} \mathcal{F} =
Rf_*\left(E \otimes_{\mathcal{O}_X}^\mathbf{L} Lf^*\mathcal{F}\right) =
Rf_*\left(E \otimes_{f^{-1}\mathcal{O}_Y}^\mathbf{L} f^{-1}\mathcal{F}\right)
$$
d'apr\`es le lemme \ref{lemma-cohomology-base-change} et Cohomologie sur les
sites, lemme \ref{sites-cohomology-lemma-variant-derived-pullback}.
Par l'hypoth\`ese (3), le complexe
$E\otimes_{f^{-1}\mathcal{O}_Y}^\mathbf{L}f^{-1}\mathcal{F}$ a ses faisceaux
de cohomologie dans un intervalle fini fix\'e, disons $[a,b]$. Son image par
$Rf_*$ a donc sa cohomologie dans $[a,\infty)$, ce qui conclut.
\end{proof}

\begin{lemma}
\label{lemma-tensor-perfect}
Soit $S$ un sch\'ema. Soit $B$ un espace alg\'ebrique noeth\'erien sur $S$.
Soit $f:X\to B$ un morphisme d'espaces alg\'ebriques localement de type fini
et quasi-s\'epar\'e. Soit $E\in D(\mathcal{O}_X)$ parfait. Soit
$\mathcal{G}^\bullet$ un complexe born\'e de $\mathcal{O}_X$-modules coh\'erents,
plats sur $B$ et dont le support est propre sur $B$. Alors
$K=Rf_*(E\otimes^\mathbf{L}_{\mathcal{O}_X}\mathcal{G}^\bullet)$ est un objet
parfait de $D(\mathcal{O}_B)$.
\end{lemma}

\begin{proof}
L'objet $K$ est parfait d'apr\`es le lemme \ref{lemma-perfect-direct-image}.
V\'erifions que les hypoth\`eses de ce lemme sont satisfaites. Localement, $E$
est isomorphe \`a un complexe fini de $\mathcal{O}_X$-modules libres de type
fini. Ainsi, localement,
$E\otimes^\mathbf{L}_{\mathcal{O}_X}\mathcal{G}^\bullet$ est isomorphe \`a un
complexe fini dont les termes sont de la forme
$$
\bigoplus\nolimits_{i = a, \ldots, b} (\mathcal{G}^i)^{\oplus r_i}
$$
pour certains entiers $a,b,r_a,\ldots,r_b$. Il en r\'esulte imm\'ediatement que
les faisceaux de cohomologie
$H^i(E\otimes^\mathbf{L}_{\mathcal{O}_X}\mathcal{G})$ sont coh\'erents.
L'hypoth\`ese sur la Tor-dimension en d\'ecoule elle aussi, puisque
$\mathcal{G}^i$ est plat sur $f^{-1}\mathcal{O}_B$.
\end{proof}

\begin{lemma}
\label{lemma-ext-perfect}
Soit $S$ un sch\'ema. Soit $B$ un espace alg\'ebrique noeth\'erien sur $S$.
Soit $f:X\to B$ un morphisme d'espaces alg\'ebriques localement de type fini
et quasi-s\'epar\'e. Soit $E\in D(\mathcal{O}_X)$ parfait. Soit
$\mathcal{G}^\bullet$ un complexe born\'e de $\mathcal{O}_X$-modules coh\'erents,
plats sur $B$ et dont le support est propre sur $B$. Alors
$K=Rf_*R\SheafHom(E,\mathcal{G})$ est un objet parfait de
$D(\mathcal{O}_B)$.
\end{lemma}

\begin{proof}
Comme $E$ est un complexe parfait, il existe un complexe parfait dual
$E^\vee$ ; voir Cohomologie sur les sites, lemme
\ref{sites-cohomology-lemma-dual-perfect-complex}. On a
$R\SheafHom(E,\mathcal{G}^\bullet)=
E^\vee\otimes^\mathbf{L}_{\mathcal{O}_X}\mathcal{G}^\bullet$.
La perfection de $K$ r\'esulte donc du lemme \ref{lemma-tensor-perfect}.
\end{proof}

\section{Une formule de projection pour Ext}
\label{section-ext}

\noindent
Le lemme \ref{lemma-compute-ext} (ou des r\'esultats analogues de la
litt\'erature) est parfois utile pour v\'erifier les propri\'et\'es d'une th\'eorie
des obstructions n\'ecessaires \`a la v\'erification de l'un des crit\`eres d'Artin
pour les foncteurs Quot, les sch\'emas de Hilbert et d'autres probl\`emes de
modules. Supposons que $f:X\to Y$ soit un morphisme propre, plat et de
pr\'esentation finie d'espaces alg\'ebriques, et que
$E\in D(\mathcal{O}_X)$ soit parfait. Le lemme affirme alors que
$$
\Ext^i_X(E, f^*\mathcal{F}) =
\Ext^i_Y((Rf_*E^\vee)^\vee, \mathcal{F})
$$
pour $\mathcal{F}$ quasi-coh\'erent sur $Y$. Sous cette forme, il ressemble \`a
une formule de projection pour Ext ; et le r\'esultat d\'ecoule en effet assez
facilement du lemme \ref{lemma-cohomology-base-change}.

\begin{lemma}
\label{lemma-compute-tensor-perfect}
Conservons les hypoth\`eses et les notations du lemme
\ref{lemma-tensor-perfect}. Il existe alors des isomorphismes fonctoriels
$$
H^i(B, K \otimes^\mathbf{L}_{\mathcal{O}_B} \mathcal{F})
\longrightarrow
H^i(X, E \otimes^\mathbf{L}_{\mathcal{O}_X}
(\mathcal{G}^\bullet \otimes_{\mathcal{O}_X} f^*\mathcal{F}))
$$
pour $\mathcal{F}$ quasi-coh\'erent sur $B$, compatibles avec les homomorphismes
cobord (voir la d\'emonstration).
\end{lemma}

\begin{proof}
On a
$$
\mathcal{G}^\bullet \otimes_{\mathcal{O}_X}^\mathbf{L} Lf^*\mathcal{F} =
\mathcal{G}^\bullet \otimes_{f^{-1}\mathcal{O}_B}^\mathbf{L} f^{-1}\mathcal{F} =
\mathcal{G}^\bullet \otimes_{f^{-1}\mathcal{O}_B} f^{-1}\mathcal{F} =
\mathcal{G}^\bullet \otimes_{\mathcal{O}_X} f^*\mathcal{F}.
$$
La premi\`ere \`egalit\'e r\'esulte de Cohomologie sur les sites, lemme
\ref{sites-cohomology-lemma-variant-derived-pullback}, la deuxi\`eme du fait que
$\mathcal{G}^n$ est un $f^{-1}\mathcal{O}_B$-module plat, et la troisi\`eme de
la d\'efinition des images inverses. On obtient donc
\begin{align*}
H^i(X, E \otimes^\mathbf{L}_{\mathcal{O}_X}
(\mathcal{G}^\bullet \otimes_{\mathcal{O}_X} f^*\mathcal{F}))
& =
H^i(X, E \otimes^\mathbf{L}_{\mathcal{O}_X} \mathcal{G}^\bullet
\otimes_{\mathcal{O}_X}^\mathbf{L} Lf^*\mathcal{F}) \\
& =
H^i(B,
Rf_*(E \otimes^\mathbf{L}_{\mathcal{O}_X} \mathcal{G}^\bullet
\otimes^\mathbf{L}_{\mathcal{O}_X} Lf^*\mathcal{F})) \\
& =
H^i(B,
Rf_*(E \otimes^\mathbf{L}_{\mathcal{O}_X} \mathcal{G}^\bullet)
\otimes^\mathbf{L}_{\mathcal{O}_B} \mathcal{F}) \\
& =
H^i(B, K \otimes^\mathbf{L}_{\mathcal{O}_B} \mathcal{F}).
\end{align*}
La premi\`ere \`egalit\'e vient de ce qui pr\'ec\`ede, la deuxi\`eme de la suite
spectrale de Leray (Cohomologie sur les sites, remarque
\ref{sites-cohomology-remark-before-Leray}), et la troisi\`eme du lemme
\ref{lemma-cohomology-base-change}. La compatibilit\'e avec les homomorphismes
cobord signifie ceci : pour toute suite exacte courte
$0\to\mathcal{F}_1\to\mathcal{F}_2\to\mathcal{F}_3\to0$, les isomorphismes
s'ins\`erent dans des diagrammes commutatifs
$$
\xymatrix{
H^i(B, K \otimes^\mathbf{L}_{\mathcal{O}_B} \mathcal{F}_3)
\ar[r] \ar[d]_\delta &
H^i(X, E \otimes^\mathbf{L}_{\mathcal{O}_X}
(\mathcal{G}^\bullet \otimes_{\mathcal{O}_X} f^*\mathcal{F}_3)) \ar[d]^\delta \\
H^{i + 1}(B, K \otimes^\mathbf{L}_{\mathcal{O}_B} \mathcal{F}_1)
\ar[r] &
H^{i + 1}(X, E \otimes^\mathbf{L}_{\mathcal{O}_X}
(\mathcal{G}^\bullet \otimes_{\mathcal{O}_X} f^*\mathcal{F}_1))
}
$$
o\`u les homomorphismes cobord proviennent du triangle distingu\'e
$$
K \otimes^\mathbf{L}_{\mathcal{O}_B} \mathcal{F}_1 \to
K \otimes^\mathbf{L}_{\mathcal{O}_B} \mathcal{F}_2 \to
K \otimes^\mathbf{L}_{\mathcal{O}_B} \mathcal{F}_3 \to
K \otimes^\mathbf{L}_{\mathcal{O}_B} \mathcal{F}_1[1]
$$
et du triangle distingu\'e de $D(\mathcal{O}_X)$ associ\'e \`a la suite exacte
courte de complexes
$$
0 \to
\mathcal{G}^\bullet \otimes_{\mathcal{O}_X} f^*\mathcal{F}_1 \to
\mathcal{G}^\bullet \otimes_{\mathcal{O}_X} f^*\mathcal{F}_2 \to
\mathcal{G}^\bullet \otimes_{\mathcal{O}_X} f^*\mathcal{F}_3 \to 0.
$$
Cette suite est exacte parce que $\mathcal{G}^n$ est plat sur $B$.
Nous omettons la v\'erification de la commutativit\'e du diagramme affich\'e.
\end{proof}

\begin{lemma}
\label{lemma-compute-ext-perfect}
Conservons les hypoth\`eses et les notations du lemme
\ref{lemma-ext-perfect}. Il existe alors des isomorphismes fonctoriels
$$
H^i(B, K \otimes^\mathbf{L}_{\mathcal{O}_B} \mathcal{F})
\longrightarrow
\Ext^i_{\mathcal{O}_X}(E,
\mathcal{G}^\bullet \otimes_{\mathcal{O}_X} f^*\mathcal{F})
$$
pour $\mathcal{F}$ quasi-coh\'erent sur $B$, compatibles avec les homomorphismes
cobord (voir la d\'emonstration).
\end{lemma}

\begin{proof}
Comme dans la d\'emonstration du lemme \ref{lemma-ext-perfect}, soit $E^\vee$
le complexe parfait dual, et rappelons que
$K=Rf_*(E^\vee\otimes_{\mathcal{O}_X}^\mathbf{L}\mathcal{G}^\bullet)$.
Comme on a aussi
$$
\Ext^i_{\mathcal{O}_X}(E,
\mathcal{G}^\bullet \otimes_{\mathcal{O}_X} f^*\mathcal{F})
=
H^i(X, E^\vee \otimes^\mathbf{L}_{\mathcal{O}_X}
(\mathcal{G}^\bullet \otimes_{\mathcal{O}_X} f^*\mathcal{F})),
$$
par construction de $E^\vee$, l'existence des isomorphismes r\'esulte du
lemme \ref{lemma-compute-tensor-perfect} appliqu\'e \`a $E^\vee$ et
$\mathcal{G}^\bullet$. La compatibilit\'e avec les homomorphismes cobord signifie
ceci : pour toute suite exacte courte
$0\to\mathcal{F}_1\to\mathcal{F}_2\to\mathcal{F}_3\to0$, les isomorphismes
s'ins\`erent dans des diagrammes commutatifs
$$
\xymatrix{
H^i(B, K \otimes^\mathbf{L}_{\mathcal{O}_B} \mathcal{F}_3)
\ar[r] \ar[d]_\delta &
\Ext^i_{\mathcal{O}_X}(E,
\mathcal{G}^\bullet \otimes_{\mathcal{O}_X} f^*\mathcal{F}_3) \ar[d]^\delta \\
H^{i + 1}(B, K \otimes^\mathbf{L}_{\mathcal{O}_B} \mathcal{F}_1)
\ar[r] &
\Ext^{i + 1}_{\mathcal{O}_X}(E,
\mathcal{G}^\bullet \otimes_{\mathcal{O}_X} f^*\mathcal{F}_1)
}
$$
o\`u les homomorphismes cobord proviennent du triangle distingu\'e
$$
K \otimes^\mathbf{L}_{\mathcal{O}_B} \mathcal{F}_1 \to
K \otimes^\mathbf{L}_{\mathcal{O}_B} \mathcal{F}_2 \to
K \otimes^\mathbf{L}_{\mathcal{O}_B} \mathcal{F}_3 \to
K \otimes^\mathbf{L}_{\mathcal{O}_B} \mathcal{F}_1[1]
$$
et du triangle distingu\'e de $D(\mathcal{O}_X)$ associ\'e \`a la suite exacte
courte de complexes
$$
0 \to
\mathcal{G}^\bullet \otimes_{\mathcal{O}_X} f^*\mathcal{F}_1 \to
\mathcal{G}^\bullet \otimes_{\mathcal{O}_X} f^*\mathcal{F}_2 \to
\mathcal{G}^\bullet \otimes_{\mathcal{O}_X} f^*\mathcal{F}_3 \to 0.
$$
Cette suite est exacte parce que $\mathcal{G}^n$ est plat sur $B$.
Nous omettons la v\'erification de la commutativit\'e du diagramme affich\'e.
\end{proof}

\begin{lemma}
\label{lemma-compute-ext}
Soit $S$ un sch\'ema. Soit $f:X\to B$ un morphisme d'espaces alg\'ebriques sur
$S$, soit $E\in D(\mathcal{O}_X)$, et soit $\mathcal{G}^\bullet$ un complexe de
$\mathcal{O}_X$-modules. Supposons que
\begin{enumerate}
\item $B$ soit noeth\'erien ;
\item $f$ soit localement de type fini et quasi-s\'epar\'e ;
\item $E\in D^-_{\textit{Coh}}(\mathcal{O}_X)$ ;
\item $\mathcal{G}^\bullet$ soit un complexe born\'e de
$\mathcal{O}_X$-modules coh\'erents, plats sur $B$ et dont le support est propre
sur $B$.
\end{enumerate}
Alors les deux assertions suivantes sont vraies :
\begin{enumerate}
\item[(A)] pour tout $m\in\mathbf{Z}$, il existe un objet parfait $K$ de
$D(\mathcal{O}_B)$ et des homomorphismes fonctoriels
$$
\alpha^i_\mathcal{F} :
\Ext^i_{\mathcal{O}_X}(E,
\mathcal{G}^\bullet \otimes_{\mathcal{O}_X} f^*\mathcal{F})
\longrightarrow
H^i(B, K \otimes^\mathbf{L}_{\mathcal{O}_B} \mathcal{F})
$$
pour $\mathcal{F}$ quasi-coh\'erent sur $B$, compatibles avec les homomorphismes
cobord (voir la d\'emonstration), tels que $\alpha^i_\mathcal{F}$ soit un
isomorphisme pour $i\leq m$ ;
\item[(B)] il existe $L\in D(\mathcal{O}_B)$ pseudo-coh\'erent et des
isomorphismes fonctoriels
$$
\Ext^i_{\mathcal{O}_B}(L, \mathcal{F}) \longrightarrow
\Ext^i_{\mathcal{O}_X}(E,
\mathcal{G}^\bullet \otimes_{\mathcal{O}_X} f^*\mathcal{F})
$$
pour $\mathcal{F}$ quasi-coh\'erent sur $B$, compatibles avec les homomorphismes
cobord.
\end{enumerate}
\end{lemma}

\begin{proof}
D\'emonstration de (A). Supposons $\mathcal{G}^i$ non nul seulement pour
$i\in[a,b]$. On peut remplacer $X$ par un voisinage ouvert quasi-compact de la
r\'eunion des supports des $\mathcal{G}^i$. On peut donc supposer $X$
noeth\'erien. Dans ce cas, $X$ et $f$ sont quasi-compacts et quasi-s\'epar\'es.
Choisissons une approximation $P\to E$ par un complexe parfait $P$ de
$(X,E,-m-1+a)$ (ce qui est possible par le th\'eor\`eme
\ref{theorem-approximation}). Le morphisme induit
$$
\Ext^i_{\mathcal{O}_X}(E,
\mathcal{G}^\bullet \otimes_{\mathcal{O}_X} f^*\mathcal{F})
\longrightarrow
\Ext^i_{\mathcal{O}_X}(P,
\mathcal{G}^\bullet \otimes_{\mathcal{O}_X} f^*\mathcal{F})
$$
est alors un isomorphisme pour $i\leq m$. En effet, le noyau, resp. le
conoyau, de ce morphisme est un quotient, resp. un sous-module, de
$$
\Ext^i_{\mathcal{O}_X}(C,
\mathcal{G}^\bullet \otimes_{\mathcal{O}_X} f^*\mathcal{F})
\quad\text{resp.}\quad
\Ext^{i + 1}_{\mathcal{O}_X}(C,
\mathcal{G}^\bullet \otimes_{\mathcal{O}_X} f^*\mathcal{F}),
$$
o\`u $C$ est le c\^one de $P\to E$. Comme les faisceaux de cohomologie de $C$
sont nuls aux degr\'es $\geq-m-1+a$, ces groupes d'Ext sont nuls pour
$i\leq m+1$ d'apr\`es Cat\'egories d\'eriv\'ees, lemme
\ref{derived-lemma-negative-exts}. On est ainsi ramen\'e au cas o\`u $E$ est un
complexe parfait, qui est le lemme \ref{lemma-compute-ext-perfect}.
La compatibilit\'e avec les homomorphismes cobord est expliqu\'ee dans la
d\'emonstration du lemme \ref{lemma-compute-ext-perfect}.

\medskip\noindent
D\'emonstration de (B). Comme dans la d\'emonstration de (A), on peut supposer
$X$ noeth\'erien. Le lemme \ref{lemma-identify-pseudo-coherent-noetherian}
montre que $E$ est pseudo-coh\'erent. D'apr\`es le lemme
\ref{lemma-pseudo-coherent-hocolim}, on peut \`ecrire
$E=\text{hocolim}\,E_n$, o\`u les $E_n$ sont parfaits et o\`u $E_n\to E$ induit
un isomorphisme sur les troncations $\tau_{\geq-n}$. Soit $E_n^\vee$ le
complexe parfait dual (Cohomologie sur les sites, lemme
\ref{sites-cohomology-lemma-dual-perfect-complex}). On obtient un syst\`eme
projectif $\ldots\to E_3^\vee\to E_2^\vee\to E_1^\vee$ d'objets parfaits.
Celui-ci donne \`a son tour un syst\`eme projectif
$$
\ldots \to K_3 \to K_2 \to K_1\quad\text{avec}\quad
K_n = Rf_*(E_n^\vee \otimes_{\mathcal{O}_X}^\mathbf{L} \mathcal{G}^\bullet)
$$
d'objets parfaits sur $B$ ; voir le lemme \ref{lemma-tensor-perfect}.
D'apr\`es le lemme \ref{lemma-compute-ext-perfect} et sa d\'emonstration, ainsi
que les arguments du paragraphe pr\'ec\'edent (avec $P=E_n$), pour tout
$\mathcal{F}$ quasi-coh\'erent sur $B$, on dispose d'homomorphismes canoniques
fonctoriels
$$
\xymatrix{
& \Ext^i_{\mathcal{O}_X}(E,
\mathcal{G}^\bullet \otimes_{\mathcal{O}_X} f^*\mathcal{F})
\ar[ld] \ar[rd] \\
H^i(B, K_{n + 1} \otimes_{\mathcal{O}_B}^\mathbf{L} \mathcal{F})
\ar[rr] & &
H^i(B, K_n \otimes_{\mathcal{O}_B}^\mathbf{L} \mathcal{F})
}
$$
qui sont des isomorphismes pour $i\leq n+a$. Soit $L_n=K_n^\vee$ le complexe
parfait dual. On voit alors que
$L_1\to L_2\to L_3\to\ldots$ est un syst\`eme inductif d'objets parfaits de
$D(\mathcal{O}_B)$ tel que, pour tout $\mathcal{F}$ quasi-coh\'erent sur $B$,
les morphismes
$$
\Ext^i_{\mathcal{O}_B}(L_{n + 1}, \mathcal{F})
\longrightarrow
\Ext^i_{\mathcal{O}_B}(L_n, \mathcal{F})
$$
sont des isomorphismes pour $i\leq n+a-1$. Il en r\'esulte que
$L_n\to L_{n+1}$ induit un isomorphisme sur les troncations
$\tau_{\geq-n-a+2}$ (indication : prendre le c\^one de $L_n\to L_{n+1}$ et
consid\'erer son dernier faisceau de cohomologie non nul). Ainsi,
$L=\text{hocolim}\,L_n$ est pseudo-coh\'erent ; voir le lemme
\ref{lemma-pseudo-coherent-hocolim}. La propri\'et\'e universelle des colimites
homotopiques donne
$\Ext^i_{\mathcal{O}_B}(L,\mathcal{F})=
\Ext^i_{\mathcal{O}_B}(L_n,\mathcal{F})$ pour $i\leq n+a-3$, ce qui ach\`eve
la d\'emonstration.
\end{proof}

\begin{remark}
\label{remark-base-change-of-L}
Le complexe pseudo-coh\'erent $L$ de la partie (B) du lemme
\ref{lemma-compute-ext} est canoniquement associ\'e \`a la situation. Par
exemple, la construction de $L$ dans (B) est compatible au changement de
base. Autrement dit, pour tout diagramme cart\'esien de sch\'emas
$$
\xymatrix{
X' \ar[r]_{g'} \ar[d]_{f'} &
X \ar[d]^f \\
Y' \ar[r]^g &
Y
}
$$
on a des isomorphismes canoniques fonctoriels
$$
\Ext^i_{\mathcal{O}_{Y'}}(Lg^*L, \mathcal{F}') \longrightarrow
\Ext^i_{\mathcal{O}_{X'}}(L(g')^*E,
(g')^*\mathcal{G}^\bullet \otimes_{\mathcal{O}_{X'}} (f')^*\mathcal{F}')
$$
pour $\mathcal{F}'$ quasi-coh\'erent sur $Y'$. Notons que nous n'utilisons
{\bf pas} l'image inverse d\'eriv\'ee de $\mathcal{G}^\bullet$ au membre de
droite. Si ce r\'esultat nous est un jour n\'ecessaire, nous en formulerons ici
un \`enonc\'e pr\'ecis et en donnerons une d\'emonstration d\'etaill\'ee.
\end{remark}

\section{Limites et cat\'egories d\'eriv\'ees}
\label{section-limits}

\noindent
Nous rassemblons dans cette section quelques r\'esultats sur la cat\'egorie
d\'eriv\'ee d'un espace alg\'ebrique qui est la limite projective d'un syst\`eme
projectif d'espaces alg\'ebriques. Plus pr\'ecis\'ement, nous nous placerons dans
la situation suivante.

\begin{situation}
\label{situation-descent}
Soit $S$ un sch\'ema. Soit $X=\lim_{i\in I}X_i$ la limite d'un syst\`eme
projectif d'espaces alg\'ebriques sur $S$, dont les morphismes de transition
affines sont $f_{i'i}:X_{i'}\to X_i$. Notons $f_i:X\to X_i$ la projection.
Nous supposons $X_i$ quasi-compact et quasi-s\'epar\'e pour tout $i\in I$.
Nous choisissons en outre un \`el\'ement $0\in I$.
\end{situation}

\begin{lemma}
\label{lemma-descend-homomorphisms}
Pla\c{c}ons-nous dans la situation \ref{situation-descent}. Soient $E_0$ et
$K_0$ des objets de $D(\mathcal{O}_{X_0})$. Posons
$E_i=Lf_{i0}^*E_0$ et $K_i=Lf_{i0}^*K_0$ pour $i\geq0$, ainsi que
$E=Lf_0^*E_0$ et $K=Lf_0^*K_0$. Alors le morphisme
$$
\colim_{i \geq 0} \Hom_{D(\mathcal{O}_{X_i})}(E_i, K_i)
\longrightarrow
\Hom_{D(\mathcal{O}_X)}(E, K)
$$
est un isomorphisme dans chacun des cas suivants :
\begin{enumerate}
\item $E_0$ est parfait et $K_0\in D_\QCoh(\mathcal{O}_{X_0})$ ;
\item $E_0$ est pseudo-coh\'erent et
$K_0\in D_\QCoh(\mathcal{O}_{X_0})$ est de Tor-dimension finie.
\end{enumerate}
\end{lemma}

\begin{proof}
Pour tout objet quasi-compact et quasi-s\'epar\'e $U_0$ de
$(X_0)_{spaces,\etale}$, consid\'erons la propri\'et\'e $P$ selon laquelle le
morphisme canonique
$$
\colim_{i \geq 0} \Hom_{D(\mathcal{O}_{U_i})}(E_i|_{U_i}, K_i|_{U_i})
\longrightarrow
\Hom_{D(\mathcal{O}_U)}(E|_U, K|_U)
$$
est un isomorphisme, o\`u $U=X\times_{X_0}U_0$ et
$U_i=X_i\times_{X_0}U_0$. Nous allons montrer, par le principe de r\'ecurrence
du lemme \ref{lemma-induction-principle}, que chaque $U_0$ v\'erifie $P$.
La condition (2) de ce lemme r\'esulte imm\'ediatement de Mayer--Vietoris pour
les Hom de la cat\'egorie d\'eriv\'ee ; voir le lemme
\ref{lemma-mayer-vietoris-hom}. Il suffit donc de d\'emontrer le lemme lorsque
$X_0$ est affine.

\medskip\noindent
Si $X_0$ est affine, le r\'esultat d\'ecoule du cas des sch\'emas ; voir
Cat\'egories d\'eriv\'ees des sch\'emas, lemme
\ref{perfect-lemma-descend-homomorphisms}. Pour le constater, utilisons
l'\'equivalence du lemme \ref{lemma-derived-quasi-coherent-small-etale-site}
et la correspondance des propri\'et\'es expliqu\'ee dans les lemmes
\ref{lemma-descend-pseudo-coherent},
\ref{lemma-descend-tor-amplitude} et
\ref{lemma-descend-perfect}.
\end{proof}

\begin{lemma}
\label{lemma-perfect-on-limit}
Dans la situation \ref{situation-descent}, la cat\'egorie des objets parfaits
de $D(\mathcal{O}_X)$ est la limite inductive des cat\'egories d'objets
parfaits de $D(\mathcal{O}_{X_i})$.
\end{lemma}

\begin{proof}
Pour tout objet quasi-compact et quasi-s\'epar\'e $U_0$ de
$(X_0)_{spaces,\etale}$, consid\'erons la propri\'et\'e $P$ selon laquelle le
foncteur
$$
\colim_{i \geq 0} D_{perf}(\mathcal{O}_{U_i})
\longrightarrow
D_{perf}(\mathcal{O}_U)
$$
est une \`equivalence, o\`u l'indice ${}_{perf}$ d\'esigne la sous-cat\'egorie
pleine des objets parfaits, et o\`u $U=X\times_{X_0}U_0$ et
$U_i=X_i\times_{X_0}U_0$. Nous allons montrer, par le principe de r\'ecurrence
du lemme \ref{lemma-induction-principle}, que chaque $U_0$ v\'erifie $P$.
D'apr\`es le lemme \ref{lemma-descend-homomorphisms}, nous savons d\'ej\`a que
le foncteur est pleinement fid\`ele. Il suffit donc d'\'etablir sa surjectivit\'e
essentielle.

\medskip\noindent
V\'erifions d'abord la condition (2) du principe de r\'ecurrence. Supposons donc
donn\'e un carr\'e distingu\'e \`el\'ementaire
$(U_0\subset X_0,V_0\to X_0)$, et supposons que $U_0$, $V_0$ et
$U_0\times_{X_0}V_0$ v\'erifient $P$. Soit $E$ un objet parfait de
$D(\mathcal{O}_X)$. On peut trouver $i\geq0$, ainsi que des objets parfaits
$E_{U,i}$ sur $U_i$ et $E_{V,i}$ sur $V_i$ dont les images inverses sur $U$
et $V$ sont respectivement isomorphes \`a $E|_U$ et $E|_V$. Notons
$$
a : E_{U, i} \to (R(X \to X_i)_*E)|_{U_i}
\quad\text{et}\quad
b : E_{V, i} \to (R(X \to X_i)_*E)|_{V_i}
$$
les morphismes adjoints aux isomorphismes
$L(U\to U_i)^*E_{U,i}\to E|_U$ et
$L(V\to V_i)^*E_{V,i}\to E|_V$. Par pleine fid\'elit\'e, quitte \`a augmenter
$i$, on peut trouver un isomorphisme
$c:E_{U,i}|_{U_i\times_{X_i}V_i}\to
E_{V,i}|_{U_i\times_{X_i}V_i}$ dont l'image inverse donne les identifications
$$
L(U \to U_i)^*E_{U, i}|_{U \times_X V} \to E|_{U \times_X V} \to
L(V \to V_i)^*E_{V, i}|_{U \times_X V}.
$$
Appliquons le lemme \ref{lemma-glue}. On obtient un objet $E_i$ sur $X_i$ et
un morphisme $d:E_i\to R(X\to X_i)_*E$ dont les restrictions sur $U_i$ et
$V_i$ sont les morphismes $a$ et $b$. Il est alors clair que $E_i$ est parfait
et que $d$ est adjoint \`a un isomorphisme $L(X\to X_i)^*E_i\to E$.

\medskip\noindent
V\'erifions enfin la condition (1) du principe de r\'ecurrence, autrement dit
le cas o\`u $X_0$ est affine. Celui-ci r\'esulte du cas des sch\'emas ; voir
Cat\'egories d\'eriv\'ees des sch\'emas, lemme
\ref{perfect-lemma-descend-perfect}. Pour le constater, utilisons
l'\'equivalence du lemme \ref{lemma-derived-quasi-coherent-small-etale-site}
et la correspondance du lemme \ref{lemma-descend-perfect}.
\end{proof}

\section{Cohomologie et changement de base, VI}
\label{section-cohomology-and-base-change-final}

\noindent
Cette dernière section sur la cohomologie et le changement de base poursuit
la discussion des sections
\ref{section-cohomology-and-base-change-perfect},
\ref{section-cohomology-base-change} et
\ref{section-producing-perfect}.
Un cas particulier facile à saisir est donné dans
la remarque \ref{remark-explain-perfect-direct-image}.

\begin{lemma}
\label{lemma-base-change-tensor-perfect}
Soit $S$ un schéma. Soit $f : X \to Y$ un morphisme de présentation finie
entre des espaces algébriques sur $S$. Soit $E \in D(\mathcal{O}_X)$ un objet
parfait. Soit $\mathcal{G}^\bullet$ un complexe borné de
$\mathcal{O}_X$-modules de présentation finie dont les termes sont plats sur
$Y$ et dont le support est propre sur $Y$. Alors
$$
K = Rf_*(E \otimes_{\mathcal{O}_X}^\mathbf{L} \mathcal{G}^\bullet)
$$
est un objet parfait de $D(\mathcal{O}_Y)$ et sa formation
commute à tout changement de base.
\end{lemma}

\begin{proof}
L'assertion concernant le changement de base est le
lemme \ref{lemma-base-change-tensor}.
Il suffit donc de montrer que $K$ est un objet parfait. Si $Y$ est
noethérien, cela résulte du
lemme \ref{lemma-tensor-perfect}.
Nous nous ramènerons à ce cas par approximation noethérienne.
Nous conseillons au lecteur de ne pas lire la suite de cette démonstration.

\medskip\noindent
La question étant locale sur $Y$, nous pouvons supposer $Y$ affine.
Écrivons $Y = \Spec(R)$. Écrivons $R = \colim R_i$ comme limite inductive
filtrante d'anneaux noethériens $R_i$. D'après Limites d'espaces, lemme
\ref{spaces-limits-lemma-descend-finite-presentation},
il existe un $i$ et un espace algébrique $X_i$ de présentation finie
sur $R_i$ dont le changement de base à $R$ est $X$. D'après
Limites d'espaces, lemme
\ref{spaces-limits-lemma-descend-modules-finite-presentation},
quitte à remplacer $i$ par un indice plus grand, nous pouvons supposer
qu'il existe un complexe borné de
$\mathcal{O}_{X_i}$-modules de présentation finie $\mathcal{G}_i^\bullet$ dont
l'image inverse sur $X$ est $\mathcal{G}^\bullet$. Quitte à remplacer $i$
par un indice plus grand,
nous pouvons supposer que $\mathcal{G}_i^n$ est plat sur $R_i$ ; voir
Limites d'espaces, lemme
\ref{spaces-limits-lemma-descend-flat}.
Quitte à remplacer $i$ par un indice plus grand, nous pouvons supposer que
le support de $\mathcal{G}_i^n$ est propre sur $R_i$ ; voir
Limites d'espaces, lemme \ref{spaces-limits-lemma-eventually-proper-support}.
Enfin, d'après le lemme \ref{lemma-perfect-on-limit},
quitte à remplacer $i$ par un indice plus grand, nous pouvons supposer
qu'il existe un objet parfait $E_i$ de $D(\mathcal{O}_{X_i})$ dont l'image
inverse sur $X$ est $E$. D'après le lemme \ref{lemma-tensor-perfect},
$K_i =
Rf_{i, *}(E_i \otimes_{\mathcal{O}_{X_i}}^\mathbf{L} \mathcal{G}_i^\bullet)$
est parfait sur $\Spec(R_i)$, où $f_i : X_i \to \Spec(R_i)$ est
le morphisme structural. D'après le résultat de changement de base
(lemme \ref{lemma-base-change-tensor}), l'image inverse de $K_i$
sur $Y = \Spec(R)$ est $K$, ce qui permet de conclure.
\end{proof}

\begin{remark}
\label{remark-explain-perfect-direct-image}
Soit $R$ un anneau. Soit $X$ un espace algébrique de présentation finie sur
$R$. Soit $\mathcal{G}$ un $\mathcal{O}_X$-module de présentation finie,
plat sur $R$ et dont le support est propre sur $R$. D'après le
lemme \ref{lemma-base-change-tensor-perfect},
il existe un complexe fini de $R$-modules projectifs de type fini
$M^\bullet$ tel que l'on ait
$$
R\Gamma(X_{R'}, \mathcal{G}_{R'}) = M^\bullet \otimes_R R'
$$
fonctoriellement en la $R$-algèbre $R'$.
\end{remark}

\begin{lemma}
\label{lemma-base-change-tensor-pseudo-coherent}
Soit $S$ un schéma.
Soit $f : X \to Y$ un morphisme de présentation finie
entre des espaces algébriques sur $S$.
Soit $E \in D(\mathcal{O}_X)$ un objet pseudo-cohérent.
Soit $\mathcal{G}^\bullet$ un complexe borné supérieurement de
$\mathcal{O}_X$-modules de présentation finie dont les termes sont
plats sur $Y$ et dont le support est propre sur $Y$. Alors
$$
K = Rf_*(E \otimes_{\mathcal{O}_X}^\mathbf{L} \mathcal{G}^\bullet)
$$
est un objet pseudo-cohérent de $D(\mathcal{O}_Y)$ et sa formation
commute à tout changement de base.
\end{lemma}

\begin{proof}
L'assertion concernant le changement de base est le
lemme \ref{lemma-base-change-tensor}.
Il suffit donc de montrer que $K$ est un objet pseudo-cohérent.
Cela résultera du lemme \ref{lemma-base-change-tensor-perfect}
grâce à une approximation par des complexes parfaits. Nous conseillons au
lecteur de ne pas lire la suite de cette démonstration.

\medskip\noindent
La question est locale pour la topologie étale sur $Y$ ; nous pouvons donc
supposer $Y$ affine. Alors $X$ est quasi-compact et quasi-séparé. De plus, il
existe un entier $N$ tel que l'image directe totale
$Rf_* : D_\QCoh(\mathcal{O}_X) \to D_\QCoh(\mathcal{O}_Y)$
soit de dimension cohomologique $N$, comme l'explique le
lemme \ref{lemma-quasi-coherence-direct-image}.
Choisissons un entier $b$ tel que $\mathcal{G}^i = 0$ pour $i > b$.
Il suffit de montrer que $K$ est $m$-pseudo-cohérent pour
tout $m$. Choisissons une approximation $P \to E$ par un complexe parfait $P$
de $(X, E, m - N - 1 - b)$. Cela est possible d'après le
théorème \ref{theorem-approximation}.
Choisissons un triangle distingué
$$
P \to E \to C \to P[1]
$$
dans $D_\QCoh(\mathcal{O}_X)$. Les faisceaux de cohomologie de $C$ sont nuls
en degrés $\geq m - N - 1 - b$. Par conséquent, les faisceaux de cohomologie
de $C \otimes^\mathbf{L} \mathcal{G}^\bullet$ sont nuls en degrés
$\geq m - N - 1$.
Ainsi, les faisceaux de cohomologie de
$Rf_*(C \otimes^\mathbf{L} \mathcal{G})$ sont nuls en degrés $\geq m - 1$.
Par conséquent,
$$
Rf_*(P \otimes^\mathbf{L} \mathcal{G}) \to
Rf_*(E \otimes^\mathbf{L} \mathcal{G})
$$
est un isomorphisme sur les faisceaux de cohomologie en degrés $\geq m$.
Supposons ensuite que $H^i(P) = 0$ pour $i > a$. Alors
$
P \otimes^\mathbf{L} \sigma_{\geq m - N - 1 - a}\mathcal{G}^\bullet
\longrightarrow
P \otimes^\mathbf{L} \mathcal{G}^\bullet
$
est un isomorphisme sur les faisceaux de cohomologie en degrés
$\geq m - N - 1$.
Nous trouvons donc de nouveau que
$$
Rf_*(P \otimes^\mathbf{L} \sigma_{\geq m - N - 1 - a}\mathcal{G}^\bullet) \to
Rf_*(P \otimes^\mathbf{L} \mathcal{G}^\bullet)
$$
est un isomorphisme sur les faisceaux de cohomologie en degrés $\geq m$.
D'après le lemme \ref{lemma-base-change-tensor-perfect}, la source
est un complexe parfait.
Nous en concluons que $K$ est $m$-pseudo-cohérent, comme voulu.
\end{proof}

\begin{lemma}
\label{lemma-flat-proper-perfect-direct-image-general}
Soit $S$ un schéma. Soit $f : X \to Y$ un morphisme propre
et de présentation finie entre des espaces algébriques sur $S$.
\begin{enumerate}
\item Soit $E \in D(\mathcal{O}_X)$ un objet parfait et supposons $f$ plat.
Alors
$Rf_*E$ est un objet parfait de $D(\mathcal{O}_Y)$ et sa formation
commute à tout changement de base.
\item Soit $\mathcal{G}$ un $\mathcal{O}_X$-module de présentation finie,
plat sur $S$. Alors $Rf_*\mathcal{G}$ est un objet parfait de
$D(\mathcal{O}_Y)$ et sa formation commute à tout changement de base.
\end{enumerate}
\end{lemma}

\begin{proof}
Ce sont des cas particuliers du
lemme \ref{lemma-base-change-tensor-perfect}, obtenus en prenant
(1) $\mathcal{G}^\bullet$ égal au complexe $\mathcal{O}_X$ placé en degré $0$
et (2) $E = \mathcal{O}_X$ et $\mathcal{G}^\bullet$ égal au complexe constitué
de $\mathcal{G}$ placé en degré $0$.
\end{proof}

\begin{lemma}
\label{lemma-flat-proper-pseudo-coherent-direct-image-general}
Soit $S$ un sch\'ema. Soit $f : X \to Y$ un morphisme propre, plat et
de pr\'esentation finie entre des espaces alg\'ebriques sur $S$.
Soit $E \in D(\mathcal{O}_X)$ un objet
pseudo-coh\'erent. Alors $Rf_*E$ est un objet pseudo-coh\'erent de
$D(\mathcal{O}_Y)$ et sa formation commute \`a tout changement de base.
\end{lemma}

\noindent
Plus g\'en\'eralement, si $f : X \to Y$ est propre et si $E$ sur $X$ est
pseudo-coh\'erent relativement \`a $Y$ (Compl\'ements sur les morphismes d'espaces, d\'efinition
\ref{spaces-more-morphisms-definition-relative-pseudo-coherence}),
alors $Rf_*E$ est pseudo-coh\'erent
(mais sa formation ne commute pas au changement de base \`a ce degr\'e de g\'en\'eralit\'e).
Pour les sch\'emas, cet \`enonc\'e est d\'emontr\'e dans \cite{Kiehl}.

\begin{proof}
C'est le cas particulier du
lemme \ref{lemma-base-change-tensor-pseudo-coherent} obtenu en prenant
$\mathcal{G} = \mathcal{O}_X$.
\end{proof}

\begin{lemma}
\label{lemma-pullback-and-limits}
Soit $R$ un anneau. Soit $X$ un espace alg\'ebrique et soit
$f : X \to \Spec(R)$ un morphisme propre, plat et
de pr\'esentation finie. Soit $(M_n)$ un syst\`eme projectif
de $R$-modules dont les morphismes de transition sont surjectifs.
Alors le morphisme canonique
$$
\mathcal{O}_X \otimes_R (\lim M_n)
\longrightarrow
\lim \mathcal{O}_X \otimes_R M_n
$$
induit un isomorphisme de la source sur l'objet obtenu en appliquant $DQ_X$
\`a la cible.
\end{lemma}

\begin{proof}
L'\'enonc\'e signifie que, pour tout objet $E$ de
$D_\QCoh(\mathcal{O}_X)$, le morphisme induit
$$
\Hom(E, \mathcal{O}_X \otimes_R (\lim M_n))
\longrightarrow
\Hom(E, \lim \mathcal{O}_X \otimes_R M_n)
$$
est un isomorphisme. Comme $D_\QCoh(\mathcal{O}_X)$ admet
un g\'en\'erateur parfait (th\'eor\`eme \ref{theorem-bondal-van-den-Bergh}),
il suffit de le v\'erifier pour $E$ parfait.
D'apr\`es le lemme \ref{lemma-Rlim-quasi-coherent}, on a
$\lim \mathcal{O}_X \otimes_R M_n = R\lim \mathcal{O}_X \otimes_R M_n$.
Le foncteur exact
$R\Hom_X(E, -) : D_\QCoh(\mathcal{O}_X) \to D(R)$,
consid\'er\'e dans Cohomologie sur les sites, section
\ref{sites-cohomology-section-global-RHom}, commute aux produits, donc aux
limites projectives d\'eriv\'ees ; par cons\'equent,
$$
R\Hom_X(E, \lim \mathcal{O}_X \otimes_R M_n) =
R\lim R\Hom_X(E, \mathcal{O}_X \otimes_R M_n)
$$
Soit $E^\vee$ le complexe parfait dual ; voir
Cohomologie sur les sites, lemme \ref{sites-cohomology-lemma-dual-perfect-complex}.
On a
$$
R\Hom_X(E, \mathcal{O}_X \otimes_R M_n) =
R\Gamma(X, E^\vee \otimes_{\mathcal{O}_X}^\mathbf{L} Lf^*M_n) =
R\Gamma(X, E^\vee) \otimes_R^\mathbf{L} M_n
$$
d'apr\`es le lemme \ref{lemma-cohomology-base-change}.
Le lemme \ref{lemma-flat-proper-perfect-direct-image-general} montre que
$R\Gamma(X, E^\vee)$ est un complexe parfait de $R$-modules.
En particulier, c'est un complexe pseudo-coh\'erent ; le
lemme \ref{more-algebra-lemma-pseudo-coherent-tensor-limit} de Compl\'ements
d'alg\`ebre donne donc
$$
R\lim R\Gamma(X, E^\vee) \otimes_R^\mathbf{L} M_n =
R\Gamma(X, E^\vee) \otimes_R^\mathbf{L} \lim M_n
$$
ce qu'il fallait d\'emontrer.
\end{proof}

\begin{lemma}
\label{lemma-perfect-enough}
Soit $A$ un anneau. Soit $X$ un espace alg\'ebrique sur $A$, quasi-compact
et quasi-s\'epar\'e. Soit $K \in D^-_\QCoh(\mathcal{O}_X)$.
Si $R\Gamma(X, E \otimes^\mathbf{L} K)$ est pseudo-coh\'erent
dans $D(A)$ pour tout $E$ parfait de $D(\mathcal{O}_X)$,
alors $R\Gamma(X, E \otimes^\mathbf{L} K)$ est pseudo-coh\'erent
dans $D(A)$ pour tout $E$ pseudo-coh\'erent de $D(\mathcal{O}_X)$.
\end{lemma}

\begin{proof}
Il existe un entier $N$ tel que
$R\Gamma(X, -) : D_\QCoh(\mathcal{O}_X) \to D(A)$
soit de dimension cohomologique $N$, comme expliqu\'e dans le
lemme \ref{lemma-quasi-coherence-direct-image}.
Soit $b \in \mathbf{Z}$ tel que $H^i(K) = 0$ pour $i > b$.
Soit $E$ pseudo-coh\'erent sur $X$.
Il suffit de montrer que $R\Gamma(X, E \otimes^\mathbf{L} K)$
est $m$-pseudo-coh\'erent pour tout $m$.
Choisissons une approximation $P \to E$ par un complexe parfait $P$
de $(X, E, m - N - 1 - b)$. C'est possible d'apr\`es le
th\'eor\`eme \ref{theorem-approximation}.
Choisissons un triangle distingu\'e
$$
P \to E \to C \to P[1]
$$
dans $D_\QCoh(\mathcal{O}_X)$. Les faisceaux de cohomologie de $C$ sont nuls
en degr\'es $\geq m - N - 1 - b$. Par suite, les faisceaux de cohomologie
de $C \otimes^\mathbf{L} K$ sont nuls en degr\'es $\geq m - N - 1$.
Les groupes de cohomologie de $R\Gamma(X, C \otimes^\mathbf{L} K)$
sont donc nuls en degr\'es $\geq m - 1$. Ainsi,
$$
R\Gamma(X, P \otimes^\mathbf{L} K) \to R\Gamma(X, E \otimes^\mathbf{L} K)
$$
induit un isomorphisme sur la cohomologie en degr\'es $\geq m$.
Par hypoth\`ese, la source est pseudo-coh\'erente.
On en d\'eduit que $R\Gamma(X, E \otimes^\mathbf{L} K)$
est $m$-pseudo-coh\'erent, comme voulu.
\end{proof}

\begin{lemma}
\label{lemma-base-change-RHom-perfect}
Soit $S$ un sch\'ema.
Soit $f : X \to Y$ un morphisme de pr\'esentation finie
entre des espaces alg\'ebriques sur $S$.
Soit $E \in D(\mathcal{O}_X)$ un objet parfait. Soit $\mathcal{G}^\bullet$
un complexe born\'e de $\mathcal{O}_X$-modules de pr\'esentation finie,
dont les termes sont plats sur $Y$ et dont le support est propre sur $Y$. Alors
$$
K = Rf_*R\SheafHom(E, \mathcal{G}^\bullet)
$$
est un objet parfait de $D(\mathcal{O}_Y)$ et sa formation
commute \`a tout changement de base.
\end{lemma}

\begin{proof}
L'assertion sur le changement de base est le lemme \ref{lemma-base-change-RHom}.
Il suffit donc de montrer que $K$ est un objet parfait. Si $Y$ est
noeth\'erien, cela r\'esulte du lemme \ref{lemma-ext-perfect}.
Nous allons nous ramener \`a ce cas par approximation noeth\'erienne.
Nous conseillons au lecteur de ne pas lire la suite de cette d\'emonstration.

\medskip\noindent
La question \`etant locale sur $Y$, on peut donc supposer $Y$ affine.
Écrivons $Y = \Spec(R)$. Écrivons $R = \colim R_i$ comme limite inductive
filtrante d'anneaux noeth\'eriens $R_i$. D'apr\`es Limites d'espaces, lemme
\ref{spaces-limits-lemma-descend-finite-presentation},
il existe un $i$ et un espace alg\'ebrique $X_i$ de pr\'esentation finie
sur $R_i$ dont le changement de base \`a $R$ est $X$. D'apr\`es
Limites d'espaces, lemme
\ref{spaces-limits-lemma-descend-modules-finite-presentation},
on peut supposer, quitte \`a remplacer $i$ par un indice plus grand,
qu'il existe un complexe born\'e de $\mathcal{O}_{X_i}$-modules de pr\'esentation
finie $\mathcal{G}_i^\bullet$ dont
l'image inverse sur $X$ est $\mathcal{G}$. Quitte \`a remplacer $i$
par un indice plus grand,
on peut supposer que $\mathcal{G}_i^n$ est plat sur $R_i$ ; voir
Limites d'espaces, lemme
\ref{spaces-limits-lemma-descend-flat}.
Quitte \`a remplacer $i$ par un indice plus grand, on peut supposer que
le support de $\mathcal{G}_i^n$ est propre sur $R_i$ ; voir
Limites d'espaces, lemme \ref{spaces-limits-lemma-eventually-proper-support}.
Enfin, d'apr\`es le lemme \ref{lemma-descend-perfect},
on peut, quitte \`a remplacer $i$ par un indice plus grand, supposer qu'il
existe un objet parfait $E_i$ de $D(\mathcal{O}_{X_i})$ dont l'image inverse
sur $X$ est $E$. En appliquant le lemme \ref{lemma-compute-ext-perfect}
\`a $X_i \to \Spec(R_i)$, $E_i$, $\mathcal{G}_i^\bullet$ et en utilisant la
propri\'et\'e de changement de base d\'ej\`a d\'emontr\'ee, on obtient le r\'esultat.
\end{proof}







\section{Complexes parfaits}
\label{section-perfect-complexes}

\noindent
Nous allons d'abord parler des lieux de saut des nombres de Betti des complexes
parfaits. Il nous faut commencer par d\'efinir les nombres de Betti.

\medskip\noindent
Soit $S$ un sch\'ema. Soit $X$ un espace alg\'ebrique sur $S$.
Soit $E$ un objet de $D(\mathcal{O}_X)$.
Soit $x\in|X|$. Nous voulons d\'efinir
$\beta_i(x)\in\{0,1,2,\ldots\}\cup\{\infty\}$.
Pour cela, choisissons dans la classe d'\'equivalence $x$ un morphisme
$f:\Spec(k)\to X$. Alors $Lf^*E$ est un objet de
$D(\Spec(k)_\etale,\mathcal{O})$.
D'apr\`es Cohomologie \'etale, lemme
\ref{etale-cohomology-lemma-all-modules-quasi-coherent}, et le th\'eor\`eme
\ref{etale-cohomology-theorem-quasi-coherent}, on a
$D(\Spec(k)_\etale,\mathcal{O})=D(k)$, la cat\'egorie d\'eriv\'ee des
$k$-espaces vectoriels. Ainsi, $Lf^*E$ est un complexe de $k$-espaces
vectoriels et l'on peut poser
$\beta_i(x)=\dim_k H^i(Lf^*E)$. On voit facilement que cette valeur ne d\'epend
pas du choix du repr\'esentant de $x$. De plus, si $X$ est un sch\'ema, cette
notion co\"incide avec celle employ\'ee dans Cat\'egories d\'eriv\'ees des
sch\'emas, section \ref{perfect-section-perfect-complexes}.

\begin{lemma}
\label{lemma-jump-loci}
Soit $S$ un sch\'ema. Soit $X$ un espace alg\'ebrique sur $S$.
Soit $E\in D(\mathcal{O}_X)$ pseudo-coh\'erent (par exemple parfait).
Pour tout $i\in\mathbf{Z}$, consid\'erons la fonction
$$
\beta_i : |X| \longrightarrow \{0, 1, 2, \ldots\}
$$
d\'efinie ci-dessus. Alors
\begin{enumerate}
\item la formation de $\beta_i$ commute \`a tout changement de base ;
\item les fonctions $\beta_i$ sont semi-continues sup\'erieurement ;
\item les ensembles de niveau de $\beta_i$ sont localement constructibles
pour la topologie \'etale.
\end{enumerate}
\end{lemma}

\begin{proof}
Choisissons un sch\'ema $U$ et un morphisme \'etale surjectif
$\varphi:U\to X$. Alors $L\varphi^*E$ est un complexe pseudo-coh\'erent sur le
sch\'ema $U$ (utiliser le lemme \ref{lemma-descend-pseudo-coherent}), et l'on
peut appliquer le r\'esultat pour les sch\'emas ; voir Cat\'egories d\'eriv\'ees
des sch\'emas, lemme \ref{perfect-lemma-jump-loci}. L'assertion (3) signifie
que les images inverses des ensembles de niveau sur $U$ sont localement
constructibles ; voir Propri\'et\'es des espaces, d\'efinition
\ref{spaces-properties-definition-locally-constructible}.
\end{proof}

\begin{lemma}
\label{lemma-jump-loci-geometric}
Soit $Y$ un sch\'ema et soit $X$ un espace alg\'ebrique sur $Y$ tel que le
morphisme structural $f:X\to Y$ soit plat, propre et de pr\'esentation finie.
Soit $\mathcal{F}$ un $\mathcal{O}_X$-module de pr\'esentation finie, plat sur
$Y$. Pour $i\in\mathbf{Z}$ fix\'e, consid\'erons la fonction
$$
\beta_i : |Y| \to \{0, 1, 2, \ldots\},\quad
y \longmapsto \dim_{\kappa(y)} H^i(X_y, \mathcal{F}_y).
$$
Alors
\begin{enumerate}
\item la formation de $\beta_i$ commute \`a tout changement de base ;
\item les fonctions $\beta_i$ sont semi-continues sup\'erieurement ;
\item les ensembles de niveau de $\beta_i$ sont localement constructibles
dans $Y$.
\end{enumerate}
\end{lemma}

\begin{proof}
D'apr\`es le th\'eor\`eme de cohomologie et changement de base (plus
pr\'ecis\'ement, le lemme
\ref{lemma-flat-proper-perfect-direct-image-general}), l'objet
$K=Rf_*\mathcal{F}$ est un objet parfait de la cat\'egorie d\'eriv\'ee de $Y$
dont la formation commute \`a tout changement de base. En particulier,
$$
H^i(X_y, \mathcal{F}_y) = H^i(K \otimes_{\mathcal{O}_Y}^\mathbf{L} \kappa(y)).
$$
Le lemme r\'esulte donc du lemme \ref{lemma-jump-loci}.
\end{proof}

\begin{lemma}
\label{lemma-chi-locally-constant}
Soit $S$ un sch\'ema. Soit $X$ un espace alg\'ebrique sur $S$.
Soit $E\in D(\mathcal{O}_X)$ parfait. La fonction
$$
\chi_E : |X| \longrightarrow \mathbf{Z},\quad
x \longmapsto \sum (-1)^i \beta_i(x)
$$
est localement constante sur $X$.
\end{lemma}

\begin{proof}
D\'emonstration omise. Indications :
on se ram\`ene au cas des sch\'emas par localisation \'etale. Voir Cat\'egories
d\'eriv\'ees des sch\'emas, lemme
\ref{perfect-lemma-chi-locally-constant}.
\end{proof}

\begin{lemma}
\label{lemma-open-where-cohomology-in-degree-i-rank-r}
Soit $S$ un sch\'ema. Soit $X$ un espace alg\'ebrique sur $S$.
Soit $E\in D(\mathcal{O}_X)$ parfait.
\'Etant donn\'es $i,r\in\mathbf{Z}$, il existe un sous-espace ouvert
$U\subset X$ caract\'eris\'e par les propri\'et\'es suivantes :
\begin{enumerate}
\item $E|_U\cong H^i(E|_U)[-i]$ et $H^i(E|_U)$ est un
$\mathcal{O}_U$-module localement libre de rang $r$ ;
\item un morphisme $f:Y\to X$ se factorise par $U$ si et seulement si
$Lf^*E$ est isomorphe \`a un module localement libre de rang $r$ plac\'e en
degr\'e $i$.
\end{enumerate}
\end{lemma}

\begin{proof}
D\'emonstration omise. Indications :
on se ram\`ene au cas des sch\'emas par localisation \'etale. Voir Cat\'egories
d\'eriv\'ees des sch\'emas, lemme
\ref{perfect-lemma-open-where-cohomology-in-degree-i-rank-r}.
\end{proof}

\begin{lemma}
\label{lemma-open-where-cohomology-in-degree-i-rank-r-geometric}
Soit $S$ un sch\'ema.
Soit $f:X\to Y$ un morphisme d'espaces alg\'ebriques sur $S$, propre, plat et
de pr\'esentation finie.
Soit $\mathcal{F}$ un $\mathcal{O}_X$-module de pr\'esentation finie, plat sur
$Y$. Fixons $i,r\in\mathbf{Z}$.
Il existe alors un sous-espace ouvert $V\subset Y$ poss\'edant la propri\'et\'e
suivante : un morphisme $T\to Y$ se factorise par $V$ si et seulement si
$Rf_{T,*}\mathcal{F}_T$ est isomorphe \`a un module localement libre de type
fini et de rang $r$, plac\'e en degr\'e $i$.
\end{lemma}

\begin{proof}
D'apr\`es le th\'eor\`eme de cohomologie et changement de base
(lemme \ref{lemma-flat-proper-perfect-direct-image-general}), l'objet
$K=Rf_*\mathcal{F}$ est un objet parfait de la cat\'egorie d\'eriv\'ee de $Y$
dont la formation commute \`a tout changement de base. Ce lemme r\'esulte donc
imm\'ediatement du lemme
\ref{lemma-open-where-cohomology-in-degree-i-rank-r}.
\end{proof}

\begin{lemma}
\label{lemma-locally-closed-where-H0-locally-free}
Soit $S$ un sch\'ema. Soit $X$ un espace alg\'ebrique sur $S$.
Soit $E\in D(\mathcal{O}_X)$ parfait, de Tor-amplitude contenue dans
$[a,b]$ pour certains $a,b\in\mathbf{Z}$.
Soit $r\geq0$.
Il existe alors un sous-espace localement ferm\'e $j:Z\to X$ caract\'eris\'e
par les propri\'et\'es suivantes :
\begin{enumerate}
\item $H^a(Lj^*E)$ est un $\mathcal{O}_Z$-module localement libre de rang $r$ ;
\item un morphisme $f:Y\to X$ se factorise par $Z$ si et seulement si, pour
tout morphisme $g:Y'\to Y$, le $\mathcal{O}_{Y'}$-module
$H^a(L(f\circ g)^*E)$ est localement libre de rang $r$.
\end{enumerate}
De plus, $j:Z\to X$ est de pr\'esentation finie et
\begin{enumerate}
\item[(3)] si $f:Y\to X$ se factorise sous la forme
$Y\xrightarrow{g}Z\to X$, alors
$H^a(Lf^*E)=g^*H^a(Lj^*E)$ ;
\item[(4)] si $\beta_a(x)\leq r$ pour tout $x\in|X|$, alors $j$ est une
immersion ferm\'ee et, pour tout $f:Y\to X$, les conditions suivantes sont
\'equivalentes :
\begin{enumerate}
\item $f:Y\to X$ se factorise par $Z$ ;
\item $H^a(Lf^*E)$ est un $\mathcal{O}_Y$-module localement libre de rang $r$ ;
\end{enumerate}
et, si $r=1$, elles sont encore \'equivalentes \`a
\begin{enumerate}
\item[(c)] $\mathcal{O}_Y\to
\SheafHom_{\mathcal{O}_Y}(H^a(Lf^*E),H^a(Lf^*E))$ est injectif.
\end{enumerate}
\end{enumerate}
\end{lemma}

\begin{proof}
D\'emonstration omise. Indications :
on se ram\`ene au cas des sch\'emas par localisation \'etale. Voir Cat\'egories
d\'eriv\'ees des sch\'emas, lemme
\ref{perfect-lemma-locally-closed-where-H0-locally-free}.
\end{proof}

\begin{lemma}
\label{lemma-proper-flat-h0}
Soit $S$ un sch\'ema.
Soit $f:X\to Y$ un morphisme d'espaces alg\'ebriques sur $S$. Supposons que
\begin{enumerate}
\item $f$ soit propre, plat et de pr\'esentation finie ;
\item pour tout morphisme $\Spec(k)\to Y$, o\`u $k$ est un corps, on ait
$k=H^0(X_k,\mathcal{O}_{X_k})$.
\end{enumerate}
Alors
\begin{enumerate}
\item[(a)] $f_*\mathcal{O}_X=\mathcal{O}_Y$, et cette \'egalit\'e subsiste apr\`es
tout changement de base ;
\item[(b)] localement pour la topologie \'etale sur $Y$, on a
$$
Rf_*\mathcal{O}_X = \mathcal{O}_Y \oplus P
$$
dans $D(\mathcal{O}_Y)$, o\`u $P$ est parfait et de Tor-amplitude dans
$[1,\infty)$.
\end{enumerate}
\end{lemma}

\begin{proof}
Il suffit de d\'emontrer (a) et (b) localement pour la topologie \'etale sur
$Y$ ; nous pouvons donc supposer, et supposerons, que $Y$ est un sch\'ema
affine. D'apr\`es le th\'eor\`eme de cohomologie et changement de base
(lemme \ref{lemma-flat-proper-perfect-direct-image-general}), le complexe
$E=Rf_*\mathcal{O}_X$ est parfait et sa formation commute \`a tout changement
de base. En particulier, pour $y\in Y$, on a
$H^0(E\otimes^\mathbf{L}\kappa(y))=
H^0(X_y,\mathcal{O}_{X_y})=\kappa(y)$.
Ainsi, $\beta_0(y)\leq1$ pour tout $y\in Y$, avec les notations du lemme
\ref{lemma-jump-loci}. Appliquons le lemme
\ref{lemma-locally-closed-where-H0-locally-free} avec $a=0$ et $r=1$.
On obtient un sous-sch\'ema ferm\'e universel $j:Z\to Y$ tel que
$H^0(Lj^*E)$ soit inversible, caract\'eris\'e par l'\'equivalence de (4)(a),
(b) et (c) du lemme. Puisque la formation de $E$ commute au changement de
base, on a
$$
Lf^*E = R\text{pr}_{1, *}\mathcal{O}_{X \times_Y X}.
$$
Le morphisme $\text{pr}_1:X\times_YX\to X$ poss\`ede une section, \`a savoir le
morphisme diagonal $\Delta$ de $X$ sur $Y$. On obtient donc des morphismes
$$
\mathcal{O}_X \longrightarrow R\text{pr}_{1, *}\mathcal{O}_{X \times_Y X}
\longrightarrow \mathcal{O}_X
$$
dans $D(\mathcal{O}_X)$, dont le compos\'e est l'identit\'e. Par cons\'equent,
$R\text{pr}_{1,*}\mathcal{O}_{X\times_YX}=\mathcal{O}_X\oplus E'$ dans
$D(\mathcal{O}_X)$. Ainsi, $\mathcal{O}_X$ est un facteur direct de
$H^0(Lf^*E)$, et l'\'equivalence de (4)(c) et (4)(a) dans le lemme pr\'ecit\'e
montre que $X\to Y$ se factorise par $Z$. Puisque $\{X\to Y\}$ est un
recouvrement fppf, on a $Z=Y$. Ainsi, $f_*\mathcal{O}_X$ est un
$\mathcal{O}_Y$-module inversible. On conclut que
$\mathcal{O}_Y\to f_*\mathcal{O}_X$ est un isomorphisme, car un morphisme
d'anneaux $A\to B$ tel que $B$ soit inversible comme $A$-module est un
isomorphisme. Les hypoth\`eses \'etant stables par changement de base, (a) est
d\'emontr\'e.

\medskip\noindent
D\'emonstration de (b). Nous avons vu ci-dessus que, pour tout $y\in Y$, le
morphisme $\mathcal{O}_Y\to H^0(E\otimes^\mathbf{L}\kappa(y))$ est surjectif.
Le lemme \ref{more-algebra-lemma-better-cut-complex-in-two} de Compl\'ements
d'alg\`ebre montre donc que, dans un voisinage ouvert de $y$, on a une
d\'ecomposition $Rf_*\mathcal{O}_X=\mathcal{O}_Y\oplus P$.
\end{proof}

\begin{lemma}
\label{lemma-proper-flat-geom-red-connected}
Soit $S$ un sch\'ema.
Soit $f:X\to Y$ un morphisme d'espaces alg\'ebriques sur $S$. Supposons que
\begin{enumerate}
\item $f$ soit propre, plat et de pr\'esentation finie ;
\item les fibres g\'eom\'etriques de $f$ soient r\'eduites et connexes.
\end{enumerate}
Alors $f_*\mathcal{O}_X=\mathcal{O}_Y$, et cette \'egalit\'e subsiste apr\`es
tout changement de base.
\end{lemma}

\begin{proof}
D'apr\`es le lemme \ref{lemma-proper-flat-h0}, il suffit de montrer que
$k=H^0(X_k,\mathcal{O}_{X_k})$ pour tout morphisme $\Spec(k)\to Y$, o\`u $k$
est un corps. Cela r\'esulte du lemme
\ref{spaces-over-fields-lemma-proper-geometrically-reduced-global-sections}
de Espaces sur un corps, et du fait que $X_k$ est g\'eom\'etriquement connexe
et g\'eom\'etriquement r\'eduit.
\end{proof}






\section{Autres applications}
\label{section-other-applications}

\noindent
Nous \'enon\c{c}ons et d\'emontrons dans cette section quelques r\'esultats qui
se d\'eduisent de la th\'eorie d\'evelopp\'ee ci-dessus.

\begin{lemma}
\label{lemma-countable-cohomology}
Soit $S$ un sch\'ema. Soit $X$ un espace alg\'ebrique quasi-compact et
quasi-s\'epar\'e sur $S$. Soit $K$ un objet de
$D_\QCoh(\mathcal{O}_X)$ tel que les faisceaux de cohomologie $H^i(K)$ aient
des ensembles d\'enombrables de sections sur les sch\'emas affines \'etales sur
$X$. Alors, pour tout morphisme \'etale quasi-compact et quasi-s\'epar\'e
$U\to X$ et tout objet parfait $E$ de $D(\mathcal{O}_X)$, les ensembles
$$
H^i(U, K \otimes^\mathbf{L} E),\quad \Ext^i(E|_U, K|_U)
$$
sont d\'enombrables.
\end{lemma}

\begin{proof}
D'apr\`es Cohomologie sur les sites, lemme
\ref{sites-cohomology-lemma-dual-perfect-complex}, il suffit de d\'emontrer le
r\'esultat pour les groupes $H^i(U,K\otimes^\mathbf{L}E)$.
Nous allons employer le principe de r\'ecurrence du lemme
\ref{lemma-induction-principle}.

\medskip\noindent
Lorsque $U=\Spec(A)$ est affine, le r\'esultat d\'ecoule du cas des sch\'emas ;
voir Cat\'egories d\'eriv\'ees des sch\'emas, lemme
\ref{perfect-lemma-countable-cohomology}.

\medskip\noindent
Pour achever la d\'emonstration, il suffit d'\'etablir ceci : si
$(U\subset W,V\to W)$ est un carr\'e distingu\'e \'el\'ementaire et si le
r\'esultat vaut pour $U$, $V$ et $U\times_WV$, alors il vaut pour $W$.
C'est une cons\'equence imm\'ediate de la suite de Mayer--Vietoris ; voir le
lemme \ref{lemma-unbounded-mayer-vietoris}.
\end{proof}

\begin{lemma}
\label{lemma-countable}
Soit $S$ un sch\'ema.
Soit $X$ un espace alg\'ebrique quasi-compact et quasi-s\'epar\'e sur $S$.
Supposons d\'enombrables les ensembles de sections de $\mathcal{O}_X$ sur les
sch\'emas affines \'etales sur $X$. Soit $K$ un objet de
$D_\QCoh(\mathcal{O}_X)$. Les conditions suivantes sont \'equivalentes :
\begin{enumerate}
\item $K=\text{hocolim}\ E_n$, o\`u les $E_n$ sont des objets parfaits de
$D(\mathcal{O}_X)$ ;
\item les faisceaux de cohomologie $H^i(K)$ ont des ensembles d\'enombrables
de sections sur les sch\'emas affines \'etales sur $X$.
\end{enumerate}
\end{lemma}

\begin{proof}
Si (1) est vraie, alors (2) l'est aussi, car la prise des faisceaux de
cohomologie commute aux colimites homotopiques (d'apr\`es Cat\'egories
d\'eriv\'ees, lemme \ref{derived-lemma-cohomology-of-hocolim}) et parce qu'un
complexe parfait est localement isomorphe \`a un complexe fini de
$\mathcal{O}_X$-modules libres de type fini ; il v\'erifie donc (2) par
l'hypoth\`ese faite sur $X$.

\medskip\noindent
Supposons (2).
Choisissons un complexe K-injectif $\mathcal{K}^\bullet$ repr\'esentant $K$.
Choisissons un g\'en\'erateur parfait $E$ de
$D_\QCoh(\mathcal{O}_X)$ et repr\'esentons-le par un complexe K-injectif
$\mathcal{I}^\bullet$. D'apr\`es le th\'eor\`eme
\ref{theorem-DQCoh-is-Ddga} et sa d\'emonstration, il existe une \'equivalence
de cat\'egories triangul\'ees
$F:D_\QCoh(\mathcal{O}_X)\to D(A,\text{d})$, o\`u $(A,\text{d})$ est
l'alg\`ebre diff\'erentielle gradu\'ee
$$
(A, \text{d}) =
\Hom_{\text{Comp}^{dg}(\mathcal{O}_X)}
(\mathcal{I}^\bullet, \mathcal{I}^\bullet),
$$
qui envoie $K$ sur le module diff\'erentiel gradu\'e
$$
M = \Hom_{\text{Comp}^{dg}(\mathcal{O}_X)}
(\mathcal{I}^\bullet, \mathcal{K}^\bullet).
$$
Notons que $H^i(A)=\Ext^i(E,E)$ et que
$H^i(M)=\Ext^i(E,K)$. De plus, puisque $F$ est une \'equivalence, $F$ et son
quasi-inverse commutent aux colimites homotopiques. Il suffit donc d'\'ecrire
$M$ comme colimite homotopique d'objets compacts de $D(A,\text{d})$.
D'apr\`es Alg\`ebre diff\'erentielle gradu\'ee, lemme
\ref{dga-lemma-countable}, il suffit que $\Ext^i(E,E)$ et $\Ext^i(E,K)$ soient
d\'enombrables pour tout $i$. Cela r\'esulte du lemme
\ref{lemma-countable-cohomology}.
\end{proof}

\begin{lemma}
\label{lemma-computing-sections-as-colim}
Soit $A$ un anneau. Soit $f:U\to X$ un morphisme plat d'espaces alg\'ebriques
de pr\'esentation finie sur $A$. Alors
\begin{enumerate}
\item il existe un syst\`eme projectif d'objets parfaits $L_n$ de
$D(\mathcal{O}_X)$ tel que
$$
R\Gamma(U, Lf^*K) = \text{hocolim}\ R\Hom_X(L_n, K)
$$
dans $D(A)$, fonctoriellement en $K\in D_\QCoh(\mathcal{O}_X)$ ;
\item il existe un syst\`eme inductif d'objets parfaits $E_n$ de
$D(\mathcal{O}_X)$ tel que
$$
R\Gamma(U, Lf^*K) = \text{hocolim}\ R\Gamma(X, E_n \otimes^\mathbf{L} K)
$$
dans $D(A)$, fonctoriellement en $K\in D_\QCoh(\mathcal{O}_X)$.
\end{enumerate}
\end{lemma}

\begin{proof}
D'apr\`es le lemme \ref{lemma-cohomology-base-change}, on a, fonctoriellement
en $K$,
$$
R\Gamma(U, Lf^*K) = R\Gamma(X, Rf_*\mathcal{O}_U \otimes^\mathbf{L} K).
$$
Notons que $R\Gamma(X,-)$ commute aux colimites homotopiques, puisqu'il commute
aux sommes directes d'apr\`es le lemme
\ref{lemma-quasi-coherence-pushforward-direct-sums}. De m\^eme,
$-\otimes^\mathbf{L}K$ commute aux colimites d\'eriv\'ees, puisqu'il commute aux
sommes directes (les sommes directes dans $D(\mathcal{O}_X)$ \'etant donn\'ees
par les sommes directes des complexes qui les repr\'esentent). Pour d\'emontrer
(2), il suffit donc d'\'ecrire
$Rf_*\mathcal{O}_U=\text{hocolim}\ E_n$ pour un syst\`eme inductif d'objets
parfaits $E_n$ de $D(\mathcal{O}_X)$. Cela fait, on obtient (1) en posant
$L_n=E_n^\vee$ ; voir Cohomologie sur les sites, lemme
\ref{sites-cohomology-lemma-dual-perfect-complex}.

\medskip\noindent
\'Ecrivons $A=\colim A_i$, avec les $A_i$ de type fini sur $\mathbf{Z}$.
D'apr\`es Limites d'espaces, lemme
\ref{spaces-limits-lemma-descend-finite-presentation}, il existe un $i$ et des
morphismes $U_i\to X_i\to\Spec(A_i)$ de pr\'esentation finie dont le changement
de base \`a $\Spec(A)$ redonne $U\to X\to\Spec(A)$.
Quitte \`a augmenter $i$, on peut supposer $f_i:U_i\to X_i$ plat ; voir
Limites d'espaces, lemme \ref{spaces-limits-lemma-descend-flat}.
D'apr\`es le lemme \ref{lemma-compare-base-change}, l'image inverse d\'eriv\'ee
de $Rf_{i,*}\mathcal{O}_{U_i}$ par $g:X\to X_i$ est \'egale \`a
$Rf_*\mathcal{O}_U$. Puisque $Lg^*$ commute aux colimites d\'eriv\'ees, il
suffit de d\'emontrer le r\'esultat pour $f_i$. Nous pouvons donc supposer $U$
et $X$ de type fini sur $\mathbf{Z}$.

\medskip\noindent
Supposons $f:U\to X$ un morphisme d'espaces alg\'ebriques de type fini sur
$\mathbf{Z}$. Pour achever la d\'emonstration, nous allons montrer que
$Rf_*\mathcal{O}_U$ est une colimite homotopique de complexes parfaits.
Appliquons pour cela le lemme \ref{lemma-countable}. Il suffit donc de montrer
que les faisceaux $R^if_*\mathcal{O}_U$ ont des ensembles d\'enombrables de
sections sur les sch\'emas affines \'etales sur $X$. Cela r\'esulte du lemme
\ref{lemma-countable-cohomology} appliqu\'e au faisceau structural.
\end{proof}









\section{La propri\'et\'e de r\'esolution}
\label{section-resolution-property}

\noindent
Cette section est, pour les espaces alg\'ebriques, l'analogue de Cat\'egories
d\'eriv\'ees des sch\'emas, section
\ref{perfect-section-resolution-property} ; nous invitons \`a lire d'abord
cette derni\`ere. On ne sait pas actuellement si tout espace alg\'ebrique lisse
et propre sur un corps poss\`ede la propri\'et\'e de r\'esolution, ou si cette
assertion est fausse. Si vous connaissez la r\'eponse, veuillez \'ecrire \`a
\href{mailto:stacks.project@gmail.com}{stacks.project@gmail.com}.

\medskip\noindent
Nous pouvons donner la d\'efinition suivante, bien qu'il ne soit gu\`ere utile
de la consid\'erer pour des espaces alg\'ebriques g\'en\'eraux.

\begin{definition}
\label{definition-resolution-property}
Soit $S$ un sch\'ema. Soit $X$ un espace alg\'ebrique sur $S$.
On dit que $X$ poss\`ede la {\it propri\'et\'e de r\'esolution} si tout
$\mathcal{O}_X$-module quasi-coh\'erent de type fini est quotient d'un
$\mathcal{O}_X$-module localement libre de type fini.
\end{definition}

\noindent
Si $X$ est un espace alg\'ebrique quasi-compact et quasi-s\'epar\'e, il suffit de
v\'erifier que tout $\mathcal{O}_X$-module de pr\'esentation finie
(automatiquement quasi-coh\'erent) est quotient d'un $\mathcal{O}_X$-module
localement libre de type fini ; voir Limites d'espaces, lemme
\ref{spaces-limits-lemma-finite-directed-colimit-surjective-maps}.
Si $X$ est un espace alg\'ebrique noeth\'erien, les modules quasi-coh\'erents de
type fini sont exactement les $\mathcal{O}_X$-modules coh\'erents ; voir
Cohomologie des espaces, lemme
\ref{spaces-cohomology-lemma-coherent-Noetherian}.

\begin{lemma}
\label{lemma-resolution-property-ample-relative}
Soit $S$ un sch\'ema.
Soit $f:X\to Y$ un morphisme d'espaces alg\'ebriques sur $S$. Supposons que
\begin{enumerate}
\item $Y$ soit quasi-compact et quasi-s\'epar\'e et poss\`ede la propri\'et\'e de
r\'esolution ;
\item il existe sur $X$ un module inversible $f$-ample (Diviseurs sur les
espaces, d\'efinition
\ref{spaces-divisors-definition-relatively-ample}).
\end{enumerate}
Alors $X$ poss\`ede la propri\'et\'e de r\'esolution.
\end{lemma}

\begin{proof}
Soit $\mathcal{F}$ un $\mathcal{O}_X$-module quasi-coh\'erent de type fini.
Soit $\mathcal{L}$ un module inversible $f$-ample.
Choisissons un sch\'ema affine $V$ et un morphisme \'etale surjectif
$V\to Y$. Posons $U=V\times_YX$. Alors $\mathcal{L}|_U$ est ample sur $U$.
D'apr\`es Propri\'et\'es, proposition
\ref{properties-proposition-characterize-ample}, il existe un nombre fini de
morphismes
$s_i:\mathcal{L}^{\otimes n_i}|_U\to\mathcal{F}|_U$ conjointement surjectifs.
Consid\'erons les $\mathcal{O}_Y$-modules quasi-coh\'erents
$$
\mathcal{H}_n =
f_*(\mathcal{F} \otimes_{\mathcal{O}_X} \mathcal{L}^{\otimes n}).
$$
On peut regarder $s_i$ comme une section sur $V$ du faisceau
$\mathcal{H}_{-n_i}$. Supposons que l'on puisse trouver des
$\mathcal{O}_Y$-modules localement libres de type fini $\mathcal{E}_i$ et des
morphismes $\mathcal{E}_i\to\mathcal{H}_{-n_i}$ dont l'image contienne $s_i$.
Alors les morphismes correspondants
$$
f^*\mathcal{E}_i
\otimes_{\mathcal{O}_X}
\mathcal{L}^{\otimes n_i} \longrightarrow \mathcal{F}
$$
sont conjointement surjectifs, ce qui d\'emontre le lemme. D'apr\`es Limites
d'espaces, lemme \ref{spaces-limits-lemma-directed-colimit-finite-type}, il
existe pour chaque $i$ un sous-module quasi-coh\'erent de type fini
$\mathcal{H}'_i\subset\mathcal{H}_{-n_i}$ qui contient la section $s_i$ sur
$V$. La propri\'et\'e de r\'esolution de $Y$ fournit alors des surjections
$\mathcal{E}_i\to\mathcal{H}'_i$, d'o\`u la conclusion.
\end{proof}

\begin{lemma}
\label{lemma-resolution-property-goes-up-affine}
Soit $S$ un sch\'ema.
Soit $f:X\to Y$ un morphisme affine ou quasi-affine d'espaces alg\'ebriques
sur $S$, avec $Y$ quasi-compact et quasi-s\'epar\'e.
Si $Y$ poss\`ede la propri\'et\'e de r\'esolution, alors $X$ la poss\`ede aussi.
\end{lemma}

\begin{proof}
D'apr\`es Diviseurs sur les espaces, lemme
\ref{spaces-divisors-lemma-quasi-affine-O-ample}, c'est un cas particulier du
lemme \ref{lemma-resolution-property-ample-relative}.
\end{proof}

\noindent
Voici un cas o\`u l'on peut faire descendre la propri\'et\'e de r\'esolution.

\begin{lemma}
\label{lemma-resolution-property-goes-down-finite-flat}
Soit $S$ un sch\'ema.
Soit $f:X\to Y$ un morphisme surjectif, fini et localement libre d'espaces
alg\'ebriques sur $S$.
Si $X$ poss\`ede la propri\'et\'e de r\'esolution, alors $Y$ la poss\`ede aussi.
\end{lemma}

\begin{proof}
La condition signifie que $f$ est affine et que $f_*\mathcal{O}_X$ est un
$\mathcal{O}_Y$-module localement libre de type fini et de rang positif.
Soit $\mathcal{G}$ un $\mathcal{O}_Y$-module quasi-coh\'erent de type fini.
Par hypoth\`ese, il existe une surjection
$\mathcal{E}\to f^*\mathcal{G}$, o\`u $\mathcal{E}$ est un
$\mathcal{O}_X$-module localement libre de type fini. Puisque $f_*$ est exact
(Cohomologie des espaces, section
\ref{spaces-cohomology-section-finite-morphisms}), on obtient une surjection
$$
f_*\mathcal{E}
\longrightarrow
f_*f^*\mathcal{G} = \mathcal{G} \otimes_{\mathcal{O}_Y} f_*\mathcal{O}_X.
$$
En prenant les duaux, on obtient une surjection
$$
f_*\mathcal{E}
\otimes_{\mathcal{O}_Y}
\SheafHom_{\mathcal{O}_Y}(f_*\mathcal{O}_X, \mathcal{O}_Y)
\longrightarrow
\mathcal{G}.
$$
Comme $f_*\mathcal{E}$ est localement libre de type fini, on conclut.
\end{proof}

\noindent
Pour d'autres r\'esultats sur la propri\'et\'e de r\'esolution des espaces
alg\'ebriques, voir Compl\'ements sur les morphismes d'espaces, section
\ref{spaces-more-morphisms-section-resolution-property}.












\section{D\'etection de la bornitude}
\label{section-detecting-boundedness}

\noindent
Nous montrons dans cette section que les g\'en\'erateurs compacts de $D_\QCoh$
d'un espace alg\'ebrique quasi-compact et quasi-s\'epar\'e, construits dans la
section \ref{section-generating}, poss\`edent une propri\'et\'e particuli\`ere.
Nous recommandons de lire d'abord cette section, tr\`es proche de celle-ci.

\begin{lemma}
\label{lemma-bounded-truncation}
Soit $S$ un sch\'ema.
Soit $X$ un espace alg\'ebrique quasi-compact et quasi-s\'epar\'e sur $S$.
Soient $P\in D_{perf}(\mathcal{O}_X)$ et
$E\in D_{\QCoh}(\mathcal{O}_X)$. Soit $a\in\mathbf{Z}$.
Les conditions suivantes sont \'equivalentes :
\begin{enumerate}
\item $\Hom_{D(\mathcal{O}_X)}(P[-i],E)=0$ pour $i\gg0$ ;
\item $\Hom_{D(\mathcal{O}_X)}(P[-i],\tau_{\geq a}E)=0$ pour $i\gg0$.
\end{enumerate}
\end{lemma}

\begin{proof}
En utilisant le triangle $\tau_{<a}E\to E\to\tau_{\geq a}E\to$, on voit
qu'il suffit, pour obtenir l'\'equivalence, de montrer que
$$
\Hom_{D(\mathcal{O}_X)}(P[-i], \tau_{< a} E) =
\Hom_{D(\mathcal{O}_X)}(P, (\tau_{< a} E)[i]) = 0
$$
pour $i\gg0$. Comme $P$ est parfait, cela r\'esulte du lemme
\ref{lemma-ext-from-perfect-into-bounded-QCoh}.
\end{proof}

\begin{lemma}
\label{lemma-bounded-below-truncation}
Soit $S$ un sch\'ema.
Soit $X$ un espace alg\'ebrique quasi-compact et quasi-s\'epar\'e sur $S$.
Soient $P\in D_{perf}(\mathcal{O}_X)$ et
$E\in D_{\QCoh}(\mathcal{O}_X)$. Soit $a\in\mathbf{Z}$.
Les conditions suivantes sont \'equivalentes :
\begin{enumerate}
\item $\Hom_{D(\mathcal{O}_X)}(P[-i],E)=0$ pour $i\ll0$ ;
\item $\Hom_{D(\mathcal{O}_X)}(P[-i],\tau_{\leq a}E)=0$ pour $i\ll0$.
\end{enumerate}
\end{lemma}

\begin{proof}
En utilisant le triangle $\tau_{\leq a}E\to E\to\tau_{>a}E\to$, on voit
qu'il suffit, pour obtenir l'\'equivalence, de montrer que
$$
\Hom_{D(\mathcal{O}_X)}(P[-i], \tau_{> a} E) =
\Hom_{D(\mathcal{O}_X)}(P, (\tau_{> a} E)[i]) = 0
$$
pour $i\ll0$. Comme $P$ est parfait, cela r\'esulte du lemme
\ref{lemma-ext-from-perfect-into-bounded-QCoh}.
\end{proof}

\begin{proposition}
\label{proposition-detecting-bounded-above}
Soit $S$ un sch\'ema.
Soit $X$ un espace alg\'ebrique quasi-compact et quasi-s\'epar\'e sur $S$.
Soit $G\in D_{perf}(\mathcal{O}_X)$ un complexe parfait qui engendre
$D_\QCoh(\mathcal{O}_X)$. Soit $E\in D_\QCoh(\mathcal{O}_X)$.
Les conditions suivantes sont \'equivalentes :
\begin{enumerate}
\item $E\in D^-_\QCoh(\mathcal{O}_X)$ ;
\item $\Hom_{D(\mathcal{O}_X)}(G[-i],E)=0$ pour $i\gg0$ ;
\item $\Ext^i_X(G,E)=0$ pour $i\gg0$ ;
\item $R\Hom_X(G,E)$ appartient \`a $D^-(\mathbf{Z})$ ;
\item $H^i(X,G^\vee\otimes_{\mathcal{O}_X}^\mathbf{L}E)=0$ pour $i\gg0$ ;
\item $R\Gamma(X,G^\vee\otimes_{\mathcal{O}_X}^\mathbf{L}E)$ appartient \`a
$D^-(\mathbf{Z})$ ;
\item pour tout objet parfait $P$ de $D(\mathcal{O}_X)$,
\begin{enumerate}
\item les assertions (2), (3) et (4) valent apr\`es remplacement de $G$ par
$P$ ;
\item $H^i(X,P\otimes_{\mathcal{O}_X}^\mathbf{L}E)=0$ pour $i\gg0$ ;
\item $R\Gamma(X,P\otimes_{\mathcal{O}_X}^\mathbf{L}E)$ appartient \`a
$D^-(\mathbf{Z})$.
\end{enumerate}
\end{enumerate}
\end{proposition}

\begin{proof}
Supposons (1). Comme
$\Hom_{D(\mathcal{O}_X)}(G[-i],E)=\Hom_{D(\mathcal{O}_X)}(G,E[i])$, ce groupe
est nul pour $i\gg0$ d'apr\`es le lemme
\ref{lemma-ext-from-perfect-into-bounded-QCoh}. Ainsi, (1) implique (2).

\medskip\noindent
Les assertions (2), (3) et (4) sont \'equivalentes d'apr\`es Cohomologie sur
les sites, section \ref{sites-cohomology-section-global-RHom}.
Les assertions (5) et (6) sont \'equivalentes puisque, par d\'efinition,
$H^i(X,-)=H^i(R\Gamma(X,-))$. Les conditions \'equivalentes (2), (3), (4)
\'equivalent aux conditions \'equivalentes (5), (6) d'apr\`es Cohomologie sur
les sites, lemme \ref{sites-cohomology-lemma-dual-perfect-complex}, et
l'\'egalit\'e $(G[-i])^\vee=G^\vee[i]$.

\medskip\noindent
Il est clair que (7) implique (2). R\'eciproquement, montrons que les conditions
\'equivalentes (2)--(6) impliquent (7). Rappelons que $G$ est un g\'en\'erateur
classique de $D_{perf}(\mathcal{O}_X)$ d'apr\`es la remarque
\ref{remark-classical-generator}. Pour $P\in D_{perf}(\mathcal{O}_X)$, notons
$T(P)$ l'assertion selon laquelle $R\Hom_X(P,E)$ appartient \`a
$D^-(\mathbf{Z})$. Il est clair que $T$ passe aux sommes directes et aux
facteurs directs, qu'elle v\'erifie la propri\'et\'e des deux sur trois pour les
triangles distingu\'es, et qu'elle est stable par d\'ecalage. La remarque
\ref{derived-remark-check-on-generator} de Cat\'egories d\'eriv\'ees montre donc
que (4) implique $T(P)$ pour tout $P\in D_{perf}(\mathcal{O}_X)$. Le m\^eme
raisonnement s'applique \`a toutes les autres propri\'et\'es ; pour (7)(b) et
(7)(c), on utilise aussi que $P\mapsto P^\vee$ est une auto-\'equivalence de
$D_{perf}(\mathcal{O}_X)$. Nous omettons ce petit d\'etail.

\medskip\noindent
Nous allons montrer que les conditions \'equivalentes (2)--(7) impliquent (1)
en employant le principe de r\'ecurrence du lemme
\ref{lemma-induction-principle}.

\medskip\noindent
Supposons d'abord $X$ affine. L'implication (2)--(7) $\Rightarrow$ (1) d\'ecoule
du cas des sch\'emas ; voir Cat\'egories d\'eriv\'ees des sch\'emas, proposition
\ref{perfect-proposition-detecting-bounded-above}.

\medskip\noindent
Supposons maintenant que $(U\subset X,j:V\to X)$ soit un carr\'e distingu\'e
\'el\'ementaire d'espaces alg\'ebriques quasi-compacts et quasi-s\'epar\'es sur
$S$, et que l'implication (2)--(7) $\Rightarrow$ (1) soit connue pour $U$, $V$
et $U\times_XV$. Il reste \`a la d\'emontrer pour $X$. Supposons donc que
$E\in D_\QCoh(\mathcal{O}_X)$ v\'erifie (2)--(7). D'apr\`es le lemme
\ref{lemma-direct-summand-of-a-restriction} et le th\'eor\`eme
\ref{theorem-bondal-van-den-Bergh}, il existe sur $X$ un complexe parfait $Q$
tel que $Q|_U$ engendre $D_\QCoh(\mathcal{O}_U)$.

\medskip\noindent
\'Ecrivons $V=\Spec(A)$. Soit $Z\subset V$ le sous-sch\'ema ferm\'e r\'eduit,
image inverse de $X\setminus U$, qui s'envoie isomorphiquement sur celui-ci
(voir la d\'efinition \ref{definition-elementary-distinguished-square}).
C'est une partie ferm\'ee r\'etrocompacte de $V$. Choisissons
$f_1,\ldots,f_r\in A$ tels que $Z=V(f_1,\ldots,f_r)$. Soit
$K\in D(\mathcal{O}_V)$ l'objet parfait correspondant au complexe de Koszul
associ\'e \`a $f_1,\ldots,f_r$ sur $A$. Puisque $K$ est \`a support dans $Z$,
son image directe $K'=Rj_*K$ est un objet parfait de $D(\mathcal{O}_X)$ dont
la restriction \`a $V$ est $K$ (voir les lemmes
\ref{lemma-pushforward-perfect} et
\ref{lemma-pushforward-with-support-in-open}). Par hypoth\`ese,
$R\Hom_{\mathcal{O}_X}(Q,E)$ et $R\Hom_{\mathcal{O}_X}(K',E)$ sont born\'es
sup\'erieurement.

\medskip\noindent
D'apr\`es le lemme \ref{lemma-pushforward-with-support-in-open}, on a
$K'=j_!K$, et donc
$$
\Hom_{D(\mathcal{O}_X)}(K'[-i], E) = \Hom_{D(\mathcal{O}_V)}(K[-i], E|_V) = 0
$$
pour $i\gg0$. On peut donc appliquer \`a $E|_V$ le lemme
\ref{perfect-lemma-orthogonal-koszul-first-variant} de Cat\'egories d\'eriv\'ees
des sch\'emas. Il existe un entier $a$ tel que
$\tau_{\geq a}(E|_V)=
\tau_{\geq a}R(U\times_XV\to V)_*(E|_{U\times_XV})$.
Alors $\tau_{\geq a}E=\tau_{\geq a}R(U\to X)_*(E|_U)$ ; on le v\'erifie en
restreignant le morphisme canonique \`a $U$ et \`a $V$. En appliquant deux fois
le lemme \ref{lemma-bounded-truncation}, on obtient
$$
\Hom_{D(\mathcal{O}_U)}(Q|_U [-i], E|_U) =
\Hom_{D(\mathcal{O}_X)}(Q[-i], R (U \to X)_* (E|_U)) = 0
$$
pour $i\gg0$. Puisque la proposition vaut pour $U$ et le g\'en\'erateur
$Q|_U$, on a $E|_U\in D^-_\QCoh(\mathcal{O}_U)$. Le foncteur
$R(U\to X)_*$ pr\'eserve $D^-_\QCoh$ d'apr\`es le lemme
\ref{lemma-quasi-coherence-direct-image} ; ainsi,
$\tau_{\geq a}E\in D^-_\QCoh(\mathcal{O}_X)$, puis
$E\in D^-_\QCoh(\mathcal{O}_X)$.
\end{proof}

\begin{proposition}
\label{proposition-detecting-bounded-below}
Soit $S$ un sch\'ema.
Soit $X$ un espace alg\'ebrique quasi-compact et quasi-s\'epar\'e sur $S$.
Soit $G\in D_{perf}(\mathcal{O}_X)$ un complexe parfait qui engendre
$D_\QCoh(\mathcal{O}_X)$. Soit $E\in D_\QCoh(\mathcal{O}_X)$.
Les conditions suivantes sont \'equivalentes :
\begin{enumerate}
\item $E\in D^+_\QCoh(\mathcal{O}_X)$ ;
\item $\Hom_{D(\mathcal{O}_X)}(G[-i],E)=0$ pour $i\ll0$ ;
\item $\Ext^i_X(G,E)=0$ pour $i\ll0$ ;
\item $R\Hom_X(G,E)$ appartient \`a $D^+(\mathbf{Z})$ ;
\item $H^i(X,G^\vee\otimes_{\mathcal{O}_X}^\mathbf{L}E)=0$ pour $i\ll0$ ;
\item $R\Gamma(X,G^\vee\otimes_{\mathcal{O}_X}^\mathbf{L}E)$ appartient \`a
$D^+(\mathbf{Z})$ ;
\item pour tout objet parfait $P$ de $D(\mathcal{O}_X)$,
\begin{enumerate}
\item les assertions (2), (3) et (4) valent apr\`es remplacement de $G$ par
$P$ ;
\item $H^i(X,P\otimes_{\mathcal{O}_X}^\mathbf{L}E)=0$ pour $i\ll0$ ;
\item $R\Gamma(X,P\otimes_{\mathcal{O}_X}^\mathbf{L}E)$ appartient \`a
$D^+(\mathbf{Z})$.
\end{enumerate}
\end{enumerate}
\end{proposition}

\begin{proof}
Supposons (1). Comme
$\Hom_{D(\mathcal{O}_X)}(G[-i],E)=\Hom_{D(\mathcal{O}_X)}(G,E[i])$, ce groupe
est nul pour $i\ll0$ d'apr\`es le lemme
\ref{lemma-ext-from-perfect-into-bounded-QCoh}. Ainsi, (1) implique (2).

\medskip\noindent
Les assertions (2), (3) et (4) sont \'equivalentes d'apr\`es Cohomologie sur
les sites, section \ref{sites-cohomology-section-global-RHom}.
Les assertions (5) et (6) sont \'equivalentes puisque, par d\'efinition,
$H^i(X,-)=H^i(R\Gamma(X,-))$. Les conditions \'equivalentes (2), (3), (4)
\'equivalent aux conditions \'equivalentes (5), (6) d'apr\`es Cohomologie sur
les sites, lemme \ref{sites-cohomology-lemma-dual-perfect-complex}, et
l'\'egalit\'e $(G[-i])^\vee=G^\vee[i]$.

\medskip\noindent
Il est clair que (7) implique (2). R\'eciproquement, montrons que les conditions
\'equivalentes (2)--(6) impliquent (7). Rappelons que $G$ est un g\'en\'erateur
classique de $D_{perf}(\mathcal{O}_X)$ d'apr\`es la remarque
\ref{remark-classical-generator}. Pour $P\in D_{perf}(\mathcal{O}_X)$, notons
$T(P)$ l'assertion selon laquelle $R\Hom_X(P,E)$ appartient \`a
$D^+(\mathbf{Z})$. Il est clair que $T$ passe aux sommes directes et aux
facteurs directs, qu'elle v\'erifie la propri\'et\'e des deux sur trois pour les
triangles distingu\'es, et qu'elle est stable par d\'ecalage. La remarque
\ref{derived-remark-check-on-generator} de Cat\'egories d\'eriv\'ees montre donc
que (4) implique $T(P)$ pour tout $P\in D_{perf}(\mathcal{O}_X)$. Le m\^eme
raisonnement s'applique \`a toutes les autres propri\'et\'es ; pour (7)(b) et
(7)(c), on utilise aussi que $P\mapsto P^\vee$ est une auto-\'equivalence de
$D_{perf}(\mathcal{O}_X)$. Nous omettons ce petit d\'etail.

\medskip\noindent
Nous allons montrer que les conditions \'equivalentes (2)--(7) impliquent (1)
en employant le principe de r\'ecurrence du lemme
\ref{lemma-induction-principle}.

\medskip\noindent
Supposons d'abord $X$ affine. L'implication (2)--(7) $\Rightarrow$ (1) d\'ecoule
du cas des sch\'emas ; voir Cat\'egories d\'eriv\'ees des sch\'emas, proposition
\ref{perfect-proposition-detecting-bounded-below}.

\medskip\noindent
Supposons maintenant que $(U\subset X,j:V\to X)$ soit un carr\'e distingu\'e
\'el\'ementaire d'espaces alg\'ebriques quasi-compacts et quasi-s\'epar\'es sur
$S$, et que l'implication (2)--(7) $\Rightarrow$ (1) soit connue pour $U$, $V$
et $U\times_XV$. Il reste \`a la d\'emontrer pour $X$. Supposons donc que
$E\in D_\QCoh(\mathcal{O}_X)$ v\'erifie (2)--(7). D'apr\`es le lemme
\ref{lemma-direct-summand-of-a-restriction} et le th\'eor\`eme
\ref{theorem-bondal-van-den-Bergh}, il existe sur $X$ un complexe parfait $Q$
tel que $Q|_U$ engendre $D_\QCoh(\mathcal{O}_U)$.

\medskip\noindent
\'Ecrivons $V=\Spec(A)$. Soit $Z\subset V$ le sous-sch\'ema ferm\'e r\'eduit,
image inverse de $X\setminus U$, qui s'envoie isomorphiquement sur celui-ci
(voir la d\'efinition \ref{definition-elementary-distinguished-square}).
C'est une partie ferm\'ee r\'etrocompacte de $V$. Choisissons
$f_1,\ldots,f_r\in A$ tels que $Z=V(f_1,\ldots,f_r)$. Soit
$K\in D(\mathcal{O}_V)$ l'objet parfait correspondant au complexe de Koszul
associ\'e \`a $f_1,\ldots,f_r$ sur $A$. Puisque $K$ est \`a support dans $Z$,
son image directe $K'=Rj_*K$ est un objet parfait de $D(\mathcal{O}_X)$ dont
la restriction \`a $V$ est $K$ (voir les lemmes
\ref{lemma-pushforward-perfect} et
\ref{lemma-pushforward-with-support-in-open}). Par hypoth\`ese,
$R\Hom_{\mathcal{O}_X}(Q,E)$ et $R\Hom_{\mathcal{O}_X}(K',E)$ sont born\'es
inf\'erieurement.

\medskip\noindent
D'apr\`es le lemme \ref{lemma-pushforward-with-support-in-open}, on a
$K'=j_!K$, et donc
$$
\Hom_{D(\mathcal{O}_X)}(K'[-i], E) = \Hom_{D(\mathcal{O}_V)}(K[-i], E|_V) = 0
$$
pour $i\ll0$. On peut donc appliquer \`a $E|_V$ le lemme
\ref{perfect-lemma-orthogonal-koszul-second-variant} de Cat\'egories d\'eriv\'ees
des sch\'emas. Il existe un entier $a$ tel que
$\tau_{\leq a}(E|_V)=
\tau_{\leq a}R(U\times_XV\to V)_*(E|_{U\times_XV})$.
Alors $\tau_{\leq a}E=\tau_{\leq a}R(U\to X)_*(E|_U)$ ; on le v\'erifie en
restreignant le morphisme canonique \`a $U$ et \`a $V$. En appliquant deux fois
le lemme \ref{lemma-bounded-below-truncation}, on obtient
$$
\Hom_{D(\mathcal{O}_U)}(Q|_U [-i], E|_U) =
\Hom_{D(\mathcal{O}_X)}(Q[-i], R (U \to X)_* (E|_U)) = 0
$$
pour $i\ll0$. Puisque la proposition vaut pour $U$ et le g\'en\'erateur
$Q|_U$, on a $E|_U\in D^+_\QCoh(\mathcal{O}_U)$. Le foncteur
$R(U\to X)_*$ pr\'eserve $D^+_\QCoh$ d'apr\`es le lemme
\ref{lemma-quasi-coherence-direct-image} ; ainsi,
$\tau_{\leq a}E\in D^+_\QCoh(\mathcal{O}_X)$, puis
$E\in D^+_\QCoh(\mathcal{O}_X)$.
\end{proof}










\section{Objets quasi-coh\'erents dans la cat\'egorie d\'eriv\'ee}
\label{section-QC}

\noindent
Soit $S$ un sch\'ema. Soit $X$ un espace alg\'ebrique sur $S$.
Rappelons que $X_{affine,\etale}$ d\'esigne la cat\'egorie des objets affines de
$X_\etale$, munie de la topologie d\'efinie par les recouvrements \'etales
standard ; voir Propri\'et\'es des espaces, d\'efinition
\ref{spaces-properties-definition-affine-etale-site}.
Rappelons aussi que le topos de $X_{affine,\etale}$ est le petit topos \'etale
de $X$ ; voir Propri\'et\'es des espaces, lemme
\ref{spaces-properties-lemma-alternative}. Le site $X_\etale$ est muni d'un
faisceau structural $\mathcal{O}_X$, dont nous notons encore
$\mathcal{O}_X$ la restriction \`a $X_{affine,\etale}$. Il existe alors une
\'equivalence de topos annel\'es
$$
(\Sh(X_{affine, \etale}), \mathcal{O}_X)
\longrightarrow
(\Sh(X_\etale), \mathcal{O}_X).
$$
Voir Descente sur les espaces, \'equation
(\ref{spaces-descent-equation-alternative-small-ringed}), ainsi que la discussion de
Descente sur les espaces, section
\ref{spaces-descent-section-alternative-quasi-coherent}.

\medskip\noindent
Dans cette section, nous notons $X_{affine}$ la cat\'egorie sous-jacente \`a
$X_{affine,\etale}$, munie de la topologie chaotique, c'est-\`a-dire telle que
les faisceaux s'identifient aux pr\'efaisceaux. En particulier, le faisceau
structural $\mathcal{O}_X$ devient aussi un faisceau sur $X_{affine}$.
On obtient un morphisme de sites annel\'es
$$
\epsilon :
(X_{affine, \etale}, \mathcal{O}_X)
\longrightarrow
(X_{affine}, \mathcal{O}_X)
$$
comme dans Cohomologie sur les sites, section
\ref{sites-cohomology-section-compare}. Nous identifierons dans cette section
$D_\QCoh(\mathcal{O}_X)$ \`a la cat\'egorie
$\mathit{QC}(X_{affine},\mathcal{O}_X)$ introduite dans Cohomologie sur les
sites, section \ref{sites-cohomology-section-modules-cohomology}.

\begin{lemma}
\label{lemma-DQCoh-alternative-small}
Dans la situation ci-dessus, il existe des \'equivalences exactes canoniques
entre les cat\'egories triangul\'ees suivantes :
\begin{enumerate}
\item $D_\QCoh(\mathcal{O}_X)$ ;
\item $D_\QCoh(X_{affine, \etale}, \mathcal{O}_X)$ ;
\item $D_\QCoh(X_{affine}, \mathcal{O}_X)$ ;
\item $\mathit{QC}(X_{affine}, \mathcal{O}_X)$.
\end{enumerate}
\end{lemma}

\begin{proof}
Si $U\to V\to X$ sont des morphismes \'etales et si $U$ et $V$ sont affines,
alors le morphisme d'anneaux
$\mathcal{O}_X(V)\to\mathcal{O}_X(U)$ est plat. L'\'equivalence de (3) et (4)
est donc un cas particulier de Cohomologie sur les sites, lemme
\ref{sites-cohomology-lemma-cartesian-flat-quasi-coherent}, dont la
d\'emonstration \'eclaire aussi l'\'enonc\'e.

\medskip\noindent
La discussion qui pr\'ec\`ede le lemme montre qu'il existe une \'equivalence de
topos annel\'es
$(\Sh(X_{affine,\etale}),\mathcal{O}_X)\to
(\Sh(X_\etale),\mathcal{O}_X)$, donc une \'equivalence entre les cat\'egories
ab\'eliennes de modules. Puisque la notion de module quasi-coh\'erent est
intrins\`eque (Modules sur les sites, lemme
\ref{sites-modules-lemma-special-locally-free}), cette \'equivalence pr\'eserve
les sous-cat\'egories de modules quasi-coh\'erents. On obtient ainsi une
\'equivalence exacte canonique entre les cat\'egories triangul\'ees de (1) et
(2).

\medskip\noindent
Pour obtenir une \'equivalence exacte entre les cat\'egories triangul\'ees de
(2) et (3), appliquons Cohomologie sur les sites, lemme
\ref{sites-cohomology-lemma-compare-topologies-derived-adequate-modules}, au
morphisme
$\epsilon:(X_{affine,\etale},\mathcal{O}_X)\to
(X_{affine},\mathcal{O}_X)$ ci-dessus. Prenons
$\mathcal{B}=\Ob(X_{affine})$ et pour
$\mathcal{A}\subset\textit{PMod}(X_{affine},\mathcal{O})$ la sous-cat\'egorie
pleine des pr\'efaisceaux $\mathcal{F}$ tels que
$\mathcal{F}(V)\otimes_{\mathcal{O}_X(V)}\mathcal{O}_X(U)	o\mathcal{F}(U)$
soit un isomorphisme. D'apr\`es Descente sur les espaces, lemme
\ref{spaces-descent-lemma-quasi-coherent-alternative-small}, les objets de
$\mathcal{A}$ sont exactement les faisceaux pour la topologie \'etale qui sont
des modules quasi-coh\'erents sur
$(X_{affine,\etale},\mathcal{O}_X)$. D'autre part, d'apr\`es Modules sur les
sites, lemme \ref{sites-modules-lemma-chaotic-quasi-coherent}, les objets de
$\mathcal{A}$ sont exactement les modules quasi-coh\'erents sur
$(X_{affine},\mathcal{O}_X)$, c'est-\`a-dire pour la topologie chaotique. Si les
hypoth\`eses du lemme
\ref{sites-cohomology-lemma-compare-topologies-derived-adequate-modules} de
Cohomologie sur les sites sont satisfaites, $\epsilon^*$ et $R\epsilon_*$
fournissent donc bien l'\'equivalence canonique entre les cat\'egories de (2) et
(3).

\medskip\noindent
Il reste \`a v\'erifier quatre conditions :
\begin{enumerate}
\item Tout objet de $\mathcal{A}$ est un faisceau pour la topologie \'etale.
Nous l'avons vu ci-dessus.
\item $\mathcal{A}$ est une sous-cat\'egorie de Serre faible de
$\textit{Mod}(X_{affine,\etale},\mathcal{O}_X)$. Nous avons vu ci-dessus que
$\mathcal{A}=\QCoh(X_{affine,\etale},\mathcal{O}_X)$ et que, par
l'\'equivalence
$\textit{Mod}(X_{affine,\etale},\mathcal{O})=
\textit{Mod}(X_\etale,\mathcal{O}_X)$, ses objets correspondent aux modules
quasi-coh\'erents sur $X$. L'assertion r\'esulte donc de Propri\'et\'es des
espaces, lemme \ref{spaces-properties-lemma-properties-quasi-coherent}, et
d'Homologie, lemme
\ref{homology-lemma-characterize-weak-serre-subcategory}.
\item Tout objet de $X_{affine}$ poss\`ede, pour la topologie chaotique, un
recouvrement dont les membres appartiennent \`a $\mathcal{B}$. Cela tient au
fait que $\mathcal{B}$ contient tous les objets.
\item Pour tout objet $U$ de $X_{affine}$ et tout $\mathcal{F}$ de
$\mathcal{A}$, on a $H^p_{Zar}(U,\mathcal{F})=0$ pour $p>0$. Cela r\'esulte de
l'annulation de la cohomologie des modules quasi-coh\'erents sur les sch\'emas
affines ; voir la discussion de Cohomologie des espaces, section
\ref{spaces-cohomology-section-higher-direct-image}, et Cohomologie des
sch\'emas, lemme
\ref{coherent-lemma-quasi-coherent-affine-cohomology-zero}.
\end{enumerate}
La d\'emonstration est achev\'ee.
\end{proof}

\begin{remark}
\label{remark-QC-compare}
Soit $S$ un sch\'ema. Soit $X$ un espace alg\'ebrique sur $S$.
Nous montrerons plus loin que
$\mathit{QC}((\textit{Aff}/X),\mathcal{O})$ est aussi canoniquement
\'equivalente \`a $D_\QCoh(\mathcal{O}_X)$. Voir Faisceaux sur les champs,
proposition \ref{stacks-sheaves-proposition-QC-compare}.
\end{remark}






\begin{multicols}{2}[\section{Autres chapitres}]
\noindent
Préliminaires
\begin{enumerate}
\item \hyperref[introduction-section-phantom]{Introduction}
\item \hyperref[conventions-section-phantom]{Conventions}
\item \hyperref[sets-section-phantom]{Théorie des ensembles}
\item \hyperref[categories-section-phantom]{Catégories}
\item \hyperref[topology-section-phantom]{Topologie}
\item \hyperref[sheaves-section-phantom]{Faisceaux sur les espaces}
\item \hyperref[sites-section-phantom]{Sites et faisceaux}
\item \hyperref[stacks-section-phantom]{Champs}
\item \hyperref[fields-section-phantom]{Corps}
\item \hyperref[algebra-section-phantom]{Algèbre commutative}
\item \hyperref[brauer-section-phantom]{Groupes de Brauer}
\item \hyperref[homology-section-phantom]{Algèbre homologique}
\item \hyperref[derived-section-phantom]{Catégories dérivées}
\item \hyperref[simplicial-section-phantom]{Méthodes simpliciales}
\item \hyperref[more-algebra-section-phantom]{Compléments d'algèbre}
\item \hyperref[smoothing-section-phantom]{Lissage des morphismes d'anneaux}
\item \hyperref[modules-section-phantom]{Faisceaux de modules}
\item \hyperref[sites-modules-section-phantom]{Modules sur les sites}
\item \hyperref[injectives-section-phantom]{Objets injectifs}
\item \hyperref[cohomology-section-phantom]{Cohomologie des faisceaux}
\item \hyperref[sites-cohomology-section-phantom]{Cohomologie sur les sites}
\item \hyperref[dga-section-phantom]{Algèbre différentielle graduée}
\item \hyperref[dpa-section-phantom]{Algèbre à puissances divisées}
\item \hyperref[sdga-section-phantom]{Faisceaux différentiels gradués}
\item \hyperref[hypercovering-section-phantom]{Hyperrecouvrements}
\end{enumerate}
Schémas
\begin{enumerate}
\setcounter{enumi}{25}
\item \hyperref[schemes-section-phantom]{Schémas}
\item \hyperref[constructions-section-phantom]{Constructions de schémas}
\item \hyperref[properties-section-phantom]{Propriétés des schémas}
\item \hyperref[morphisms-section-phantom]{Morphismes de schémas}
\item \hyperref[coherent-section-phantom]{Cohomologie des schémas}
\item \hyperref[divisors-section-phantom]{Diviseurs}
\item \hyperref[limits-section-phantom]{Limites de schémas}
\item \hyperref[varieties-section-phantom]{Variétés}
\item \hyperref[topologies-section-phantom]{Topologies sur les schémas}
\item \hyperref[descent-section-phantom]{Descente}
\item \hyperref[perfect-section-phantom]{Catégories dérivées de schémas}
\item \hyperref[more-morphisms-section-phantom]{Compléments sur les morphismes}
\item \hyperref[flat-section-phantom]{Compléments sur la platitude}
\item \hyperref[groupoids-section-phantom]{Schémas en groupoïdes}
\item \hyperref[more-groupoids-section-phantom]{Compléments sur les schémas en groupoïdes}
\item \hyperref[etale-section-phantom]{Morphismes étales de schémas}
\end{enumerate}
Thèmes de théorie des schémas
\begin{enumerate}
\setcounter{enumi}{41}
\item \hyperref[chow-section-phantom]{Homologie de Chow}
\item \hyperref[intersection-section-phantom]{Théorie de l'intersection}
\item \hyperref[pic-section-phantom]{Schémas de Picard des courbes}
\item \hyperref[weil-section-phantom]{Théories cohomologiques de Weil}
\item \hyperref[adequate-section-phantom]{Modules adéquats}
\item \hyperref[dualizing-section-phantom]{Complexes dualisants}
\item \hyperref[duality-section-phantom]{Dualité pour les schémas}
\item \hyperref[discriminant-section-phantom]{Discriminants et différentes}
\item \hyperref[derham-section-phantom]{Cohomologie de de Rham}
\item \hyperref[local-cohomology-section-phantom]{Cohomologie locale}
\item \hyperref[algebraization-section-phantom]{Géométrie algébrique et formelle}
\item \hyperref[curves-section-phantom]{Courbes algébriques}
\item \hyperref[resolve-section-phantom]{Résolution des surfaces}
\item \hyperref[models-section-phantom]{Réduction semi-stable}
\item \hyperref[functors-section-phantom]{Foncteurs et morphismes}
\item \hyperref[equiv-section-phantom]{Catégories dérivées de variétés}
\item \hyperref[pione-section-phantom]{Groupes fondamentaux de schémas}
\item \hyperref[etale-cohomology-section-phantom]{Cohomologie étale}
\item \hyperref[crystalline-section-phantom]{Cohomologie cristalline}
\item \hyperref[proetale-section-phantom]{Cohomologie pro-étale}
\item \hyperref[relative-cycles-section-phantom]{Cycles relatifs}
\item \hyperref[more-etale-section-phantom]{Compléments de cohomologie étale}
\item \hyperref[trace-section-phantom]{La formule des traces}
\end{enumerate}
Espaces algébriques
\begin{enumerate}
\setcounter{enumi}{64}
\item \hyperref[spaces-section-phantom]{Espaces algébriques}
\item \hyperref[spaces-properties-section-phantom]{Propriétés des espaces algébriques}
\item \hyperref[spaces-morphisms-section-phantom]{Morphismes d'espaces algébriques}
\item \hyperref[decent-spaces-section-phantom]{Espaces algébriques décents}
\item \hyperref[spaces-cohomology-section-phantom]{Cohomologie des espaces algébriques}
\item \hyperref[spaces-limits-section-phantom]{Limites d'espaces algébriques}
\item \hyperref[spaces-divisors-section-phantom]{Diviseurs sur les espaces algébriques}
\item \hyperref[spaces-over-fields-section-phantom]{Espaces algébriques sur des corps}
\item \hyperref[spaces-topologies-section-phantom]{Topologies sur les espaces algébriques}
\item \hyperref[spaces-descent-section-phantom]{Descente et espaces algébriques}
\item \hyperref[spaces-perfect-section-phantom]{Catégories dérivées d'espaces}
\item \hyperref[spaces-more-morphisms-section-phantom]{Compléments sur les morphismes d'espaces}
\item \hyperref[spaces-flat-section-phantom]{Platitude sur les espaces algébriques}
\item \hyperref[spaces-groupoids-section-phantom]{Groupoïdes dans les espaces algébriques}
\item \hyperref[spaces-more-groupoids-section-phantom]{Compléments sur les groupoïdes dans les espaces}
\item \hyperref[bootstrap-section-phantom]{Amorçage}
\item \hyperref[spaces-pushouts-section-phantom]{Sommes amalgamées d'espaces algébriques}
\end{enumerate}
Thèmes de géométrie
\begin{enumerate}
\setcounter{enumi}{81}
\item \hyperref[spaces-chow-section-phantom]{Groupes de Chow des espaces}
\item \hyperref[groupoids-quotients-section-phantom]{Quotients de groupoïdes}
\item \hyperref[spaces-more-cohomology-section-phantom]{Compléments sur la cohomologie des espaces}
\item \hyperref[spaces-simplicial-section-phantom]{Espaces simpliciaux}
\item \hyperref[spaces-duality-section-phantom]{Dualité pour les espaces}
\item \hyperref[formal-spaces-section-phantom]{Espaces algébriques formels}
\item \hyperref[restricted-section-phantom]{Algébrisation des espaces formels}
\item \hyperref[spaces-resolve-section-phantom]{Résolution des surfaces revisitée}
\end{enumerate}
Théorie des déformations
\begin{enumerate}
\setcounter{enumi}{89}
\item \hyperref[formal-defos-section-phantom]{Théorie formelle des déformations}
\item \hyperref[defos-section-phantom]{Théorie des déformations}
\item \hyperref[cotangent-section-phantom]{Le complexe cotangent}
\item \hyperref[examples-defos-section-phantom]{Problèmes de déformation}
\end{enumerate}
Champs algébriques
\begin{enumerate}
\setcounter{enumi}{93}
\item \hyperref[algebraic-section-phantom]{Champs algébriques}
\item \hyperref[examples-stacks-section-phantom]{Exemples de champs}
\item \hyperref[stacks-sheaves-section-phantom]{Faisceaux sur les champs algébriques}
\item \hyperref[criteria-section-phantom]{Critères de représentabilité}
\item \hyperref[artin-section-phantom]{Axiomes d'Artin}
\item \hyperref[quot-section-phantom]{Espaces de Quot et de Hilbert}
\item \hyperref[stacks-properties-section-phantom]{Propriétés des champs algébriques}
\item \hyperref[stacks-morphisms-section-phantom]{Morphismes de champs algébriques}
\item \hyperref[stacks-limits-section-phantom]{Limites de champs algébriques}
\item \hyperref[stacks-cohomology-section-phantom]{Cohomologie des champs algébriques}
\item \hyperref[stacks-perfect-section-phantom]{Catégories dérivées de champs}
\item \hyperref[stacks-introduction-section-phantom]{Introduction aux champs algébriques}
\item \hyperref[stacks-more-morphisms-section-phantom]{Compléments sur les morphismes de champs}
\item \hyperref[stacks-geometry-section-phantom]{La géométrie des champs}
\end{enumerate}
Thèmes de théorie des espaces de modules
\begin{enumerate}
\setcounter{enumi}{107}
\item \hyperref[moduli-section-phantom]{Champs de modules}
\item \hyperref[moduli-curves-section-phantom]{Espaces de modules de courbes}
\end{enumerate}
Divers
\begin{enumerate}
\setcounter{enumi}{109}
\item \hyperref[examples-section-phantom]{Exemples}
\item \hyperref[exercises-section-phantom]{Exercices}
\item \hyperref[guide-section-phantom]{Guide bibliographique}
\item \hyperref[desirables-section-phantom]{Desiderata}
\item \hyperref[coding-section-phantom]{Style de programmation}
\item \hyperref[obsolete-section-phantom]{Obsolète}
\item \hyperref[fdl-section-phantom]{Licence de documentation libre GNU}
\item \hyperref[index-section-phantom]{Index généré automatiquement}
\end{enumerate}
\end{multicols}


\bibliography{my}
\bibliographystyle{amsalpha}

\end{document}
